% !TEX program = xelatex
% 中文译本：Tryggvason, Scardovelli and Zaleski,
% Direct Numerical Simulations of Gas-Liquid Multiphase Flows (2011).
\documentclass[11pt,a4paper,openany,fontset=none]{ctexbook}
\usepackage[margin=2.55cm,headheight=15pt]{geometry}
\usepackage{microtype}
\usepackage{fancyhdr}
\usepackage{hyperref}
\usepackage{xurl}
\usepackage{titlesec}
\usepackage{amsmath,amssymb,amsfonts,bm,mathtools}
\usepackage[export]{adjustbox}
\usepackage{enumitem}
\usepackage{graphicx}
\usepackage{caption}
\usepackage{longtable,booktabs,array}
\usepackage{needspace}
\setmainfont[AutoFakeBold=2,AutoFakeSlant=.2]{Arial Unicode MS}
\setCJKmainfont{Songti SC}
\setCJKsansfont{Songti SC}
\setCJKmonofont{Songti SC}
\setlength{\parindent}{2em}
\setlength{\parskip}{0.62em}
\setlength{\emergencystretch}{3em}
\linespread{1.32}\selectfont
\xeCJKsetup{CJKglue=\hskip 0.08em plus 0.04em minus 0.03em}
\setlist[itemize]{leftmargin=2.8em,itemsep=0.32em,topsep=0.48em}
\titleformat{\chapter}[display]{\bfseries\Huge\centering}{第\thechapter 章}{0.75em}{}
\titleformat{\section}{\large\bfseries}{}{0pt}{}
\titleformat{\subsection}{\normalsize\bfseries}{}{0pt}{}
\titlespacing*{\section}{0pt}{1.45em}{0.72em}
\titlespacing*{\subsection}{0pt}{1.20em}{0.58em}
\pagestyle{fancy}
\fancyhf{}
\fancyhead[C]{\small\nouppercase{\leftmark}}
\fancyfoot[C]{\thepage}
\renewcommand{\headrulewidth}{0.3pt}
\hypersetup{hypertexnames=false,colorlinks=true,linkcolor=black,urlcolor=blue,
  pdftitle={气液多相流直接数值模拟},
  pdfauthor={Grétar Tryggvason, Ruben Scardovelli, Stéphane Zaleski}}
\graphicspath{{./}}
\title{气液多相流直接数值模拟\\[0.65em]\large 中文译文}
\author{Grétar Tryggvason、Ruben Scardovelli、Stéphane Zaleski 著\\
\small AI 辅助翻译与排版}
\date{原著：Cambridge University Press，2011}
\clubpenalty=10000
\widowpenalty=10000
\displaywidowpenalty=10000
\raggedbottom
\allowdisplaybreaks[1]
\setlength{\abovedisplayskip}{0.82\baselineskip plus 0.2\baselineskip minus 0.1\baselineskip}
\setlength{\belowdisplayskip}{0.82\baselineskip plus 0.2\baselineskip minus 0.1\baselineskip}
\setlength{\abovedisplayshortskip}{0.58\baselineskip}
\setlength{\belowdisplayshortskip}{0.58\baselineskip}
\newenvironment{sourceblock}
  {\par\addvspace{0.65\baselineskip}\noindent\begin{minipage}{\linewidth}\centering}
  {\end{minipage}\par\addvspace{0.48\baselineskip}}
% 参数依次为：原 PDF 页号、裁片左下与右上坐标（bp）、最大高度比例。
\newcommand{\sourcecrop}[6]{%
  \includegraphics[page=#1,pagebox=mediabox,viewport={#2bp #3bp #4bp #5bp},clip,
    width=0.94\textwidth,max height=#6\textheight]{Tryggvason_Book.pdf}%
}
\newcommand{\eqnote}[1]{\par\smallskip\noindent\textit{式中：}#1\par}

\makeatletter
\renewcommand{\@pnumwidth}{2.2em}
\renewcommand{\@tocrmarg}{3em}
\makeatother

\begin{document}
\frontmatter
\maketitle
\chapter*{译本说明}
\markboth{译本说明}{译本说明}
本译文按正常教材阅读顺序重新排版，不与原 PDF 逐页对齐，也不保留原书页码目录。
章节另起页，页眉显示当前章名；正文使用较宽松的字号、字距、行距和段间距。

说明性文字由 AI 按多相流直接数值模拟、界面捕捉与界面追踪领域的通行术语翻译并统一。
短式和中等长度公式重写为 LaTeX；只有实验照片、数值结果图、界面重构示意和极复杂推导
才从原 PDF 紧裁保留，并配中文图意或符号说明。裁片作为不可拆分版块处理，避免跨页、侵入
页眉或产生孤立题注。译文供学习与检索使用；涉及公式、物性参数和模型细节时建议核对原著。

\begingroup\small\tableofcontents\endgroup

\chapter*{常用符号与中文释义}
\addcontentsline{toc}{chapter}{常用符号与中文释义}
\markboth{常用符号与中文释义}{常用符号与中文释义}
\begingroup\small
\renewcommand{\arraystretch}{1.20}
\begin{longtable}{>{\centering\arraybackslash}p{0.22\textwidth}p{0.68\textwidth}}
\toprule 符号 & 中文释义 \\ \midrule
\endfirsthead
\toprule 符号 & 中文释义 \\ \midrule
\endhead
$\rho,\mu,\nu$ & 密度、动力黏度和运动黏度 \\
$\bm u=(u,v,w)$；$p$ & 速度矢量及其分量；压力 \\
$\bm x$；$t$ & 位置矢量；时间 \\
$\Gamma$；$\bm n$；$\bm t$ & 流体界面；界面单位法向矢量；单位切向矢量 \\
$\kappa$；$\sigma$ & 界面曲率；表面张力系数 \\
$H$；$\delta_\Gamma$ & Heaviside 指示函数；集中在界面上的 Dirac 分布 \\
$C$ 或 $\chi$ & 颜色函数、体积分数或相指示函数，具体定义依章节而定 \\
$\alpha$ & VOF 体积分数 \\
$\phi$ & 水平集函数或相场序参量，依上下文而定 \\
$[\![q]\!]$ & 量 $q$ 穿过界面的跳跃 \\
$D/Dt$ & 随体导数（物质导数） \\
$Re,We,Ca,Bo,Oh$ & Reynolds、Weber、Capillary、Bond 与 Ohnesorge 数 \\
$Fr,Pe,Pr$ & Froude、Péclet 与 Prandtl 数 \\
$\Delta x,\Delta t$ & 网格间距与时间步长 \\
$\mathcal S,\mathcal V$ & 界面面积（曲线长度）与控制体体积 \\
$\dot m$；$L$ & 相变质量通量；潜热 \\
\bottomrule
\end{longtable}
\endgroup

\clearpage
\makeatletter\@mainmattertrue\makeatother
\pagenumbering{arabic}
% ===== 译文分片 1（已合并）=====

% Tryggvason 等《气液多相流直接数值模拟》中文译稿，第 10--85 页
% 本文件为可嵌入主文档的正文片段；图号、公式号由主文档统一生成。

% [原PDF p10]
\chapter*{前言}
\addcontentsline{toc}{chapter}{前言}
\markboth{前言}{前言}

科学进步通常由一连串相互促进的事件构成：一个领域的突破为另一个领域创造机会，后者又推动第三个领域前进。高速、高性能计算机的发展，催生了多相流直接数值模拟（direct numerical simulation，DNS）的新方法，使研究者能够细致考察多相流并加深对其机理的认识。多相流 DNS 的源头可追溯至 20 世纪 60 年代初计算流体力学诞生之时，但这一领域真正迅速发展，主要还是近十余年的事情。本书作者有幸参与了这些方法的早期开拓，也较早把 DNS 用于相当复杂的多相流问题；同时，我们通过参加学术会议、访问实验室以及编辑《计算物理杂志》和《国际多相流杂志》等期刊，持续关注同行的进展。我们认为当前状况可以概括为两点：

\begin{itemize}
  \item 各种多相流 DNS 路线在表面上有所不同，但它们之间的共同点远多于差异；
  \item 随着方法越来越精细、研究问题越来越复杂，新研究者进入这一领域所要跨越的门槛也在不断升高。
\end{itemize}

本书试图同时回应这两个问题。

互不相溶的两种或多种流体天然包含尖锐界面，因此，如何为这种流动建立数值方法已成为一个活跃方向，多个研究团队都取得了显著进展。本书主要讨论两类方法：流体体积法（volume of fluid，VOF）和界面追踪法（front tracking）。这一选择既反映作者自身的研究背景，也因为两类方法都已相当成功，并借助 DNS 揭示了许多重要的多相流动力学现象。更重要的是，各类界面方法的共性足够强；掌握本书内容后，读者通常也能较容易地转向水平集法、相场法等其他界面捕捉方法。

多相流 DNS 的最终目标是理解流动行为及其性质。数值方法的发展固然重要，但最有价值的成果往往来自方法的实际应用。因此，本书安排了若干章节，介绍 DNS 如何用于认识具体系统，以及目前已经获得了哪些物理认识。所选案例不可避免带有作者偏好，因为我们主要讨论自己最熟悉的工作；不过，这些章节的意义并不限于相应专题。我们确信，现有方法已经足够成熟，可以用来探索大量复杂流动。把成熟方法应用于适合的问题，并为更复杂的流动发展新方法，正是多相流 DNS 当前最令人振奋的方向之一。

% [原PDF p11]
本书不仅介绍界面数值方法，也始终把多相流 DNS 的物理目标放在中心位置：真正的目的，是理解这类流动的行为和性质。方法发展固然重要，最重要的回报却来自实际应用。因此书中安排若干章节，说明怎样用 DNS 理解具体系统，以及截至写作时已经获得了哪些认识。案例难免偏向作者本人最熟悉的研究，但其价值并不局限于所讨论的专题。作者相信，这些方法已经成熟到足以探索许多复杂流动；把已有方法用于适合的问题，同时针对末章所述更复杂物理发展新方法，是当前最令人振奋的近期方向之一。作者还强调，不同界面捕捉路线的共同基础很强，理解本书重点介绍的 VOF 与界面追踪后，转向水平集或相场等方法通常不会十分困难。

我们的工作得益于许多同事和朋友的努力。首先，我们要感谢学生、博士后研究人员和访问学者；他们对本书所介绍的工作作出了大量而重要的贡献。GT 感谢他的学生 D. Yu、M. Song、S. O. Unverdi、E. Ervin、M. R. Nobari、C. H. H. Chang、Y.-J. Jan、S. Nas、M. Saeed、A. Esmaeeli、F. Tounsi、D. Juric、N. C. Suresh、J. Han、J. Che、B. Bunner、N. Al-Rawahi、W. Tauber、M. Stock、S. Biswas 和 S. Thomas 博士，也感谢访问学者与博士后 S. Homma、J. Wells、A. Fernandez 和 J. Lu。RS 感谢他的学生 E. Aulisa、L. Campioni、A. Cervone 和 V. Marra 博士，以及合作者 S. Manservisi、P. Yecko 和 G. Zanetti。SZ 感谢他的学生 B. Lafaurie、F.-X. Keller、J. Li、D. Gueyffier、S. Popinet、A. Leboissetier、L. Duchemin、O. Devauchelle、A. Bagu\'e 和 G. Agbaglah 博士，以及合作者、访问学者和博士后 G. Zanetti、A. Nadim、J.-M. Fullana、C. Josserand、P. Yecko、M. Renardy、Y. Renardy、E. Lopez-Pages、T. Boeck、P. Ray、D. Fuster、G. Tomar 和 J. Hoepffner。SZ 还要感谢他所信赖的朋友与导师 Y. Pomeau、D. H. Rothman 和 E. A. Spiegel，感谢他们不可估量的建议。

我们还感谢 Victoria Tsengu\'e Ingoba 女士两次通读全书，并指出许多排印错误和其他问题。当然，书中任何遗留的错误、遗漏或含混之处都只能由作者负责。

最后但同样重要的是，我们感谢家人始终如一的支持，感谢他们接受我们长时间工作，也容忍那些有时看起来颇难理解的事务优先次序。

\begin{flushright}
Gr\'etar Tryggvason\quad Ruben Scardovelli\quad St\'ephane Zaleski
\end{flushright}

% [原PDF p12]
\chapter{引言}

气液多相流在自然界和人类工程活动中都具有基础地位。日常生活中，我们接触液体时几乎总会遇到自由表面，例如饮水、清洗、漂洗和烹饪。工业中更是如此：常规与核电站普遍依靠沸腾传热；钢铁冶炼、钢包冶金以及铝、铜的二次精炼常利用气泡驱动循环；大量能源来自液体燃料，而液体必须先蒸发才能燃烧，因此通常需要雾化成大量小液滴，以获得巨大的界面面积。除阻力、管路压降和气体燃料混合等少数问题外，几乎所有工业流体过程都以某种形式涉及多相流。

某些系统的一相是固体，例如浆体和流化床；但大量问题由液相与气相组成。在自然气液流动中，雨滴最容易使人联想到多相流，而气泡和液滴还控制着海洋与大气之间的热质交换，并参与火山爆发。生命体本身也可以看成庞大而复杂的多相系统。

理解气液多相流动力学具有重要的工程与科学意义，但这类问题在数学上素以困难著称。两相中的流动方程高度非线性，而且相界面的位置通常也是未知量，必须作为解的一部分同时确定。因此，严格解析解只存在于最简单的情形，例如 Stokes 流中气泡或液滴的定常运动、线性无黏波以及气泡和液滴的小振幅振荡。实验研究同样不易：许多实际流动的长度尺度很小、时间尺度很短，而且多数区域缺乏光学可视条件。正因如此，自 20 世纪 50 年代末至 60 年代初 CFD 兴起以来，多相流研究者就一直期待依靠控制方程的数值解。

本书所谓 DNS，是指对包含非平凡多尺度结构的非定常流动，在足够细的网格上直接求解控制方程，使连续介质意义下的全部时间与空间尺度得到分辨。这一定义介于两种极端用法之间：一类作者只把完全分辨湍流称为 DNS，另一类则把所有不含湍流模型的流动计算都称为 DNS。在雾化研究中，有人把未完全分辨且不含湍流模型的模拟称为 LES；本书仍倾向于称其为 DNS，尤其是在持续通过网格加密检验结果的情况下。

多相流 DNS 落后于均相流并不奇怪，但近十余年的积累已开始产生显著回报。现在，人们能够直接模拟许多相当复杂的系统，并得到实验或理论很难以其他方式获取的信息。本书旨在评述这一领域的现状、回顾其发展路径，并为研究更复杂的多相流建立基础。

\section{多相流实例}

本书关注数值模拟，但从若干真实系统出发更有助于理解目标问题的范围。

% [原PDF p13]
气泡存在于大量工业过程：在沸腾传热中把蒸气带离热表面，在化工设备中分散气体并搅动液体，还会显著改变海洋中的声传播。设计气泡流系统时，需要知道大量气泡的群体上升速度如何随空隙率和粒径分布变化、气泡如何弥散以及它们如何搅动连续相。图 1.1 所示为静止液池中大量浮力气泡上升的情形。一个典型小型气泡柱中，平均气泡直径约为 $2.2\,\mathrm{mm}$，空隙率约为 $0.75\%$，平均上升速度约为 $0.27\,\mathrm{m\,s^{-1}}$；由于尺寸并不一致，各气泡通常以不同速度上升。

\begin{sourceblock}
\captionof{figure}{大量浮力气泡在原本静止的液池中上升。平均气泡直径约为 $2.2\,\mathrm{mm}$，空隙率约为 $0.75\%$。引自 Br\"oder 与 Sommerfeld（2007），经许可转载。}
\end{sourceblock}

% [原PDF p14]
燃烧、涂覆、喷漆、灌溉和增湿等过程需要把液体射流雾化。能否预测雾化速率、所得液滴粒径分布以及液滴速度，直接关系到这些过程的成功设计。图 1.2 的同轴水-气射流实验中，直径 $8\,\mathrm{mm}$ 的液体喷口以 $0.6\,\mathrm{m\,s^{-1}}$ 射出水柱；为加速破碎，射流被注入速度为 $35\,\mathrm{m\,s^{-1}}$ 的同向空气流。最初，空气与液体之间的剪切产生尺度较大的轴对称波；当这些波向下游移动时，空气从波峰拉出细长液丝，液丝随后因毛细不稳定性断裂成液滴。更详细的讨论见 Marmottant 与 Villermaux（2001）以及 Villermaux 等（2004）。

\begin{sourceblock}
\captionof{figure}{同轴水射流和空气射流雾化实验的照片。引自 Villermaux 等（2004），经许可转载。}
\end{sourceblock}

% [原PDF p15]
液滴撞击固体或液体表面通常会飞溅，并且往往会显著扰乱表面。例如，雨滴落地经常造成土壤侵蚀；但在淬火和喷雾冷却中，液滴撞击也能增强传热，降雨还常显著加强海洋表层的混合。图 1.3 中，直径约 $3\,\mathrm{mm}$ 的液滴从接近 $0.5\,\mathrm{m}$ 的高度释放，撞击一层深度略大于一个液滴直径的液体。撞击产生液体凹坑和边缘液冠，液冠经常进一步破碎成液滴；凹坑坍缩时，有时还会把气泡困入液体之中。

\begin{sourceblock}
\captionof{figure}{液滴撞击自由表面时形成的飞溅。图片由 Rochester Institute of Technology 的 A. Davidhazy 提供，经许可转载。}
\end{sourceblock}

% [原PDF p16]
气泡既可由向液池注气或自由表面卷吸形成，例如波浪破碎时的卷吸；也常由液体发生液-汽相变产生。这种相变往往在固体表面成核：把液体加热到饱和温度以上会发生沸腾，把压力降到蒸气压以下则会发生空化。图 1.4 显示浸在水中的水翼在最大厚度附近发生大范围空化。水翼弦长为 $7.6\,\mathrm{cm}$，水流以 $13.7\,\mathrm{m\,s^{-1}}$ 从左向右运动；流体绕过前缘时速度升高，压力下降到足以使液体“沸腾”。蒸气泡移到水翼后部的高压区后会坍缩。然而，溶解在液体中的残余气体在气泡存在期间会扩散进入其中，因此蒸气凝结后仍留下可见痕迹。

\begin{sourceblock}
\captionof{figure}{水翼最大厚度附近的大范围空化，流动方向从左向右。除相变引起的体积变化外，气泡的可压缩性往往也很重要，见第 11.1.5 节。引自 Kermeen（1956），经许可转载。}
\end{sourceblock}

多相系统的一相也常为固体。气体或液体中的固体颗粒悬浮很常见，“多相流”一词有时还扩展到流体穿过或流过复杂静止固体的情形，例如填充床、多孔介质、森林和城市。固-气或固-液系统与气液系统的主要差别通常在于前者的界面保持自身形状，虽然固体位置可以改变。但也有例外：柔性固体能响应流动而变形；在凝固或侵蚀过程中，边界也会演化，有时会形成与气液系统中同样曲折复杂的形状。金属合金凝固时，溶质最初被固相排斥，这会造成成分过冷并使凝固前沿失稳。富溶质相最终也会凝固，但其成分与最先凝固的材料大不相同。所得微观组织的尺度、形状和成分决定材料性质，而这些特征通常对各项工艺参数十分敏感。图 1.5 给出经金相制样和腐蚀、使相界及界面显现出来的 Al--Si 合金显微照片。浅灰色相几乎是纯铝，最先凝固；但成分过冷使它形成尺度为几十微米量级的枝晶结构。

\begin{sourceblock}
\captionof{figure}{铝--硅合金的微观组织。图片由 Worcester Polytechnic Institute 的 D. Apelian 提供，经许可转载。}
\end{sourceblock}

% [原PDF p17]
生命系统也提供了丰富的多相流实例。在“身体”尺度，血流中悬浮的血细胞、肺部流动中的气溶胶、器官乃至整个个体的运动，都是显而易见的例子。更复杂的系统，例如鸟群穿过空气以及鱼群穿过水体，同样属于多相流。图 1.6 显示大量黄尾副绯鲤共同游动并协调动作。自推进物体会形成产生推力的尾迹，因此其群体动力学很可能与上升或下沉物体的群体运动显著不同。理解单条鱼的运动以及大鱼群的集体运动，不仅可能有助于种群控制和捕捞管理，也可能启发机械游动和飞行装置的构造。

\begin{sourceblock}
\captionof{figure}{夏威夷群岛西北部附近的黄尾副绯鲤（\textit{Mulloidichthys flavolineatus}）鱼群。图片来自 NOAA Photo Library。}
\end{sourceblock}

\section{计算建模}

多流体流动指两种不同流体的流动，多相流则指同一物质不同相之间的流动；它们的计算研究几乎与常密度流动计算一样古老。与常密度流动一样，人们尝试了许多不同路线，并采用了多种简化。本节力求简明而较完整地回顾模拟多流体流动的主要工作。我们不打算逐篇列举所有论文，但希望不遗漏主要发展。

多流体流动（两种不同流体）和多相流动（同一物质的不同相）的计算历史，几乎与常密度流动计算一样悠久。研究者尝试过多种路线并采用不同程度的简化。下面概述主要发展脉络。

\subsection{\texorpdfstring{简单流动：$Re=0$ 与 $Re\to\infty$}{简单流动：Re=0 与 Re 趋于无穷}}

% [原PDF p18]

当黏度极大或极小，即 Reynolds 数趋于两个极限时，可分别忽略惯性得到 Stokes 流，或忽略黏性得到无黏流。无黏流通常还要假设除流体界面外流动无旋。对于不可变形球形分散相，两种极限下都可把控制方程约化为粒子位置的耦合常微分方程组。Stokes 流方面，Brady 及合作者系统研究了剪切流中颗粒悬浮体；无黏流方面，周期区域中的球形气泡模拟显示，气泡容易形成水平团簇，尤其在气泡速度方差较小时。

对可变形边界，Stokes 流和无黏势流都可采用边界积分法。Birkhoff（1954）以无黏、无旋两流体界面追踪初始重流体位于轻流体上方时的 Rayleigh--Taylor 不稳定性，是最早的尝试之一；此后，这一问题和这种方法都成为多相流模拟的经典内容。Longuet-Higgins 与 Cokelet（1976）提出水波边界积分法并用于研究破碎波，这项工作影响极大，随后出现了大量非常成功的扩展和应用，特别是 Vinje 与 Brevig（1981）、Baker 等（1982）以及 Schultz 等（1994）的水波研究。其他应用包括 Rayleigh--Taylor 不稳定性的演化（Baker 等，1980）、空化泡的生长与坍缩（Blake 与 Gibson，1981；Robinson 等，2001）、气泡与自由表面聚并所引起的气泡和液滴生成（Oguz 与 Prosperetti，1990；Boulton-Stone 与 Blake，1993）、孔口处的气泡和液滴形成（Oguz 与 Prosperetti，1993）以及涡旋流与自由表面的相互作用（Yu 与 Tryggvason，1990）等。

所有无黏流边界积分法（若积分按单元进行，也称边界元法）的共同思想，都是通过积分 Bernoulli 型方程随时间推进表面奇点的强度。表面奇点给出一个速度分量，Green 第二恒等式再给出另一个速度分量，于是可以推进表面位置。选择不同的表面奇点便产生大量不同方法：有的只适用于自由表面，有的适用于两流体问题；不同实现又进一步增加了可能性。详尽讨论及较新的进展可参见 Hou 等（2001）。尽管改进仍在继续，新应用也不断出现，但在 Longuet-Higgins 和 Cokelet 的论文发表三十多年后，二维无黏流边界积分技术已经相当成熟。全三维计算仍然少见，例如 Chahine 与 Duraiswami（1992）计算了少量无黏空化泡的相互作用，Xue 等（2001）模拟了三维破碎波。势流假设在短时瞬态流动中取得了许多引人注目的成功，但其内在局限也很多：缺乏小尺度耗散机制使模型容易形成奇性，而没有耗散通常又使它不适合预测任何系统的长期演化。

% [原PDF p19]
当黏度极大或极小时，以 Reynolds 数衡量的两个极限会大幅简化流动描述：Stokes 极限完全忽略惯性，无黏极限完全忽略黏性；无黏流通常还要假设除流体界面外处处无旋。最成功的对象是不可变形球形分散相，因为两个极限都可把控制方程约化为粒子位置的耦合常微分方程组。Stokes 流方面，Brady 及其合作者系统研究了剪切流颗粒悬浮体，早期工作综述见 Brady 与 Bossis（1988）；无黏流方面，Sangani 与 Didwania（1993）和 Smereka（1993）模拟周期盒中的球形气泡，发现气泡倾向形成水平团簇，气泡速度方差较小时尤其明显。

可变形边界可由 Stokes 流或无黏势流的边界积分技术处理。Birkhoff（1954）追踪初始重流体位于轻流体上方的无黏、无旋界面，成为 Rayleigh--Taylor 数值研究的早期范例。Longuet-Higgins 与 Cokelet（1976）的水波方法则引发破碎波、空化泡、自由表面聚并和涡旋--界面相互作用等大量后续工作。所有无黏边界积分法都通过 Bernoulli 型方程推进表面奇点强度：奇点给出一个速度分量，Green 第二恒等式给出另一个分量，二者共同推进界面。

\begin{sourceblock}
\captionof{figure}{线性剪切流中液滴破裂的 Stokes 流模拟。数值结果后方几乎不可见的细线，是从实验照片上描出的液滴轮廓。引自 Cristini 等（1998），经许可转载。}
\end{sourceblock}

% [原PDF p20]
势流界面问题能够改写成边界积分，关键在于势方程的线性性。在相反的 Stokes 极限中，惯性可忽略，流动完全受黏性耗散支配，Navier--Stokes 方程同样变为线性，控制方程也可化成移动表面上的积分方程。非定常两流体 Stokes 问题的边界积分模拟始于 Youngren 与 Acrivos（1976）和 Rallison 与 Acrivos（1978）；他们分别模拟了伸长流中气泡和液滴的变形。随后许多作者研究了不同问题。Pozrikidis 及其合作者考察了液滴悬浮体的多个方面，其中较早的是 Zhou 与 Pozrikidis（1993）对通道中少量二维液滴悬浮体的研究。Loewenberg 与 Hinch（1996）以及 Zinchenko 与 Davis（2000）完成了全三维悬浮体模拟。Pozrikidis（1992）的专著详细说明了这种方法，Pozrikidis（2001）则较完整地总结了各种应用。图 1.7 给出强黏液滴在线性剪切流中破裂的计算；所用方法在液滴变形时自适应加密表面网格，从而持续分辨局部细颈。

% [原PDF p21]
Hele--Shaw 槽中的二维两流体流动也适合边界积分法。虽然流动完全由黏性支配，但远离界面处仍是势流。界面可由无黏流所用的奇点表示（de Josselin de Jong，1960），不过奇点强度的演化方程不同。Tryggvason 与 Aref（1983，1985）用这种表示研究了 Saffman--Taylor 不稳定性，即低黏流体置换高黏流体时界面的失稳演化。他们在固定网格上求法向速度，而不是使用 Green 定理；后来多位作者改用 Green 定理建立 Hele--Shaw 槽界面的边界积分法，例如 DeGregoria 与 Schwartz（1985）、Meiburg 与 Homsy（1988），以及 Hou 等（2001）的综述。

“简单流动”项下还应提到这样一种固体颗粒运动模拟：忽略流体运动，动力学只由颗粒惯性和颗粒间碰撞支配。许多研究者追踪了大量仅在相互碰撞时发生作用的颗粒；此时同样只需解粒子运动的常微分方程组，通常称为“颗粒动力学”。早期讨论见 Louge（1994），较新的介绍可参见 P\"oschel 与 Schwage（2005）。这些方法模拟某些固--气多相流非常成功，但适用范围很窄。若把由势函数相互作用的粒子也包括在内，行星或星系的引力相互作用模拟以及分子动力学在形式上都可归入这一类。全面讨论将大幅扩大本书范围，感兴趣者可分别查阅 Schlick（2002）的分子模拟文献和 Hockney 与 Eastwood（1981）的天体物理及其他粒子体系文献。

\subsection{有限 Reynolds 数流动}

% [原PDF p22]

中等 Reynolds 数下必须求解完整 Navier--Stokes 方程。约在 Birkhoff 用边界积分模拟无黏 Rayleigh--Taylor 问题十年以后，Harlow 团队在 Los Alamos 提出了标记-单元（marker-and-cell，MAC）方法。MAC 以分布在流体区域内的标记粒子辨认流体，同时在覆盖流体区和空区的规则网格上求解控制方程。Harlow 与 Welch（1965）首次介绍该方法，并在那篇论文中给出两个所谓溃坝问题的算例。随后很快出现一系列工作：Harlow 与 Welch（1966）研究 Rayleigh--Taylor 问题，见图 1.8；Harlow 与 Shannon（1967）研究液滴撞击液面的飞溅。原始 MAC 只处理自由表面，因此域中只有一种流体，需要在表面施加边界条件，而域内其余区域完全被动。Los Alamos 团队很快认识到同一方法也可用于两流体问题。Daly（1969b）计算有限密度比 Rayleigh--Taylor 不稳定性，Daly 与 Pracht（1968）研究密度流的初始运动，随后 Daly（1969a）加入表面张力并再次计算 Rayleigh--Taylor 不稳定性。MAC 很快吸引了一批追随者：Chan 与 Street（1970）将其用于自由表面波，Foote（1973，1975）分别模拟轴对称液滴振荡和液滴撞击刚性壁面，Chapman 与 Plesset（1972）以及 Mitchell 与 Hammitt（1973）追踪空化泡坍缩。

\begin{sourceblock}
\captionof{figure}{多相流计算研究的早期开端：用两流体 MAC 方法计算非线性 Rayleigh--Taylor 不稳定性的演化。引自 Daly（1969b），经许可转载。}
\end{sourceblock}

% [原PDF p23]
Los Alamos 团队在 20 世纪 60 年代和 70 年代初计算了许多不同问题，Harlow 还在《Scientific American》上介绍了基本思想（Harlow 与 Fromm，1965），但这一新工具的巨大潜力当时大体没有引起流体力学研究界的广泛兴趣。MAC 原本就是为多流体问题设计的，名称中的 M 正是“标记”的首字母；但它还首次成功用原始变量，即速度和压力，求解 Navier--Stokes 方程，并创新性地采用交错网格。今天，人们常把交错网格配合投影时间推进的方法统称为 MAC 型方法，见第 3 章。

下一代多流体方法从 MAC 逐渐演化而来。Harlow 与 Welch（1965）已经指出标记粒子会带来误差；在 Los Alamos 团队探索的众多算法思想中，用标记函数取代粒子很快成为最受欢迎的选择，于是流体体积法（VOF）诞生。VOF 最早在 Hirt 与 Nichols（1981）的同行评审期刊论文中系统讨论，但其起源更早（DeBar，1974；Noh 与 Woodward，1976）。单元平均标记函数的平流会产生数值扩散，见第 4 章；VOF 通过在单元内“重构”界面，使当前单元尚未充满前标记量不会进入下一个单元。一维实现几乎不费力；早期 VOF 在水平平流时简单假定界面是竖直平面，在竖直平流时则假定为水平平面。这种粗糙重构常产生大量所谓“漂浮碎屑”——从主界面脱落的非物理小液滴——从而降低计算精度。Youngs（1982）、Ashgriz 与 Poo（1991）等后来引入更复杂的重构，用可任意定向的二维直线或三维平面尽可能拟合界面。方法虽然复杂得多，却显著改善了标记函数的平流。

% [原PDF p24]
即使每个单元内采用更高阶的界面表示，准确计算表面张力仍是一个主要困难。Daly（1969b）用 MAC 研究表面张力对 Rayleigh--Taylor 不稳定性的影响时，为此引入了显式表面标记；然而，VOF 的出发点恰恰是彻底摆脱任何表面标记，因此表面张力必须直接由标记函数求得。Brackbill 等（1992）解决了这一问题：曲率以及表面张力可以通过标记函数梯度的单位法向之离散散度计算，并把界面力转化为局域体力，这就是连续表面力（CSF）法。Lafaurie 等（1994）又发展了该方法的“守恒”形式。VOF 随后由许多作者从不同方面扩展：界面重构方法不断改进（Rider 与 Kothe，1998；Scardovelli 与 Zaleski，2000；Aulisa 等，2007），表面张力算法持续完善（Renardy 与 Renardy，2002；Popinet，2009），动量方程平流格式和压力方程求解器也更先进（例如 Rudman，1997）；Chen 等（1997a）还使用子单元使界面尽可能保持尖锐。今天，VOF 已被广泛应用，许多商业软件也用它追踪界面和自由表面。图 1.9 是现代 VOF 对液滴撞击自由表面飞溅的计算实例；后续章节将详细讨论 VOF。

\begin{sourceblock}
\captionof{figure}{采用先进 VOF 方法计算的液滴飞溅。引自 Rieber 与 Frohn（1999），经 Elsevier 许可转载。}
\end{sourceblock}

% [原PDF p25]
20 世纪 90 年代初，由 MAC 和 VOF 的基本思想衍生出多种固定网格新路线。Unverdi 与 Tryggvason（1992）提出多流体界面追踪法，用相互连接的标记点显式表示界面。与标记点散布在整个流体区域的 MAC 不同，这些点只存在于界面上，并用于推进密度、黏度等物性以及计算表面张力，其余计算仍在固定网格上完成。

与此同时，固定网格界面追踪还发展出水平集法、CIP 法和相场法。Osher 与 Sethian（1988）提出水平集法；其首次用于流体界面追踪似乎见于 Sussman 等（1994）和 Chang 等（1996），他们模拟了二维气泡上升和液滴下落；Sussman 与 Smereka（1997）随后用轴对称版本研究气泡和液滴。VOF 平流不连续的标记函数，而水平集法平流连续的水平集函数，并把其零等值面认作界面。密度、黏度等物性则通过由水平集构造的标记函数得到；该标记函数在两流体间有一个平滑过渡区，因而比界面限制在一个网格间距内的 VOF 更光滑。不过，这种映射要求水平集函数在界面附近保持固定形状。Sussman 等（1994）为此引入重初始化，使水平集在任何时刻都近似等于点到界面的最短距离；这一步是水平集法成功应用于流体动力学模拟的关键。表面张力仍按 Brackbill 等（1992）的 CSF 思路求取。早期水平集法质量守恒较差，随后出现多项改进：Sussman 等（1998）以及 Sussman 与 Fatemi（1999）改善质量守恒，Sussman 等（1999）把水平集追踪与自适应网格结合，Sussman 与 Puckett（2000）则发展混合 VOF/水平集法。

% [原PDF p26]
CIP（constrained interpolated propagation，约束插值传播）方法由 Takewaki 等（1985）提出，尤其受到日本研究者欢迎，并被用于多种多相问题。CIP 用三次多项式描述两种流体间的过渡，同时更新标记函数及其导数以平流界面；除两流体问题外，它还用于包含漂浮固体等更复杂的应用，见 Yabe 等（2001）。

相场法则修改控制方程，使不同流体之间平滑过渡区的动力学满足热力学一致性。在数值实现中，过渡区比真实系统厚得多；修改后的净效果相当于引入“反扩散”项，使界面维持得相当尖锐。相场法和水平集法表面上有很多相似之处，但基本思想截然不同：水平集的平滑相界完全是为数值方便而人为引入；相场法的过渡区在物理上真实存在，只是出于数值原因被显著加厚。对于人为加厚的界面，保持正确热力学条件是否比从完全尖锐界面出发更有优势，在本书写作时尚无定论。其关键缺点是界面的传播和性质对过渡区内动力学非常敏感，所以必须充分分辨该区域。对于本可用尖锐界面很好描述的互不相溶流体，这一要求比其他“单流体”方法更严格。相场最初用于凝固建模（Kobayashi，1992，1993），并已广泛用于这类模拟；除含流动凝固（Beckermann 等，1999；Tonhardt 与 Amberg，1998）外，其在流体动力学中的应用仍较少（Jacqmin，1999；Jamet 等，2001）。相场法最有吸引力的场合，是必须考虑小尺度物理而又难以在尖锐界面极限下表示这些物理的问题。

% [原PDF p27]
上述“单流体”方法均以一套控制方程描述两种流体，流体运动大多在规则结构网格上计算。各种方法的主要区别在于标记函数如何平流以及表面张力如何求取；数值界面厚度从 VOF 的一个单元到水平集和某些界面追踪法的数个单元不等。一旦物性标记函数构造完成，流动计算的其余部分基本不再依赖具体的界面平流方案。尽管这些方法极为成功，其精度通常仍有限，因此近来出现了多种尝试，希望保留它们的大部分优点，同时把界面处理为“完全尖锐”。这类尝试可追溯到 Glimm 及其合作者（Glimm 等，1981；Glimm 与 McBryan，1985；Chern 等，1986）：他们在界面附近局部修改网格，使界面与单元边界重合；也受到复杂物体无黏流模拟中较新切割单元方法的影响（Quirk，1994；Powell，1998）。

现代尖锐界面技术包括幽灵流体法、浸入界面法和 Udaykumar 等（2001）的方法。Fedkiw 等（1999）提出的幽灵流体法仍通过平流水平集函数标记界面；但当界面附近的差分模板包含界面另一侧节点时，给这些节点赋予外推得到的虚构值，Glimm 等（2001）讨论了多种外推方式。Lee 与 LeVeque（2003）的浸入界面法则把跨界面的已知跳跃显式写入界面附近的有限差分近似。跳跃条件简单时这样做并不困难，复杂情形则会迅速变得繁琐，所以 Lee 与 LeVeque 将其方法限制在两流体黏度相同的情况。Udaykumar 等（2001）在规则网格上合并界面附近单元，并用多项式拟合求界面上的场值，从而表示复杂固体边界；它与无黏可压缩流的切割单元法（Powell，1998）有关，迄今主要用于固--液问题，包括凝固（Yang 与 Udaykumar，2005），但原则上同样可以用于多流体问题。三维推广见 Marella 等（2005）。

% [原PDF p28]
传统单流体方法除允许密度和黏度变化外，基本无需在界面附近修改流动求解器；上述尖锐界面法都要局部修改基本离散格式，因此算法复杂得多，但也很可能提高精度。这在控制参数处于极端范围时可能很重要，例如两流体物性差异极大、黏度很低，或界面物理非常复杂时。不过，直接实现单流体控制方程并配以先进界面或标记函数平流格式，已经在大量问题上证明有效，其简洁性很可能保证它仍会得到广泛使用。

除提高单流体路线的精度外，许多研究者还把基本方法扩展到比两种互不相溶液体更复杂的问题。已有方法把受污染界面、传质和化学反应、电流变效应、沸腾、凝固以及固体同自由表面或流体界面的相互作用纳入模拟。第 11 章将在全书末尾简要回顾这些高级应用。

% [原PDF p29]
贴体结构网格方法让一条网格线始终与流体界面对齐，而远离界面的节点由保角映射或椭圆网格生成技术决定。Ryskin 与 Leal（1984）用这种方法计算浮力作用下可变形轴对称气泡的定常上升，当时忽略气泡内部流体；Dandy 与 Leal（1989）以及 Kang 与 Leal（1987）随后把它推广到两流体和非定常问题。这类方法对单粒子定常运动或自由表面中等变形相当有效，但网格拓扑预先固定，难以容纳意外的界面连接与断裂。全三维结果较少，Takagi 等（1997）是一个例子；对于多个粒子的三维非定常运动，贴体结构网格通常不会成为首选。

在单流体路线发展的同时，研究者也探索了其他技术。Hirt 等（1970）较早使用结构化贴体 Lagrangian 网格：采用逻辑矩形结构网格，但网格点随流体速度运动，整个网格因而随流动变形。这适合界面拓扑较简单且不会出现意外界面构型的问题。相关做法是令一条网格线与流体界面对齐，而远离界面的网格通过保角映射或更先进的椭圆网格生成方法获得。Ryskin 与 Leal（1984）用它计算浮力作用下可变形轴对称气泡的定常上升，并忽略气泡内部流体；Dandy 与 Leal（1989）以及 Kang 与 Leal（1987）又将其扩展到两流体和非定常问题。多位作者用这种路线研究单粒子的定常运动、自由表面的中等变形等较简单问题；全三维模拟仍少见，Takagi 等（1997）是一个例外。对于多个粒子的三维非定常运动等复杂问题，这种方法大概不会成为首选。

% [原PDF p30]
用网格线连续表示流体界面的更一般方法，是采用完全非结构网格。这样可以按需插入和删除节点，并重塑畸变单元。早期应用中整个网格随流体速度运动；现代做法则只移动界面点，或令内部节点以不同于流体速度的速度运动，从而在减小网格畸变的同时保持足够分辨率。这一大类方法很多，这里只列举少数例子：Oran 与 Boris（1987）模拟二维液滴破裂，Shopov 等（1990）研究浮力气泡的初始变形，Feng 等（1994，1995）和 Hu（1996）计算多个粒子的二维非定常运动，Fukai 等（1995）追踪单个轴对称液滴撞壁。尽管路线看似复杂，Johnson 与 Tezduyar（1997）以及 Hu 等（2001）却对多个球形颗粒的三维非定常运动给出了十分出色的结果。图 1.10 展示 Hu 等任意 Lagrangian--Eulerian（ALE）方法的一个算例：两个固体球最初沿同一直线下落；后球被前球遮蔽，受到的阻力较小，于是追上并“亲吻”前球；共线构型不稳定，两球随即“翻滚”，翻滚后又彼此漂开。

\begin{sourceblock}
\captionof{figure}{两个下落球体的相互作用。由左至右给出四个不同时刻：追赶、接触、翻滚并最终分离。引自 Hu 等（2001），经 Elsevier 许可转载。}
\end{sourceblock}

格子 Boltzmann 方法（LBM）是有限 Reynolds 数多相流数值方法中较新的成员。现在已经清楚，LBM 能获得与传统方法相当的精度；但它是否如有时所宣称的那样显著更快或更简单，仍不明确。有关讨论可见 Shan 与 Chen（1993）和 Sankaranarayanan 等（2002）；LBM 与 Unverdi--Tryggvason 界面追踪法的结果比较见 Sankaranarayanan 等（2003）。这类方法很可能会长期存在，但本书不再专门讨论 LBM，读者可参阅 Rothman 与 Zaleski（1997）以及 Prosperetti 与 Tryggvason（2007）的第 6 章。

\section{展望}

% [原PDF p31]

近十余年来，多相流 DNS 已取得长足进步。现在可以对数百个气泡、液滴或颗粒组成的分散流进行足够长时间的准确模拟，从而得到可靠的统计量；凝固、蒸发等更复杂流动的方法也在快速成熟，大规模沸腾和凝固模拟已逐渐可行。

不过，多相流 DNS 用于研究与工程设计仍处于早期阶段。若能直接预测搅混湍动沸腾气泡流、射流破碎并形成蒸发液滴等复杂过程，将会根本改变我们对许多重要工业流动的认识。目前多数多相过程的控制仍主要依靠经验和直觉，而相关产业规模巨大；即使喷雾器、气泡柱注气器或管内破滴构件的效率只提高少量，也会产生可观收益。对空间飞行器等实验代价很高的系统，可靠预测还可显著降低设计成本。随着固定式或自由运动的微机电器件操控微尺度流动逐渐成为现实，预测这种操控的影响也会更加关键。

近期的发展重点预计将从简单两流体问题转向复杂物理，包括变表面张力、热场和电场耦合、相变、传质及化学反应。耦合系统不仅方程更多，时间与长度尺度跨度也更大，因此刚性问题的隐式时间积分和自适应网格会越来越重要。完全分辨的 DNS 与对最小尺度建模的模拟之间，界限也会逐渐模糊：例如用亚网格模型处理雾化中的薄膜，把极小液滴视为点粒子，或把连续介质模拟同接触线微观动力学、凝固界面动力学和薄火焰反应模型耦合。非 Newton 流体中对微观结构的模型化也属于这一方向。

\enlargethispage{2\baselineskip}
除算法外，大规模模拟的“人的因素”同样关键。三维界面处理、动态自适应和并行连续介质求解器的复杂度，已经使完整软件开发难以由单个研究生独立完成。降低知识门槛、让新研究者能够进入多相流 DNS，可能和提高数值方法的效率与精度同等重要。本书正希望为此提供帮助。

\chapter{含界面流动的流体力学}

% [原PDF p32]

本章给出尖锐界面分隔互不相溶流体或不同相时的多相流控制方程。首先以标准方式推导无界面流动方程；随后讨论移动界面的数学表示，以及耦合界面两侧方程所需的跳跃条件；最后引入“单流体表述”，即在整个流场统一书写的方程中把界面表示为奇异分布。单流体形式是本书后续数值方法的基础。

\section{基本假设}

推导以三个一般原则为基础：连续介质假设、尖锐界面假设和忽略显式分子间作用力。

真实流体由原子或分子组成。若在边长极小的盒子中测量单位体积质量，结果会剧烈涨落；随着测量尺度 $\ell$ 增大，密度逐渐平滑，并可由连续函数 $\rho$ 很好地近似，详细讨论见 Batchelor（1970）。常温液体在几十纳米以上通常已经满足连续介质假设，其中 $1\,\mathrm{nm}=10^{-9}\,\mathrm m$。稀薄气体中，物质离散性可能影响更大的长度尺度；此时分子两次碰撞之间的平均距离，即平均自由程 $\ell_{\mathrm{mfp}}$，是关键尺度。当 $\ell\gg\ell_{\mathrm{mfp}}$ 时，气体服从 Navier--Stokes 方程。

% [原PDF p33]
追踪大量单个分子足够长时间、再求取有意义平均量的分子模拟表明，只需出乎意料地少的分子，流体就会表现为连续介质。例如 Koplik 等（1988）发现，在真实压力和温度下，通道中数百个分子所形成的 Poiseuille 流已经与连续介质理论预测相符。

多相流还采用尖锐界面假设。界面可以分隔空气与水、油与醋等不同流体，也可分隔固-液、汽-液等不同热力学相。状态方程、密度、黏度和导热系数等物性通常跨界面突变。真实过渡发生在极小尺度上，因此在连续介质尺度可把界面视为厚度为零。

本书不考虑带电流体中的远程电磁力。van der Waals 力等分子间作用对界面物理很重要，但这里主要保留其宏观效应，即毛细作用或表面张力，并把它表示为集中在尖锐界面上的应力。若要同时覆盖纳米至厘米尺度，每一方向可能需要约 $10^7$ 个网格点；即使采用自适应网格，这在目前也远超实际能力。除本章先给出一般形式外，本书主要研究不可压缩流。

\section{基本方程}

质量、动量和能量守恒给出流动方程；还需用本构关系规定具体流体的性质。以下仅考虑 Newton 流体。

\subsection{质量守恒}

% [原PDF p34]
取空间中固定控制体 $V$，边界为 $S$，外法向为 $\bm n$。质量既不能产生也不能消灭，因此控制体内质量只能因物质穿过边界而改变。穿过面积元 $\mathrm dS$ 的外向质量流率为 $\rho\bm u\cdot\bm n\,\mathrm dS$，所以
\begin{equation}
 \frac{\mathrm d}{\mathrm dt}\int_V\rho\,\mathrm dV
 =-\oint_S\rho\bm u\cdot\bm n\,\mathrm dS .
 \tag{2.1}
\end{equation}
\begin{sourceblock}
\captionof{figure}{固定控制体 $V$，其边界曲面记为 $S$。}
\end{sourceblock}
固定控制体允许把时间导数移入积分号，再用散度定理得
\begin{equation}
 \int_V\left[\frac{\partial\rho}{\partial t}+\nabla\cdot(\rho\bm u)\right]\mathrm dV=0.
 \tag{2.2}
\end{equation}
由于 $V$ 任意，局部连续方程为
\begin{equation}
 \frac{\partial\rho}{\partial t}+\nabla\cdot(\rho\bm u)=0. \tag{2.3}
\end{equation}
引入物质导数
\begin{equation}
 \frac{D(\,\cdot\,)}{Dt}=\frac{\partial(\,\cdot\,)}{\partial t}+\bm u\cdot\nabla(\,\cdot\,), \tag{2.4}
\end{equation}
即可写成对流形式
\begin{equation}
 \frac{D\rho}{Dt}=-\rho\nabla\cdot\bm u . \tag{2.5}
\end{equation}
这说明，流体质点的密度只有在流体发生压缩或膨胀时才会变化。

\subsection{动量守恒}

% [原PDF p35]
物质质点密度只有在流体压缩或膨胀，即 $\nabla\cdot\bm u\ne0$ 时才会变化。动量守恒则说，固定控制体内动量变化率等于边界动量净流入加全部外力。式（2.6）右端第一项是穿过边界的动量通量，第二项是体力总和；最常见的体力是重力 $\bm f=\rho\bm g$。最后一项为表面力，构造对称应力张量 $\bm T$，使 $\bm n\cdot\bm T\,\mathrm dS$ 恰为法向为 $\bm n$ 的面积元所受力。对任意小控制体应用同样论证得到局部守恒式。并矢积 $\bm a\bm b$ 的 $ij$ 分量为 $a_i b_j$，故 $\bm u\bm u$ 的分量是 $u_i u_j$；利用连续方程消去质量散度项，便得到适用于任意连续介质的 Cauchy 运动方程。对水、油、空气等 Newton 流体，应力被假定为应变率的线性函数，压力给出各向同性法向部分，黏度给出变形率所对应的偏应力。


固定控制体内动量的变化率，等于边界动量通量之负值加上体力和表面力：
\begin{equation}
 \frac{\mathrm d}{\mathrm dt}\int_V\rho\bm u\,\mathrm dV
 =-\oint_S\rho\bm u(\bm u\cdot\bm n)\,\mathrm dS
 +\int_V\bm f\,\mathrm dV+\oint_S\bm n\cdot\bm T\,\mathrm dS . \tag{2.6}
\end{equation}
其中 $\bm f$ 为单位体积体力（重力时 $\bm f=\rho\bm g$），$\bm T$ 为对称 Cauchy 应力张量，使 $\bm n\cdot\bm T\,\mathrm dS$ 表示面积元上的作用力。局部形式为
\begin{equation}
 \frac{\partial(\rho\bm u)}{\partial t}=-\nabla\cdot(\rho\bm u\bm u)+\bm f+\nabla\cdot\bm T . \tag{2.7}
\end{equation}
这里 $\bm a\bm b$ 表示并矢积，分量为 $a_i b_j$。利用
\begin{equation}
 \nabla\cdot(\rho\bm u\bm u)=\rho\bm u\cdot\nabla\bm u+\bm u\,\nabla\cdot(\rho\bm u), \tag{2.8}
\end{equation}
以及连续方程，得到 Cauchy 运动方程
\begin{equation}
 \rho\frac{D\bm u}{Dt}=\bm f+\nabla\cdot\bm T . \tag{2.9}
\end{equation}
Newton 流体的应力与应变率线性相关：
\begin{align}
 \bm T&=(-p+\lambda\nabla\cdot\bm u)\bm I+2\mu\bm S, \tag{2.10}\\
 S_{ij}&=\frac12\left(\frac{\partial u_i}{\partial x_j}+\frac{\partial u_j}{\partial x_i}\right),\qquad
 \bm S=\frac12\left(\nabla\bm u+\nabla\bm u^{T}\right). \tag{2.11}
\end{align}
$\lambda$ 是第二黏性系数；采用 Stokes 假设时 $\lambda=-2\mu/3$。有些文献改用体积黏度 $\kappa=2\mu/3+\lambda$，Stokes 假设等价于体积黏度为零。代入得一般 Newton 流体的 Navier--Stokes 方程
\begin{equation}
 \rho\frac{D\bm u}{Dt}=\bm f-\nabla p+\nabla(\lambda\nabla\cdot\bm u)+\nabla\cdot(2\mu\bm S). \tag{2.12}
\end{equation}

\subsection{能量守恒}

% [原PDF p36]
$\lambda$ 是第二黏性系数；Stokes 假设给出 $\lambda=-2\mu/3$。有些文献改用体积黏度 $\kappa_b=2\mu/3+\lambda$，因此 Stokes 假设也可表述为体积黏度为零。将 Newton 本构代入 Cauchy 方程就得到一般 Navier--Stokes 方程。

能量守恒的积分式逐项具有清楚的物理意义：左端是控制体内内能与动能之和的变化率；右端第一项是内能和动能随质量穿过边界的净通量，第二项是体力做功，第三项是压力和黏性剪切等边界应力做功，最后一项是热流出。把速度与动量方程点乘可得机械能方程；用它消去总能量式中的对应项，再将积分定律局部化，便得到内能的对流方程。张量双点积定义为 $\bm A:\bm B=\sum_i\sum_j A_{ij}B_{ji}$。为封闭能量方程，还必须给出热流、状态方程和输运系数的本构关系。


令 $e$ 为单位质量内能，$\bm q$ 为热流密度，则总内能与动能满足
\begin{align}
 \frac{\mathrm d}{\mathrm dt}\int_V\rho\left(e+\frac12u^2\right)\mathrm dV
 ={}&-\oint_S\rho\left(e+\frac12u^2\right)\bm u\cdot\bm n\,\mathrm dS
 +\int_V\bm u\cdot\bm f\,\mathrm dV \notag\\
 &+\oint_S\bm n\cdot(\bm u\cdot\bm T)\,\mathrm dS
 -\oint_S\bm q\cdot\bm n\,\mathrm dS . \tag{2.13}
\end{align}
速度与动量方程点乘得到机械能方程。消去相应项后，能量方程的对流形式为
\begin{equation}
 \rho\frac{De}{Dt}-\bm T:\nabla\bm u+\nabla\cdot\bm q=0, \tag{2.15}
\end{equation}
其中 $\bm A:\bm B=\sum_i\sum_j A_{ij}B_{ji}$。采用 Fourier 定律
\begin{equation}
 \bm q=-k\nabla T, \tag{2.16}
\end{equation}
并忽略辐射传热，可写成
\begin{equation}
 \rho\frac{De}{Dt}+p\nabla\cdot\bm u=\Phi+\nabla\cdot(k\nabla T),\qquad
 \Phi=\lambda(\nabla\cdot\bm u)^2+2\mu\bm S:\bm S\ge0 . \tag{2.17}
\end{equation}
% [原PDF p37]
Fourier 定律 $\bm q=-k\nabla T$ 规定热流与温度梯度成正比，$k$ 是导热系数。假定 Newton 流体且忽略辐射传热后，能量方程中的耗散函数 $\Phi=\lambda(\nabla\cdot\bm u)^2+2\mu\bm S:\bm S$ 表示机械功转化为热的速率，必定非负。封闭方程还需要状态方程 $p=p(e,\rho)$，并给出 $\mu$、$\lambda$ 和 $k$ 随状态的变化。

原书方框 2.1--2.3 分别汇总积分形式、守恒微分形式和对流（非守恒）形式。积分形式直接对应有限体积法；守恒微分形式适合有限差分，并通常产生与有限体积本质相同的离散；重新组合后的物质导数形式则是教科书和理论分析中最常见的写法。不可压缩流是重要特例：当特征速度远小于声速，即 $Ma=U/c_s\ll1$ 时，每个流体质点沿轨迹保持密度不变，由连续方程进一步得到速度无散。

一般还需要状态方程，例如
\begin{equation}
 p=p(e,\rho), \tag{2.18}
\end{equation}
以及输运系数 $\mu$、$\lambda$、$k$ 随流体状态变化的关系。

控制方程在原书的方框 2.1--2.3 中分别汇总为三种等价形式。积分形式直接把守恒原理施加于小控制体，虽然不便于解析计算，却是有限体积法的起点：
\begin{align*}
\frac{\mathrm d}{\mathrm dt}\int_V\rho\,\mathrm dV
\oint_S\rho\bm u\cdot\bm n\,\mathrm dS&=0,\\
\frac{\mathrm d}{\mathrm dt}\int_V\rho\bm u\,\mathrm dV
\oint_S\rho\bm u(\bm u\cdot\bm n)\,\mathrm dS
\oint_S\bm n\cdot\bm T\,(-\mathrm dS)&=\int_V\bm f\,\mathrm dV,
\end{align*}
而由积分定律在一点成立得到的守恒微分形式适于有限差分法，所得离散通常与有限体积法本质相同。把各项重新组合便得到对流形式或非守恒形式，即质量、动量和内能的物质导数表达；教科书中最常见的正是这种形式。

% [原PDF p38]
\begin{sourceblock}
\captionof{table}{流体运动方程的三种等价表述：方框 2.1 为积分形式，方框 2.2 为守恒微分形式，方框 2.3 为对流或非守恒形式。通量项可与应力项合并写入同一表面积分。}
\end{sourceblock}

\subsection{不可压缩流}

若特征速度 $U$ 远小于声速 $c_s$，即 $Ma=U/c_s\ll1$，每个流体质点的密度沿轨迹不变：
\begin{equation}
 \frac{D\rho}{Dt}=\frac{\partial\rho}{\partial t}+\bm u\cdot\nabla\rho=0, \tag{2.19}
\end{equation}
于是质量守恒化为
\begin{equation}
 \nabla\cdot\bm u=0,\qquad \oint_S\bm u\cdot\bm n\,\mathrm dS=0. \tag{2.20--2.21}
\end{equation}
% [原PDF p39]
不可压缩并不要求密度处处相同，只要求每个物质质点的密度保持不变。不同质点可以具有不同密度；这些质点随流动平移时，固定空间点的密度仍可随时间变化，此时必须用式（2.19）更新密度场。只有密度在整个区域内都相同时，才无需这一步。

压力在不可压缩流中不是由密度和温度等状态变量决定，而是作为约束乘子，取恰当数值使速度场无散。可以把这一过程理解为：先不含压力地预测速度，再求出强制不可压缩所需的压力，并把预测速度投影到无散函数空间。若物性不依赖温度，求速度和压力无需同时求能量方程。

物质导数可写成若干解析等价形式：
\begin{equation}
 \frac{D\bm u}{Dt}=\frac{\partial\bm u}{\partial t}+\nabla\cdot(\bm u\bm u)
 =\frac{\partial\bm u}{\partial t}+\bm u\cdot\nabla\bm u
 =\frac{\partial\bm u}{\partial t}+\frac12\nabla(\bm u\cdot\bm u)-\bm u\times(\nabla\times\bm u). \tag{2.22}
\end{equation}
它们的连续形式相同，但离散后略有差别。对常密度、常黏度流体，$\nu=\mu/\rho$，动量方程化为
\begin{equation}
 \frac{\partial\bm u}{\partial t}+\bm u\cdot\nabla\bm u
 =-\frac{\nabla p}{\rho}+\frac{\bm f}{\rho}+\nu\nabla^2\bm u . \tag{2.23}
\end{equation}

\subsection{边界条件}

% [原PDF p40]
正确实现边界条件是流动数值模拟的主要困难之一，尽管物理边界条件原则上定义清楚。黏性不可压缩流在固壁必须黏附，因此法向和切向速度都等于壁面速度；无黏流没有黏性剪应力，可以沿壁面自由滑移，只需法向速度与壁面一致。在极小通道等情形，切向速度可采用 Navier 滑移条件。参数 $\beta$ 具有长度量纲：若把近壁速度线性外推到固体内部，它会在距壁面 $\beta$ 处降为零，所以 $\beta$ 称为滑移长度。

入口和出口是另一类困难。入口速度通常可给定，真实出口却没有简单通用条件。实验风洞也必须兼顾均匀入口和不干扰上游的出口；实验可以延长风洞，数值计算照做却会消耗大量网格。计算边界条件的目标，是在足够远处截断区域而让上游“感觉”不到截断；速度的零梯度等变化率条件只有在出口已远离关注区时才合理。速度条件确定后，压力边界条件应由不可压缩方程推导，而不是随意另行指定。


黏性不可压缩流在固壁满足无滑移条件
\begin{equation}
 \bm u=\bm U_{\mathrm{wall}}, \tag{2.24}
\end{equation}
同时约束法向和切向速度。无黏流只要求法向速度等于壁面法向速度。微通道等情况下可对切向速度 $u_t$ 使用 Navier 滑移条件
\begin{equation}
 u_t-U_{\mathrm{wall}}=\beta\frac{\partial u_t}{\partial n}, \tag{2.25}
\end{equation}
其中 $\beta$ 是具有长度量纲的滑移长度；把近壁速度线性外推到固体内部，速度在距壁面 $\beta$ 处为零。

% [原PDF p41]
\begin{sourceblock}
\captionof{figure}{某些模型允许流体在固体表面滑移。把切向速度线性外推到固体内部，它在距壁面 $\beta$ 处为零；$\beta$ 称为滑移长度。处理两种流体与固体相交的三相线（接触线）时，这一模型很有用。}
\end{sourceblock}

我们经常只模拟较大充液区域的一部分，此时必须给定入口和出口边界条件。入口通常给定速度；出口条件则困难得多，而且没有简单的普适解。实验人员设计风洞时也面临类似困难：入口速度要均匀，出口又应尽量少影响上游。计算中实现均匀入口通常比实验容易，但实验可把风洞做得足够长，而计算若照做会需要大量网格点。数值边界的目标，是截断计算域而让它“看起来”仍然很长。若出口离关注区足够远，可对速度施加零梯度等变化率条件；速度边界条件给定后，再由方程导出入口和出口的压力条件。

在理论研究中，周期边界很有吸引力，因为无需处理往往棘手的入口和出口条件。若 $x$ 方向周期为 $L$，则 $\bm u(x,y,z)=\bm u(x+L,y,z)$；压力和其他变量同理，$y$ 或 $z$ 方向周期时也类似。若重力方向周期，为避免整个流体均匀加速，需要加入大小为 $\rho_{\mathrm{av}}\bm g$ 的平衡体力，其中 $\rho_{\mathrm{av}}$ 是全域平均密度。该力产生静水压力梯度，但它保持的是域内总动量而不是平均速度，因此与具有刚性底壁的容器并不完全等价。

\section{界面的描述与定义}

% [原PDF p42]
准确跟随分隔不同流体或不同相的可变形界面，是可靠预测多相流的核心。界面位置和运动有多种描述。最简单的是二维用一个坐标的单值函数、三维用两个坐标的单值函数表示。薄液层可写成 $y=h(x)$ 或 $z=h(x,y)$；这种高度函数在薄膜中很方便，但只要界面发生悬垂就会多值，必须采用更一般的参数化。

任意二维界面可引入沿界面的参数 $u$，写成 $\bm x(u)=(x(u),y(u))$。对闭曲线，$\mathrm d\bm x/\mathrm du$ 是切向量。参数并非必须为弧长；若选弧长 $s$，则 $\mathrm ds^2=\mathrm dx^2+\mathrm dy^2$，$\bm t=\mathrm d\bm x/\mathrm ds$ 自动成为单位切向。法向与切向正交，方向可按问题方便选取；闭曲线常取外法向，而本书通常规定它从参考流体 1 指向流体 2。


准确预测多相流的核心是跟随可变形相界面。若界面可由单值函数描述，则二维和三维分别可写成
\begin{equation}
 y=h(x,t),\qquad z=h(x,y,t). \tag{2.26}
\end{equation}
这对薄膜非常方便，但出现悬垂时 $h$ 会多值。图 2.3 正说明了这种局限；尽管薄膜问题中高度函数常是最方便的描述，一般界面仍需要更精细的表示。

\begin{sourceblock}
\captionof{figure}{用 $y=h(x)$ 参数化界面是最简单的方式；但曲线出现悬垂时，高度函数 $h$ 会成为多值函数，必须改用其他表示。}
\end{sourceblock}

% [原PDF p43]
任意二维界面可先用任意坐标 $u$ 参数化；最简单时再取弧长 $s$ 为参数，表示为
\begin{equation}
 \bm x(s)=(x(s),y(s)). \tag{2.27}
\end{equation}
若 $s$ 取弧长，则 $\mathrm ds^2=\mathrm dx^2+\mathrm dy^2$，单位切向量 $\bm t=\mathrm d\bm x/\mathrm ds$，$|\bm t|=1$，单位法向 $\bm n$ 满足 $\bm t\cdot\bm n=0$。法向朝向原则上可任意选择；闭曲线通常令其朝外，本书通常把流体 1 作为参考相，法向从流体 1 指向流体 2。例如空气--水混合物若以水为参考流体，则法向从水指向空气。曲线按逆时针方向参数化且 $\bm t=(t_1,t_2)$ 时，朝外法向为 $\bm n=(t_2,-t_1)$。

\begin{sourceblock}
\captionof{figure}{界面的参数化。坐标 $u$ 沿界面变化；单位切向 $\bm t$ 沿参数方向，单位法向 $\bm n$ 从流体 1 指向流体 2。}
\end{sourceblock}

单位切向长度恒为一，沿曲线移动时 $\bm t$ 的变化只能沿法向；其对弧长的变化率大小就是曲率。采用相应定向时，Frenet--Serret 关系为
\begin{equation}
 \frac{\mathrm d\bm t}{\mathrm ds}=\kappa\bm n,\qquad
 \frac{\mathrm d\bm n}{\mathrm ds}=-\kappa\bm t,\qquad
 \kappa=\bm n\cdot\frac{\mathrm d\bm t}{\mathrm ds}=-\bm t\cdot\frac{\mathrm d\bm n}{\mathrm ds}. \tag{2.28--2.30}
\end{equation}
曲率符号依赖法向约定；曲线向法向一侧弯曲时曲率为正。按本书约定，外法向圆周半径为 $R$ 时 $\kappa=-1/R$。附录 A 会以分量形式更详细地讨论二维曲线。

% [原PDF p44]
曲率符号完全取决于法向的定向；若曲线向法向一侧折弯，则按本书约定曲率为正。二维曲线的分量公式见附录 A。三维界面必须用两个独立参数 $u,v$ 表示。由于三维曲面几何本身内容很多，本书只在正文引用所需结果，较系统的推导放入附录 A。

从参数曲面 $\bm x(u,v)$ 分别对 $u$ 和 $v$ 求导，得到两个切向量 $\bm x_u$、$\bm x_v$；它们的叉积垂直于曲面，归一化后就是单位法向。除了位置和法向，计算中还频繁使用两条独立切向和平均曲率。曲面元边界上的向量 $\bm p$ 位于切平面内、垂直于边缘，它在后续表面张力积分中扮演“沿边缘拉力方向”的角色。平均曲率既可写成法向的负曲面散度，也可把法向延拓到界面外后取通常的三维散度。

\begin{sourceblock}
\captionof{figure}{曲率符号取决于法向的定向；曲线向法向一侧折弯时，曲率为正。}
\end{sourceblock}

三维曲面需要两个独立参数。图 2.6 给出由 $u,v$ 参数化的小曲面元；三维曲面几何内容繁多，附录 A 汇集了本书所需的结论：
\begin{equation}
 \bm x(u,v)=(x(u,v),y(u,v),z(u,v)),\quad
 \bm x_u=\frac{\partial\bm x}{\partial u},\quad
 \bm x_v=\frac{\partial\bm x}{\partial v}. \tag{2.31--2.32}
\end{equation}
两个偏导均切于曲面，单位法向为
\begin{equation}
 \bm n=\frac{\bm x_u\times\bm x_v}{|\bm x_u\times\bm x_v|}. \tag{2.33}
\end{equation}
% [原PDF p45]
平均曲率可由曲面散度或延拓到界面外的法向场求得：
\begin{equation}
 \kappa=-\nabla_S\cdot\bm n=-\nabla\cdot\bm n. \tag{2.34--2.35}
\end{equation}
若面积为 $\delta A$ 的微小曲面元边界上，$\bm p=\bm t\times\bm n$ 是垂直于边界曲线但位于切平面内的向量，则
\begin{equation}
 \kappa\bm n=\lim_{\delta A\to0}\frac1{\delta A}\oint\bm p\,\mathrm dl. \tag{2.36}
\end{equation}

其中 $\bm p=\bm t\times\bm n$，$\bm t$ 沿曲面元边界，所以 $\bm p$ 位于曲面切平面内并垂直于边界；它可以理解为对曲面元边缘的“拉力”方向。

\begin{sourceblock}
\captionof{figure}{由 $u,v$ 参数化的小曲面元。$\bm p$ 垂直于面元边缘但位于曲面切平面内。}
\end{sourceblock}

也可用全域定义的标记函数隐式描述界面。特征函数为
\begin{equation}
 H(\bm x)=\begin{cases}1,&\bm x\text{ 位于流体 1 内},\\0,&\bm x\text{ 位于流体 2 内},\end{cases} \tag{2.37}
\end{equation}
界面就是 $H$ 的跳变位置。二维中还可写成闭曲线所围区域上的积分
\begin{equation}
H(x,y)=\int_V\delta(x-x')\delta(y-y')\,\mathrm dV'. \tag{2.38}
\end{equation}
点在轮廓内部时积分中存在与它重合的 $(x',y')$，故值为一；点在外部时没有这样的点，故值为零。其梯度是界面上的奇异分布：
\begin{equation}
 \nabla H=-\delta_S\bm n. \tag{2.40}
\end{equation}

\begin{sourceblock}
\captionof{figure}{用阶跃函数辨认特定流体所占区域；界面就是 $H$ 从 1 跃变到 0 的位置。}
\end{sourceblock}

% [原PDF p46]
阶跃函数 $H$ 的性质在附录 B 进一步讨论。二维时，$\nabla H$ 可写为沿界面的 Dirac 分布积分；在与界面对齐的局部坐标中，它等于 $-\delta(n)\bm n$。定义曲面分布 $\delta_S(\bm x)=\delta(n)$，便得到紧凑关系 $\nabla H=-\delta_S\bm n$。这个公式把全域标记函数的跃变与界面上的法向集中分布直接联系起来，是单流体表述的基础。

另一条路线是光滑水平集函数 $F$：指定某一等值面，例如 $F=0$，作为界面，界面两侧分别有 $F<0$ 和 $F>0$。若参考流体 1 位于闭曲面内部并取 $F>0$，则外法向为 $-\nabla F/|\nabla F|$；再取其散度即可得到曲率。这种把界面认作特定等值面的做法正是水平集方法的几何表示。

还可用光滑函数 $F$ 的零等值面表示界面。若 $F>0$ 位于参考相内部，则
\begin{equation}
 \bm n=-\frac{\nabla F}{|\nabla F|},\qquad
 \kappa=\nabla\cdot\left(\frac{\nabla F}{|\nabla F|}\right). \tag{2.41--2.42}
\end{equation}
这正是水平集法采用的表示。

\begin{sourceblock}
\captionof{figure}{用水平集函数辨认界面：界面是 $F=0$ 的等值面。}
\end{sourceblock}

% [原PDF p47]
参数曲面每一点的运动由法向速度 $V(u,v,t)$ 决定。它可以就是流体速度的法向分量，也可因蒸发或凝结而与流体不同。推进几何界面无需指定切向速度：切向运动只会重新分布参数点，不会改变界面集合本身。虽然从物理上说，构成界面的流体质点一般同时具有切向和法向速度，但几何演化只依赖后者。

从水平集观点可清楚看出这一点。若 $F$ 随流体运动，则物质导数 $DF/Dt=0$；把 $\nabla F$ 写成法向与梯度模的乘积后，输运方程化为 $\partial_tF-V|\nabla F|=0$，其中 $V=\bm u\cdot\bm n$，切向速度完全不出现。这种等值面输运形式也为第 4 章的界面平流方法奠定基础。

界面运动由每一点的法向速度决定；该速度可以等于流体法向速度，也可因蒸发或凝结而不同：
\begin{equation}
 V(u,v,t)=\bm n\cdot\frac{\partial\bm x}{\partial t}. \tag{2.43}
\end{equation}
切向速度只改变界面参数点的分布，不改变几何形状。若 $F$ 随流体输运，则
\begin{equation}
 \frac{DF}{Dt}=\frac{\partial F}{\partial t}+\bm u\cdot\nabla F=0,\qquad
 \frac{\partial F}{\partial t}-V|\nabla F|=0,\quad V=\bm u\cdot\bm n. \tag{2.44--2.45}
\end{equation}

\section{含界面流动的守恒关系}

% [原PDF p48]

考虑两种互不相溶不可压缩流体。每一相内部都满足通常的 Navier--Stokes 方程；下面由跨越零厚度界面的薄控制体导出两侧解的耦合条件。

\subsection{质量守恒与速度条件}

图 2.9 所示薄控制体 $\delta V$ 跨过界面 $S$，边界为 $\delta S$；令其厚度趋于零，所以控制体内部不能积累质量。设界面两侧速度分别为 $\bm u_1$、$\bm u_2$，界面法向速度为 $V$，法向 $\bm n$ 从流体 1 指向流体 2。相对随界面运动的控制体，入流法向速度为 $\bm u_1\cdot\bm n-V$，出流法向速度为 $\bm u_2\cdot\bm n-V$，因此跨界面质量流率 $\dot m$ 满足 Rankine--Hugoniot 条件
\begin{equation}
 \rho_1(\bm u_1\cdot\bm n-V)=\rho_2(\bm u_2\cdot\bm n-V)=\dot m. \tag{2.46}
\end{equation}
无相变时 $\dot m=0$，故
\begin{equation}
 V=\bm u_1\cdot\bm n=\bm u_2\cdot\bm n. \tag{2.47}
\end{equation}
质量守恒本身不约束切向速度。无黏流可在界面滑移；正常条件下的黏性流体则像固壁无滑移一样，实验上没有有限滑移。若把尖锐界面视为有限厚度过渡层的零厚度极限，有限黏度下的切向速度跳跃还会造成无限大应力。因此无相变的不可压缩黏性两流体满足
\begin{equation}
 [\bm u]_S=\bm0,\qquad [a]_S\equiv a_2-a_1. \tag{2.48}
\end{equation}

\begin{sourceblock}
\captionof{figure}{包含一部分界面 $S$ 的薄控制体 $\delta V$。控制体随界面法向速度 $V$ 运动；令厚度趋于零，因此其中不发生积累。}
\end{sourceblock}

\subsection{表面张力}

% [原PDF p49]

从分子观点看，界面处的热力学状态并非自由能最低状态，因此单位面积具有过剩能量：
\begin{equation}
 \mathrm dE_\sigma=\sigma\,\mathrm dS, \tag{2.49}
\end{equation}
其中 $\sigma$ 为表面张力系数。从力学观点看，表面张力是作用于曲面内任意线元、垂直于该线元且位于切平面内的单位长度拉力。若 $\bm p$ 为这个方向的单位向量，则单位长度“拉力”为 $\sigma\bm p$。两种观点等价：拉伸曲面必然要对其做功。若沿 $x$ 方向拉伸，而曲面在另一方向宽度为 $L$，则面积增量 $\mathrm dS=L\,\mathrm dx$，所做功 $\sigma L\,\mathrm dx=\sigma\,\mathrm dS$，恰好等于界面自由能的增加。

仿照三维连续介质应力张量的引入，可把小曲面元边缘上的力写成 $\bm T^S_\sigma\cdot\bm p$。由于
\begin{equation}
 \bm T^S_\sigma\cdot\bm p=\sigma\bm p=\sigma\bm I_S\cdot\bm p, \tag{2.50}
\end{equation}
可知 $\bm T^S_\sigma=\sigma\bm I_S$；$\bm I_S$ 是三维恒等张量在切平面上的投影。

设 $\bm I_S=\bm I-\bm n\bm n$ 为切平面投影张量，表面张力张量为
\begin{equation}
 \bm T_\sigma^S=\sigma\bm I_S=\sigma(\bm I-\bm n\bm n)
 =\sigma(\bm t_1\bm t_1+\bm t_2\bm t_2). \tag{2.51--2.52}
\end{equation}
对由闭合曲线 $C$ 围成的曲面片 $S$，边界拉力的合力可转为曲面散度积分：
\begin{equation}
 \delta\bm F_\sigma=\oint_C\bm T_\sigma^S\cdot\bm p\,\mathrm dl
 =\int_S\nabla_S\cdot\bm T_\sigma^S\,\mathrm dS. \tag{2.53}
\end{equation}
令曲面片缩小到一点，得到单位界面面积上的表面力
\begin{equation}
 \bm f_\sigma=\nabla_S\cdot\bm T_\sigma^S
 =\sigma\kappa\bm n+\nabla_S\sigma. \tag{2.55}
\end{equation}
第一项是法向毛细力，第二项是由表面张力梯度产生的切向 Marangoni 力；当 $\sigma$ 为常数时后者为零。

\subsection{跨界面动量守恒}

% [原PDF p50]

对随界面移动的薄控制体应用动量守恒：
\begin{equation}
 0=-\oint_{\delta S}\rho\bm u(\bm u\cdot\bm n-V)\,\mathrm dS
 +\oint_{\delta S}\bm n\cdot\bm T\,\mathrm dS
 +\int_S\bm f_\sigma\,\mathrm dS. \tag{2.56}
\end{equation}
前两项沿薄控制体边界积分，最后一项沿界面部分积分。无相变不可压缩流中，控制体以流体法向速度移动，第一项为零。厚度趋于零后，边界与界面重合；先在一侧、再在另一侧积分应力便给出界面两侧应力跳跃。由于结论对包含界面的任意控制体均成立，于是逐点有
\begin{equation}
 -[\bm T]_S\cdot\bm n=\sigma\kappa\bm n+\nabla_S\sigma. \tag{2.57}
\end{equation}
将其分解为法向和切向，并与速度连续条件合并，可写为
\begin{align}
 [\bm u]_S&=\bm0,\\
 -[-p+2\mu\bm n\cdot\bm S\cdot\bm n]_S&=\sigma\kappa,\\
 -[2\mu\bm t^{(k)}\cdot\bm S\cdot\bm n]_S&=\bm t^{(k)}\cdot\nabla_S\sigma. \tag{2.57a--c}
\end{align}
再加上每一相内部的 $\nabla\cdot\bm u=0$ 与 Navier--Stokes 方程，以及外部壁面、入口和出口条件，就构成无蒸发/凝结的不可压缩 Newton 两相流完整问题。

\subsection{自由表面流}

% [原PDF p51]

自由表面是两流体问题的极限情形。例如空气-水流中，有时可把气相压力 $p_{\mathrm{free}}$ 视为常数或仅随时间变化，并忽略气相黏性应力。两相跳跃条件便成为液体表面的边界条件。常表面张力时，切向应力为零：
\begin{equation}
 \left.\bm t^{(k)}\cdot\bm S\cdot\bm n\right|_S=0. \tag{2.58}
\end{equation}
这一切向应力条件是纯运动学的，完全不含物性，因此有一个有趣后果。涡量向量定义为 $\nabla\times\bm u$；二维中只有 $z$ 分量非零，记作 $\omega=(\nabla\times\bm u)\cdot\bm e_z$。二维自由表面上的涡量满足
\begin{equation}
 \omega=2\kappa q,\qquad q=|\bm u|. \tag{2.59}
\end{equation}
因此高速且高曲率处会出现强涡量，也正是无旋假设偏离最明显的位置。密度、黏度比很大的两相流近似自由表面流，也会在高曲率区域形成特征涡量斑块。

\section{单流体表述}

% [原PDF p52]

上一节在每一相分别写方程，并以跳跃条件耦合。另一种做法是在整个区域只写一套控制方程，把各相视为物性在界面突变的同一种“流体”。界面上的附加作用以 Dirac $\delta$ 分布加入方程。这些奇异项正是跳跃条件的对应物，两种表述在数学上等价。

由于解可以跨界面不连续，方程应在弱意义下理解，即只要求其积分形式成立；也可以允许解包含阶跃函数、$\delta$ 函数等广义函数，本书采用后一观点。固定网格多相算法多以这种单流体表述为起点。其推导与第 2.2 节完全相同，唯一差别是必须把表面张力作为体力加入动量方程。总质量守恒没有额外源项；发生化学反应时，各组分守恒方程可有局域源；相变流动中，第 11 章还会看到相界面处自然出现局域体积源。

% [原PDF p53]
对包含界面的控制体，表面张力是 $\bm f_\sigma$ 在控制体所截界面 $S$ 上的积分。利用表面分布 $\delta_S$，就像用散度定理把表面积分化成体积分一样，可将控制体内的界面力改写为
\begin{equation}
 \int_S\bm f_\sigma\,\mathrm dS=\int_V\bm f_\sigma\delta_S\,\mathrm dV. \tag{2.60}
\end{equation}
不可压缩 Newton 流体的单流体 Navier--Stokes 方程为
\begin{equation}
 \rho\frac{\partial\bm u}{\partial t}+\rho\nabla\cdot(\bm u\bm u)
 =-\nabla p+\bm f+\nabla\cdot\{\mu[\nabla\bm u+(\nabla\bm u)^T]\}
 +\bm f_\sigma\delta_S. \tag{2.61}
\end{equation}
即使 $\rho$ 和 $\mu$ 跨界面不连续，该式仍在全流场成立。常表面张力时
\begin{equation}
 \bm f_\sigma\delta_S=\sigma\kappa\bm n\delta_S. \tag{2.62}
\end{equation}
也可把它表示为界面上的三维 Dirac 分布积分：
\begin{equation}
\bm f_\sigma\delta_S=\int_S\bm f_\sigma\delta(\bm x-\bm x_s)\,\mathrm dA,
\qquad
\sigma\kappa\bm n\delta_S=\int_S\sigma\kappa\bm n\delta(\bm x-\bm x_s)\,\mathrm dA. \tag{2.63}
\end{equation}
其中 $\delta(\bm x-\bm x_s)=\delta(x-x_s)\delta(y-y_s)\delta(z-z_s)$。将任意常数 $\Delta p_0$ 加到 $\sigma\kappa$ 上，只会相应改变界面压强跳跃；因为两侧流动由压强梯度而非绝对压强驱动，所以速度场不变。Jan（1994）曾利用这一自由度，从气泡内部压力中减去平均值，以减小洁净或受污染轴对称和二维气泡运动计算中的界面压力跳跃。

\section{无量纲数}

% [原PDF p54]

多流体/多相流受多种无量纲数控制。以参考流体物性、长度 $L$ 和速度 $U$ 定义 Reynolds 数
\begin{equation}
 Re=\frac{\rho LU}{\mu}=\frac{LU}{\nu}, \tag{2.64}
\end{equation}
衡量惯性力与黏性力之比。若以 $[(\Delta\rho/\rho)gL]^{1/2}$ 为速度尺度，$Re^2$ 变为 Galileo 数
\begin{equation}
 N=\frac{g\rho\Delta\rho L^3}{\mu^2}. \tag{2.65}
\end{equation}
$N$ 大时重力主要与惯性和湍流耗散平衡，$N$ 小时重力与黏性力平衡。

惯性与表面张力之比为 Weber 数，黏性应力与毛细应力之比为毛细数：
\begin{equation}
 We=\frac{\rho LU^2}{\sigma},\qquad Ca=\frac{\mu U}{\sigma}. \tag{2.66--2.67}
\end{equation}
消去 $U$ 得 Ohnesorge 数
\begin{equation}
 Oh=\frac{\mu}{\sqrt{\rho\sigma L}},\qquad La=Oh^{-2}. \tag{2.68}
\end{equation}
$Ca$ 和 $Oh$ 都比较黏性与毛细作用，看起来似乎重复；区别在于前者用特征速度估计两类力，例如推进接触线，后者则用特征长度，例如振荡液滴或气泡的半径。E\"otv\"os（Bond）数比较重力与毛细力：
\begin{equation}
 Eo=\frac{\Delta\rho gL^2}{\sigma}. \tag{2.69}
\end{equation}
Morton 数
\begin{equation}
 Mo=\frac{\Delta\rho g\mu^4}{\rho^2\sigma^3} \tag{2.70}
\end{equation}
在给定 $g$ 后只依赖流体物性。在地球标准重力下，若气液密度差并不小，Morton 数也可用来比较黏性与毛细作用。气--水或气--液态金属流动的 $Mo$ 往往极小，表示表面张力相对黏性很强，也使这些流动特别难以计算。

% [原PDF p55]
惯性与重力竞争时，例如表面波问题，常用 Froude 数
\begin{equation}
 Fr=\frac{U^2}{gL}. \tag{2.71}
\end{equation}
若还有传热、传质、相变和化学反应，则需引入更多无量纲数。

\section{薄膜、分子间力与接触线}

多相流的拓扑变化，例如两个液滴聚并为一个、一个液滴破裂成两个，发生在薄膜破裂或细液丝断裂时。细丝是两者中容易处理得多的一种：有充分理由相信 Navier--Stokes 方程能描述其直径如何在有限时间内降为零，而且模拟表明整体结果对具体断裂处理相对不敏感。薄膜则必须加入额外物理。严格说来，分子效应无法纳入连续介质和尖锐界面框架；但略大于微观尺度的介观尺度仍可能由修正的连续模型描述。本节只作最基础的介绍，破裂问题将在第 9 章进一步讨论。

\subsection{分离压与界面间作用力}

% [原PDF p56]

当两界面间距小于数百纳米时，分子间力不可忽略。常把界面单位面积力写成
\begin{equation}
 \bm f_I=-Ah^{-3}\bm n, \tag{2.72}
\end{equation}
其中 $A$ 为 Hamaker 常数，$h$ 为两界面距离，$\bm n$ 指向远离另一界面的方向。自由薄膜和固壁附近薄膜都可能受到这种力。

\begin{sourceblock}
\captionof{figure}{厚度为 $h$ 的薄膜可能受到分子间力，使两个界面出现有效吸引或排斥。}
\end{sourceblock}

\begin{sourceblock}
\captionof{figure}{固壁附近的液膜；按原书符号约定，正的分离压会导致薄膜破裂。}
\end{sourceblock}

表面力于是成为
\begin{equation}
 \bm f_\sigma\delta_S=\left(-Ah^{-3}\bm n+\sigma\kappa\bm n+\nabla_S\sigma\right)\delta_S. \tag{2.73}
\end{equation}
自由表面的法向应力条件改为
\begin{equation}
 (-p+2\mu\bm n\cdot\bm S\cdot\bm n)|_S=-p_{\mathrm{free}}+\sigma\kappa-Ah^{-3}. \tag{2.74}
\end{equation}
静止时
\begin{equation}
 p=p_{\mathrm{free}}-\sigma\kappa+Ah^{-3}. \tag{2.75}
\end{equation}
因此 $\Pi(h)=-Ah^{-3}$ 称为分离压。$A>0$ 时 $\Pi<0$，两界面相吸；$A<0$ 时界面相斥。$A$ 的数值依流体而异，$10^{-20}\,\mathrm J$ 量级很常见。众所周知，$A>0$ 的纯水薄膜会破裂，$A<0$ 的含表面活性剂水膜则可能稳定。正 $A$ 会升高狭窄区域的压力，把流体从那里排出，于是出现灾变失稳并使薄膜破裂。原则上可把式（2.74）直接与宏观 Navier--Stokes 方程联立，但所需尺度范围极大，文献中极少有完全分辨的实现。

\subsection{接触线静力学与动力学}

% [原PDF p57]

两流体界面可在固壁处形成接触线。液-气、液-固和气-固单位面积表面能分别为 $\sigma$、$\sigma_{ls}$ 和 $\sigma_{gs}$。把接触点沿固面移动微小距离并平衡表面能，得 Young 关系
\begin{equation}
 \sigma_{gs}=\sigma_{ls}+\sigma\cos\theta_{\mathrm{eq}}, \tag{2.76}
\end{equation}
其中 $\theta_{\mathrm{eq}}$ 为平衡接触角。若实际角度大于平衡角，接触线向前移动在能量上更有利；小于平衡角时方向相反。不过实际表面通常有接触角迟滞：在后退角与前进角之间的一段范围内，接触线会被表面非均匀性钉扎而不移动。例如窗玻璃上滑动的液滴可能突然停住，此时接触线就钉扎在前进角和后退角之间的某个平衡位置。

\begin{sourceblock}
\captionof{figure}{部分润湿情形 $0<\theta<\pi$ 中向前推进的接触线。}
\end{sourceblock}

接触线移动时，外力和表面能的释放由黏性耗散平衡。困难在于接触线附近耗散奇异。小角度液体楔形区可写成
\begin{equation}
 h(x,t)=\theta(x_c-x), \tag{2.77}
\end{equation}
若接触线以速度 $U$ 前进，固壁无滑移而界面附近速度量级为 $U$，故
\begin{equation}
 \frac{\partial u}{\partial y}\sim\frac{U}{h}. \tag{2.78}
\end{equation}
距接触线 $a$ 范围内的耗散为
\begin{equation}
 \varepsilon=\int_{x_c}^{x_c+a}\int_0^{h(x,t)}\mu\left(\frac{\partial u}{\partial y}\right)^2\mathrm dy\,\mathrm dx
 \sim\int_0^a\frac{\mu U^2}{\theta x}\,\mathrm dx, \tag{2.79--2.80}
\end{equation}
在 $x\to0$ 时对数发散。这就是经典接触线悖论：连续介质无滑移模型预言任何接触线运动都耗散无限能量。采用滑移长度 $\beta$ 的部分滑移条件可截断发散：
\begin{equation}
 \varepsilon=\frac{\mu U^2}{\theta}\ln\frac{a}{\beta}. \tag{2.81}
\end{equation}

\section{补充说明}

% [原PDF p58]
本专著的首要目的不是重新编写一部界面流体力学教科书。需要进一步推导的读者可参阅 Batchelor（1970），其中对自由表面和流体界面有较多讨论；界面几何可参阅 Weatherburn（1927），连续介质与张量分析还可参阅 Aris（1962）。式（2.49）的界面本构已经允许表面张力因污染物或温度而变化，但没有包含曲面内剪切或面积膨胀引起的界面黏性应力。若界面带有不可忽略的表面流变性，可以引入更复杂的二维本构方程，相关理论见 Scriven（1960）、Aris（1962）和 Edwards 等（1961）。这些扩展并不改变质量和动量守恒的结构，却会在切向应力跳跃条件中增加新的曲面散度项。


本书并不以推导流体界面控制方程、甚至也不以作教科书式展示为首要目的。需要更系统背景的读者可参阅 Batchelor（1970）的基础流体力学专著，其中对界面问题有较多篇幅；界面几何可参阅 Weatherburn（1927），也可参阅 Aris（1962）。式（2.49）的本构假设允许表面张力因污染物或温度变化而改变，但我们没有包括表面剪切和面积膨胀所产生的应力。更复杂的界面本构可以纳入这些效应，相关讨论见 Scriven（1960）、Aris（1962）以及 Edwards 等（1961）。

% [原PDF p59]
分子间力的模型依赖界面极性、电荷以及界面上是否存在大分子，因此式（2.72）并不是唯一选择。其他形式的分离压以及薄膜动力学的一般介绍可参阅 Oron 等（1997）。一般薄膜方程常写为
\begin{equation}
 \frac{\partial h}{\partial t}=\frac{\partial}{\partial x}\left[f(h)\frac{\partial h}{\partial x}
 +g(h)\frac{\partial^3h}{\partial x^3}+k(h)\right], \tag{2.82}
\end{equation}
其中 $f(h)$、$g(h)$ 和 $k(h)$ 的函数形式取决于具体物理。对于固体基底上的 Hamaker 定律和表面张力，即式（2.74）描述的情形，薄膜会发生破裂，并形成厚度 $h$ 按 $t^{1/5}$ 趋于零的自相似解（Lister 与 Zhang，1999）。

式（2.25）的滑移长度条件最早就出现在 Navier（1823）提出黏性流方程的论文中。Huh 与 Scriven（1971）用它消除接触线耗散悖论后，这一条件重新受到关注。接触线悖论详见 Dussan（1979）；静态和动态润湿的一般问题还可参阅 de Gennes（1985）、de Gennes 等（2003）、Oron 等（1997）、Pomeau（2002）以及 Bonn 等（2009）。

实验常把接触线速度与距尖端 $a$ 处的表观接触角联系起来：
\begin{equation}
 \theta(a)=f(Ca,a/\beta,\ldots). \tag{2.83}
\end{equation}
其中 $\beta$ 是某种微观长度，例如式（2.25）的滑移长度；函数 $f$ 的理论形式仍有争议。完全润湿、平衡角为零且角度较小时，多种近似与理论给出
\begin{equation}
 \theta(a)\simeq[9Ca\ln(a/\beta)]^{1/3}. \tag{2.84}
\end{equation}
对数依赖可理解为源于式（2.81）的对数奇性。该式由 Navier 滑移条件论证，但还存在其他微观理论。Shikmurzaev（1997）提出广义 Navier 滑移模型，并指出 Navier 滑移虽能消除耗散奇性，却不能消除压力奇性；Hervet 与 de Gennes（1984）及 de Gennes（1985）引入分子间力；Pomeau（2002）讨论相场（亦称 Cahn--Hilliard、二阶梯度或弥散界面）模型；Blake 与 Shikmurzaev（2002）提出界面弛豫理论。较新的综述见 Bonn 等（2009）。

% [原PDF p60]

数值计算通常只能在靠壁第一层单元规定接触角，网格尺度 $\Delta x$ 是方法能够“看到”的最小距离，因此有效数值边界条件应具有
\begin{equation}
 \theta_{\mathrm{num}}=f(Ca,\Delta x/\beta,\ldots). \tag{2.85}
\end{equation}
Afkhami 等（2009）讨论过这种可视为数值“有效边界条件”的关系。不过，它并未得到普遍采用；许多研究者直接固定与网格无关的动态接触角，或尝试把网格尺度取得小于滑移长度以回避网格依赖。固定一个与网格无关的动态接触角，或让网格尺度小于滑移长度，都会改变网格收敛性，使用时必须明确相应微观假设。

\chapter{Navier--Stokes 方程的数值解}

% [原PDF p61]
本章把单流体方程的流动求解部分单独拿出来讨论。由于所有相共享同一速度--压力方程组，单相不可压缩流的大量成熟技术都可以复用；多相性主要通过空间变化且可能跳跃的密度、黏度以及集中在界面上的表面力进入。为了把这些内容与界面几何算法清楚分开，本章暂不讨论指示函数如何平流，也不讨论曲率和表面张力如何计算，而是假定相应场已经给定。这样，时间推进、空间通量、压力投影和边界处理的数值一致性就能单独说明，后续再同第 4--7 章的界面模块组合。

第 2 章的单流体表述用一套 Navier--Stokes 方程描述所有流体，因此可直接借用单相流数值方法。新增困难有两个：密度和黏度通常随相而变，必须构造辨识各相的指示函数；还要在界面上计算表面张力。指示函数平流将在第 4--6 章讨论，表面张力算法见第 7 章。本章只讨论变密度、变黏度 Navier--Stokes 方程的求解，采用规则 Cartesian 网格上的有限体积法，并针对本书低速流动假设不可压缩。

任一非定常 Navier--Stokes 求解器都必须决定：（i）存放各离散近似的网格点如何排列；（ii）速度场如何作时间积分；（iii）对流项与黏性项如何离散；（iv）如何求解由不可压缩约束产生的压力方程；（v）如何实现边界条件。这些任务都有多种完成方式；下面给出的路线已广泛用于多相流，能得到合理准确而稳健的数值方法。

把动量方程记为
\begin{equation}
 \rho\frac{\partial\bm u}{\partial t}+\rho\bm A=-\nabla p+\bm D+\bm f, \tag{3.1}
\end{equation}
其中 $\bm A=\nabla\cdot(\bm u\bm u)$ 为对流项，$\bm D=\nabla\cdot\{\mu[\nabla\bm u+(\nabla\bm u)^T]\}$ 为扩散项，$\bm f$ 包含重力和表面张力等其余作用。对流项也可采用式（2.22）的其他等价形式。

% [原PDF p62]
不可压缩流的压力不是独立演化的热力学变量，而是在每个时间步调整到使新速度无散。投影算法先忽略压力梯度求临时速度，再用压力梯度校正；校正步在数学上就是把临时速度投影到无散函数空间。早期 MAC 程序很可能已经采用这种实现，Chorin（1968）和 Yanenko（1971）则明确提出了这一概念。

实际计算通常要求二阶或更高时间精度，但一阶显式格式最便于说明。上标 $n$、$n+1$ 表示时间步起点和终点，$\Delta t$ 是步长，$\bm A_h$ 与 $\bm D_h$ 是对流和扩散项的离散近似，$\nabla_h$ 同时表示离散梯度或散度。动量式在压力已知时能推进速度，连续式却只给出步末速度必须满足的约束，因此需要安排正确的求解顺序。

\section{时间积分}

不可压缩流中，每个时间步结束时都必须调整压力，使速度无散。最常用的是投影法：先忽略压力梯度得到临时速度，该速度通常并不无散；再加入恰当压力梯度进行修正。第二步在形式上把速度投影到无散速度场空间，名称即由此而来。早期不可压缩 MAC 程序很可能已经按这种方式实现，但作为明确概念，投影法由 Chorin（1968）和 Yanenko（1971）提出。

先以一阶显式 Euler 离散说明。令上标 $n$、$n+1$ 分别表示步初和步末，$\nabla_h$、$\bm A_h$、$\bm D_h$ 为离散算子，则
\begin{equation}
 \frac{\bm u^{n+1}-\bm u^n}{\Delta t}+\bm A_h^n
 =\frac1{\rho^n}\left(-\nabla_h p+\bm D_h^n+\bm f^n\right), \tag{3.2}
\end{equation}
并要求
\begin{equation}
 \nabla_h\cdot\bm u^{n+1}=0. \tag{3.3}
\end{equation}
% [原PDF p63]
预测步忽略压力，利用已知的旧时刻速度、物性、体力和离散对流、扩散项求 $\bm u^*$；投影步再加入压力梯度得到 $\bm u^{n+1}$。两步相加严格恢复原动量更新。为求压力，对投影式取散度并用 $\nabla_h\cdot\bm u^{n+1}=0$ 消去新速度，得到变系数 Poisson 方程。压力求得后回代投影式即可获得新速度。

由于上述方法完全显式，步长必须满足稳定性限制。二阶中心扩散在二维要求 $\mu\Delta t/(\rho h^2)\le1/4$，三维右端为 $1/6$，其中 $h$ 是最小网格间距。二阶中心对流配显式向前时间推进在无黏极限本身不稳定，只有黏性项能稳定它，因此还会出现不含网格尺度、却与速度和黏度有关的限制；高 Reynolds 数时这一限制可能十分苛刻。

式（3.2）含速度时间导数，压力已知时可求下一步速度；连续方程却只是新速度必须满足的条件，没有显式压力方程。为在所需信息都已知的顺序下求解，投影法将动量更新拆成预测步
\begin{equation}
 \frac{\bm u^*-\bm u^n}{\Delta t}=-\bm A_h^n+\frac1{\rho^n}(\bm D_h^n+\bm f^n), \tag{3.4}
\end{equation}
和投影步
\begin{equation}
 \frac{\bm u^{n+1}-\bm u^*}{\Delta t}=-\frac1{\rho^n}\nabla_h p. \tag{3.5}
\end{equation}
对后式取散度并用无散条件，得变系数压力 Poisson 方程
\begin{equation}
 \nabla_h\cdot\left(\frac1{\rho^n}\nabla_h p\right)=\frac1{\Delta t}\nabla_h\cdot\bm u^*. \tag{3.6}
\end{equation}
求得 $p$ 后，用式（3.5）得到新时刻速度。

完全显式格式受稳定性限制。二维二阶中心扩散要求
\begin{equation}
 \frac{\mu\Delta t}{\rho h^2}\le\frac14, \tag{3.7}
\end{equation}
三维右端为 $1/6$。二阶中心对流与显式向前时间推进在无黏极限不稳定，黏性稳定作用还要求
\begin{equation}
 (\bm u\cdot\bm u)\frac{\rho\Delta t}{\mu}\le2. \tag{3.8}
\end{equation}
% [原PDF p64]
对流项采用显式而稳定的时空离散时，通常受 CFL 条件限制：一个时间步内物质点的位移不能超过一个网格间距。多维非分裂格式的条件可能更严格；若高级对流格式采用方向分裂，依次完成各坐标方向的平流，通常可沿用一维限制。

一阶投影只具一阶时间精度。最简单的二阶扩展是预测--校正：先用一阶显式法求临时速度，再用临时状态估算步末右端，最后以梯形规则平均步初和临时右端。编程上也可连续做两个一阶步，再平均 $\bm u^n$ 与 $\bm u^{n+2}$；这与上述格式等价，能很方便地把现有一阶程序升级到二阶。其稳定范围通常仍受一阶方法相同限制，低 Reynolds 数时显式黏性步长过小，应改用隐式黏性处理。

稳定的显式对流格式通常服从 Courant--Friedrichs--Lewy（CFL）条件
\begin{equation}
 \frac{u_{\max}\Delta t}{h}\le1, \tag{3.9}
\end{equation}
即一个时间步内物质点不应穿过一个以上网格间距。多维算子分裂逐方向推进时常可沿用一维限制。

二阶预测-校正可先用上述一阶方法得到临时速度，再以梯形规则
\begin{equation}
 \frac{\bm u^{n+1}-\bm u^n}{\Delta t}=\frac12(\mathrm{RHS}^n+\mathrm{RHS}^{\mathrm{tmp}}) \tag{3.10}
\end{equation}
校正。编程时也可连续做两个一阶步后对相应状态平均。低 Reynolds 数下显式黏性限制过严，应隐式处理扩散项。

% [原PDF p65]
二阶 Navier--Stokes 时间推进方法很多，最常见的一类把黏性项隐式处理而把对流项显式处理。下面采用扩展到变物性流动的 Kim--Moin（1985）二步投影法：用 Adams--Bashforth 显式处理对流，用 Crank--Nicolson 型方式半隐式处理黏性：
\begin{equation}
 \frac{\bm u^*-\bm u^n}{\Delta t}
 =-\left(\frac32\bm A_h^n-\frac12\bm A_h^{n-1}\right)
 +\frac1{2\rho^n}(\bm D_h^n+\bm D_h^*), \tag{3.11}
\end{equation}
随后用类似压力的变量 $\Phi$ 投影：
\begin{equation}
 \frac{\bm u^{n+1}-\bm u^*}{\Delta t}=-\frac{\nabla\Phi}{\rho^n}. \tag{3.12}
\end{equation}
$\Phi$ 与真实压力满足
\begin{equation}
 \frac12(\bm D^*-\bm D^{n+1})-\nabla\Phi=-\nabla p. \tag{3.13}
\end{equation}
Kim 与 Moin 使用交错网格和二阶空间差分，并用方向分裂求解隐式方程（3.11）。投影法并非唯一选择；Chorin（1967）的人工可压缩法、Patankar（1980）的 SIMPLE、Issa（1985）的 PISO 等几乎所有均相 Navier--Stokes 方法，在单流体表述下都可推广到多相流。

强毛细流还受毛细波稳定性限制。波数 $k$ 的毛细波相速为
\begin{equation}
 c=\left(\frac{\sigma k}{\rho_1+\rho_2}\right)^{1/2}. \tag{3.14}
\end{equation}
最短可分辨波取 $k=\pi/h$，令其一步不跨过一个单元，得
\begin{equation}
 \Delta t\le\left[\frac{(\rho_1+\rho_2)h^3}{\pi\sigma}\right]^{1/2}. \tag{3.15}
\end{equation}

% [原PDF p66]
有限体积离散把质量和动量守恒直接施加于小控制体。可以从方框 2.1 的积分守恒式出发，也可把方框 2.4 的微分方程在控制体上积分，两条路线得到相同结构。控制体中心速度以体平均近似；二维矩形控制体体积为 $V=\Delta x\Delta y$。对流和黏性项的体积分用散度定理化为控制体边界上的通量积分，再用数值求积得到 $\bm A_h$、$\bm D_h$。压力梯度同样转为压力乘法向的表面积分，体力则保留体积平均。连续条件在新时刻直接要求各面的法向体积流率代数和为零。这种写法使每个面通量由相邻控制体共享，因而自然具有离散守恒性。

\section{空间离散}

有限体积法把质量和动量守恒施加于小控制体。中心速度用体平均近似：
\begin{equation}
 \bm u_c=\frac1V\int_V\bm u(\bm x)\,\mathrm dV. \tag{3.16}
\end{equation}
用散度定理，对流项和黏性项的控制体平均分别为
\begin{align}
 \bm A_c&=\frac1V\int_V\nabla\cdot(\bm u\bm u)\,\mathrm dV
 =\frac1V\oint_S\bm u(\bm u\cdot\bm n)\,\mathrm dS, \tag{3.17}\\
 \bm D_c&=\frac1V\int_V\nabla\cdot\bm T_v\,\mathrm dV
 =\frac1V\oint_S\bm n\cdot\bm T_v\,\mathrm dS, \tag{3.18}\\
 \bm T_v&=\mu[\nabla_h\bm u+(\nabla_h\bm u)^T]. \tag{3.19}
\end{align}
压力梯度和体力平均为
\begin{equation}
 (\nabla p)_c=\frac1V\oint_Sp\bm n\,\mathrm dS,
 \qquad \bm f_c=\frac1V\int_V\bm f\,\mathrm dV, \tag{3.20--3.21}
\end{equation}
连续方程则直接以积分形式离散：
\begin{equation}
 \oint_S\bm u^{n+1}\cdot\bm n\,\mathrm dS=0. \tag{3.22}
\end{equation}

% [原PDF p67]
在导出离散式前必须决定控制体形状和各变量相互位置。规则 Cartesian 网格上可把所有量放在同一节点，这称为同位网格；但经验表明，不可压缩流通常采用每类变量位于不同子网格的交错 MAC 网格更简单，直线网格尤其如此。为避免表述繁琐，以下限于二维；推广至三维并不困难。

Harlow 与 Welch（1965）为 MAC 方法引入交错网格，此后它一直广泛用于 CFD。二维中，压力 $p_{i,j}$、密度和物性存于压力控制体中心；水平速度 $u_{i+1/2,j}$ 存于左右面中心；竖直速度 $v_{i,j+1/2}$ 存于上下表面中心。不可压缩流中，压力的作用是迫使速度散度为零：若控制体有净流入就应升高压力，有净流出则应降低压力。这样计算压力控制体净流量时，所需速度恰好就在相应表面。

\begin{sourceblock}
\captionof{figure}{标准交错 MAC 网格的记号。压力位于粗线所示控制体中心；水平速度 $u$ 位于左右边中点，竖直速度 $v$ 位于上下边中点。}
\end{sourceblock}

% [原PDF p68]
水平速度控制体相对压力控制体向右移动半格，竖直速度控制体向上移动半格；半整数下标正表示压力控制体边界上的存储位置。实际程序的数组仍用整数下标，只需记住各数组的几何原点不同；密度与其他物性通常存放在压力节点。

离散连续方程从每个压力控制体的边界积分得到。左右两个竖直面贡献水平速度通量，上下两个水平面贡献竖直速度通量；用中点规则计算各面积分，再除以 $\Delta x\Delta y$，就得到式（3.25）的中心散度。它也可以直接由离散无散条件推出。关键是四个面速度恰好存放在需要求通量的位置，因此无须额外插值，压力控制体的质量守恒关系非常紧凑。

水平速度网格相对压力网格向右平移半个单元，竖直速度网格向上平移半个单元；因此通常用 $(i,j)$ 标记压力，用 $(i+1/2,j)$ 和 $(i,j+1/2)$ 分别标记 $u$ 与 $v$。实际程序中各数组当然仍用整数下标，只是物理位置互相错开。密度及其他物性通常存于压力节点。

\begin{sourceblock}
\captionof{figure}{$u$、$v$ 速度分量的控制体分别相对压力控制体向右、向上平移半个网格。半整数下标表示压力控制体边界上的变量。}
\end{sourceblock}

对压力控制体各面用中点积分，二维离散连续方程为
\begin{equation}
 \frac{u_{i+1/2,j}^{n+1}-u_{i-1/2,j}^{n+1}}{\Delta x}
 +\frac{v_{i,j+1/2}^{n+1}-v_{i,j-1/2}^{n+1}}{\Delta y}=0. \tag{3.25}
\end{equation}
% [原PDF p69]
预测速度方程在各自速度控制体中心计算。体力可分为普通体力 $\bm f_b$ 与界面表面张力 $\bm f_\sigma$；扩散和表面力除以相邻压力节点密度的算术平均，再同显式对流项一起推进临时速度。投影校正使用相邻压力点差分，并使用同一个面密度，这种一致性使压力工作与离散质量通量相匹配。凡某个量没有直接存储在所需位置，例如面中心密度，都以相邻值线性平均。

交错网格的三项优势分别是有效压力梯度分辨率较高、守恒通量容易构造，以及压力--速度耦合紧密。最后一点最重要：同位网格的中心连续式可能只连接隔点，从而允许棋盘格式的高频压力或速度仍满足离散连续条件。Rhie--Chow（1983）等插值能修正同位网格，但交错网格从变量布局上直接避免了这一解耦。

预测速度可写为
\begin{align}
 u_{i+1/2,j}^*&=u_{i+1/2,j}^n+\Delta t\left[-(A_x)^n_{i+1/2,j}+(f_{bx})^n_{i+1/2,j}
 +\frac{(D_x)^n_{i+1/2,j}+(f_{\sigma x})^n_{i+1/2,j}}{(\rho_{i+1,j}^n+\rho_{i,j}^n)/2}\right], \tag{3.26}\\
 v_{i,j+1/2}^*&=v_{i,j+1/2}^n+\Delta t\left[-(A_y)^n_{i,j+1/2}+(f_{by})^n_{i,j+1/2}
 +\frac{(D_y)^n_{i,j+1/2}+(f_{\sigma y})^n_{i,j+1/2}}{(\rho_{i,j+1}^n+\rho_{i,j}^n)/2}\right]. \tag{3.27}
\end{align}
投影校正为
\begin{align}
 u_{i+1/2,j}^{n+1}&=u_{i+1/2,j}^*-
 \frac{\Delta t}{(\rho_{i+1,j}^n+\rho_{i,j}^n)/2}\frac{p_{i+1,j}-p_{i,j}}{\Delta x}, \tag{3.28}\\
 v_{i,j+1/2}^{n+1}&=v_{i,j+1/2}^*-
 \frac{\Delta t}{(\rho_{i,j+1}^n+\rho_{i,j}^n)/2}\frac{p_{i,j+1}-p_{i,j}}{\Delta y}. \tag{3.29}
\end{align}
未在所需位置定义的物性用线性插值，即取相邻值平均。交错网格广泛使用有三点原因。第一是精度：压力梯度由相邻压力点之差计算，而同位网格若直接按对应速度位置处理常跨越两个 $\Delta x$ 或 $\Delta y$；不过其他导数仍可能用较远点，所以它不等价于把整个同位网格加密一倍。第二是守恒格式容易构造。第三且可能最重要的是变量耦合紧密。同位网格的连续方程可能只连接隔点，因而速度即使剧烈振荡仍能精确满足连续条件；虽然 Rhie--Chow（1983）等方法可克服这一问题，交错布置仍更直接。

% [原PDF p70]
\section{对流项离散}

原始 MAC 方法对空间项采用中心差分，时间上采用一阶显式投影；后来的实现常对对流项使用一阶迎风或供体单元格式。有限体积中，对流通量的两个分量由速度控制体边界积分得到：
\begin{align}
 (A_x)_{i+1/2,j}&=\frac{(uu)_{i+1,j}-(uu)_{i,j}}{\Delta x}
 +\frac{(uv)_{i+1/2,j+1/2}-(uv)_{i+1/2,j-1/2}}{\Delta y}, \tag{3.30}\\
 (A_y)_{i,j+1/2}&=\frac{(uv)_{i+1/2,j+1/2}-(uv)_{i-1/2,j+1/2}}{\Delta x}
 +\frac{(vv)_{i,j+1}-(vv)_{i,j}}{\Delta y}. \tag{3.31}
\end{align}
各格式的区别主要在于如何从控制体中心值构造边界速度与动量通量。

二阶中心格式以相邻值线性插值得到面速度，再计算通量。它在完全分辨流动中通常比非中心二阶格式更准确，却有两个严重缺点：未充分分辨的陡梯度会产生非物理振荡；与显式向前时间推进组合时，无黏极限无条件不稳定，只能依靠黏性稳定，因此高 Reynolds 数下时间步往往过小。

最简单的稳健选择是一阶迎风，例如
\begin{equation}
 u_{i,j}=\begin{cases}
 u_{i-1/2,j},&\tfrac12(u_{i-1/2,j}+u_{i+1/2,j})>0,\\
 u_{i+1/2,j},&\tfrac12(u_{i-1/2,j}+u_{i+1/2,j})<0.
 \end{cases} \tag{3.34}
\end{equation}
% [原PDF p71]
它在零黏度时也稳定，但只有一阶精度并产生过量数值扩散，第 4 章会看到这一点。QUICK（quadratic upstream interpolation for convective kinematics，对流运动学二次迎风插值）格式由 Leonard（1979）提出，使用三阶迎风偏置多项式，例如正向流动时
\begin{equation}
 u_{i,j}=\frac18(3u_{i+1/2,j}+6u_{i-1/2,j}-u_{i-3/2,j}), \tag{3.35}
\end{equation}
反向流动取镜像模板。QUICK 在足够陡的梯度处仍可能出现小幅振荡，但比中心格式稳健、比一阶迎风准确。

更稳健的 ENO（essentially non-oscillatory，本质无振荡）格式由 Shu 与 Osher（1989）提出。二阶 ENO 从迎风侧中心值外推到面：
\begin{equation}
 u_{i,j}=\begin{cases}
 u_{i-1/2,j}+s_i\Delta x/2,&\bar u_i>0,\\
 u_{i+1/2,j}-s_{i+1}\Delta x/2,&\bar u_i<0,
 \end{cases} \tag{3.36}
\end{equation}
其中 $\bar u_i=(u_{i-1/2,j}+u_{i+1/2,j})/2$，斜率 $s_i$ 在
\begin{equation}
 s_i^+=\frac{u_{i+1/2,j}-u_{i-1/2,j}}{\Delta x},\qquad
 s_i^-=\frac{u_{i-1/2,j}-u_{i-3/2,j}}{\Delta x} \tag{3.37}
\end{equation}
中选择绝对值较小者，这与 Sweby（1984）的 MINMOD 限制器一致。Bell--Colella--Glaz（BCG；Bell 等，1989；另见 Martin，1998）格式是相关变体，Popinet（2003，2008，2009）的 Gerris 软件就采用它，本书若干模拟也由 Gerris 完成。QUICK 或 ENO 的面速度既可代入守恒离散，也常用于非守恒对流项。两者都显著增强低黏流动计算的稳健性，在黏性足以使中心格式可靠时又能保持相近精度。

% [原PDF p72]
二阶 ENO 面值从迎风控制体中心外推。斜率在相邻两段差分中选择绝对值较小者，这与 MINMOD 限制器一致。Bell--Colella--Glaz 格式是相关变体，Gerris 求解器采用了它。QUICK 和 ENO 得到的面速度既能用于守恒通量，也常用于非守恒对流离散；它们在低黏度下显著提高稳健性，而在中心差分本来可靠的黏度范围内仍保持相近精度。

黏性项仍按有限体积通量计算。对水平速度控制体，$x$ 方向需要左右面的 $T^{v,xx}$，$y$ 方向需要上下边的 $T^{v,xy}$；竖直速度控制体则需要 $T^{v,xy}$ 和 $T^{v,yy}$。面通量采用中点规则，应力中的速度导数采用标准二阶中心差分，同时保留黏度在空间上的变化，故每个应力都在其自然几何位置乘以局部黏度。

\section{黏性项}

黏性通量同样在速度控制体边界上用中点规则近似。二维分量为
\begin{align}
 (D_x)_{i+1/2,j}&=\frac{T^{v,xx}_{i+1,j}-T^{v,xx}_{i,j}}{\Delta x}
 +\frac{T^{v,xy}_{i+1/2,j+1/2}-T^{v,xy}_{i+1/2,j-1/2}}{\Delta y}, \tag{3.38}\\
 (D_y)_{i,j+1/2}&=\frac{T^{v,xy}_{i+1/2,j+1/2}-T^{v,xy}_{i-1/2,j+1/2}}{\Delta x}
 +\frac{T^{v,yy}_{i,j+1}-T^{v,yy}_{i,j}}{\Delta y}. \tag{3.39}
\end{align}
应力中的导数采用二阶中心差分，同时保留空间变黏度。例如
\begin{align}
 T^{v,xx}_{i,j}&=2\mu_{i,j}\frac{u_{i+1/2,j}-u_{i-1/2,j}}{\Delta x},\\
 T^{v,xy}_{i+1/2,j+1/2}&=\mu_{i+1/2,j+1/2}\left[
 \frac{u_{i+1/2,j+1}-u_{i+1/2,j}}{\Delta y}
 +\frac{v_{i+1,j+1/2}-v_{i,j+1/2}}{\Delta x}\right],\\
 T^{v,yy}_{i,j}&=2\mu_{i,j}\frac{v_{i,j+1/2}-v_{i,j-1/2}}{\Delta y}. \tag{3.40}
\end{align}

图 3.3 示意水平速度控制体以及计算 $T^{v,xx}$、$T^{v,xy}$ 的位置；界面切过控制体，灰色区域表示流体 1。

\begin{sourceblock}
\captionof{figure}{交错网格上的黏性应力计算。图中画出 $u_{i+1/2,j}$ 的控制体；$T^{xx}$ 与 $T^{xy}$ 分别需在相应面中心和角点求取，灰色表示流体 1。}
\end{sourceblock}

% [原PDF p73]
式（3.40）允许黏度随空间变化。若两相内黏度各为常数而在界面突变，精确连续表示是 $\mu(\bm x)=\mu_1H+\mu_2(1-H)$。数值方法用近似 Heaviside 的标记函数代替 $H$：VOF 在压力节点以体积分数 $C_{i,j}$ 作算术混合；界面追踪或水平集通常先构造平滑指示函数，再由它设置黏度。黏度差小时这一做法通常良好；差异很大时却会产生误差，因为界面应力条件要求速度导数发生相应跳跃，而跨界面的二阶中心导数并不一致。

提高精度的一种方案，是只使用界面同侧数据作单边速度导数；已有少数实现。另一种是借鉴 Patankar（1980）处理热扩散以及 Coward 等（1997）处理黏性剪切流的思路，对跨界面扩散系数采用调和平均。调和平均的依据可由水平界面两相平行剪切流严格导出。

若两相黏度为常数但界面处突变，连续表示为
\begin{equation}
 \mu(\bm x)=\mu_1H(\bm x)+\mu_2[1-H(\bm x)]. \tag{3.41}
\end{equation}
VOF 中可用体积分数 $C_{i,j}$ 代替 $H$：
\begin{equation}
 \mu_{i,j}=\mu_1C_{i,j}+\mu_2(1-C_{i,j}). \tag{3.42}
\end{equation}
小黏度比时算术混合通常足够；黏度差大时，界面条件要求速度梯度跳跃，跨界面的中心导数便不一致。可采用界面单侧差分，或在适当方向使用调和平均。

% [原PDF p74]
为说明调和平均，考虑图 3.4 中水平界面分隔的两相平行剪切流。界面既与网格方向对齐，也与流向平行；剪应力在界面连续且为常数。若界面距下方速度节点为 $h$，相邻节点间距为 $\Delta y$，则
\begin{equation}
 \tau=\mu_2\frac{u_{j+1}-u_I}{\Delta y-h}=\mu_1\frac{u_I-u_j}{h}. \tag{3.45}
\end{equation}
消去界面速度 $u_I$，跨界面单元的等效黏度为
\begin{equation}
 \mu_{i+1/2,j+1/2}=\left(\frac{\phi}{\mu_1}+\frac{1-\phi}{\mu_2}\right)^{-1}, \tag{3.48}
\end{equation}
其中 $\phi=h/\Delta y$ 为相 1 在该控制体中的比例。对 VOF，可由邻近色函数估计 $\phi=C_j+C_{j+1}-1/2$；只有 $1/2\le C_j+C_{j+1}\le3/2$ 时界面才位于这两层之间，其他情况应使用单相公式。但调和平均并非普适：若界面改为竖直，算术平均反而精确；一般取向界面的一致二阶处理更复杂，本书写作时尚无公开的通用二阶方法。

\begin{sourceblock}
\captionof{figure}{水平界面分隔的两相平行剪切流。黏度跳跃使速度梯度不连续，而剪应力连续。}
\end{sourceblock}

实际剪切流中界面常与流向平行，调和平均通常更准确，但稳健性可能较差。大物性比时，算术平均偏向大黏度，相当于把有效界面推向低黏侧，形成一层较黏的保护区，可能抑制短波物理或数值失稳；调和平均则偏向小黏度。隐式黏性离散只需把相应速度时间层从 $n$ 改为 $n+1$，但线性系统求解也随之成为时间步的一部分。

% [原PDF p75]
调和平均的推导还说明了为什么必须关注界面相对于网格和流向的几何位置。在图 3.4 的平行剪切流中，剪应力不仅跨界面连续，而且沿整个横向为常数；因此每一相中的速度梯度分别与其黏度成反比。把跨界面的一整格简单当成均匀介质时，正确的等效系数必须使两段速度降相加后仍给出同一剪应力，这正产生“厚度除以黏度”相加的调和形式，而不是黏度本身的算术相加。

实际剪切流的界面往往沿流向排列，所以调和平均通常更准确；然而它偏向较小黏度，可能比算术平均更不稳健。算术平均偏向较大黏度，相当于把有效界面向低黏一侧移动，在真实界面周围形成一层较黏的数值保护区，可能压制短波物理失稳或数值失稳。若黏性项改为隐式计算，离散结构不变，只需把相关速度由时间层 $n$ 改为 $n+1$，同时求解相应线性系统。

消去界面速度 $u_I$ 后，相邻两速度节点之差等于两段速度降之和；每段速度降是恒定剪应力乘以该相厚度再除以相应黏度。因而跨界面控制体的等效黏度是以相占比为权的调和平均。若 $\phi=h/\Delta y$ 表示流体 1 所占比例，则 $\mu_{\mathrm{eff}}=(\phi/\mu_1+(1-\phi)/\mu_2)^{-1}$。

VOF 的 $C_{i,j}$ 定义在压力节点，可用二阶近似 $\phi=C_j+C_{j+1}-1/2$；只有 $1/2\le C_j+C_{j+1}\le3/2$ 时，界面才确实位于两层之间，否则应使用单相公式。调和平均不是万能办法：若界面从水平改为竖直，算术平均反而精确；一般取向的一致二阶处理在原书写作时尚无公开通用方案。剪切流中界面常与流向对齐，所以调和平均更准确，却可能不如偏向大黏度的算术平均稳健。

\section{压力方程}

把投影校正式（3.28）、（3.29）代入离散连续方程（3.25），得到压力的变系数五点 Poisson 方程：
\begin{align}
 &\frac1{\Delta x^2}\left[
 \frac{p_{i+1,j}-p_{i,j}}{\rho_{i+1,j}^n+\rho_{i,j}^n}
 -\frac{p_{i,j}-p_{i-1,j}}{\rho_{i,j}^n+\rho_{i-1,j}^n}\right] \notag\\
 &+\frac1{\Delta y^2}\left[
 \frac{p_{i,j+1}-p_{i,j}}{\rho_{i,j+1}^n+\rho_{i,j}^n}
 -\frac{p_{i,j}-p_{i,j-1}}{\rho_{i,j}^n+\rho_{i,j-1}^n}\right] \notag\\
 &=\frac1{2\Delta t}\left[
 \frac{u_{i+1/2,j}^*-u_{i-1/2,j}^*}{\Delta x}
 +\frac{v_{i,j+1/2}^*-v_{i,j-1/2}^*}{\Delta y}\right]. \tag{3.49}
\end{align}
% [原PDF p76]
分母中的因子 2 已与右端写法一致。交错网格使边界方程的推导很直接。若图 3.5 的左边界穿过 $(i-1/2,j)$，法向速度给定为 $U_{b,j}$，边界相邻压力控制体的连续方程为
\begin{equation}
 \frac{u_{i+1/2,j}^{n+1}-U_{b,j}}{\Delta x}
 +\frac{v_{i,j+1/2}^{n+1}-v_{i,j-1/2}^{n+1}}{\Delta y}=0. \tag{3.50}
\end{equation}
只对上、下、右三个未知面速度代入投影校正式，即得到边界相邻压力节点的方程；不必人为另加压力边界条件，仅用不可压缩条件就自然得到正确的边界方程。角点和其他边界同理。

\begin{sourceblock}
\captionof{figure}{竖直边界附近压力方程的推导。边界一侧给定法向速度，压力控制体的连续方程只对其余三个未知面速度作投影校正。}
\end{sourceblock}

由于边界上不指定压力绝对值，Poisson 方程只确定到任意常数。可规定全域平均压力或某一点压力；很多迭代法也会由初值平均自动固定这一常数。若所有边界都是入口或不可穿透壁面，还必须保证总边界流量 $\sum\bm U_b\cdot\bm n=0$，否则压力迭代不可能收敛。

% [原PDF p77]
式（3.51）适用于固壁和指定入口速度的边界；出口另见第 3.7 节。代码开发和小规模测试时，逐次超松弛（SOR）最方便。把每个节点的 $p_{i,j}$ 从方程中解出，用上一迭代近似计算新值，再取
\begin{equation}
 p_{i,j}^{\alpha+1}=\beta\,p_{i,j}^{\mathrm{GS}}+(1-\beta)p_{i,j}^{\alpha},\qquad 1<\beta<2, \tag{3.52}
\end{equation}
其中 $p^{\mathrm{GS}}$ 为以最新邻点值计算的 Gauss--Seidel 更新。$\beta=1.2$--$1.5$ 往往是稳定性与收敛速度的良好折中。初始时间可令压力为零，随后以前一时间步压力为初猜。向量机或并行实现常用红黑 SOR。

SOR 的主要优点是简单，缺点是收敛慢，而压力求解通常是不可压缩流模拟最耗时的部分。它适合代码开发或偶尔求解的小问题，生产计算应采用快速 Poisson、循环约化、多重网格或 Krylov 子空间方法。20 世纪 80 年代，离散 FFT 或循环约化的快速 Poisson 求解器非常流行，部分原因是 FISHPAK（Swarztrauber 与 Sweet，1979）提供了可自由使用的成熟例程；但它们只适合简单边界和可分离常系数椭圆算子，不能直接处理式（3.49）的变密度系数。

% [原PDF p78]
SOR 的主要优点是简单，却通常收敛过慢；压力求解又往往是不可压缩模拟最耗时的部分，因此生产计算需要更先进的椭圆求解器。离散 FFT 或循环约化的快速 Poisson 法在 20 世纪 80 年代很流行，FISHPAK 提供了成熟例程，但它只适合简单边界和可分离的常系数算子，无法直接处理变密度压力方程。

多重网格的思想可由把 Poisson 迭代视作扩散方程的显式伪时间推进来理解。最大稳定伪时间步正比于 $\Delta x^2$，细网格被迫取小步；与此同时，高波数误差衰减很快，长期残留的是低波数误差。继续在细网格上迭代，就等于用小步去分辨已经消失的高波数。多重网格先在细网格平滑高频误差，再把残差转移到粗网格，以较大步长消除低频误差；逐级到最粗网格求解后，再把修正返回细网格并平滑网格转移误差。实际的 V 循环或 W 循环会更频繁地在粗细网格间往返。

多重网格法在 20 世纪 80 年代后期得到系统发展（Briggs，1987；Wesseling，1992），许多研究者用它求解多流体压力方程。其核心是不同波数误差应在不同网格上消除。可以把 Poisson 方程迭代看成扩散方程沿显式“伪时间”推进到稳态；最大稳定步长正比于 $\Delta x^2$。细网格因此只能取小步，而误差的 Fourier 分量衰减率正比于波数平方：高波数误差很快消失，真正控制收敛的低波数误差却衰减缓慢。多重网格先在细网格迭代到高波数误差衰减，再把近似解或残差转移到粗网格，在那里用较大步长消除低波数误差；继续转到更粗网格，直到能直接求解，再逐级把修正插值回细网格，并消除网格转移带来的误差。实际常采用 V 循环、W 循环等更频繁的网格间转移策略。

% [原PDF p79]
极大密度比会使多重网格更难收敛。另一选择是 GMRES、BiCGSTAB 等非定常 Krylov 方法，常配合多重网格预条件。Krylov 方法只需矩阵--向量乘，以当前残差构造搜索方向；预条件器把原算子替换成更易求解、同时保留主要性质的近似算子。早期方法包括共轭梯度，较新的方法多以 GMRES 为基础。GMRES 每次迭代要利用所有先前迭代结果，存储和工作量随迭代数增加，但已成功用于多相计算（Francois 等，2006）。对于这里系数跨界面突变的 Poisson 方程，目前并不存在适用于所有问题的唯一最优方法。van der Vorst（1992，2003）的 BiCGSTAB 已由 Gerlach 等（2006）等用于多相流；Takahira 等（2004）报告其可处理约 $1:10^4$ 的密度比。

压力求解难度总体随密度跳跃增加。低松弛 SOR 通常仍能收敛，但会很慢；高级方法也可能失败。指示函数中的小误差若造成负密度，往往会立即破坏收敛，必须从物理上保证密度插值有界。某些气泡问题在气液密度比足够大后对其精确值不敏感，可适度降低密度比以加速计算；液滴受空气动力、撞击或破波等问题则必须使用真实密度比。

% [原PDF p80]
\section{速度边界条件}

交错网格使压力和法向速度边界处理简洁，却使切向速度略复杂。更新靠近边界的切向速度时需要域外节点，因此通常设置幽灵点。以图 3.6 的左侧固壁为例，壁面切向速度 $v_{\mathrm{bdry},j+1/2}$ 位于内点 $v_{i,j+1/2}$ 与幽灵点 $v_{i-1,j+1/2}$ 中间。把幽灵点当作常规点，并要求线性插值给出指定壁面切向速度：

\begin{sourceblock}
\captionof{figure}{用计算域外幽灵点施加壁面切向速度。选择幽灵值，使内点与幽灵点的线性插值等于指定壁面速度。}
\end{sourceblock}
\begin{equation}
 v_{\mathrm{bdry},j+1/2}=\frac12(v_{i,j+1/2}+v_{i-1,j+1/2}), \tag{3.53}
\end{equation}
故无滑移条件对应
\begin{equation}
 v_{i-1,j+1/2}=2v_{\mathrm{bdry},j+1/2}-v_{i,j+1/2}. \tag{3.54}
\end{equation}
完全滑移要求切向速度法向导数为零，中心差分下只需镜像同值：
\begin{equation}
 v_{i-1,j+1/2}=v_{i,j+1/2}. \tag{3.55}
\end{equation}

% [原PDF p81]
\section{出口边界条件}

实际常只模拟流场中关注的一段区域。入口速度往往可明确给定；出口则应让流体尽可能自然地离开，且少扰动上游。只有在物理上确知出口速度分布时，才能把它逐点指定；多数情况下这种条件过强。

截断计算域隐含假设：边界下游的流动对上游影响很小，通常意味着出口附近已较平滑、均匀。若矩形域右边界上的流线垂直于边界，则 $v=0$；连续方程又给出 $\partial u/\partial x=0$，离散为
\begin{equation}
 u_{i+1/2,j}=u_{i-1/2,j}. \tag{3.56}
\end{equation}
配合常压力和零切向速度，这是较“柔和”的出口条件，不强迫 $u$ 在截面上均匀，甚至允许局部回流。但若出口速度仍有显著横向变化，边界位置会影响解，必须加长计算域或采用更合适的对流/无反射出口条件。

周期边界是另一常用选择。若 $x$ 方向周期为 $L$，任意变量 $a$ 满足 $a(x,y,z)=a(x\pm mL,y,z)$。若物理域含 $N$ 个单元且左右各一层幽灵点，则左幽灵点取右侧对应内点值，右幽灵点取左侧对应内点值。

% [原PDF p82]
\section{自适应网格加密}

多相流和单相流一样，分辨率需求常随位置相差巨大。按需求分配计算资源可显著节省成本。曲线结构网格能够把网格线推向需要加密的位置，但节点总数固定，可实现的自适应有限，只适合较简单问题。

当前主流是对规则结构网格进行局部加密：二维把一个单元分为四个，三维分为八个；可嵌入一整块细网格，也可逐单元划分。后一种通常用四叉树或八叉树组织网格并在层级间传递信息，每个子单元还能递归细分。这类技术统称自适应网格加密（AMR）。局部加密与多重网格有内在联系：都在不同尺度的网格之间限制和延拓变量，但 AMR 让细网格只存在于确有需要的空间区域。

AMR 已与多种界面算法结合：包括界面追踪，自适应水平集和 VOF，以及通用 AMR 框架；应用例子涉及生物细胞、爆炸气泡、浮升气泡和射流雾化等问题。多相 AMR 除单相流已有的粗细网格守恒与插值问题外，还必须确保界面位置、体积分数、物性跳跃和表面张力在层级之间彼此一致。

% [原PDF p83]
Gerris 求解器的一个射流雾化算例使用 VOF 与八叉树 AMR；在液气密度比 $\rho_l/\rho_g=28$、$We_l=11600$、$Re_l=5800$，射流直径为区域宽度 $1/12$ 时，七级加密把均匀细网格所需的一亿多个单元减少到约三百万个。多相 AMR 的基本难点与单相 AMR 相似，但还要确保界面重构、物性和表面张力在粗细网格间一致传递。

\section{小结}

本章给出了一套求解变密度、变黏度 Navier--Stokes 方程的标准数值方法。把它与后续章节的相界面、密度和黏度推进算法以及表面张力计算结合，即可模拟相当广泛的多流体流动。格式的精度、稳健性和效率仍可通过高阶对流、隐式扩散、自适应网格和高级压力求解器等方式改进。

尽管方程看似复杂，基本实现并不一定很长。许多研究者已经公开了求解器和代码片段，复用经过验证的离散与线性求解模块，可显著缩短开发时间；但必须核对其守恒形式、变量位置、物性插值与界面算法是否彼此一致。

% [原PDF p84]
把本章的变密度、变黏度 Navier--Stokes 求解器同后续章节的界面平流、密度和黏度设置以及表面张力算法结合，就得到能够较准确模拟广泛多流体问题的方法。精度、稳健性和效率仍可通过高阶对流、隐式扩散、局部加密和更强线性求解器继续改善。尽管公式看起来复杂，基本代码未必很长；许多研究者已经实现这些方程，并在互联网上公开完整程序或代码片段，以经验证实现为起点可显著缩短开发周期。

守恒形式和非守恒形式的差异可由一个带尖锐密度界面的一维匀速流说明。忽略压力梯度和黏性，只需推进守恒动量与质量。假设速度和密度同位存储，动量用时间向前、空间中心格式。这个中心格式本身实际上不稳定，但这不影响它用于说明离散一致性的论证。若初速度处处为常数 $U$，动量更新能否保持常速度，完全取决于密度更新是否使用与动量完全相同的数值通量。

\section{补记：守恒形式与非守恒形式}

为理解两种动量离散的差别，考虑一维匀速流动，密度在尖锐界面处跳跃，压力梯度和黏性项均为零。守恒动量和质量方程为
\begin{equation}
 \frac{\partial(\rho u)}{\partial t}+\frac{\partial(\rho u^2)}{\partial x}=0,
 \qquad
 \frac{\partial\rho}{\partial t}+\frac{\partial(\rho u)}{\partial x}=0. \tag{3.57--3.58}
\end{equation}
在同位网格上，对动量用时间向前、空间中心格式：
\begin{equation}
 u_i^{n+1}=\frac1{\rho_i^{n+1}}\left[\rho_i^nu_i^n-
 \frac{\Delta t}{2h}\bigl((\rho u^2)_{i+1}^n-(\rho u^2)_{i-1}^n\bigr)\right]. \tag{3.59}
\end{equation}
若初速度处处为常数 $U$，并且密度也用完全相同的中心通量离散
\begin{equation}
 \rho_i^{n+1}=\rho_i^n-\frac{U\Delta t}{2h}(\rho_{i+1}^n-\rho_{i-1}^n), \tag{3.61}
\end{equation}
则动量更新中的密度因子严格抵消，$u_i^{n+1}=U$。

% [原PDF p85]
若密度改用迎风格式
\begin{equation}
 \rho_i^{n+1}=\rho_i^n-\frac{U\Delta t}{h}(\rho_i^n-\rho_{i-1}^n), \tag{3.62}
\end{equation}
而动量仍用中心通量，抵消不再成立，界面附近会产生小但非零的速度误差；密度比越大，误差越突出。利用质量方程，动量方程可写为非守恒形式
\begin{equation}
 \frac{\partial u}{\partial t}+\frac{\partial(u^2)}{\partial x}=0. \tag{3.63}
\end{equation}
任意一致离散都会保持常速度。

因此，当密度由标记函数构造，而标记平流与动量通量并不一致时，非守恒速度形式往往更稳妥。另一做法是先用与动量相同的通量推进密度以保证一致性，再在时间步末用更精细的界面算法重新构造密度。关键不是简单选择“守恒”或“非守恒”标签，而是使质量、动量和界面输运在离散层面协调一致。
% ===== 译文分片 2（已合并）=====

\chapter{流体界面的平流}
\markboth{第四章\quad 流体界面的平流}{第四章\quad 流体界面的平流}

% [原PDF p86]
在这种单场描述中，流场所有位置都采用同一组控制方程，因此必须另行辨认当地是哪一种流体。标记函数只承担身份识别，不把它解释成新的物理性质；流体运动和界面位移以后，身份场也必须以同样的速度更新。准确更新看似只是求解线性平流方程，实际上却是多相流计算最关键、也最容易破坏精度的步骤之一。本章先说明直接平流间断标记函数为何困难，再比较为此发展出的不同策略。VOF 是其中历史最久而仍广泛使用的一类；水平集、相场和 CIP 同样直接推进网格函数。前沿追踪则推进相连标记点，随后从界面位置重建标记函数。两类方法还可配合不同的表面张力离散，这一问题留到第七章讨论。

在固定网格上以一组方程求解整个流场时，必须用某种量区分不同流体。通常引入在各流体中取不同值的\textbf{标记函数}。对于不可压缩流，密度等物性有时也能充当标记量；本章则把标记函数仅看成流体身份的编码。界面随流体运动以后，标记函数必须同步更新。这个步骤既决定多相流计算能否成功，也远比表面上困难。

VOF 是直接平流标记函数的经典方法，后来又发展出水平集、相场与 CIP 等方法。另一条路线不直接搬运网格标量，而是用标记点显式追踪两种流体的边界，再由界面位置重建标记函数。前者统称\textbf{界面捕捉}，后者称\textbf{前沿追踪}。

% [原PDF p87]
各种近似并非可以任意互换：$C$ 精确保留控制体平均，适于守恒通量；$I$ 有连续过渡，适于给固定网格设置物性和奇异源项；$F$ 的零集和导数适于几何计算，却不天然保存相体积。后续算法的差异很大程度上来自选择哪一种量作为主变量。

流体质点跨越计算单元时仍属于原来的相，故 $H$ 的物质导数严格为零。这里的 $H_i$ 是第 $i$ 相的特征函数；两相时 $H_1+H_2=1$，只存一个函数即可。颜色函数不是点值，而是特征函数在控制体内的体积平均，因此界面单元的分数值直接给出参考相占据的面积或体积比例。指示函数 $I$ 则可理解为 $H$ 与平滑核 $G$ 的卷积；实际程序既可先平滑 $C$，也可从显式界面直接构造 $I$。此外还可使用与阶跃外形无关的光滑函数 $F$，只以 $F=0$ 定位界面。为避免记号混淆，本书在相场节用 $c$，在 CIP 节用 $f$。

\section{记号}

为判断位置 $\bm x$ 是否属于第 $i$ 种流体，定义 Heaviside 特征函数
\begin{equation}
H_i(\bm x)=
\begin{cases}
1,&\bm x\text{ 位于流体 }i\text{ 中},\\
0,&\bm x\text{ 不位于流体 }i\text{ 中}.
\end{cases}
\tag{4.1}
\end{equation}
对于两流体体系，$H_2=1-H_1$，因而只需使用 $H=H_1$。流体质点在运动中保持身份不变，所以
\begin{equation}
\frac{DH}{Dt}=\frac{\partial H}{\partial t}+\bm u\cdot\nabla H=0.
\tag{4.2}
\end{equation}
已知 $H$ 后即可重建密度、黏度等物性，再求解速度场。

实际计算必须使用 $H$ 的数值近似。颜色函数 $C$ 定义为单元内 $H$ 的平均值；二维矩形单元中
\begin{equation}
C_{i,j}=\frac{1}{\Delta x\Delta y}\int_V H(x,y)\,dx\,dy.
\tag{4.3}
\end{equation}
远离界面的单元为满单元 $C=1$ 或空单元 $C=0$；界面单元满足 $0<C<1$。也可使用光滑指示函数
\begin{equation}
I(x,y)=\iint G(x-x',y-y')H(x',y')\,dx'\,dy',
\tag{4.4}
\end{equation}
其中 $G$ 是平滑核。另一种表示完全不要求标记函数形状接近阶跃函数，而使用在两相中异号的光滑函数 $F$，以零等值面 $F=0$ 定位界面；若 $|F|$ 等于到界面的距离，$F$ 就是符号距离函数。

% [原PDF p88]
图 4.1 的一维剖面清楚展示这些表示的差别：$H(x)$ 在真实界面处由 1 突跳为 0；颜色函数 $C_j$ 除含界面的单元 $j$ 外都与 $H$ 相同；平滑的 $I(x)$ 需要经过若干单元才回到常值；符号距离型水平集满足 $|F|=d$，其中 $d$ 为至界面的距离。严格说来，$C_j$ 属于单元而非空间点，不过后文为了叙述方便，仍把它说成由各单元体积平均近似的连续颜色函数。乍看之下，用给定速度推动一段分段常值函数应当很简单；然而它恰是计算科学中的困难问题。一维情形虽有欺骗性的简单外观，却足以暴露数值扩散和振荡的根源，并说明各种界面方法的动机。

\begin{quote}\small 图 4.1：网格上界面的不同表示。实线为阶跃函数 $H$，阴影为颜色函数 $C_j$，虚线为光滑指示函数 $I$，点划线为以零等值线标记界面的水平集函数。\end{quote}
这些表示承担不同任务：$H$ 精确描述尖锐跳跃，$C$ 保存单元体积，$I$ 适合把奇异界面量分布到固定网格，而 $F$ 的导数则便于计算法向和曲率。

\section{颜色函数的平流}

% [原PDF p89]
这里把通量写成 $F=UC$，并假定 CFL 条件 $U\Delta t\le\Delta x$，保证一个时间步内信息至多穿过一个单元。守恒更新虽保证全域总量不变，却不能单凭守恒避免界面涂抹；必须同时恢复单元内满、空区域的相对位置。

设速度 $U>0$，单元 $j$ 的左、右面分别记为 $j-1/2$ 与 $j+1/2$。式（4.6）表明，一个时间步中 $C_j$ 的改变量完全由两面累计通量之差决定，真正的难点是确定面上的 $C_{j\pm1/2}$。若把每个单元内部都取为其平均值，面值就是迎风单元平均，并得到式（4.7）。图 4.2 取 200 个等宽控制体、物质段宽 $20\Delta x$、时间步 $0.5\Delta x/U$；精确解只应平移，而迎风解在移动一至两个自身长度后已经严重摊平。错误不在守恒通量积分，而在单元内常值假设。若单元最初半满，时间步只够界面抵达右面，则它应先变满，右邻单元仍为空；平均值通量却会过早把材料送出。

先考察一维、不可压缩且速度恒为 $U>0$ 的情形。守恒平流方程为
\begin{equation}
\frac{\partial C}{\partial t}+\frac{\partial F}{\partial x}=0,
\qquad F=UC.
\tag{4.5}
\end{equation}
在单元 $j$ 上积分得到
\begin{equation}
C_j^{n+1}=C_j^n-\frac{1}{\Delta x}
\int_{t^n}^{t^{n+1}}\left(F_{j+1/2}-F_{j-1/2}\right)dt.
\tag{4.6}
\end{equation}
也就是说，单元内颜色函数的改变等于流入量减流出量。若把每个单元内的 $C$ 都看成常数，并满足 $U\Delta t\leq\Delta x$，便得到一阶迎风格式
\begin{equation}
C_j^{n+1}=C_j^n-\lambda(C_j^n-C_{j-1}^n),
\qquad \lambda=\frac{U\Delta t}{\Delta x}.
\tag{4.7}
\end{equation}
该格式的通量本身并没有算错，问题在于以单元平均值代表单元内真实分布。假设界面单元半满，而一个时间步只够界面移到右边界，那么该单元应先填满，右侧单元仍为空；常值重构却会立刻把一部分 $C$ 输送出去，导致界面不可逆地扩散。

% [原PDF p90]
消除涂抹的自然想法，是在每个单元内重构能辨认“左满右空”的高阶分布。若用斜率 $s_j$ 的线性函数，时间 $\Delta t$ 内流入量为 $[C_{j-1}^n+(\Delta x-U\Delta t)s_{j-1}/2]U\Delta t$，流出量为 $[C_j^n+(\Delta x-U\Delta t)s_j/2]U\Delta t$，相减便得到式（4.8）。左右差分斜率 $s_j^\pm$ 再按参数 $\kappa$ 加权；不同 $\kappa$ 给出不同二阶格式，$\kappa=1$ 对应 Lax--Wendroff。几何图中，深灰区域是尚未重构时的单元平均分布，浅灰区域表示利用斜率完成推进、但还未重新平均的新解。这样的处理显著减少扩散，却没有自动解决间断附近的非单调性。

\begin{quote}\small 图 4.2：一阶迎风法输运一段物质。黑线精确解保持形状向下游移动，灰线数值解迅速过度扩散；区域长 1、共 200 个等距控制体，团块宽 $20\Delta x$，$\Delta t=0.5\Delta x$，给出 $t=0,140\Delta t,280\Delta t$。\end{quote}
若在单元内用斜率 $s_j$ 作线性重构，几何积分给出
\begin{equation}
C_j^{n+1}=(1-\lambda)C_j^n+\lambda C_{j-1}^n
-\frac{\Delta x}{2}(1-\lambda)\lambda(s_j-s_{j-1}),
\tag{4.8}
\end{equation}
其中
\begin{align}
s_j^+&=\frac{C_{j+1}^n-C_j^n}{\Delta x},&
s_j^-&=\frac{C_j^n-C_{j-1}^n}{\Delta x},\tag{4.9}\\
s_j&=\frac{1-\kappa}{2}s_j^-+\frac{1+\kappa}{2}s_j^+.
\tag{4.10}
\end{align}
$\kappa=1$ 对应 Lax--Wendroff 型选择。高阶重构减弱扩散，却会在跳跃附近产生过冲和欠冲；再提高多项式阶数只会引入更多振荡。这就是间断平流的基本矛盾：一阶格式耗散过强，高阶线性格式则具有色散振荡。

% [原PDF p91]
对 $\kappa=-1,0,1$ 的三种选择，图 4.4 都显示界面两侧出现正负振荡；采用抛物线等更高次函数只会产生更多振荡。因此，间断平流的普遍困难是：一阶方法过度耗散，而线性高阶方法在跳跃附近产生色散波。现代补救方法大体相同：光滑区保持高阶，在间断邻域限制重构或加入人工黏性以消除振荡，代价是局部形式阶数下降。人工黏性最初由 von Neumann 用于气体动力学激波；后来 van Leer 的高阶 Godunov 方法和 Boris--Book 的通量校正输运赋予它现代形式。对被动标记函数，这些成熟激波捕捉技术在短时间内有效，长时间却仍常不理想。

\begin{quote}\small 图 4.3：高阶格式平流颜色函数。深灰色表示把单元值取为平均值的初始解；由邻单元构造斜率后，可按几何关系精确计算通量；浅灰色为重新平均前的新时刻解。\end{quote}
现代激波捕捉格式通常在光滑区使用高阶离散，在间断附近通过限制器或人工黏性抑制振荡。对被动标记函数而言，这仍不完全理想：可压缩激波受特征线汇聚而具有自发变陡机制，颜色函数的接触间断却没有这种恢复能力。一旦界面被数值扩散，流动方程不会自动把它重新压尖，因此长时间平流需要专门的界面算法。

% [原PDF p92]
激波之所以通常比接触间断更容易保持尖锐，是因为特征线汇入激波；轻微的数值平滑会被动力学重新压缩。被动标记没有这种变陡机制，控制方程无法区分真正的零厚度界面和已被摊宽的过渡层。故保持单调的高分辨率激波格式虽能短时输运颜色函数，长时间演化仍需采用专门的几何或重构方法。VOF 的出发点正是重新审视颜色函数的含义：$C$ 不是任意标量，而是阶跃特征函数的单元平均。利用满、空及界面单元的几何约束，可以避免把所有单元平均都当成均匀分布。

\begin{quote}\small 图 4.4：高阶方法输运物质段。黑实线为精确解，三种深浅不同的灰线对应不同 $\kappa$；时刻及其余参数与图 4.2 相同。\end{quote}
\section{体积分数方法（VOF）}

% [原PDF p93]
一维精确几何通量还依赖知道参考相位于界面的哪一侧，这一信息可由相邻满、空单元确定。二维仅有面积分数不足以决定一条线的角度与截距，因此高维方法必须从邻域补充法向信息，这正是重构误差出现的根本原因。

一维 VOF 几乎没有歧义：除一个界面单元外，颜色函数只能为 0 或 1；只要知道满区位于哪一侧，分数 $C_j$ 就立即给出界面距该侧面的距离。例如 $C_{j-1}=1,C_j=1/3,C_{j+1}=0$ 时，界面距左面 $\Delta x/3$。对于 $U\ge0$，左面始终有满相流入；右面通量在界面抵达前为零，抵达后才等于 $U$，从而得到分段式（4.11）。图 4.2 和 4.4 中的精确平移就是用这一公式计算。高维推广则不再唯一。原始 SLIC 采用方向分裂：做 $x$ 扫掠时以竖线划分满空，做 $y$ 扫掠时改用横线，三维还需第三次扫掠。

\begin{quote}\small 图 4.5：VOF 的一维平流。由界面单元 $j$ 的颜色函数值以及满单元位于哪一侧，可以确定界面位置并精确计算通量。\end{quote}
在一维中，VOF 的核心几乎直接来自几何。除界面单元外，$C_j$ 只能是 0 或 1；若 $C_{j-1}=1$、$C_j=1/3$、$C_{j+1}=0$，界面就在单元左边界右侧 $\Delta x/3$ 处。于是可精确判断一个时间步内界面是否到达右侧面，并计算通量：
\begin{equation}
\int_{t}^{t+\Delta t}F_{j+1/2}\,dt=
\begin{cases}
0,&\Delta t\leq (1-C_j)\Delta x/U,\\
(C_j-1)\Delta x+U\Delta t,&\Delta t>(1-C_j)\Delta x/U.
\end{cases}
\tag{4.11}
\end{equation}

% [原PDF p94]
Hirt--Nichols 方法仍只允许界面平行于坐标轴，但用邻近 $C$ 估算法向，并在横、竖两种取向间选择与法向更一致的一种；同一取向用于两个方向的平流。虽然这种构造比每次扫掠都改变界面方向更自然，Rudman 的测试表明其精度并未显著超过原始 SLIC。两种低阶重构不仅扭曲界面，还会产生脱离主体的非物理较大碎片，文献中称为 floatsam 和 jetsam。图 4.6 对同一真实界面依次给出 SLIC、Hirt--Nichols 和 PLIC 的近似，说明关键改进不是单纯换通量公式，而是更准确地恢复混合单元内界面方向。

\begin{quote}\small 图 4.6：二维平流的 VOF 界面重构：（a）原界面；（b）原始 SLIC；（c）Hirt--Nichols 重构；（d）PLIC 重构。改绘自 Rudman（1997）。\end{quote}
高维困难在于：一个分数值不能单独确定界面在单元中的方向。早期 SLIC 方法采用方向分裂；做 $x$ 向平流时用竖直线把界面单元分成满、空两部分，做 $y$ 向平流时再改用水平线。Hirt--Nichols 方法先由邻近 $C$ 估算法向，再从水平或竖直两种取向中择一。两者都会使界面明显畸变，并可能产生脱离主体的非物理碎片。

% [原PDF p95]
PLIC 线段跨单元不连续并不意味着物理界面真的断开；它只是单元局部积分表示。细化网格时，相邻线段的端点错位应随截断误差减小。若法向估计不收敛，即使每个单元面积严格正确，整体轮廓仍可能呈锯齿或产生错误小滴。

PLIC 由 DeBar 和 Youngs 提出，它允许每个界面单元中的直线段相对坐标轴任意取向。先用中心及邻单元的体积分数求法向，再沿法向移动直线，直至切出的面积恰与 $C_j$ 相符；随后用几何交叠计算单元间通量。因而，PLIC 的整体平流误差高度依赖法向估计。二维直线可穿过相邻的两条边或相对的两条边，基本构型有限；三维平面与立方体相交的情况多得多，实现复杂度显著增加。重构线段在单元边界处不要求连续，但对缓慢变化的光滑界面，错位通常很小；曲率变化剧烈时误差较大。下一章将系统讨论现代 VOF，后文也会展示实际多相流计算。

PLIC（分段线性界面计算）把每个混合单元中的界面表示成任意方向的直线段。法向由中心单元与邻单元的 $C$ 共同确定，再沿法向移动直线，使其切出的面积恰好等于 $C$ 所规定的相体积。随后以几何方式计算相邻单元间的输运体积。PLIC 的精度主要取决于法向和截距重构；二维中只有少数相交构型，三维平面切割多面体则复杂得多。各单元线段通常不要求在边界处连续，但对光滑界面，其错位远小于 SLIC。

% [原PDF p96]
相连标记点用于黏性流界面的历史可追溯到 Daly 的 MAC 计算，他用这些点求 Rayleigh--Taylor 不稳定性中的表面张力；Hyman 随后用标记点追踪接触间断。Glimm 的“前沿追踪”原本服务于双曲守恒律，界面由精确或近似 Riemann 解推进，标记点也携带间断两侧状态。本书把这一名称推广到所有以相连标记点显式表示流体界面的算法。二维前沿由线元连接，三维通常用三角网格。由于曲面会拉伸和压缩，必须动态增删标记、维护连接和表面物理量。非结构表面网格已有成熟生成与管理技术，而且前沿只承载较简单的信息，质量要求低于直接求解全套流体方程的体网格。

\begin{quote}\small 图 4.7：前沿追踪用相互连接的标记粒子表示流体界面，粒子以从固定网格插值得到的流体速度推进。\end{quote}
\section{前沿追踪}

另一种做法是用随流体运动的连接标记点直接表示界面。二维界面为一条折线或高阶曲线，三维界面通常采用非结构三角网格。前沿追踪必须回答四个问题：前沿如何存储；拉伸、压缩时怎样增删标记点；界面量如何与求解流动方程的固定网格交换；液滴合并或破碎时如何改变连接关系。

% [原PDF p97]
光滑分布法的代价是把原本零厚度的密度、黏度跳跃扩展到几个网格；核函数和前沿点密度必须匹配，否则会出现网格印记。其优势是固定网格流体求解器保持规则，复杂几何的处理集中在低维前沿和双向插值运算中。

各类前沿追踪的主要区别，在于移动前沿如何同固定流体网格交换信息，以及界面邻近的变量怎样更新。最简单的办法把尖锐界面视为有限宽的光滑过渡层，把控制方程中的奇异界面源项近似为固定网格上的光滑分布。这一思想与 particle-in-cell、vortex-in-cell 和 Peskin 浸入边界方法一脉相承。此时前沿的主要作用是近乎无误差地推进标记函数；流动求解器只需允许物性随空间变化，不必改变界面附近的差分格式。Unverdi--Tryggvason 的有限雷诺数方法还用同一前沿更新标记函数并施加表面张力，第六章将详细说明。

若把尖锐界面在固定网格上平滑成有限厚过渡层，则界面源项可以用光滑分布表示，固定网格离散无需在界面邻域改变。Unverdi--Tryggvason 方法正是用前沿更新标记函数并施加表面张力。若坚持保持尖锐跳跃，可在界面附近外推场量生成幽灵值，或把跳跃条件直接嵌入差分模板；幽灵流体法和浸入界面法属于这一类。还可采用切割单元，让局部控制体边界贴合界面，但需处理过小控制体。

% [原PDF p98]
切割单元法可把跳跃条件直接放在控制体边界上，概念上接近每相使用独立贴体网格，但绝大多数背景网格仍不变。很小的切割控制体会造成稳定性和时间步限制，所以通常需要合并。显式方法对拓扑事件的“非自动性”也正提供了加入亚网格薄膜模型的接口。

若不希望平滑界面，就必须修改前沿附近的离散。Glimm 方法把一侧场量外推到另一侧，形成幽灵点，使常规差分模板仍可使用；Fedkiw 等人的幽灵流体法采用类似思想，但用水平集定位界面。浸入界面法则把跳跃条件直接写进靠近前沿的梯度近似。另一条路线是局部改造固定网格，使网格线与界面重合；这会形成不规则切割控制体，并常需合并过小单元。显式前沿不会自动完成合并和断裂。虽然必须改连接增加了编程量，但也允许对薄膜是否破裂作物理控制；当膜厚仅一个网格时，自动连通未必正确，这一问题将在第六章末再讨论。

前沿追踪不会自动合并邻近界面，这常被视为缺点，却也提供了物理控制：当液膜厚度只有一个网格时，是否应立即破裂并无普适答案。显式前沿可把合并准则、薄膜排液模型或亚网格物理明确写入算法，而不是由数值分辨率暗中决定。

% [原PDF p99]
符号约定必须一致：本书法向指向 $F$ 减小的一侧，所以式（4.13）带负号；若另一本资料把法向定义为 $\nabla F/|\nabla F|$，相应 $V_n$ 与曲率符号也会改变。数值实现比较公式时应先统一这一点。

水平集函数在一相为正、另一相为负，零等值线即相界面。只有法向速度改变零等值线位置，故由 $\bm n=-\nabla F/|\nabla F|$ 和 $V_n=\bm u\cdot\bm n$ 可把平流式改写为式（4.14）。由于 $F$ 光滑，可用 ENO 等高阶双曲格式推进。密度与黏度等物性并不直接由 $F$ 的数值大小给出，而通过式（4.15）的光滑 Heaviside 映射生成指示函数；经验参数 $\alpha\approx3$ 时，过渡区约六个网格。对 $I(F)$ 求梯度又得到界面 delta 分布 $\bm\delta=(dI/dF)\nabla F$，可用来分布表面源项。

\section{水平集方法}

水平集方法用光滑函数 $F(\bm x,t)$ 标记两相，界面为 $F=0$。其平流方程为
\begin{equation}
\frac{\partial F}{\partial t}+\bm u\cdot\nabla F=0.
\tag{4.12}
\end{equation}
若法向指向 $F$ 减小的一侧，则
\begin{equation}
\bm n=-\frac{\nabla F}{|\nabla F|},
\tag{4.13}
\end{equation}
且只需法向速度 $V_n=\bm u\cdot\bm n$ 即可推进零等值面：
\begin{equation}
\frac{\partial F}{\partial t}-V_n|\nabla F|=0.
\tag{4.14}
\end{equation}
由于 $F$ 光滑，可采用 ENO 等高阶双曲格式。为了恢复物性，常把 $F$ 映射为有限厚度的指示函数
\begin{equation}
I(F)=
\begin{cases}
0,&F<-\alpha\Delta x,\\
\dfrac12\left[1+\dfrac{F}{\alpha\Delta x}
+\dfrac1\pi\sin\!\left(\dfrac{\pi F}{\alpha\Delta x}\right)\right],
&|F|\leq\alpha\Delta x,\\
1,&F>\alpha\Delta x,
\end{cases}
\tag{4.15}
\end{equation}
并由
\begin{equation}
\bm\delta=\nabla I=\frac{dI}{dF}\nabla F
\tag{4.16}
\end{equation}
构造光滑界面 delta 分布。$\alpha$ 常取 3，过渡层约为六个单元。

% [原PDF p100]
普通流动平流会改变 $F$ 的斜率：界面受拉伸处梯度变陡，受压缩处梯度变缓。Sussman 等人因而在每个物理时间步后引入伪时间方程（4.17），其稳态满足 $|\nabla F|=1$，把 $F$ 恢复为符号距离。慢界面可间隔若干步重初始化，快速变形时往往每步都做；但这一操作可能使零等值线人工移动并造成质量损失。Peng 等人提出式（4.18）的光滑符号函数，并只在零等值面窄带内推进几个伪时间步。若速度只在界面上已知，例如凝固前沿问题，还须解式（4.19）把每个速度分量沿法线延拓，使其在法向近似常值。

纯平流会使 $|\nabla F|$ 在拉伸区增大、压缩区减小，因此需要重初始化，使 $F$ 恢复为符号距离函数：
\begin{equation}
\frac{\partial F}{\partial\tau}+\operatorname{sgn}(F_0)(|\nabla F|-1)=0.
\tag{4.17}
\end{equation}
其中 $\tau$ 是伪时间，稳态满足 $|\nabla F|=1$。较平滑的符号函数可取
\begin{equation}
S(F)=\frac{F}{\sqrt{F^2+|\nabla F|^2\Delta x^2}}.
\tag{4.18}
\end{equation}
重初始化只需在零等值面附近窄带内推进若干伪时间步，但它可能引起人工界面位移，成为质量误差的重要来源。

有些问题只在界面上给定速度，此时需把速度沿法向延拓到窄带内。对每个速度分量 $q$ 可解
\begin{equation}
\frac{\partial q}{\partial\tau}+S(F)\bm n\cdot\nabla q=0,
\tag{4.19}
\end{equation}
得到沿法线近似不变的延拓场。水平集的原始吸引力是实现简洁，但为了改善守恒性，往往还要加入质量校正、VOF 耦合（CLSVOF）或幽灵流体处理，复杂度也随之增加。

% [原PDF p101]
原始水平集方法吸引人的地方是简单：流体力学计算仅多解一个偏微分方程，不需要 PLIC 的几何切割，也不新增显式前沿对象。早期实现却有明显质量守恒问题；加入重初始化、质量校正和更精确的界面条件后，复杂度逐渐接近其他方法。Sussman 等人提出守恒改进，Sussman--Puckett 构造了 VOF/水平集耦合的 CLSVOF，Fedkiw 等人则发展幽灵流体法。相场方法采用不同哲学：承认界面有有限厚度，并用热力学一致的守恒律控制它。数值过渡层远厚于真实分子界面，相场在凝固模拟中应用广泛，在流体动力学中相对较少。

\section{相场方法}

相场方法承认界面具有有限厚度，并以热力学一致的自由能模型保持过渡层宽度。实际网格上的界面远厚于真实分子尺度，因此其优势主要在于统一描述扩散界面和某些微尺度现象，而不是逐点再现真实厚度。

Jacqmin 的等密度近似模型可写为
\begin{align}
\rho_0\frac{D\bm u}{Dt}&=-\nabla S+\mu\nabla^2\bm u-c\nabla\phi+\rho(c)\bm g,
\tag{4.20}\\
\frac{Dc}{Dt}&=\kappa_d\nabla^2\phi,
\tag{4.21}\\
\phi&=\beta\frac{d\psi(c)}{dc}-\alpha\nabla^2c,
\qquad \psi(c)=\left(c+\frac12\right)^2\left(c-\frac12\right)^2.
\tag{4.22}
\end{align}
这里 $S$ 是保证不可压缩性的压力型变量，$\kappa_d$ 是相场迁移率；$\alpha$ 控制梯度能，$\beta$ 控制体自由能。界面厚度量级为 $\sqrt{\alpha/\beta}$，表面张力与 $\sqrt{\alpha\beta}$ 成正比。对上述势函数，$\sigma=\sqrt{\alpha\beta}/18$，90\% 的相场变化宽度约为 $4.164\sqrt{\alpha/\beta}$。给定物理表面张力和数值界面厚度后，系数便受约束。

扩散项阻止拉伸使界面无限变薄，双阱势产生的反扩散作用又防止界面过宽。相场法与水平集法都在固定网格上使用光滑过渡区，但出发点不同；相场模型原则上可描述接触线等尖锐界面假设不适用的现象。

% [原PDF p102]
这些参数并非彼此独立：指定物理表面张力与希望解析的数值界面宽度后，$\alpha$、$\beta$ 已受两个关系约束，迁移率还必须足以恢复界面而不过度扩散。Jacqmin 的四阶空间离散正是为了用两至三个网格仍解析这一平衡。后续研究把相场用于相变、热毛细和接触线问题，说明其热力学结构在微尺度现象中可能比纯几何平滑更有意义。

相场动量方程中的 $-c\nabla\phi$ 是界面毛细力，$S$ 是用于实施不可压缩约束的压力型变量，$\kappa_d$ 为迁移率。双阱势在 $c=\pm1/2$ 有两个稳定极小值，梯度能惩罚过陡变化，两者平衡形成稳定的有限厚界面。网格细化时应按厚度和表面张力重新缩放参数，不能把数值系数固定不变。

相场变量 $c$ 的平流方程加入非线性扩散，形成带对流的 Cahn--Hilliard 方程。扩散部分阻止拉伸把界面拉得过薄，反扩散部分又限制压缩时界面过厚；Navier--Stokes 方程同时加入由自由能导出的界面力。Jacqmin 在 Boussinesq 近似、两相黏度和质量扩散系数相同的情形求解式（4.20）--（4.22）。$\alpha$ 与 $\beta$ 分别控制梯度能和体自由能，界面厚度为 $O(\sqrt{\alpha/\beta})$，表面张力与 $\sqrt{\alpha\beta}$ 成比例。对给定双阱势，$\sigma=\sqrt{\alpha\beta}/18$，90\% 变化宽度约 $4.164\sqrt{\alpha/\beta}$；给定 $\sigma$ 和取为两三个网格的厚度后，系数基本确定。

\section{CIP 方法}

CIP 方法同时平流函数值及其导数，以三次多项式降低色散误差。对常速度一维问题
\begin{align}
\frac{\partial f}{\partial t}+u\frac{\partial f}{\partial x}&=0,
\tag{4.23}\\
\frac{\partial g}{\partial t}+u\frac{\partial g}{\partial x}&=0,
\qquad g=\frac{\partial f}{\partial x}.
\tag{4.24}
\end{align}
% [原PDF p103]
CIP 系数由四个端点条件 $P(0),P(\Delta x),P'(0),P'(\Delta x)$ 唯一确定。回溯坐标的具体写法取决于局部原点；原书令多项式区间起点为上游节点，因此式（4.26）与“$\xi=u\Delta t$”的表述要结合坐标定义理解。函数值与导数必须采用同一多项式推进。

Jacqmin 还要求迁移率按界面厚度的适当幂次缩放，并采用四阶方法使两三个网格宽的界面保持稳定。相场已用于 Boussinesq Rayleigh--Taylor、不等温相变、Hele--Shaw 热毛细不稳定性及接触线动力学。其数值表现与其他固定网格单场方法相近；虽然平滑区外观类似水平集，两者的模型假设不同，相场原则上可研究锐界面假设不适用的微尺度接触线现象。CIP 的目标则是通过同时推进函数及其导数降低色散。名称从 cubic interpolated pseudo-particle 演变为 cubic interpolated propagation，后来称 constrained interpolation profile，但缩写始终不变。

当 $u>0$ 时，在节点 $j-1$ 与 $j$ 之间构造 $P(x)=a_jx^3+b_jx^2+c_jx+d_j$，取 $x_{j-1}=0$、$x_j=\Delta x$，由两端的 $f$ 与 $g$ 得
\begin{align}
a_j&=\frac{2(f_{j-1}-f_j)}{\Delta x^3}+\frac{g_{j-1}+g_j}{\Delta x^2},\\
b_j&=\frac{3(f_j-f_{j-1})}{\Delta x^2}-\frac{g_j+2g_{j-1}}{\Delta x},\\
c_j&=g_{j-1},\qquad d_j=f_{j-1}.
\tag{4.25}
\end{align}
令 $\xi=\Delta x-u\Delta t$ 表示回溯点相对左端的坐标，则
\begin{align}
f_j^{n+1}&=a_j\xi^3+b_j\xi^2+g_{j-1}^n\xi+f_{j-1}^n,\\
g_j^{n+1}&=3a_j\xi^2+2b_j\xi+g_{j-1}^n.
\tag{4.26}
\end{align}
原书以方形脉冲平流说明：CIP 能很好保持间断，但仍有小振荡。多维可采用方向分裂，也可构造多维多项式并按局部速度整体平移。

% [原PDF p104]
常速度下，$f$ 与 $g=\partial f/\partial x$ 都保持形状平移；若速度非均匀，$g$ 方程可拆成纯平流和速度梯度源项两步处理。对 $u>0$，用上游节点 $j-1$ 和节点 $j$ 的 $f,g$ 唯一确定三次多项式；对 $u<0$ 则改用 $j,j+1$。将该多项式回溯 $u\Delta t$ 即得新时刻的函数和导数。方形脉冲试验中，初始 $g$ 由中心差分生成；CIP 几乎保持间断，但仍出现小振荡，这一表现源于其极低色散，而非专为间断设计。多维既可方向分裂，也可建立多维多项式并按局部速度整体平移。本章由此说明，准确、尖锐且守恒地推进流体身份场需要专门算法。

\begin{quote}\small 图 4.8：CIP 方法的平流结果；全部参数与图 4.2 相同。\end{quote}
\section{小结}

固定网格上的统一方程需要一个准确的界面表示。直接平流离散标记函数时，既要保持界面尖锐，又要避免非物理振荡并维持守恒。VOF 通过几何重构和相体积通量满足这一要求；水平集以光滑隐式曲面换取几何便利；相场以自由能控制扩散界面；CIP 同时推进函数及其导数；前沿追踪则把界面本身作为计算对象。后续两章分别详细讨论现代 VOF 与 Unverdi--Tryggvason 型前沿追踪。

% [原PDF p105]
在后续章节中，本书集中讨论两条路线：现代 VOF，以及 Unverdi--Tryggvason 前沿追踪。VOF 是直接推进标记函数中应用最广的方法；早期版本精度有限，但经过 PLIC、几何通量和守恒改进以后，解质量已可与较新的替代方案相比。水平集等方法最初因形式简洁而流行，不过为提高守恒和界面条件精度所加入的步骤也会增加复杂度。总体而言，固定网格上最先进的直接平流方法已接近早期前沿追踪；若以相同计算投入比较，显式界面标记往往仍更准确，但实现、数据结构和拓扑维护更加复杂。

从比较的角度看，直接在固定网格上平流标记函数的先进方法已经能够达到早期前沿追踪实现的解质量；不过，在相当的计算投入下，显式界面标记通常仍能给出更准确的几何位置。它的代价是程序更复杂，因此在可预见的时期内，固定网格界面捕捉方法仍会更普及。现代 VOF 与 Unverdi--Tryggvason 型前沿追踪。

\chapter{体积分数方法}
\markboth{第五章\quad 体积分数方法}{第五章\quad 体积分数方法}

% [原PDF p106]
第四章列举了多种界面定位和推进方法，本章专门讨论 VOF。它早期主要用于只有一个显式流体的自由表面问题，如今已常规用于真正的两流体流动。标记量是参考相占据一个计算单元的面积或体积比例。已提出的 VOF 变体有数十种，选择时必须在几何精度、实现简易性和体积守恒之间权衡。本章限定于两相或自由表面，并先研究边长 $h=\Delta x=\Delta y$ 的二维 Cartesian 方格。满单元、空单元和混合单元分别满足 $C=1,0$ 和 $0<C<1$，混合单元带定位了界面。

体积分数（VOF）方法最初主要用于自由表面流，如今也广泛用于两流体计算。它把参考相在计算单元中所占的比例作为标记函数。文献中存在许多 VOF 变体，选择算法时通常要在精度、实现复杂度与体积守恒之间取舍。

\section{基本性质}

颜色函数 $C$ 是特征函数 $H$ 的离散单元平均。以下主要讨论边长 $h=\Delta x=\Delta y$ 的二维 Cartesian 方格：满单元 $C=1$，空单元 $C=0$，混合单元 $0<C<1$ 确定界面所在的过渡带。低阶 VOF 不显式确定界面位置；从几何上看，相当于在每个混合单元内采用平行于坐标轴的线段，单元边界上的错位通常为 $O(h)$。PLIC 则用垂直于
\[
\bm m=-\nabla C
\]
的任意方向线段重构界面，其中法向约定为从参考相指向外部。各单元线段仍可不连续，但错位随网格间距和局部曲率而减小。

% [原PDF p107]
低阶 VOF 并不显式给出过渡带中的界面位置；其几何解释是每个混合单元用平行于某一坐标轴的线段，跨单元跳跃通常为 $O(h)$。标准高阶 PLIC 则用垂直于局部梯度 $\bm m=-\nabla C$ 的线段，本文约定法向从参考相指向外部；读者应注意部分文献采用相反符号。线段仍不必跨单元连接，但错位会随 $h$ 和局部曲率变化而显著缩小。一个 VOF 步骤先由邻域体积分数重构界面形状，关键是求法向；再在已知速度场中推进该重构，本质是交换相邻单元之间的参考相体积。

\begin{quote}\small 图 5.1：VOF 的基本原理：（a）一段界面及各单元颜色函数；（b）SLIC 重构，线段平行于坐标轴；（c）PLIC 重构，相邻单元线段不要求连接，法向 $\bm m=-\nabla C$ 指向参考相外部。\end{quote}
一个典型 VOF 时间步包含两步：
\begin{enumerate}[label=(\roman*)]
\item 根据本单元及其邻域的体积分数重构界面；PLIC 的关键是求局部法向；
\item 在给定速度场中平流重构界面，本质是计算穿过相邻单元公共面的参考相体积。
\end{enumerate}
% [原PDF p108]
VOF 自然适合拓扑改变，界面重连或破碎无需显式调整连接关系；二维到三维 Cartesian 网格的推广也相对直接。$C$ 是固定网格上的静态矩阵，算法只依赖邻单元，因而易于域分解并行。最具辨识度的优点是质量守恒。把不可压缩标记方程积分到单元边界，再对时间积分，得到式（5.3）的四面面积通量。全域求和时每个内部通量在相邻单元中一正一负而抵消，故总面积守恒。但该证明隐含两个条件：更新不得产生 $C>1$ 或 $C<0$，也不能把同一块材料穿过单元面计算两次；后文将多次回到这两个一致性问题。

\begin{quote}\small 图 5.2：用面积与体积分数成比例的阴影方块理解重构问题：（a）对应光滑圆弧的良好颜色函数分布；（b）较混乱的分布。\end{quote}
VOF 的数据结构固定、局部且易于域分解并行；界面合并与破碎也隐含在体积分数场中。最突出的性质是守恒。不可压缩流中
\begin{equation}
\frac{\partial H}{\partial t}+\nabla\cdot(\bm uH)=0.
\tag{5.1}
\end{equation}
在单元 $(i,j)$ 上积分得
\begin{equation}
h^2\frac{\partial C_{i,j}}{\partial t}
+\int_{\Gamma}\bm u\cdot\bm n\,H\,dl=0,
\tag{5.2}
\end{equation}
再对一个时间步积分：
\begin{align}
h^2(C_{i,j}^{n+1}-C_{i,j}^{n})={}&
-(\Phi^n_{x:i+1/2,j}-\Phi^n_{x:i-1/2,j})\notag\\
&-(\Phi^n_{y:i,j+1/2}-\Phi^n_{y:i,j-1/2}).
\tag{5.3}
\end{align}
对所有单元求和，内部面通量成对抵消，因此适当边界条件下
\begin{equation}
\sum_{i,j}C_{i,j}^{n+1}=\sum_{i,j}C_{i,j}^{n}.
\tag{5.4}
\end{equation}
这一结论还要求算法不重复输运同一块面积，也不产生 $C>1$ 或 $C<0$。

% [原PDF p109]
PLIC 重构分两步：先从中心单元及邻域的 $C$ 求非单位法向 $\bm m$，再以 $m_xx+m_yy=\alpha$ 表示直线，并调节 $\alpha$ 使切割面积等于 $h^2C_{i,j}$。这是由面积反求曲线的逆几何问题。为判断收敛阶，可类比梯形积分：对二阶导数连续的函数，分段线性积分的误差为 $O(f''(\xi)h^2)$，任意直线因二阶导数为零而应被精确再现。因此，不能精确重构任意直线的算法不可能是完整的二阶方法。圆弧重构理论上可达三阶，但三维代价高；局部二次信息更适合用在表面张力等环节。

\section{界面重构}

PLIC 重构首先由邻域 $C$ 求法向 $\bm m$，再写界面直线
\begin{equation}
\bm m\cdot\bm x=m_xx+m_yy=\alpha.
\tag{5.5}
\end{equation}
沿法向调节 $\alpha$，直到直线切出的参考相面积等于 $h^2C_{i,j}$。

\subsection{重构方法的收敛阶}

重构是几何积分的逆问题：已知曲线下的面积，反求曲线。梯形积分用分段线性函数逼近具有连续二阶导数的 $f(x)$，面积截断误差为 $O(f''(\xi)h^2)$，因而任意直线都应被精确再现。由此可把“能否精确重构任意直线”作为二阶 PLIC 的必要基准；不能做到这一点的方法阶数低于二阶。用圆弧重构理论上可达到三阶，但三维实现代价很大，通常不如在表面张力计算中局部使用二次信息更实用。

% [原PDF p110]
法向估计可以是 $\nabla C$ 的有限差分，也可由某个误差泛函最小化；本节方法都用中心单元周围 $3\times3$ 模板。Youngs 法先在四个角点以相邻四个体积分数求梯度，再平均到单元中心。右上角的两个分量由式（5.7）、（5.8）给出，四角平均即式（5.9）。该方法由 Youngs 和 Li 独立发展，计算快速，但一般不能精确再现任意直线。例如在边长 $h$ 的 $3\times3$ 方块中考察 $y=2x/3+h$，直接代入即可看到误差；数值直线测试也会揭示这一弱点。

\begin{quote}\small 图 5.3：一条直线穿过两个网格单元。梯形规则的分段线性近似跨单元连续；VOF--PLIC 线段跨单元不连续但满足面积守恒。点状区域分别表示 PLIC 与梯形规则的面积误差。\end{quote}
\subsection{界面单位法向的求取}

法向估计可直接近似 $\nabla C$，也可通过误差最小化得到。以下算法都以中心混合单元周围的 $3\times3$ 模板为基础。

\subsubsection{Youngs 有限差分法}

Youngs 方法取
\begin{equation}
\bm m=-\nabla_h C.
\tag{5.6}
\end{equation}
先在中心单元四个角点计算梯度。例如右上角
\begin{align}
m_{x:i+1/2,j+1/2}&=-\frac{C_{i+1,j+1}+C_{i+1,j}-C_{i,j+1}-C_{i,j}}{2h},
\tag{5.7}\\
m_{y:i+1/2,j+1/2}&=-\frac{C_{i+1,j+1}-C_{i+1,j}+C_{i,j+1}-C_{i,j}}{2h}.
\tag{5.8}
\end{align}
再平均四角的向量：
\begin{equation}
\bm m_{i,j}=\frac14(\bm m_{i+1/2,j+1/2}+\bm m_{i+1/2,j-1/2}
+\bm m_{i-1/2,j+1/2}+\bm m_{i-1/2,j-1/2}).
\tag{5.9}
\end{equation}
该格式速度快、鲁棒，但不能精确再现所有斜率的直线，渐近上通常只有一阶。

% [原PDF p111]
中心列高度 $hy_{i-1}=h^2\sum_{k=-1}^{1}C_{i-1,j+k}$ 的量纲处理表明，实际高度是三单元面积和除以列宽。由高度求斜率时，$m_y$ 的绝对大小通常被归一化，只需另存其符号；宽度表示则相反。这些细节决定法向方向是否从参考相正确指向外部。

中心列法把三行体积分数沿竖直方向求和形成高度函数 $y=f(x)$，或沿水平方向求和形成宽度函数 $x=g(y)$。只有当直线穿过一列的两条竖边时，该列高度才恰落在真实直线上；宽度表示同理。列求和会丢失哪一种流体在上方的信息，所以必须另外由 $C$ 的中心差分确定 $m_y$ 的符号。高度的中心斜率为式（5.11），宽度的中心斜率为式（5.13）。当界面穿过模板的相对两边时，其中一个表示给出正确斜率；若穿过相邻边，单纯中心差分一般不正确，需要更宽模板或前、后向差分。

\subsubsection{中心列差分法}

把同一列三个单元的体积分数相加，可构造高度函数 $y=f(x)$；按行相加则构造宽度函数 $x=g(y)$。若界面穿过列的两个竖直侧面，列高落在真实直线上。用中心差分表示
\begin{equation}
\operatorname{sgn}(m_y)y=-m_xx+\alpha',
\tag{5.10}
\end{equation}
并取
\begin{equation}
m_{xc}=-\frac{y_{i+1}-y_{i-1}}{2h}
=-\frac12\sum_{k=-1}^{1}(C_{i+1,j+k}-C_{i-1,j+k}).
\tag{5.11}
\end{equation}
类似地，宽度表示为
\begin{equation}
\operatorname{sgn}(m_x)x=-m_yy+\alpha'',
\tag{5.12}
\end{equation}
其中心差分为
\begin{equation}
m_{yc}=-\frac{x_{j+1}-x_{j-1}}{2h}
=-\frac12\sum_{k=-1}^{1}(C_{i+k,j+1}-C_{i+k,j-1}).
\tag{5.13}
\end{equation}
% [原PDF p112]
式（5.14）的“最小”用于在高度图和宽度图两个不同表示之间选择；式（5.15）的“最大”用于同一表示下的多个偏心差分候选。二者不能混用。后者修正中心列在界面穿相邻边时对高度的系统高估，从而恢复任意直线的可能性。

由于界面方向事先未知，必须在高度与宽度候选中选择。对 $45^\circ$ 直线，两种斜率绝对值都为 1；较平坦的高度图满足 $|m_x|=1/|m_y|<1$，较陡时则宽度图斜率更小。因此，对两个方向候选先取绝对值较小者。若同一方向还有中心、前向或后向等多个估计，穿相邻边时中心高度可能高估而使斜率绝对值偏小，此时应在同方向候选中取绝对值较大者。两个简单规则足以精确恢复直线，但对曲线未必最高效。ELVIRA 因而不只按斜率大小判断，而通过整个邻域的预测面积误差选择候选。

对穿过相对边的直线，选绝对值较小的斜率
\begin{equation}
|m^*|=\min(|m_{xc}|,|m_{yc}|)
\tag{5.14}
\end{equation}
可得到正确表示；若同一方向已有多个候选斜率，在相邻边相交构型中应取绝对值较大者
\begin{equation}
|m^*|=\max(|m_{x1}|,|m_{x2}|).
\tag{5.15}
\end{equation}
仅用中心差分仍不能精确重构全部直线，但比 Youngs 方法误差小。

% [原PDF p113]
ELVIRA 对高度计算后向、中心、前向三种斜率，对宽度再计算三种，共六个法向候选。对每个候选，先在中心单元按面积守恒求 $\alpha_n$，再把同一直线延伸到整个 $3\times3$ 模板，预测九个单元的 $\widetilde C_{k,l}$。式（5.17）和（5.18）分别用 $L_2$ 与 $L_\infty$ 型面积误差衡量预测与真实体积分数的差别，最后取误差最小者。任意直线至少贯穿模板中两行或两列，因此六个候选中至少一个准确。矩形大小不等时需更宽模板；三维预测体积的次数很多，速度成为 ELVIRA 的主要缺点。

\begin{quote}\small 图 5.4：中心列法。（a）按列累加得到高度 $y_{i-1},y_{i+1}$；（b）按行累加得到宽度 $x_{j-1},x_{j+1}$；（c）应采用偏心差分高度 $y_i,y_{i+1}$。\end{quote}
\subsubsection{ELVIRA 方法}

ELVIRA 同时计算高度函数的后向、中心和前向斜率，并对宽度函数重复一次，共有六个候选。例如
\begin{align}
m_{xb}&=-\frac{y_i-y_{i-1}}{h}
=-\sum_{k=-1}^{1}(C_{i,j+k}-C_{i-1,j+k}),\\
m_{xf}&=-\frac{y_{i+1}-y_i}{h}
=-\sum_{k=-1}^{1}(C_{i+1,j+k}-C_{i,j+k}).
\tag{5.16}
\end{align}
对每个候选法向 $\bm m_n$，先在中心单元用面积守恒确定 $\alpha_n$，再把该直线延伸到整个 $3\times3$ 模板，得到预测体积分数 $\widetilde C_{k,l}$。候选误差可用
\begin{align}
E_2(n)&=h^2\sum_{k=i-1}^{i+1}\sum_{l=j-1}^{j+1}
[\widetilde C_{k,l}(\bm m_n)-C_{k,l}]^2,
\tag{5.17}\\
E_\infty(n)&=h^2\max_{k,l}|\widetilde C_{k,l}(\bm m_n)-C_{k,l}|.
\tag{5.18}
\end{align}
选择误差最小者。由于任意直线至少贯穿两个连续行或列，总有一个候选能精确重构直线；代价是反复计算预测面积，三维尤其昂贵。

% [原PDF p114]
样点选在线段两端和中心，而不是只取中心，是为了让拟合点尽可能落在真实曲线两侧，避免直线整体偏向一边。拟合只更新斜率、再用面积守恒重算截距，保证几何误差下降的同时不破坏本单元给定体积分数。

最小二乘方法先用 Youngs 或中心列法快速完成全域预重构，同时估计每个单元的 $\alpha$。然后在每条预重构线段的两端和中心取点；好的直线近似通常在一个单元内与真实曲线相交两次，因此样点会分布在界面两侧。围绕中心单元汇总 $n_p$ 个点后，根据预估斜率选择把 $y$ 对 $x$ 或 $x$ 对 $y$ 拟合，最小化式（5.19）之一。权重可随距中心单元的距离改变，本书测试取 1。只解关于斜率的 $2\times2$ 线性系统，截距仍由中心单元面积守恒决定；反复迭代可精确恢复直线，实际通常一次拟合已足够。

\subsubsection{最小二乘拟合法}

最小二乘方法先用 Youngs 或中心列差分作快速预重构，并在每条初始线段的两端和中点取样。围绕中心单元汇集 $n_p$ 个界面点 $(x_i,y_i)$ 后，根据预估斜率选择最合适的图形表示，最小化
\begin{equation}
I_1=\sum_{i=1}^{n_p}\omega_i(y_i+mx_i+\alpha)^2,
\quad\text{或}\quad
I_2=\sum_{i=1}^{n_p}\omega_i(x_i+my_i+\alpha)^2.
\tag{5.19}
\end{equation}
这里只对斜率求解小型线性系统，中心单元的 $\alpha$ 仍由面积守恒确定。迭代使用拟合结果可在少数迭代内精确恢复直线；多数实际计算做一次拟合已经足够。

% [原PDF p115]
LVIRA 不把法向限制在六个候选，而在法向空间搜索面积误差的局部或全局最小。López 等人先用 Youngs 重构并按弧长排列线段中点，再用三次样条插值 $x(s),y(s)$，通过微分得到更光滑的法向。Scardovelli--Zaleski 还在高曲率区每单元使用两条连续线段。Youngs 和中心列法可直接推广到三维；三维 ELVIRA 需要 72 至 144 个候选，并通常采用 $5^3$ 模板，因为 $3^3$ 单元不足以用差分恢复任意平面。最小二乘法也已推广到三维。求得法向后，下一步就是确定面积守恒截距 $\alpha$。

\begin{quote}\small 图 5.5：最小二乘重构示意。（a）在预重构线段上选择端点和中点；（b）以这些点拟合虚线，只示出 $3\times3$ 模板的一部分。\end{quote}
\subsubsection{其他方法及三维扩展}

LVIRA 不把候选限制为六个，而在法向空间搜索面积误差的局部或全局极小值。另有方法先用 Youngs 重构，按弧长排列各线段中点，再用三次样条求连续法向；高曲率区也可在一个单元中用两条连续线段。三维 Youngs 和中心列法可直接推广；三维 ELVIRA 候选法向可达 72--144 个，并通常需要 $5\times5\times5$ 模板，因为 $3\times3\times3$ 不足以对任意平面计算所有法向分量。最小二乘拟合同样可以推广到三维三角面片。

\subsection{\texorpdfstring{截距 $\alpha$ 的确定}{截距 alpha 的确定}}

% [原PDF p116]
对一般 $\Delta x\times\Delta y$ 矩形，可猜测 $\alpha$、构造切割多边形 $ABFGD$、用顶点鞋带公式求 $A(\alpha)$，再以求根法使 $A=\Delta x\Delta yC$。Cartesian 对称性允许无迭代反演：先关于坐标轴镜像，使 $m_x,m_y>0$，切割面积等于大三角形 $AEH$ 减去越过右边和上边的两个相似三角形，得到式（5.22）。界面穿相邻边时它是 $\alpha$ 的二次函数，穿相对边时是线性函数；该函数连续且严格单调，所以逆解唯一而稳定。三维长方体的直接切割和逆关系在附录 C 讨论。

\begin{quote}\small 图 5.6：“切割面积”为矩形单元 $ABCD$ 内、直线 $EH$ 下方的灰色区域。\end{quote}
给定 $\bm m$ 后，要解
\begin{equation}
g(\alpha)=A(\alpha)-\Delta x\Delta y C=0,
\tag{5.20}
\end{equation}
其中 $A(\alpha)$ 是矩形被直线切出的参考相面积。可按逆时针排列多边形顶点并使用鞋带公式
\begin{equation}
A=\frac12\sum_{k=1}^{n}(x_ky_{k+1}-x_{k+1}y_k),
\qquad (x_{n+1},y_{n+1})=(x_1,y_1).
\tag{5.21}
\end{equation}
再用求根算法调节 $\alpha$。Cartesian 单元的对称性还允许直接反演。经镜像使 $m_x,m_y>0$ 后，
\begin{equation}
A(\alpha)=\frac{\alpha^2-F_2(\alpha-m_x\Delta x)-F_2(\alpha-m_y\Delta y)}{2m_xm_y},
\tag{5.22}
\end{equation}
其中 $F_2(z)=z^2$（$z>0$），否则为零。相邻边相交时它对 $\alpha$ 为二次函数，相对边相交时为线性函数，而且连续、严格单调，因而可稳定求逆。

% [原PDF p117]
误差 $E$ 不是单纯比较线段端点，而是积分两特征函数之差，因此能正确计量线段交叉、方向不同以及局部面积补偿。随机平移和转动后再平均，可避免某条界面恰好沿网格线时产生零误差而掩盖算法的一般表现。

重构测试固定界面不做平流，从而排除时间误差。设精确特征函数 $H$，网格尺度 $h$ 上的重构为 $\widetilde H_h$；式（5.23）的 $L_1$ 误差就是二者不同区域的面积。归一化 $B=1$ 时为绝对面积误差，取精确界面长度 $L$ 时可看成平均偏移宽度，取流体总面积时成为无量纲相对误差。类似地可定义 $L_2$ 和 $L_\infty$。两级网格误差按式（5.24）给出观测阶；它与 $B$ 无关，却会受几何形状、位置、方向和网格尺度影响。真正的渐近阶只在 $h\to0$ 时出现，数值上只能由逐级细化外推。

\section{重构方法测试}

对静止界面，设精确特征函数为 $H$，网格尺度 $h$ 上重构得到 $\widetilde H_h$。$L^1$ 几何误差可定义为
\begin{equation}
E=\frac1B\iint|H-\widetilde H_h|\,dx\,dy,
\tag{5.23}
\end{equation}
$B=1$ 时是所有混合单元中真实与重构界面之间的面积和；$B$ 取界面长度时，$E$ 可理解为平均法向位移。由两级网格估计收敛阶
\begin{equation}
\mathcal O=\frac{\ln(E_h/E_{h/2})}{\ln 2}.
\tag{5.24}
\end{equation}

% [原PDF p118]
第一个试验随机生成 1000 条直线：在单位方形中心附近随机取线上一点，并随机选择界面与水平轴的夹角。误差通过精确求真实和重构界面的交点并解析积分得到，多次随机平均避免偶然与网格对齐导致的异常小误差。表 5.1 显示 Youngs 误差随加密近似一阶；中心列法同样约一阶但常数更小；以中心列为初值的一次线性拟合也约一阶下降，不过绝对误差已显著更低。ELVIRA 对任意直线在舍入精度内完全重构。若继续迭代最小二乘拟合，每次误差下降数个数量级，第二次通常已低于机器精度。

随机平移和转动的直线测试表明：ELVIRA 能在舍入误差内精确恢复任意直线；Youngs 与中心列方法大体为一阶，后者误差更小；以中心列为初值的一次最小二乘拟合误差已很低，第二次迭代可降到机器精度。% [原PDF p119]
第二个试验采用 $x^2/a^2+y^2/b^2=1$ 的椭圆，其中 $a^2=0.12,b^2=0.02$，最大与最小曲率半径之比约 15；随机平移中心并转动长轴，约需 150 个样本才能稳定到三位有效数字。表 5.2 表明 Youngs 在极粗网格上颇有竞争力，但渐近回到一阶；中心列法在实用分辨率内介于一、二阶；ELVIRA 与一次线性拟合清楚地趋向二阶。随后进入平流问题：PLIC 先形成不连续分段直线，再用几何端点推进或面通量推进。把一维格式解释成平面对自身的映射，可证明更新后始终有 $0\le C\le1$。

对高偏心率椭圆测试，Youngs 在极粗网格上表现尚可，但渐近阶仍接近一阶；中心列法在实用分辨率内介于一、二阶；ELVIRA 与一次最小二乘拟合清楚地呈二阶收敛。

\section{界面平流}

PLIC 先用不连续的分段直线重构界面，再把它推进到下一时刻。几何方法直接推进线段端点；通量方法则计算穿过单元面的参考相面积。将一维推进写成平面到自身的映射，能清楚说明算法如何同时保证 $0\leq C\leq1$。

% [原PDF p120]
无量纲化后，局部面速度 $u_{i\pm1/2,j}$ 已含 $\Delta t/h$，所以映射式中不再显式出现时间步和网格间距。恢复有量纲变量时必须重新乘回这些尺度。映射保持直线性的前提是单元内速度采用线性插值。

方向分裂二维平流依次做 $x$、$y$ 扫掠；为降低方向偏置，奇数步可先 $x$ 后 $y$，偶数步反序。出单元显式映射从质点方程 $dx/dt=u(x)$ 出发，用一阶显式式（5.25）。以网格宽和时间步无量纲化后，速度就是 CFL 数，单元成为单位方格，局部原点位于左下角。面心速度在单元内按式（5.27）线性插值，把旧坐标映为式（5.28）。系数 $b=1+u_{i+1/2,j}-u_{i-1/2,j}$ 表示横向伸缩，常数项表示平移；这一定义记作 $T_x^E$。

\subsection{几何一维线性映射}

二维方向分裂按顺序完成 $x$、$y$ 两次一维平流；为降低不对称性，奇数步可先 $x$ 后 $y$，偶数步反序。三维需要三个扫掠，并循环改变方向次序。

\subsubsection{出单元显式映射}

质点满足 $dx/dt=u(x)$。一阶显式推进为
\begin{equation}
x(t^{n+1})=x(t^n)+u[x(t^n)]\Delta t.
\tag{5.25}
\end{equation}
以 $h$ 和 $\Delta t$ 无量纲化后，速度分量就是局部 CFL 数。单元内对面心速度作线性插值
\begin{equation}
x(t^{n+1})=x(t^n)+u[x(t^n)].
\tag{5.26}
\end{equation}
\begin{equation}
u(x)=u_{i-1/2,j}(1-x)+u_{i+1/2,j}x,
\tag{5.27}
\end{equation}
得到出单元映射
\begin{equation}
T_x^E:\quad
\begin{cases}
x'=bx+u_{i-1/2,j},\\
y'=y,
\end{cases}
\qquad b=1+u_{i+1/2,j}-u_{i-1/2,j}.
\tag{5.28}
\end{equation}
% [原PDF p121]
像矩形铺砌性质来自相邻单元共享同一面速度：一个单元的右边界和右邻单元的左边界被映到完全相同的位置。若两侧使用不一致速度或独立限制，矩形之间便会重叠或留缝，有界性证明也随之失效。

线性 $T_x^E$ 把每个方格 $\Sigma$ 变为矩形 $\Gamma_x$，直线界面仍是直线。变形后把参考相重新裁切并投影到原网格，中心单元体积分数由左、自身、右三个像矩形的交叠面积 $D+E+F$ 得到。三维中这些贡献是平面切割直六面体形成的多面体，体积按附录 C 计算。所有像矩形恰好铺满平面，不重叠也不留空，因此既不会漏输运，也不会重复输运，更新天然满足有界性。一个二维显式步骤需要在 $x$ 扫掠后更新 $C^*$ 并重新重构，再做 $y$ 扫掠；三维则要三次重构和三次一维推进。

单位方格被映为平移并拉伸或压缩的矩形
\begin{equation}
\Gamma_x=T_x^E(\Sigma).
\tag{5.29}
\end{equation}
线性映射保持直线为直线，故可把变形后的参考相面积重新投影到原网格。中心单元的新体积分数来自自身及左右相邻单元的三个面积贡献
\begin{equation}
C_{i,j}^{n+1}=D+E+F.
\tag{5.30}
\end{equation}
% [原PDF p122]
出单元映射的几何图说明，像矩形既可被压缩，也可膨胀，并可相对原方格向左或向右移。中心单元接收的 $D,E,F$ 都是非负真实面积，它们合起来不超过单位面积，这是 $0\le C\le1$ 的直接几何理由。二维完整流程为：在 $t^n$ 重构界面；以 $T_x^E$ 移动每个切割单元的两个界面端点；按附录 C 求面积并形成 $C^*$；从 $C^*$ 再重构；最后以 $T_y^E$ 得到 $C^{n+1}$。这种两次显式映射虽有界，却有一个主要缺点：对二维无散流也不能精确守恒总面积，必须同另一种一维映射组合。

\begin{quote}\small 图 5.7：（a）水平出单元映射 $T_x^E$ 把方格 $\Sigma$ 变成矩形 $\Gamma_x$；（b）三个连续单元映成三个矩形，并向原网格投影，给出中心单元的 $D,E,F$。\end{quote}

\begin{quote}\small 图 5.8：$T_x^E$ 的所有像矩形铺砌平面；它们可相对原方格伸缩并左右平移。\end{quote}
所有像矩形完整铺砌平面，既不重叠也不留缝，所以 $0\leq C\leq1$。二维显式分裂步骤是：重构；以 $T_x^E$ 推进并得到
\begin{align}
C^*_{i,j}&=D_x+E_x+F_x,\tag{5.31}\\
C^{n+1}_{i,j}&=D_y+E_y+F_y.\tag{5.32}
\end{align}
在两次式之间从 $C^*$ 再重构。它保持有界性，却不严格守恒面积。

% [原PDF p123]
隐式一词只指质点位置在速度插值中的取值点，并不意味着需要求解全局隐式方程；单元内线性速度使映射可解析反演。参数 $a$ 必须为正，实际由 CFL 和局部形变限制保证，否则原像方向翻转，几何解释失效。

到单元映射改用一阶隐式轨迹式（5.33）：新位置等于旧位置加在新位置处取值、但仍属于旧时间层的速度。结合单元内线性速度可解出式（5.34），其中 $a=[1-u_{i+1/2,j}+u_{i-1/2,j}]^{-1}$。与把源方格向外映射不同，现在先找目标方格 $\Sigma$ 在 $T_x^I$ 下的原像矩形 $\Gamma_x$，再统计该矩形中会进入目标单元的参考相面积。原像矩形同样完整铺砌计算域，所以不会丢失或重复面积；然而连续做 $x$、$y$ 两个隐式映射仍不保证二维总面积守恒。

\subsubsection{到单元隐式映射}

改用
\begin{equation}
x(t^{n+1})=x(t^n)+u[x(t^{n+1})]\Delta t
\tag{5.33}
\end{equation}
可得
\begin{equation}
T_x^I:\quad
\begin{cases}
x'=ax+au_{i-1/2,j},\\
y'=y,
\end{cases}
\qquad a=\frac{1}{1-u_{i+1/2,j}+u_{i-1/2,j}}.
\tag{5.34}
\end{equation}
这里从目标方格 $\Sigma$ 反求其原像矩形 $\Gamma_x$，满足
\begin{equation}
\Sigma=T_x^I(\Gamma_x),
\tag{5.35}
\end{equation}
并把 $\Gamma_x$ 内的参考相映入目标单元：
\begin{equation}
C_{i,j}^{n+1}=D+E+F.
\tag{5.36}
\end{equation}
原像矩形同样铺砌全域，因而不会漏算或重复计算面积；但连续采用两个隐式方向映射也不自动守恒二维面积。

% [原PDF p124]
守恒性质来自隐式与显式映射的特定组合。先以 $T_x^I$ 把原像矩形映到方格，面积乘 $a$；再以 $T_y^E$ 把方格映到竖向矩形，面积乘 $b=1+v_{i,j+1/2}-v_{i,j-1/2}$。组合倍率即式（5.40）。MAC 方格上的离散无散条件使分子和分母相等，因此 $ab=1$，整个步骤严格保存面积，同时继承两个铺砌映射的有界性。为了减少不对称，应在相邻时间步交替 $T_y^ET_x^I$ 与 $T_x^ET_y^I$。若把显式和隐式次序错误交换，$a,b$ 会引用不同单元的速度，离散散度关系不再能推出守恒。

\begin{quote}\small 图 5.9：（a）到单元映射 $T_x^I$ 用单元面速度把原像矩形 $\Gamma_x$ 映到方格 $\Sigma$；（b）三个连续单元在 $\Gamma_x$ 内的贡献映成中心单元中的 $D,E,F$。\end{quote}
\subsubsection{组合线性映射}

先用到单元映射 $T_x^I$，再用出单元映射 $T_y^E$：
\begin{align}
C^*_{i,j}&=D_x+E_x+F_x,\tag{5.38}\\
C^{n+1}_{i,j}&=D_y+E_y+F_y.\tag{5.39}
\end{align}
\begin{equation}
T_y^E:\quad x'=x,\qquad y'=by+v_{i,j-1/2},
\quad b=1+v_{i,j+1/2}-v_{i,j-1/2}.
\tag{5.37}
\end{equation}
组合映射的面积倍率是
\begin{equation}
ab=\frac{1+v_{i,j+1/2}-v_{i,j-1/2}}
{1-u_{i+1/2,j}+u_{i-1/2,j}}.
\tag{5.40}
\end{equation}
交错 MAC 网格上的离散不可压缩条件
\begin{equation}
u_{i+1/2,j}-u_{i-1/2,j}+v_{i,j+1/2}-v_{i,j-1/2}=0
\tag{5.41}
\end{equation}
给出 $ab=1$，故组合 $T_y^ET_x^I$ 严格守恒面积且保持 $0\leq C\leq1$。

% [原PDF p125]
六个二维子速度分量只有五个独立约束，剩余自由度允许从已知边界面开始逐列或逐层构造。只要每个子场在其坐标平面上满足离散无散条件，二维组合映射的 Jacobian 就严格为一。三维实际步进还要轮换三组平面与方向次序，使某一坐标轴不会长期承担相同角色；否则即使总质量守恒，也可能留下可测的方向分裂误差。

三维速度分解中的自由度可用来满足给定边界或简化递推。每个子场只在一个坐标平面内非零并保持无散，使二维守恒映射可以复用；最终三个子步之和近似完整三维运动。分解和扫掠次序仍可能引入各向异性，需要周期轮换。

三维推广可把 $\bm u=(u,v,w)$ 分解为三个二维无散场 $(u_1,v_1,0)$、$(u_2,0,w_2)$、$(0,v_3,w_3)$，并让三者之和为原速度。六个分量受三个无散条件和三个和式约束，其中只有五个独立，剩余自由度可由边界起逐步确定。每个二维无散场再采用守恒组合映射。传统术语把这些几何映射称为 Lagrangian 或 Eulerian：前者用插值速度推进粒子或界面点，再投影回固定网格；后者直接更新固定单元的场量，通过单元面交换面积通量。两者差异由图 5.10 的端点推进和通量矩形清楚表现。

下一步应交替使用 $T_x^ET_y^I$。次序不能任意交换，因为错误组合会从不同单元取速度，无法利用同一离散散度条件。三维扩展可把速度拆成三个二维无散场，分别应用守恒组合映射。

% [原PDF p126]
方向分裂时，即使完整多维速度无散，单个方向仍可压缩或膨胀材料。将 $\partial_tH+\nabla\cdot(\bm uH)=H\nabla\cdot\bm u=0$ 限制到 $x$ 扫掠并对单元积分，右端必须保留平均的 $C\,\partial u/\partial x$。时间积分、中心差分速度梯度并无量纲化后得到式（5.44）。该式中的 $\widetilde C$ 有多种选法，但选定后，面通量必须与压缩项采用一致的几何解释；否则会发生重复输运、漏输运或有界性破坏。

\begin{quote}\small 图 5.10：（a）Lagrangian 方法以面速度插值推进界面端点 $A,B$；（b）Eulerian 方法把右侧面总通量视为宽 $u^+\Delta t$ 的矩形，阴影部分为参考相通量。\end{quote}
\subsection{相关的一维平流格式}

几何映射也可用传统的 Lagrangian/Eulerian 语言解释。Lagrangian 方法推进界面点，再投影回固定网格；Eulerian 方法直接在相邻单元之间交换面积通量。方向分裂的一维体积分数方程必须保留该方向上的压缩或膨胀：
\begin{equation}
\frac{\partial H}{\partial t}+\nabla\cdot(\bm uH)=H\nabla\cdot\bm u=0.
\tag{5.42}
\end{equation}
对单元积分可写为
\begin{equation}
h^2\frac{\partial C_{i,j}}{\partial t}
+\int_\Gamma \bm u\cdot\bm n\,H\,dl
=h^2C_{i,j}\frac{\partial u}{\partial x}.
\tag{5.43}
\end{equation}
\begin{equation}
C_{ij}^{n+1}=C_{ij}^n-\widetilde\Phi_{i+1/2,j}+\widetilde\Phi_{i-1/2,j}
+\widetilde C_{ij}(u_{i+1/2,j}-u_{i-1/2,j}).
\tag{5.44}
\end{equation}
即使总体速度无散，单个方向的速度梯度一般也不为零。

% [原PDF p127]
显式式中的压缩系数 $b$ 与隐式式中的 $a$ 互为不同时间离散下的 Jacobian。把通量区域与错误的 Jacobian 搭配，会把几何一致性误差表现为 wisps。因而“先映射界面再裁切通量”不是实现细节，而是保持有界性的必要条件。

若取 $\widetilde C=C^n$，得到显式 Lagrangian 式（5.45）；其中 $bC^n$ 正是出单元映射中像矩形的中心面积 $E$，两面通量必须对应图 5.7 的 $D,F$。若仍从未映射界面直接算通量，会出现小过冲、欠冲、满相内部小孔和外部 wisps。这些小碎屑与低阶 SLIC 的 $O(h)$ 界面错位所造成的大块 floatsam、jetsam 不同。若取 $\widetilde C=C^{n+1}$，则得到隐式 Eulerian 式（5.46），系数 $a$ 等于到单元映射压缩率；可先求原像矩形内三份贡献，再统一乘 $a$，无须显式搬动所有线段。

取 $\widetilde C=C^n$ 得显式 Lagrangian 格式
\begin{equation}
C_{ij}^{n+1}=bC_{ij}^{n}-\widetilde\Phi_{i+1/2,j}+\widetilde\Phi_{i-1/2,j},
\tag{5.45}
\end{equation}
其三个面积正是出单元映射的 $D,E,F$。若通量不是在映射后的界面上几何计算，就会出现小的过冲、欠冲、孔洞和“细丝状碎屑（wisps）”。这不同于 SLIC 的 $O(h)$ 界面错位造成的更大块漂浮物。

取 $\widetilde C=C^{n+1}$ 得隐式 Eulerian 格式
\begin{equation}
C_{ij}^{n+1}=a\left(C_{ij}^n-\widetilde\Phi_{i+1/2,j}+\widetilde\Phi_{i-1/2,j}\right),
\tag{5.46}
\end{equation}
等价于到单元映射。可以先在原像矩形内求三个体积分数贡献，再统一乘压缩系数 $a$，不必显式移动所有线段。

\subsection{非分裂方法}

% [原PDF p128]
代数总量守恒与物理可接受性是两个独立条件：一个单元出现负体积分数时，其他单元会有等量超额，故全域和仍正确。若随后简单截断，才把这对误差转化为真正的质量损失。设计非分裂算法必须从几何上排除重复区域。

最朴素的非分裂 Eulerian 法只重构一次，同时计算四个面流出的参考相面积，并直接使用式（5.3）。内部面通量仍成对抵消，因此全局质量在代数上守恒；但局部可能不满足 $0\le C\le1$。若参考相只占单元右上角面积 $\delta=u_{i+1/2,j}v_{i,j+1/2}$，水平、竖直通量会把同一面积各输出一次；没有左、下流入时，新值变成 $-\delta$。即使流入足以避免负值，重复输运仍然存在。此外，材料只能进入正交邻单元而不能直接进入对角单元，平移或刚体旋转会迅速生成锯齿并可能导致数值破碎。

\subsubsection{朴素的非分裂 Eulerian 平流}

最直接的非分裂方案从同一次重构同时计算四个面的参考相通量，再按式（5.3）更新。它形式上严格守恒质量，却可能违反有界性：若参考相只占据单元右上角的一小块面积 $\delta$，且水平、竖直出流各自都把这块面积计为 $\delta$，同一材料就被输运两次，得到 $C^{n+1}=-\delta$。同时，没有材料直接进入对角单元，平移或刚体转动也会迅速产生界面锯齿，进而诱发数值破碎。

% [原PDF p129]
将 $T_x^I$ 和 $T_y^E$ 合并为二维线性映射 $\Pi_{xy}$，可在两次方向操作之间省去重构。直线仍映为直线，一个目标方格的原像最多涉及三个相邻界面线段，所以几何裁切比方向分裂复杂但仍局部。每个时间步只需一次 PLIC 重构。为了不让同一坐标方向永远充当隐式方向，需要交替 $\Pi_{xy}$ 与 $\Pi_{yx}$；还可在同一步对二者作线性组合以增强对称性。整个网格可理解为：一组水平矩形先映成方格铺砌，再映成竖直矩形铺砌，刚体旋转示例说明该解释。

\begin{quote}\small 图 5.11：并发分裂平流中，右上角矩形面积同时向右和向上输运，却没有直接流入对角单元，因而发生重复计数。\end{quote}
\subsubsection{非分裂几何线性映射}

不在两个方向之间重构，可把组合写成一个映射
\begin{equation}
\Pi_{xy}=T_x^I T_y^E:\quad
\begin{cases}
x'=ax+au_{i-1/2,j},\\
y'=by+v_{i,j-1/2}.
\end{cases}
\tag{5.47}
\end{equation}
它仍把直线映成直线，一个时间步只需一次界面重构，但几何裁切稍复杂。为避免固定偏向某一坐标方向，应在 $\Pi_{xy}$ 与对偶映射 $\Pi_{yx}$ 间交替；也可在同一步内对二者作对称组合。

% [原PDF p130]
非分裂 Lagrangian 方法沿特征线解 $d\bm x/dt=\bm u(\bm x,t)$。可以推进整个单元的四个顶点，也可以只推进参考相多边形的顶点；局部速度由固定网格插值。轨迹积分可采用二阶预测--校正或更高阶 Runge--Kutta。推进后，必须把变形多边形与固定网格求交；图 5.14 中四块阴影区域分别贡献给四个接受单元，其顶点坐标先由多边形裁切得到，再用式（5.21）算面积。汇总所有供体单元的贡献后得到下一时刻 $C$。理论上无散流保持面积，实际却有若干独立误差源。

\begin{quote}\small 图 5.12：非分裂几何法：（a）初态；（b）$\Gamma_x$ 映到单位方格 $\Sigma$；（c）$\Sigma$ 再映到 $\Gamma_y$。\end{quote}

\begin{quote}\small 图 5.13：所有 $\Pi_{xy}$ 映射把水平矩形铺砌先变成方格网，再变成竖直矩形铺砌；示例为 $u=\omega y,\ v=-\omega x$ 的刚体旋转。\end{quote}
\subsubsection{非分裂 Lagrangian 方法}

Lagrangian 方法沿特征线求 $d\bm x/dt=\bm u(\bm x,t)$，可用预测--校正或 Runge--Kutta 推进单元顶点或参考相多边形顶点。随后求变形多边形与固定网格各单元的交叠面积，并用式（5.21）分配体积。% [原PDF p131]
第一，流动求解器给出的离散速度本身可能不严格无散；第二，从网格点插值到移动顶点的连续速度通常不保持散度为零；第三，特征线时间积分有截断误差；第四，高于一阶的点推进会把原直线段变为连续但分段可微的曲线，速度导数和曲线斜率在跨越单元边界时跳变。因此，变形单元面积一般不等于初始方格面积。可按某些经验准则把面积盈亏重新分给各交叠区，但这并不自动保证每个接受单元仍有 $0\le C\le1$，往往还需另一层局部再分配。

\begin{quote}\small 图 5.14：非分裂 Lagrangian 平流。单元四顶点沿特征线推进，四块阴影面积分别贡献给固定网格的四个单元。\end{quote}
理论上无散流应保持面积，实际仍会因以下因素失守恒：流动求解器的离散速度并非严格无散；插值速度不保持无散；轨迹时间积分有误差；高阶推进使直线段变成跨网格导数不连续的分段曲线。面积余量可再分配，但不保证同时保持 $0\leq C\leq1$。

% [原PDF p132]
非分裂 Eulerian 法不推进供体单元，而构造在 $\Delta t$ 内穿过每个面的“流束多边形”。Rider--Kothe 用 MAC 面心速度把分裂算法的矩形改成梯形：上方增加由横向速度决定的三角形，下方减去另一个三角形。这样可反映多维流动，但梯形面积不一定等于面总通量，而且相邻流束可能重叠。其离散更新于是加入净散度项，并在最后对越界体积分数局部再分配。图 5.15 还比较了 López 方法，它先插值顶点速度，再按守恒面积移动远端边。

\begin{quote}\small 图 5.15：非分裂 Eulerian 的两种面通量多边形。（a）Rider--Kothe 用 MAC 面心速度构造梯形 $1$--$2$--$3$--$4$；（b）López 等人先把速度插值至顶点 1、2，再由面积守恒和速度空间变化确定远端边 3--4。\end{quote}
\subsubsection{非分裂 Eulerian 方法}

Eulerian 非分裂算法构造一个在 $\Delta t$ 内穿过单元面的流束多边形。Rider--Kothe 方法用面心速度把分裂格式的矩形扩展为梯形，以考虑横向流动；但其梯形面积未必等于面通量，而且不同面的流束仍可能重叠。该方法因而在离散方程中加入净散度修正，并对越界的 $C$ 作局部再分配。

L\'opez 等人的方案先把速度插值到面两端点，以局部轨迹确定梯形侧边；再沿既定方向移动远端边，使梯形面积严格等于 $u_{i+1/2,j}\Delta t\,h$。在离散速度无散时，这种构造不会重复输运，所有流入面积与流出面积严格平衡。

% [原PDF p133]
再分配的搜索顺序会影响局部界面形状，即便总质量相同。超量优先传给最接近满且尚未满的邻单元，可减少扩散范围；若邻域容量不足则继续扩展。频繁触发再分配通常说明原通量构造不一致，不应把后处理当成主算法。

López 梯形的两条侧边由面端点速度轨迹确定，远端边斜率还包含速度空间变化；随后沿该方向平移远端边，直至多边形面积严格等于 $u_{i+1/2,j}\Delta t h$。若离散速度无散，该构造既不重复输运，也使全域流入与流出相等。各种不一致仍可能造成 $C<0$ 或 $C>1$。直接截断 $C\leftarrow\min(1,\max(C,0))$ 最简单，却破坏质量。Harvie--Fletcher 再分配把超量传给邻域中尚未满且 $C$ 最大的单元，必要时继续搜索；负值用对称过程补足，从而保持总量。

\subsubsection{再分配算法}

小幅越界最简单的处理是截断
\begin{equation}
C_{i,j}\leftarrow\min(1,\max(C_{i,j},0)).
\end{equation}
但截断破坏全局守恒。若要求精确守恒，可把 $C_{i,j}>1$ 的超量 $\varepsilon=C_{i,j}-1$ 传给邻域中尚未填满且 $C$ 最大的单元，必要时继续向下一邻单元传递；$C<0$ 时反向处理缺量。此类后处理应被看成平流离散不一致的修复，而不是几何通量计算的替代。

% [原PDF p134]
重构和平流的基本验证让流体区域作匀速平移或刚体旋转，理想情况下形状保持不变。更苛刻的非均匀涡量场会强烈拉伸并折叠界面，甚至使其接近破碎或聚并。本书采用初始圆在 $t=T$ 回到起点的可逆试验，以式（5.48）计算初末体积分数的 $L_1$ 几何差。后两个试验统一采用 ELVIRA 法向重构和严格守恒的非分裂几何平流。平移路径特意由多段不同方向的直线组成，以消除始终与网格轴对齐带来的有利或不利偏差。

\begin{quote}\small 图 5.16：平移试验。三个圆沿连续分段直线路径运动，图中还画出 $80\times80$ 网格的一部分。\end{quote}
\section{重构与平流方法测试}

测试通常让液体区域作平移、刚体转动或可逆强变形，精确解应保持形状并返回初始位置。几何误差取
\begin{equation}
E=\sum_{i,j}h^2|C_{i,j}^{T}-C_{i,j}^{0}|.
\tag{5.48}
\end{equation}

\subsection{平移测试}

在单位方域的 $80\times80$ 网格上，半径满足 $r/h=2.5,4,5.5$ 的三个圆沿分段直线路径平移。ELVIRA 加非分裂几何映射在 CFL$=1$ 时几乎只是把体积分数阵列平移一个单元，误差可为零；% [原PDF p135]
单位方域划分为 $80^2$ 单元，三个圆的 $r/h$ 分别为 2.5、4 和 5.5，初始圆心为 $(0.5,0.9)$。每段路径速度恒定，全程使用由最大速度分量定义的固定 CFL。CFL$=1$ 时沿某方向每步恰平移一个单元，$C$ 阵列不变而只改索引，理论误差为零；CFL$=0.1$ 时走同一距离需十倍重构，每次误差积累，小圆最终明显失真。误差随步数增长后趋于饱和，因为相邻时刻界面已很接近。结果提示：单线段 PLIC 要把曲率视为解析良好，局部曲率半径至少约四个网格间距。

\begin{quote}\small 图 5.17：平移试验。（a）$r/h=2.5,4,5.5$ 的三个圆的初始重构；（b）CFL$=1$（实线段）与 CFL$=0.1$（虚线段）的最终重构。\end{quote}
CFL 降至 0.1 后，同一距离需要更多次重构，误差逐步积累，小圆尤其明显。实验表明，单混合单元一线段的 VOF 要可靠解析曲面，局部曲率半径最好至少约为四个网格间距。固定 CFL 下，细化网格时收敛阶由接近一阶逐渐趋向二阶。

% [原PDF p136]
流函数在单元顶点定义，面速度由相邻顶点差分，因而离散旋度和散度关系与 MAC 网格严格兼容。时间调制在 $t=T/2$ 变号，使后一半轨迹理论上逐点逆转；最终误差因此同时衡量数值不可逆性和界面重构损失。

平移误差曲线随 $1/\mathrm{CFL}$ 增大而上升，并在固定 CFL 下由粗网格近一阶逐渐趋向二阶。方盒涡旋则同时检验重构和通量。流函数式（5.49）在单位方域产生顺时针无散旋转，顶点上的 $\psi$ 经中心差分得到 MAC 面速度，离散散度达到机器精度。时间因子 $\cos(\pi t/T)$ 使流动在 $T/2$ 后反向。半径 0.15、圆心 $(0.5,0.75)$ 的圆在半周期被卷成长螺旋丝，到 $T$ 应返回原圆。

\begin{quote}\small 图 5.18：采用 ELVIRA 与非分裂几何线性映射时，平移误差随 CFL 数和网格分辨率的变化。\end{quote}
\subsection{方盒涡旋测试}

可逆涡旋由流函数
\begin{equation}
\psi(x,y,t)=\frac1\pi\cos\left(\frac{\pi t}{T}\right)
\sin^2(\pi x)\sin^2(\pi y)
\tag{5.49}
\end{equation}
生成，$u=\partial\psi/\partial y$、$v=-\partial\psi/\partial x$。% [原PDF p137]
“精确”界面由大量有序标记点用四阶 Runge--Kutta 沿流线推进并连接得到。$T=2$、初始 CFL$=1$ 时，$32^2$ 网格在半周期的细丝头尾处无法解析曲率半径与网格相当的区域；这些位置分别映射为最终圆的下部和上部，也是误差集中的地方。提高到 $128^2$ 后，明显误差只剩在尾端很小的区域。表 5.3 给出 $32^2,64^2,128^2$ 三个分辨率的误差，观测阶约 2.02 和 2.20，说明此组数据呈二阶收敛。

\begin{quote}\small 图 5.19：方盒涡旋速度场。\end{quote}

\begin{quote}\small 图 5.20：$T=2$、初始 CFL$=1$ 的方盒涡旋试验。左图 $t=T/2$、右图 $t=T$；点线为标记点给出的“精确”界面，虚线段与实线段分别为 $32^2$、$128^2$ 网格重构。\end{quote}
初始圆半径为 0.15，圆心位于 $(0.5,0.75)$；在 $T/2$ 时被拉成长螺旋细丝，随后速度反向，于 $T$ 返回原位。% [原PDF p138]
把周期增大为 $T=8$ 会让同一流场作用更久，半周期细丝厚度达到一个网格甚至更薄。每个混合单元只有一条线段的重构无法保存这种亚网格结构，界面逐渐被破坏；算法把微小碎屑收拢为圆形小滴，效果类似人为数值表面张力。图 5.21 对 $32^2$ 和 $128^2$ 网格比较半周期与终态：细网格能更久保持长丝，但仍受最小可解析尺度限制。这一试验提醒，质量严格守恒并不等于几何信息没有丢失；当局部特征低于网格尺度，拓扑和曲率仍会被数值方法重塑。

\begin{quote}\small 表 5.3：$T=2$ 方盒涡旋的几何误差与收敛阶：$32^2$ 网格 $E=2.520\times10^{-3}$；$64^2$ 网格 $E=6.224\times10^{-4}$、$\mathcal O=2.018$；$128^2$ 网格 $E=1.351\times10^{-4}$、$\mathcal O=2.203$。\end{quote}

\begin{quote}\small 图 5.21：$T=8$、初始 CFL$=1$ 的方盒涡旋。绘图约定同图 5.20；强拉伸使细丝达到或薄于网格尺度。\end{quote}
$T=2$ 时，ELVIRA 与非分裂几何平流在 $32^2$ 到 $128^2$ 网格上呈约二阶几何收敛。$T=8$ 时细丝可薄于一个单元，单线段重构不能再解析其头尾；算法会把碎屑聚成小圆滴，表现得像一种数值表面张力。

% [原PDF p139]
颜色函数和水平集分别携带互补信息：$C$ 的单元积分提供严格体积，$F$ 的光滑梯度提供连续法向与曲率。若两者独立推进而不校正，误差会使零等值面与 PLIC 界面逐渐分离；因此每步必须在重构和重初始化两个阶段闭合信息回路。加权拟合中的 $H_\varepsilon'(F)$ 只在窄带非零，也避免远场水平集误差影响法向。

式（5.51）的光滑 delta 权重只让距界面约 $\varepsilon$ 的点显著参与拟合，远处大幅度的水平集值不会支配最小二乘。面积守恒决定 $c$ 后，$a,b$ 主要确定方向；这样 VOF 的体积约束和水平集的光滑法向被明确分工。

混合方法试图同时利用 VOF 的守恒和其他表示的几何优点。CLSVOF 对颜色函数 $C$ 与水平集 $F$ 都求解守恒平流式（5.50）；$C$ 用几何方法，$F$ 用高阶迎风或加权 ENO。两者在重构阶段交换信息：最直接取 $\bm m=-\nabla F$；Sussman--Puckett 还在 $3\times3$ 模板中拟合线性 $\widetilde F=a(x-x_i)+b(y-y_j)+c$。截距 $c$ 由中心单元面积守恒确定，只对 $a,b$ 最小化式（5.51），光滑 delta 权重把拟合集中在界面附近。

\section{混合方法}

VOF 可与水平集结合，分别利用前者的守恒性和后者的光滑几何。不可压缩流中两者都可按
\begin{equation}
\frac{\partial G}{\partial t}+\nabla\cdot(\bm uG)=0
\tag{5.50}
\end{equation}
推进，其中 $G=C$ 时用几何通量，$G=F$ 时用高阶迎风或 WENO。耦合发生在两个环节：重构时直接用 $\bm m=-\nabla F$ 给 PLIC 法向；重构后再把 $F$ 初始化为到 PLIC 线段的符号距离。

还可在 $3\times3$ 模板中拟合线性水平集 $\widetilde F=a(x-x_i)+b(y-y_j)+c$，以面积守恒确定 $c$，再最小化
\begin{equation}
E_{i,j}(a,b)=\sum_{k=i-1}^{i+1}\sum_{l=j-1}^{j+1}
\omega_{k,l}H_\varepsilon'(F_{k,l})
[F_{k,l}-\widetilde F(x_k,y_l)]^2.
\tag{5.51}
\end{equation}
这里 $H_\varepsilon'$ 是宽度 $\varepsilon\simeq\sqrt2h$ 的光滑 delta。% [原PDF p140]
重初始化窄带外取饱和值，是因为流体物性和界面几何只依赖零等值面附近的 $F$；无须在全域维护精确距离。计算端点距离和单元边界距离也防止相邻 PLIC 线段间的小缝在 $F$ 中变成不连续或伪零点。

第二次耦合发生在重构以后：用 PLIC 界面重新初始化 $F$ 为符号距离。若点到直线的垂足在线段内部，距离取法向线段 $AD$；若垂足在段外，还须比较到两个端点的距离。由于各单元 PLIC 线段通常不连续，为使 $F$ 连续还需把适当的单元边界段纳入距离候选。精确距离只在界面周围厚度 $K\le8h$ 的窄带中计算，带外可截为 $(1+K/2)h$。这种耦合也已推广到自适应非结构三角网格：$F$ 用有限元推进，$C$ 用非分裂 Lagrangian 几何法推进。

\begin{quote}\small 图 5.22：（a）点 $A$ 的水平集值由其到线段 1--2 的有符号距离 $AD$ 初始化，若垂足在段外还需比较两端点；（b）为保证 $F$ 连续，还须考虑部分单元边界。\end{quote}
符号距离的计算既要考虑点到线段内部的垂距，也要考虑到端点和必要的单元边界段，才能使 $F$ 连续；通常只在厚度不超过约 $8h$ 的窄带中精确初始化，带外截为常值。

% [原PDF p141]
非结构 CLSVOF 在每个移动三角单元中心以周围共享顶点的单元值作二次最小二乘拟合，再由 $\nabla\widetilde F$ 给法向。另一类混合把 VOF 与显式前沿结合。Aulisa 等人的前沿由网格交点标记和面积守恒标记组成；每个切割单元内参考相形成按逆时针排列的多边形。推进时，单元顶点用四阶 Runge--Kutta 沿流线移动，内部标记依据保存的归一化局部坐标作双线性插值。变形多边形与固定网格求交后，把不同阴影区域的面积分配给接受单元。

\begin{quote}\small 图 5.23：（a）整单元的 Lagrangian 平流及参考相向邻单元再分配；（b）面积守恒标记 $E,F$ 替换内点，同时保持两个多边形面积相同。\end{quote}
VOF 也可与前沿追踪耦合：显式前沿由网格交点标记和面积守恒标记连接而成，Lagrangian 推进单元顶点后，以双线性插值更新内部标记，再将变形多边形面积分配给接受单元。标记点聚集或稀疏时，沿前沿重新分布，并沿法向微调新点，使重新连接前后包围面积不变。% [原PDF p142]
运动中标记点会聚集或拉开，故每步都沿界面重新均匀分布，并在需要时增删。新面积守恒点先沿线段均匀放置，再沿法向微调，使替换前后的多边形面积相等。早期版本平流满单元和混合单元，$C$ 与 Lagrangian 多边形耦合很强，但多次相交和同单元标记聚集需额外约简。改进版只在界面窄带内保存面积，速度更快，且能局部增删而不误删交点；最后按逆时针收集界面标记与单元角点计算 $C$。该方法已推广到三维 Cartesian 和非结构四边形网格。

\begin{quote}\small 图 5.24：（a）界面平流后，仅在界面附近窄带内再分配面积；（b）新面积守恒标记替换内点并保持包围面积，随后按逆时针收集标记与单元角点更新颜色函数。\end{quote}
改进版本只在界面周围窄带内实施面积修正，速度更快，也避免不必要地删除交点标记。此类算法一旦把连续显式前沿作为核心数据结构，实质上已经接近下一章的前沿追踪方法。

% [原PDF p143]
图 5.25 在 $T=16$ 的逆时针方盒涡旋中展示混合前沿：即使 $32^2$ 粗网格也保持连续拓扑，$128^2$ 网格能更细致地保存长丝。随着界面位置主要由连续标记点链表示、体积分数只在再分配后重建，这种“混合 VOF”事实上已成为前沿追踪方法。它自然引出下一章：如何组织标记点和连接、如何在变形中重构前沿，以及怎样同固定流体网格可靠交换信息。

\begin{quote}\small 图 5.25：$32^2$（左）和 $128^2$（右）网格上的方盒涡旋；周期 $T=16$，旋转方向为逆时针。\end{quote}
方盒涡旋的混合前沿试验进一步表明：在 $32^2$ 与 $128^2$ 网格、周期 $T=16$ 的逆时针流动中，面积守恒标记能在强拉伸后仍维持连续前沿。随着耦合算法逐步把连续显式前沿作为主表示，它在算法分类上也自然过渡为前沿追踪。数据结构，实质上已经接近下一章的前沿追踪方法。

\chapter{标记点的平流：前沿追踪}
\markboth{第六章\quad 标记点的平流：前沿追踪}{第六章\quad 标记点的平流：前沿追踪}

% [原PDF p144]
前沿追踪不直接平流网格标记函数，而用随流体运动的相连标记点表示两相边界。具体实现不同会形成不同算法，但都必须处理四个方面。第一是前沿数据结构；“前沿”不仅是点坐标，还包括点之间的连接、单元方向以及界面物理属性。二维几乎任何结构都能工作，三维数据结构却可能决定算法成败。第二是维护受拉伸、压缩和弯曲的前沿，要动态增删对象、重新连接，并在必要时平滑。第三是前沿与求解 Navier--Stokes 方程的固定网格交换速度、物性和表面张力。第四是液滴聚并或破碎时的拓扑改变。

前沿追踪不用网格上的颜色函数直接搬运界面，而是用相互连接、随流体运动的标记点表示两相边界。本章的“前沿”指表示界面的全部计算对象：标记点、连接单元、邻接关系以及界面物理属性。

% [原PDF p145]
除光滑耦合外，尖锐前沿法还包括跨界面外推形成幽灵值，以及在前沿附近修改固定网格。前者保持常规模板但需准确跳跃关系，后者把界面变成控制体边界却要管理不规则小单元。本章选择光滑耦合，是为集中说明 Unverdi--Tryggvason 型方法。

本章采用光滑界面耦合：把前沿代表的奇异量分布到固定网格的有限宽区域，而不是在界面附近修改差分、构造幽灵值或重划切割网格。一个计算同时含两个网格：规则体网格求流动，低一维前沿网格追踪界面并计算表面张力。图 6.1 的浮升气泡中，气泡面清楚显示非结构三角剖分，背景剖面则是规则速度网格。三维前沿必须能高效访问角点、邻接面和界面属性；拉伸后要加密，压缩后要粗化。拓扑改变长期被视为显式追踪最大难题，本章最后会讨论检测和重连的主要原则。

一个完整方法必须处理四类问题：合适的数据结构；界面拉伸和压缩时的动态加密、粗化与平滑；移动前沿和固定流体网格之间的信息交换；液滴合并、液丝夹断等拓扑改变。二维前沿的管理相对简单，三维曲面通常必须采用非结构三角网格。固定网格用来求解 Navier--Stokes 方程，低一维的前沿网格负责准确平流界面和计算表面张力。

\section{前沿的结构}

% [原PDF p146]
二维界面可直接存为有序点数组 $\bm x_i=(x_i,y_i)$。点的数组次序就是沿前沿的物理次序，因此点 $i$ 的邻点为 $i\pm1$，距离、切向和曲率的局部计算都很直接。若同时存在多条界面，可以为每条曲线使用独立数组，也可在一个数组中保存每条曲线的起止索引。为了求切向或在两点之间插点，仍必须假设两标记点间的曲线形状。点很密或原型程序中可用直线段，但线性插点会带来面积误差和曲率波动，表面张力较强时尤其不理想，因此常采用局部高阶插值。

\begin{quote}\small 图 6.1：浮升气泡的前沿追踪计算。气泡表面由非结构三角网格追踪，流体方程则在规则结构网格上求解。\end{quote}
\subsection{结构化二维前沿}

二维闭合界面可用有序标记点数组表示。只需存储
\begin{equation}
\bm x_i=(x_i,y_i),\qquad i=1,\ldots,N,
\tag{6.1}
\end{equation}
点 $i$ 的两个邻点即 $i\pm1$。多条界面可以分别存成数组，也可额外记录每条曲线的首尾索引。

% [原PDF p147]
Legendre/Lagrange 型局部插值只依赖附近少量点。令参数 $u$ 在标记点上取连续整数，式（6.2）以 $n+1$ 个点给出 $n$ 次曲线。四点时得到式（6.3）的三次多项式；代入 $u=3/2$ 可求点 1、2 之间的曲线中点，微分则给出该处未归一化切向。上标“3”表示三阶插值，不是空间维数。需要注意，式（6.5）的 $\bm t^{(3)}$ 尚未除以长度。图 6.2 中点列的参数顺序使上述运算完全局部；只要前沿不发生强烈拉伸或压缩，有序数组实现简洁而高效。

\begin{quote}\small 图 6.2：二维前沿用一串有序物质点表示的布局。\end{quote}
为了求切向或插入新点，必须假定标记点之间的曲线形状。线性插值简单，但插点时可能导致面积误差，并在高表面张力问题中引入曲率波动。局部 Legendre 插值更平滑。设局部参数 $u$ 在各标记点处取整数，则通过 $n+1$ 个点的多项式为
\begin{equation}
\bm x^{(n)}(u)=\sum_{i=0}^{n}
\left[\prod_{\substack{j=0\\j\ne i}}^{n}\frac{u-u_j}{u_i-u_j}\right]\bm x_i.
\tag{6.2}
\end{equation}
对 $u=0,1,2,3$ 的四点，三次式为
\begin{align}
\bm x^{(3)}(u)={}&-\frac16(u-1)(u-2)(u-3)\bm x_0
+\frac12u(u-2)(u-3)\bm x_1\notag\\
&-\frac12u(u-1)(u-3)\bm x_2
+\frac16u(u-1)(u-2)\bm x_3.
\tag{6.3}
\end{align}
点 1 和 2 之间的中点及其未归一化切向为
\begin{align}
\bm x^{(3)}(3/2)&=\frac1{16}[-\bm x_0+9(\bm x_1+\bm x_2)-\bm x_3],
\tag{6.4}\\
\bm t^{(3)}(3/2)&=\left.\frac{\partial\bm x}{\partial u}\right|_{3/2}
=\frac1{24}[-\bm x_0+27(\bm x_1-\bm x_2)+\bm x_3].
\tag{6.5}
\end{align}

% [原PDF p148]
全局等弧长重分布会移动所有点，即使局部间距原本合适；局部阈值法只改问题区域，扰动较小。另一方面，全局重分布能避免长期积累的点密度不均。实际选择取决于界面变形强度、是否需要点携带历史变量以及面积守恒要求。

界面变形会使一部分点间距过大、另一部分过密。间距过大会丢失几何精度，过密既浪费存储，也可能在前沿分辨率远高于速度网格时产生小尺度波纹。因此可把点按弧长重新均匀分布，或只在超过最大/最小阈值时局部插入、合并。真实弧长计算较繁，本书用相邻点弦长累加的式（6.6）；闭合曲线总长为 $s_N$，目标间距 $\Delta s=s_N/N$，新点取插值曲线 $\bm x(i\Delta s)$。局部插点或合点时采用式（6.4）的曲线中点，比简单算术平均更少扰动面积和曲率。有序数组中批量增删还需要移动后续索引。

界面变形后，一些点间距过大，另一些点又过密。可把点按弧长均匀重分布，或仅在超过阈值时局部增删。弧长常用弦长近似
\begin{equation}
s_i=\sum_{j=1}^{i-1}
\sqrt{(x_{j+1}-x_j)^2+(y_{j+1}-y_j)^2}.
\tag{6.6}
\end{equation}
闭合曲线的目标间距为 $\Delta s=s_N/N$，新点由插值曲线 $\bm x(i\Delta s)$ 给出。局部加点时，把新点放在式（6.4）确定的曲线中点；删点时把过近两点合并到同一位置，比简单坐标平均更能保持面积和曲率。

% [原PDF p149]
当点数较少而精度要求很高，可用分段三次样条 $\bm x(s)=\bm a_is^3+\bm b_is^2+\bm c_is+\bm d_i$，要求函数及一、二阶导数在标记点连续。以弧长近似作参数，$x(s)$ 与 $y(s)$ 的系数分别由三对角系统求解。样条能产生很高精度，但主要适用于二维和较简单拓扑。三维复杂曲面不可能用单个有序数组有效表示，除非始终近似平面；通用做法是非结构三角前沿，由点和三角单元共同组成。为了直观说明数据结构，作者先用二维“点 + 两点线元”示例，再推广到三维三角形。

\begin{quote}\small 图 6.3：二维前沿的数据结构。点只保存坐标；单元携带前沿结构信息，包括角点和相邻单元指针。\end{quote}
需要更高精度时可采用分段三次样条
\begin{equation}
\bm x(s)=\bm a_i s^3+\bm b_i s^2+\bm c_i s+\bm d_i,
\tag{6.7}
\end{equation}
并要求函数及一、二阶导数在标记点连续。样条适合二维高精度问题，但三维复杂曲面更常使用局部三角网格插值。

% [原PDF p150]
非结构二维前沿中，点对象只存坐标，线元对象则保存两个角点指针、前后邻接线元、方向、表面张力、标记函数跳跃和其他局部界面量。方向定义界面的内外侧，同一物理界面所有单元必须一致。运行时频繁增删点和单元，最方便的存储是链表：每个对象除自身数据外还含下一对象指针；不用的数组槽也组成空闲链表。C/C++ 结构体可把这些量集中保存，其他语言也可用多个平行数组实现。维护双向链表虽非必须，却能显著简化删除操作。

\begin{quote}\small 图 6.4：二维前沿单元的链表连接。图中含两条界面和一组幽灵单元；链表次序原则上与单元的物理次序无关。\end{quote}
\subsection{非结构前沿}

三维界面除近似平面外，不适合用有序数组。通用表示是“点 + 三角单元”：点只保存坐标；单元保存角点指针、共享边的邻单元指针、表面张力系数、标记函数跳跃等界面信息。每个单元带一致的方向，从而定义界面的“内侧”和“外侧”。

% [原PDF p151]
链表次序与前沿的物理次序无须一致；任意指定首对象，沿 next 指针访问预存数量的对象即可。多条液滴界面通常也无需分成不同计算结构，即使表面张力系数不同，只要各单元携带局部物性；只有处理方式根本不同或出于计算组织需要时才另设前沿。界面与外边界相交时常加入域外幽灵点，并以幽灵单元连接壁面点。幽灵对象单独成链，但壁面处真实单元把幽灵单元视作邻居；调节幽灵点位置可令壁面切向或接触角满足边界条件。三维同理，只是一个面单元含三个角点和三个邻面。

\begin{quote}\small 图 6.5：三维前沿结构。点只保存坐标；三角单元保存角点与相邻单元指针等全部拓扑信息。\end{quote}
由于计算中要频繁增删对象，点和单元适合存入双向链表。每个对象除自身数据外，还含前、后对象指针；另建空闲对象链表，分配和回收都无需移动大数组。多条物理界面未必需要不同数据结构，只要单元携带各自物性即可；若界面与外边界相交，可在域外加幽灵点和幽灵单元，通过调节幽灵点位置施加指定接触角或切向条件。

% [原PDF p152]
重心坐标不依赖全局 Cartesian 方向，适合任意朝向三角面。式（6.9）使用邻面外点，是为了获得跨边的二阶几何信息；若相邻面缺失，例如接触边界，则必须用幽灵面或适当的一侧插值补齐模板。

三维每个三角形的角点与边按从界面外侧观察的统一顺时针或逆时针次序编号。需要在三角形内部插值时，以中心三角形的三个点和跨三条边的三个邻点构造重心坐标 $(u,v,w)$，满足 $u+v+w=1$；三条边分别由 $u=0,v=0,w=0$ 表示。式（6.9）给出六点二次插值，可用于坐标或任意标量、向量界面量。它在中心三角形的顶点保持给定值，并利用邻面外点提供曲率信息。局部插值的准确性决定新点放置、面心位置和几何导数的平滑程度。

\begin{quote}\small 图 6.6：三维三角单元的重心坐标；单元 $i$ 的三条边分别由 $u=0,v=0,w=0$ 给出。\end{quote}
三维每个三角单元由三个点构成，并与三个相邻三角形共享边。为用中心三角形及其三个邻三角形上的共六个点插值，可引入重心坐标 $(u,v,w)$，满足
\begin{equation}
u+v+w=1.
\tag{6.8}
\end{equation}
若中心三角形顶点为 1、2、3，对边邻三角形的外顶点为 4、5、6，则任意标量或向量 $p$ 的二次局部插值为
\begin{align}
p(u,v,w)={}&\frac12(1-u)[-up_5+(1-v)p_3+(1-w)p_2]\notag\\
&+\frac12(1-v)[(1-u)p_3-vp_6+(1-w)p_1]\notag\\
&+\frac12(1-w)[(1-u)p_2+(1-v)p_1-wp_4].
\tag{6.9}
\end{align}
% [原PDF p153]
把式（6.9）代入 $u=v=w=1/3$ 可得中心三角形形心的二次近似；代入 $(1/2,1/2,0)$ 得边 1--2 的曲线中点。二者不仅使用本三角形三个角点，也用相邻三角形的外点修正，因此比简单算术平均更贴近光滑曲面。二维增删示意见图 6.7：过长线元在曲线中点插点后分为两段，过短线元则把两端合并并删除。三维同样要维护点、边和面邻接；加密或粗化改变局部拓扑以后，必须同步更新所有相邻指针和界面方向。

\begin{quote}\small 图 6.7：二维前沿单元的增删。（a）在长单元中点插点并分裂；（b）合并短单元两端点并删除该单元。\end{quote}
中心三角形形心和边 1--2 的中点分别近似为
\begin{align}
\bm x(1/3,1/3,1/3)&=\frac19[4(\bm x_1+\bm x_2+\bm x_3)
-(\bm x_4+\bm x_5+\bm x_6)],
\tag{6.10}\\
\bm x(1/2,1/2,0)&=\frac12(\bm x_1+\bm x_2)+\frac14\bm x_3
-\frac18(\bm x_5+\bm x_6).
\tag{6.11}
\end{align}

% [原PDF p154]
归一化宽高比在等边三角形取 1，细长退化三角形则迅速增大，因此比只看最长边更能识别差形单元。边长阈值控制分辨率，形状阈值控制插值和曲率计算的条件数；两者都必须检查。

本节所谓“重构前沿”是在物理拓扑不变的前提下调整分辨率和单元形状。拉伸区如果不加点，几何误差会增长；压缩区若不删点，对象数持续增加且可能形成网格尺度以下的波纹。二维只需分裂长线元或塌缩短线元。三维分裂一条过长边会增加一个点和两个三角形；塌缩短边会删除一个点和两个三角形；翻转公共边则不改变对象数，只改变连接以改善三角形形状。由于前沿只追踪几何并携带较少物理量，不直接离散完整流体方程，其网格质量要求低于体网格，维护成本通常也远小于流动求解。

\section{前沿的拓扑重构}

这里的“重构”指在不改变物理拓扑的前提下，维护三角化或点列分辨率。界面受流动拉伸时需加点，压缩时应删点，以免点数无谓增长或出现由过度分辨引起的小尺度波纹。前沿网格主要承载几何和界面量，并不直接离散全部守恒方程，所以网格质量要求可比体网格简单。

% [原PDF p155]
翻边并非总是改善几何。若长窄三角形沿尖锐脊线排列，翻转公共长边可能让新边横跨脊线，把原本平滑的折线变成锯齿；因此有时应保留原连接，或删除另一面。实际中先分裂长边、再塌缩短边，常会自然实现所需连接调整，不必频繁翻边。二维通常使每个固定网格间距包含 2--4 个前沿线元。三维可用周长平方除以面积、再除等边三角形值 $12\sqrt3$ 的归一化宽高比。经验阈值是最短边约 $\Delta x/3$、最长边约 $\Delta x$、宽高比不超过 $3/2$。宜先全局标记不良单元，再统一修改并视需要第二次扫描。

\begin{quote}\small 图 6.8：三维前沿重构。分裂长边会增加一个点和两个单元；塌缩短边会删除一个点和两个单元；翻边重塑不改变对象总数。\end{quote}
二维中，长线元通过插入一个曲线中点分裂为两个，短线元则把端点合并后删除。三维中，加密通常分裂过长边：增加一个点、两个三角形；粗化则塌缩过短边：删除一个点、两个三角形。还可通过翻边改善形状而不改变点、面数。但翻边不总有利，例如狭窄脊线上沿脊方向排列的长窄单元被翻转后，会把光滑脊变成锯齿；很多时候先细化、后粗化已经能够隐式完成所需的连接调整。

二维可直接按线元长度设最小和最大阈值，通常让每个固定网格间距内含 2--4 个前沿线元。三维除边长外，还可用归一化宽高比衡量三角形形状：周长平方除以面积，再以等边三角形的对应值 $12\sqrt3$ 归一化。经验上，最短边约取固定网格间距的 $1/3$、最长边取一个网格间距、最大归一化宽高比取 $3/2$，常能取得良好结果。

增删元素会改变邻域，因此宜先扫描整个前沿，标记所有不良单元，再统一修改；若一次修改数量很多，可再扫描一轮。前沿管理通常只占多相流计算总耗时的一小部分，第二轮检查的代价很低。

% [原PDF p156]
定位公式中的 $N_x/L_x$ 取决于网格点与单元数的编号约定；实际代码若有 $N_x$ 个单元、$N_x+1$ 个点，应相应调整。书中公式采用给定的周期规则网格约定，边界点还要避免 floor 产生越界索引。

固定网格把速度传给前沿，前沿又把物性跳跃与表面张力传回网格。这里采用光滑耦合，避免修改界面附近的流体离散。第一步是把前沿对象定位到固定网格。规则域长 $L_x$、含 $N_x$ 点且 $i=0$ 对应 $x=0$ 时，式（6.12）用 floor 运算给出 $x_f$ 左侧网格索引。其他方向完全类似。对逻辑矩形曲线网格，可先把物理坐标映到规则参数空间，再套用同一索引方法。前沿点到邻近网格的定位必须非常高效，因为插值、分布、标记函数重构和接触检测都会重复使用它。

\section{前沿与固定网格之间的通信}

前沿速度来自固定网格，而密度、黏度和表面张力等信息又要由前沿反馈到网格。以下讨论光滑界面耦合：把尖锐界面量扩展为固定网格上有限宽的分布，流动离散模板不必在界面附近改变。

\subsection{把前沿定位到固定网格}

对长度 $L_x$、含 $N_x$ 个网格点的规则网格，若 $i=0$ 对应 $x=0$，前沿点 $x_f$ 左侧的网格索引为
\begin{equation}
i=\operatorname{FLOOR}\!\left(x_f\frac{N_x}{L_x}\right).
\tag{6.12}
\end{equation}
% [原PDF p157]
局部候选链表本质上是空间哈希。前沿移动以后，只需清空并重新填充受影响的格子，或每步线性重建；其成本与前沿对象数成正比。拓扑检测也可复用同一结构，只比较相邻格子中的线元或三角面。

周期域中前沿坐标不必每次折回基本周期；可对 $x_f$ 取 $\operatorname{MOD}(x_f,L_x)$ 后再求 floor，得到等价的固定网格单元。计算跨周期相邻点的距离或插值时，需要加减一个周期以避免把近邻误判为远点。从一个前沿点找到附近网格点是 $O(1)$，反向从每个网格点遍历全部 $N_f$ 前沿点则是 $O(N_f)$，大量查询会很昂贵。更好的方法是遍历前沿一次，为每个固定网格单元建立局部前沿对象短链表；之后只在本单元和相邻单元的候选中搜索，既降低成本，也便于并行域分解。

周期域中，前沿无需人为限制在固定网格采用的那一个周期，只需用
\begin{equation}
i=\operatorname{FLOOR}\!\left[
\operatorname{MOD}(x_f,L_x)\frac{N_x}{L_x}\right]
\tag{6.13}
\end{equation}
找到等价网格位置。周期曲线跨域端点之间计算距离或插值时，再补上周期偏移。逻辑矩形曲线网格也可先映射到直角参数域，再用同样公式。

从前沿点找最近网格点很容易，反向从网格点找最近前沿点则昂贵。逐点搜索 $N_f$ 个前沿点是 $O(N_f)$；若很多网格点都要查询，先遍历前沿，为每个网格单元建立短的局部候选链表，再在局部列表中搜索会高效得多。

% [原PDF p158]
识别邻近网格点后，网格量按式（6.14）插值到前沿位置。权重必须非负并归一，使插值不超出邻域极值；前沿点恰落在网格点上时，还应恢复原网格值。相反方向的操作，是把定义在前沿上的线密度或面密度分布到固定网格。图 6.9 表明同一组邻近点可用于双向通信。采用光滑界面以后，密度、黏度和奇异表面力都表现为若干单元宽的过渡区；核函数的支撑宽度和平滑阶数直接控制网格印记、界面厚度与计算量。

\begin{quote}\small 图 6.9：前沿与固定网格。前沿标记由单元相连，信息在前沿点与附近固定网格点之间双向传递。\end{quote}
\subsection{插值与平滑}

固定网格上的量可插值到前沿位置 $l$：
\begin{equation}
\phi_f^l=\sum_{i,j}w_{i,j}^l\phi_{i,j},
\tag{6.14}
\end{equation}
权重应保证插值值不超出邻域范围，前沿点恰落在网格点上时能恢复原值，并满足
\begin{equation}
\sum_{i,j}w_{i,j}^l=1.
\tag{6.15}
\end{equation}
% [原PDF p159]
若分布核没有严格归一，前沿总力会随标记点相对网格的位置周期性变化，形成非物理寄生振荡。用同一权重作插值和分布还能满足功率一致性：网格力对网格速度所做功与前沿力对插值速度所做功相匹配。

前沿到网格的分布必须保持积分总量：一段前沿上的 $\phi_f\Delta s$ 应等于覆盖网格体积中 $\phi\Delta v$ 的和。二维离散式（6.17）因此包含 $\Delta s_l/(\Delta x\Delta y)$；三维相应使用三角面面积除以单元体积。通常使用同一核完成网格到前沿的插值和前沿到网格的分布，使两者形成离散伴随。若界面携带表面张力或标量跳跃，这一守恒关系保证总界面力或总源项不因平滑而改变。支撑内权重还应满足零阶矩条件式（6.15），必要时满足更高矩条件以降低平移误差。

前沿上的面密度（二维为线密度）量向固定网格平滑时，还必须保存积分总量：
\begin{equation}
\int_{\Delta s}\phi_f(s)\,ds=\int_{\Delta v}\phi(\bm x)\,dv.
\tag{6.16}
\end{equation}
二维离散式为
\begin{equation}
\phi_{i,j}=\sum_l\phi_f^l w_{i,j}^l\frac{\Delta s_l}{\Delta x\Delta y}.
\tag{6.17}
\end{equation}
通常用同一核函数完成“网格到前沿”的插值和“前沿到网格”的分布，使两个操作互为离散伴随。

% [原PDF p160]
多维权重常写成一维函数的张量积，例如二维 $w=d(r_x)d(r_y)$。最紧致的面积权重 $d(r)=1-|r|$（$|r|<1$）只覆盖二维四点、三维八点，边界实现简单且多数问题足够。更光滑的余弦核、Peskin 四点核或三次 B 样条具有连续导数并减弱网格印记，但三维支撑可扩大到 27 或更多网格点。图 6.10 的几何面积解释说明，前沿点在一个单元内移动时，四个权重按与对角子矩形面积成比例连续变化。选择核函数必须在紧致性、光滑性和界面过渡宽度之间折中。

\begin{quote}\small 图 6.10：面积权重既用于把网格量插值到前沿，也用于把前沿信息平滑分布到网格。\end{quote}
多维权重常写成一维核的乘积：
\begin{equation}
w_{i,j}^l=d(r_x)d(r_y),
\tag{6.18}
\end{equation}
其中 $r_x,r_y$ 是以网格间距归一化的坐标距离。最紧致的面积权重为
\begin{equation}
d(r)=
\begin{cases}
1-r,&0<r<1,\\
1+r,&-1<r<0,\\
0,&|r|\geq1.
\end{cases}
\tag{6.19}
\end{equation}
它在二维用四点、三维用八点，效率高且边界处理简单。Peskin 型余弦核、四点光滑核和三次 B 样条具有更光滑的导数，但三维支撑可扩展到 27 个点。支撑过窄会增强网格印记，支撑过宽又增加成本并加厚过渡层；面积权重在多数问题中已足够，强表面张力或对平滑性敏感时可选更光滑的核。

% [原PDF p161]
表 6.1 汇集了若干常用的一维平滑核：线性面积权重支撑最窄；余弦和四点 Peskin 核在支撑边界更光滑；三次样条也提供连续导数。多维时按各坐标方向乘积。完成速度插值后，最简单的前沿推进是显式 Euler 式（6.20），实际计算通常采用预测--校正或 Runge--Kutta 以降低时间误差。与 VOF 面通量的代数守恒不同，点轨迹推进不自动保证闭合前沿包围面积守恒；即便固定网格速度离散无散，普通插值后的连续速度也未必严格无散。

\section{前沿的平流}

已由插值得到每个前沿点的速度后，就可积分其轨迹。一阶显式 Euler 更新为
\begin{equation}
\bm x_f^{n+1}=\bm x_f^n+\bm v_f^n\Delta t.
\tag{6.20}
\end{equation}
实际通常使用高阶 Runge--Kutta 或预测--校正格式。与 VOF 通过公共面通量严格保存单元平均不同，标记点轨迹推进本身不保证包围体积守恒。误差来自速度插值以及时间积分；即便固定网格上的离散速度无散，插值后的连续速度也不一定无散。可构造保持散度的插值，但会使压力方程更复杂。

对长时间、多液滴且单滴分辨率较低的计算，累计体积漂移可能不可接受。这时可隔若干时间步对液滴尺度作很小的全局校正；由于每次修正很小，其动力学影响通常可忽略，但应同时记录校正量以监测算法误差。

% [原PDF p162]
可构造保持散度的速度插值，但通常会增加压力求解或插值结构的复杂度。对长时间、多液滴且每滴分辨率较低的模拟，微小面积误差会积累；可隔若干步对每个液滴作极小的尺度校正，同时记录修正量作为误差诊断。图 6.11 的可逆涡旋采用与 VOF 方盒试验相同的时空结构：圆在半周期被拉成长丝，随后速度反向。原书以 $32^2$ 速度网格、$T=8$、$\Delta t=0.005$ 和二阶预测--校正得到几乎重合的初末前沿；同等形状精度下，显式前沿所需固定网格明显粗于 VOF。

\begin{quote}\small 图 6.11：前沿平流精度试验。圆形前沿在涡旋流中变形，随后流动反向，使前沿返回初始形状。\end{quote}
可逆方盒涡旋是常用验证。速度场为
\begin{align}
u&=2\cos(\pi t/T)\sin^2(\pi x)\sin(\pi y)\cos(\pi y),\\
v&=-2\cos(\pi t/T)\sin^2(\pi y)\sin(\pi x)\cos(\pi x).
\end{align}
半径 0.15、圆心 $(0.5,0.75)$ 的标记点圆在 $T/2$ 达到最大拉伸，随后流向反转，并应在 $T$ 回到起始圆。原书使用 $32^2$ 速度网格、$T=8$、$\Delta t=0.005$ 和二阶预测--校正，返回界面与初始界面几乎重合。相同精度下，显式前沿所需固定网格明显粗于前章的 VOF 测试，体现了标记点对几何位置的高分辨能力。

% [原PDF p163]
前沿只推进界面位置，密度、黏度等网格物性没有直接被平流，所以每个时间步都要从新界面重设。通常先构造在各相中为常数的标记函数 $I(\bm x)$，再按 $I$ 插值得到物性。三种路线分别是：按网格点到界面的最短法向距离更新附近点；把界面上的 $I$ 跳跃梯度分布到网格并积分恢复 $I$；直接由前沿切割单元的几何位置求体积分数。第一种最直观，但当两段界面同时位于同一对网格点之间或界面折回时，最终值取决于处理哪一段，薄膜中这种歧义很常见，因此本章重点讨论后两种通用构造。

\section{标记函数的构造}

前沿追踪只推进两相边界，密度和黏度等网格物性必须在每个时间步由界面位置重新设置。通常先构造在各流体中为常数的标记函数 $I(\bm x)$，再由 $I$ 插值得到物性。有三条主要路线：
\begin{enumerate}[label=(\roman*)]
\item 由到界面的最短法向距离直接设置附近网格点；
\item 把标记函数在界面上的跳跃梯度分布到固定网格，再积分恢复 $I$；
\item 根据前沿与单元的交线直接计算单元体积分数。
\end{enumerate}
第一种方法在两段界面落在同一对网格点之间时会产生歧义，折叠界面和薄膜中很常见，因此通用算法更偏向后两种。

% [原PDF p164]
界面表示标记函数的跳跃；在分布意义下，网格梯度等于沿前沿积分 $\Delta I\bm n\delta(\bm x-\bm x_f)$。相距很近、法向相反的两段界面会使分布后的梯度自然抵消，无须判断先穿过哪一段。离散式（6.22）把每个前沿线元的跳跃、法向、长度和核权重累加到固定网格。不同梯度分量可置于不同的交错位置，例如 $\partial_xI$ 在竖直面心、$\partial_yI$ 在水平面心。随后若已知左邻点的 $I$，可用式（6.23）直接跨面积分得到当前点。

\begin{quote}\small 图 6.12：在交错网格上由标记函数梯度构造标记函数。\end{quote}
\subsection{由梯度构造标记函数}

界面代表 $I$ 的跳跃，固定网格上的陡峭梯度可写成
\begin{equation}
\nabla I=\int_\Gamma \Delta I\,\bm n\,
\delta(\bm x-\bm x_f)\,ds.
\tag{6.21}
\end{equation}
若两条方向相反的界面很近，它们的网格梯度会自然抵消，因而不必判断先遇到哪一条。用上一节的分布权重得到
\begin{equation}
(\nabla I)_{i,j}=\sum_l(\Delta I\,\bm n)_l
w_{i,j}^l\frac{\Delta s_l}{\Delta x\Delta y}.
\tag{6.22}
\end{equation}
% [原PDF p165]
把四个单向积分结果平均，不只是为了形式对称，也相当于寻找一个标量场，使其梯度在最小二乘意义下最接近由前沿分布得到的向量场。离散前沿梯度未必严格无旋，所以逐路径积分会给不同答案；Poisson 方程取了所有路径的一致投影。求解时必须由窄带外原有的相标记或一个固定参考点消除常数零空间，并保持界面两侧跳跃符号与法向方向一致。

右端散度使用半网格梯度值，而不是再对已经重建的 $I$ 求差分；这样保留了前沿提供的跳跃信息。Poisson 解只确定到一个常数，窄带边界或已知相内值提供该常数。若在周期全域求解，还需显式固定均值。

沿一条网格线反复使用式（6.23）需要先找到已知标记值的起点，而且结果略依赖积分路径；跨越多个点后误差会积累，使界面另一侧的常值偏离目标。更对称的做法从左右上下四邻点分别积分并平均，整理后得到式（6.24）。左端是标准二阶中心 Laplacian，右端是从前沿直接分布到半网格点的梯度之散度，因此连续形式为 $\nabla^2I=\nabla\cdot(\nabla I)$。这一 Poisson 问题可用标准方法求解，却没有必要遍历整个域，因为旧标记函数只有前沿扫过的单元及其邻域发生改变。

两个梯度分量可以放在交错网格的相应面心上。若已知左邻点的标记值，最直接的积分是
\begin{equation}
I_{i,j}=I_{i-1,j}+\Delta x
\left(\frac{\partial I}{\partial x}\right)_{i-1/2,j}.
\tag{6.23}
\end{equation}
但逐线积分依赖路径，并会累积误差。把从四个相邻方向积分得到的结果平均，可整理成离散 Poisson 方程
\begin{align}
&\frac{I_{i+1,j}+I_{i-1,j}-2I_{i,j}}{\Delta x^2}
+\frac{I_{i,j+1}+I_{i,j-1}-2I_{i,j}}{\Delta y^2}\notag\\
&=\frac{(\partial_x I)_{i+1/2,j}-(\partial_x I)_{i-1/2,j}}{\Delta x}
+\frac{(\partial_y I)_{i,j+1/2}-(\partial_y I)_{i,j-1/2}}{\Delta y},
\tag{6.24}
\end{align}
即
\begin{equation}
\nabla^2 I=\nabla\cdot(\nabla I).
\tag{6.25}
\end{equation}

% [原PDF p166]
窄带实现分三步：遍历前沿，标记它穿过的固定网格单元；再扫描网格，把这些单元及一圈邻居串成单向链表；最后只沿链表迭代式（6.24）。界面每步移动不超过一个单元，带宽数格即可，迭代通常很快。靠近外边界时，平均式中不包含指向墙外的方向。圆滴示例先把 $x$、$y$ 梯度分布到 $32^2$ 网格，再解 Poisson 恢复 $I$。若前沿点太稀，紧致核无法均匀覆盖所有网格点，会产生轻微过冲、欠冲；若在全域求解，还可能出现平均值漂移，必要时可后处理修正。

\begin{quote}\small 图 6.13：在前沿附近窄带内求解式（6.24）。黑点接收了平滑后的梯度；开圆点通常也纳入窄带，以容纳前沿运动。\end{quote}
标记函数只在界面扫过的单元中改变，无需在全域求解。实现时先遍历前沿，标记包含界面的单元及其一圈邻居；再把这些单元组织成单向链表；最后只在窄带内迭代求解式（6.24）。带宽仅数个单元，收敛通常很快。靠近固壁时，可从平均式中去掉指向墙外的方向。若前沿点过稀，离散 delta 核不能均匀覆盖网格，重建 $I$ 会出现轻微过冲和欠冲；若在全域解 Poisson 方程，还可能产生均值漂移，必要时可在后处理阶段校正。

% [原PDF p167]
直接由前沿构造体积分数的目标，是允许一个单元被任意多段界面以任意方向穿过，并尽量不依赖邻单元构型。若界面与单元边界围成参考相多边形，体积分数就是多边形面积除以单元面积。二维且每单元只穿越一次时实现较直接；多次穿越、非结构前沿以及三维都会显著增加交点排序和边界补段的复杂度。为了统一面积计算，引入满足 $\nabla\cdot\bm A=1$ 的向量场，并由散度定理把区域积分变成闭合边界上的线积分，如式（6.26）。

\begin{quote}\small 图 6.14：由上方的 $x$、$y$ 梯度场生成下方的标记函数。\end{quote}
\subsection{由前沿位置构造体积分数}

另一种选择是直接计算前沿切过单元后形成的参考相多边形面积。二维中，若每个单元只被界面穿过一次，这一过程较直接；多次穿越、非结构前沿和三维情形则需要可靠排序所有交点和单元边界片段。

对闭合区域 $\Omega$，选取满足 $\nabla\cdot\bm A=1$ 的向量场，由散度定理
\begin{equation}
\operatorname{Area}(\Omega)=\int_\Omega dv
=\oint_{\partial\Omega}\bm A\cdot\bm n\,ds.
\tag{6.26}
\end{equation}
取 $\bm A=(x,0)$、$(0,y)$ 或 $(x,y)/2$，并用
$\bm n=(-\partial_s y,\partial_s x)$，得到等价形式
\begin{equation}
A=\oint y\frac{dx}{ds}\,ds
=-\oint x\frac{dy}{ds}\,ds
=\frac12\oint\left(y\frac{dx}{ds}-x\frac{dy}{ds}\right)ds.
\tag{6.27}
\end{equation}
% [原PDF p168]
有符号面积公式要求多边形顶点按一致方向排列；顺、逆时针会改变符号。包含多个不相连参考相区域或孔洞时，应分别按外环和内环方向求和。多次交叉问题的核心正是建立这些有向闭合环，而非式（6.28）本身。

可选择 $\bm A=(x,0)$、$(0,y)$ 或 $(x,y)/2$。配合法向 $\bm n=(-\partial_sy,\partial_sx)$，得到式（6.27）的三种等价面积线积分。若边界由直线元连接，中点规则对每段积分是精确的，最终成为式（6.28）的顶点求和：第一式累加每段下方的有符号面积，第二式累加左侧面积，第三式取二者平均。这些公式与鞋带公式等价。真正困难不在积分，而是识别前沿何时进入、离开单元，并决定哪些单元边界段应补入闭合多边形；图 6.15 显示多次交叉时排序很容易变复杂。

\begin{quote}\small 图 6.15：切割单元内的标记函数构造；已知前沿位置后可按闭合多边形面积求单元体积分数。\end{quote}
若边界由直线元连接，积分可精确写成
\begin{align}
A&=\frac12\sum_i(y_{i+1}+y_i)(x_{i+1}-x_i)\notag\\
&=-\frac12\sum_i(x_{i+1}+x_i)(y_{i+1}-y_i)\notag\\
&=\frac14\sum_i\left[(y_{i+1}+y_i)(x_{i+1}-x_i)
-(x_{i+1}+x_i)(y_{i+1}-y_i)\right].
\tag{6.28}
\end{align}
第一式是各线元下方面积之和，第二式是左侧面积之和，第三式取两者平均。

% [原PDF p169]
对三维三角前沿，与六个单元面的精确求交和排序更复杂。一种近似是把每个三角形细分成许多小片，把每个小片完全归入其形心所在的单元；子片越小，界面单元及其“下游”单元的累积体积误差越小。得到体积分数后，密度等物性通常作线性插值，黏度比很大时可采用第三章讨论的调和平均。随后考虑拓扑：显式连接的前沿不会因距离小于一格而自动合并，这既增加代码复杂度，也避免把网格分辨率误当成物理破裂准则。真实聚并常由中间薄膜排液控制。

算法难点不在面积公式，而在确定前沿以何种顺序进出单元，并补齐属于闭合多边形的单元边界段。三维还要处理三角面片与六个单元面的相交。一个简单近似是把每个三角面细分成许多小片，并把每个小片归入其形心所在单元；细分越密，局部及其下游单元的体积误差越小。构造出体积分数后，密度等通常按线性关系设置，黏度跨大比值变化时可使用调和平均。

\section{前沿拓扑的改变}

显式连接的前沿不会自动合并或断裂，必须主动修改拓扑。这增加了代码复杂度，却也避免把“相距不到一个网格”误当成普适物理准则。界面捕捉方法通常在两界面靠近到约一个网格时自动连接，而真实液滴聚并受中间薄膜排液速率和短程作用力控制，结果可能完全不同。

拓扑事件可分为两类：
\begin{itemize}
\item \textbf{液膜破裂}：液滴接近液滴或固壁，中间液体被挤出；膜薄到数百埃量级后，未包含在普通连续介质方程中的短程吸引力可触发破裂。
\item \textbf{液丝夹断}：细长流体柱受 Rayleigh 不稳定性发展，局部半径持续减小，表面张力最终把它夹成两段。液丝可能在有限时间趋于零半径，其演化比薄膜破裂更接近 Navier--Stokes 连续模型本身。
\end{itemize}

% [原PDF p170]
拓扑改变可分两类。液膜破裂发生在液滴接近另一滴或固壁时，中间流体被挤出，膜薄至数百埃后由连续模型通常未包含的短程吸引力触发。液丝断裂则由 Rayleigh 不稳定性使细圆柱局部半径不断减小，表面张力最终夹断；它可能在有限时间形成奇异性，更接近 Navier--Stokes 自身描述。程序要先检测候选近邻，再执行连接修改。暴力比较 $N$ 个前沿单元形心是 $O(N^2)$；把空间分盒，只比较同盒和邻盒对象可大幅降本。也可重建 $I$ 后再插值回前沿，若值显著偏离 0.5，就表明相反法向的另一段前沿很近。

\begin{quote}\small 图 6.16：两个液滴的偏心碰撞。拓扑改变算法允许液滴先聚并、随后再次破碎；时间从左向右，第一帧箭头为初速度。图由 S. Mortazavi 博士提供。\end{quote}
程序首先要找到可能发生事件的前沿区域，再执行重新连接。最朴素的检测计算任意两前沿单元形心距离，$N$ 个单元代价为 $O(N^2)$；把空间划分成小盒，只比较同盒及邻盒中的单元，可显著降本。另一种检测利用近距离平行界面的标记函数梯度相互抵消：先由前沿重建网格标记函数，再把它插值回前沿；若在 $I=0$ 与 $I=1$ 的两相体系中，界面点得到的值显著偏离 0.5，说明附近可能存在另一段前沿。

识别近邻后，还要决定是否真正改变拓扑。由于薄膜厚度或液丝直径通常已经低于网格分辨率，除非使用亚网格排液与破裂模型，否则数值解无法可靠给出物理临界条件。许多前沿追踪计算因而沿用“距离小于一个网格”作为操作准则；显式前沿的优势是可以系统改变该阈值，检验宏观结果对拓扑模型的敏感性。

% [原PDF p171]
检测到近邻以后，是否真正改变拓扑仍需物理判据；但薄膜厚度或液丝直径此时往往已低于网格，除非有亚网格排液/破裂模型，否则无法准确判断。多数计算沿用“距离小于一个网格”触发，显式前沿至少允许改变阈值并做敏感性分析。二维点链重连较容易，三维三角网格困难得多。Nobari--Tryggvason 实现了碰撞液滴聚并但不能一般夹断；Du 等人在 FronTier 中发展更通用算法；Shin--Juric 则周期性从网格标记函数等值面重新生成前沿。第六章注释还给出本质二维的径向 blob 核式（6.29）、（6.30），以及由前沿构造体积分数的后续工作。

二维点链的重新连接较容易，三维三角网格则复杂得多。已有算法可处理两滴碰撞聚并，也可在通用前沿库中完成合并和夹断。另一种策略是不直接修改三角连接，而周期性地从网格标记函数的某个等值面重新生成前沿；这简化拓扑处理，但会引入重建误差。三维碰撞模拟已能表现液滴先聚并、再拉伸破碎的全过程，不过通用、简单而高效的拓扑重构仍是前沿追踪中的开放问题。

\section{补充说明}

插值核不一定要写成一维函数的张量积，也可以采用本质上多维的径向“团块”函数。涡方法中常见的两种二维正则化核为
\begin{align}
\phi(r)&=\frac1\pi(2-r^2)e^{-r^2},
\tag{6.29}\\
\phi(r)&=\frac1{6\pi}(24-36r^2+12r^4-r^6)e^{-r^2}.
\tag{6.30}
\end{align}
它们直接按径向距离分布奇异量，旋转对称性优于张量积核。由前沿构造体积分数、以及 VOF--前沿追踪混合算法仍在持续发展；选择实现时应同时考察几何精度、守恒、有界性、并行局部性和三维拓扑处理成本。
% ===== 译文分片 3（已合并）=====

\chapter{表面张力}

% [原PDF p172]
本章所谓“准确”，不仅指某一时刻得到正确的 Laplace 压差，还要求界面长期输运后仍能保持压力与毛细力平衡。在单流体方程中，原来只作用于零厚度界面的力必须转换成固定网格可接受的体积力；因此，界面表示、曲率估计和离散 Delta 函数三者不能彼此独立选择。VOF 与水平集法从标记函数出发，前沿追踪法则先在离散界面上求力；两条路线最终都要把力送到速度所在的网格位置。

对于需要长时间追踪互不相溶流体界面的数值方法，准确计算表面张力几乎是最关键的环节之一。在单流体表述中，运动方程在固定网格上求解，表面张力作为集中于界面附近的体积力加入离散 Navier--Stokes 方程。界面若由 VOF 颜色函数或水平集函数直接输运，表面力通常由标记函数构造；界面若由离散标记点显式追踪，则可以先在移动前沿上计算力，再传递到固定网格。本章依次说明两类处理，并以静止液滴、寄生流和毛细振荡检验其精度与稳健性。

\section{由标记函数计算表面张力}

颜色函数 $C$ 或水平集函数 $F$ 用来标记界面位置。由这类函数构造表面张力的标准方法称为连续表面力法（continuous surface force, CSF）；若将表面张力写成应力散度并直接离散，则得到严格守恒动量的连续表面应力法（continuous surface stress, CSS）。

\subsection{连续表面力法}

% [原PDF p173]
下标 $h$ 表示有限差分或其他数值近似。将 $\delta_S$ 换成 $|\nabla_hC|$ 后，曲率 $\kappa_h$ 和法向 $\boldsymbol n_h$ 仍须由同一离散界面求得。对静止圆滴，连续平衡式为 $-\nabla p+\boldsymbol f_\sigma\delta_S=0$。若人为给定精确常曲率 $\kappa_0$，则两个坐标方向上的中心差分都可逐格积分，说明未平滑 CSF 至少在这一特殊情形能够构造严格的离散平衡。
先考虑表面张力系数 $\sigma$ 为常数的情形。加入动量方程的奇异表面力为
\begin{equation}
  \boldsymbol f_\sigma\,\delta_S
  =\sigma\kappa\boldsymbol n\,\delta_S,
  \tag{7.1}
\end{equation}
其中 $\kappa$ 为曲率，$\boldsymbol n$ 为单位法向，$\delta_S$ 为支承在界面 $S$ 上的表面 Delta 分布。因为颜色函数近似 Heaviside 函数 $H$，且
\begin{equation}
  \nabla H=-\delta_S\boldsymbol n,
  \tag{7.2}
\end{equation}
CSF 用 $|\nabla_h C|$ 近似 $\delta_S$，在网格节点 $(i,j)$ 上写成
\begin{equation}
  \boldsymbol f_{i,j}=\sigma\kappa_h|\nabla_h C|\boldsymbol n_h,
  \qquad
  \boldsymbol n_h=-\frac{\nabla_hC}{|\nabla_hC|},
  \tag{7.3--7.4}
\end{equation}
从而得到未平滑形式
\begin{equation}
  \boldsymbol f_{i,j}=-\sigma\kappa_h\nabla_hC.
  \tag{7.5}
\end{equation}

% [原PDF p174]
式（7.10）中 $p_2$ 是 $C=0$ 相的压力；液滴内部 $C=1$，外法向约定使内部参考相的曲率为负，故所得压力跃迁与 Laplace 定律一致。三维推导完全相同。若以密度代替颜色函数，$\rho=\rho_2-[\rho]_SH$，从而 $\nabla\rho=-[\rho]_S\nabla H$；这种写法可直接利用流场中已有的密度数组。MAC 网格实现先在各面构造法向分量，再用四个相邻面的平均补齐切向分量。
这个看似朴素的离散有一个重要优点：对曲率为常数 $\kappa_0$ 的球形液滴或圆柱，在无重力静止状态下，压力梯度与表面张力可在离散层面精确平衡。离散平衡式
\begin{equation}
  -\nabla_hp-\sigma\kappa_0\nabla_hC=0
  \tag{7.7}
\end{equation}
可直接积分为
\begin{equation}
  p_{i,j}=p_2-\sigma\kappa_0 C_{i,j},
  \tag{7.10}
\end{equation}
这正是 Laplace 压差的网格表示。由曲率--散度关系 $\kappa=-\nabla\cdot\boldsymbol n$，常用实现为
\begin{equation}
  \boldsymbol f_{i,j}=\sigma(\nabla_h\cdot\boldsymbol n_h)\nabla_hC,
  \qquad
  \boldsymbol n_h=-\frac{\nabla_h I}{|\nabla_h I|},
  \tag{7.11--7.12}
\end{equation}
其中 $I$ 可以是任意界面指示函数。也可直接用密度标记两相。若
$\rho=\rho_1H+\rho_2(1-H)$，且 $[\rho]_S=\rho_2-\rho_1$，则
\begin{equation}
  \boldsymbol f_{i,j}
  =-\sigma\frac{\nabla_h\rho}{[\rho]_S}
  (\nabla_h\cdot\boldsymbol n_h).
  \tag{7.15}
\end{equation}

% [原PDF p175]
非单位法向的模在远离界面处严格为零，所以式（7.18）不能在那里直接使用。加入 $\varepsilon_{\rm tiny}$ 只是编程层面的除零保护；更好的实现应先判定梯度是否有意义。曲率在单元中心由面法向的散度得到，再平均回相应速度面。式（7.23）和（7.24）分别给出 $x$、$y$ 动量方程所需的面力，保证颜色函数差分、曲率平均及速度布置彼此对齐。
在 MAC 交错网格上，先于单元面计算非单位法向。例如
\begin{equation}
 m_{x,i+1/2,j}=-\frac{C_{i+1,j}-C_{i,j}}{\Delta x},
 \qquad
 m_{y,i,j+1/2}=-\frac{C_{i,j+1}-C_{i,j}}{\Delta y},
 \tag{7.16}
\end{equation}
缺失的切向分量由相邻面平均得到，再归一化为
\begin{equation}
  \boldsymbol n_{i+1/2,j}
  =\frac{\boldsymbol m_{i+1/2,j}}
  {|\boldsymbol m_{i+1/2,j}|+\varepsilon_{\rm tiny}}.
  \tag{7.21}
\end{equation}
微小量只用于避免远离界面处的零除；更稳妥的程序应显式区分 $\nabla C\ne0$ 与 $\nabla C=0$ 的单元。曲率取单位法向的离散散度：
\begin{equation}
 \kappa_{i,j}=-\left[
 \frac{n_{x,i+1/2,j}-n_{x,i-1/2,j}}{\Delta x}
 +\frac{n_{y,i,j+1/2}-n_{y,i,j-1/2}}{\Delta y}
 \right].
 \tag{7.22}
\end{equation}
随后将曲率平均到相应速度面，得到
\begin{align}
 f_{x,i+1/2,j}&=-\frac{\sigma}{2\Delta x}
 (C_{i+1,j}-C_{i,j})(\kappa_{i+1,j}+\kappa_{i,j}),\\
 f_{y,i,j+1/2}&=-\frac{\sigma}{2\Delta y}
 (C_{i,j+1}-C_{i,j})(\kappa_{i,j+1}+\kappa_{i,j}).
 \tag{7.23--7.24}
\end{align}

% [原PDF p176]
具体地，角点梯度为
\begin{align}
(C_x)_{i+1/2,j+1/2}&=\frac{C_{i+1,j}-C_{i,j}+C_{i+1,j+1}-C_{i,j+1}}{2\Delta x},\\
(C_y)_{i+1/2,j+1/2}&=\frac{C_{i,j+1}-C_{i,j}+C_{i+1,j+1}-C_{i+1,j}}{2\Delta y}.
\tag{7.27--7.28}
\end{align}
其余角点循环平移即可。$T_{xx}$ 与 $T_{yy}$ 由围绕单元中心的四个角点平均，$T_{xy}$ 不作此平均，因为 $x$ 面上 $\partial_yT_{xy}$ 与 $y$ 面上 $\partial_xT_{xy}$ 都正好需要角点值。该布置不是任意插值：它使每个控制体面的应力只计算一次，相邻单元使用相反符号，从而把“连续表面应力”的守恒性质保留到 MAC 离散层面。与 CSF 相比，CSS 无需先从噪声法向再求一次散度曲率，但静态压力场不再能写成 $p_2-\sigma\kappa_0C$ 的简单格点解。
具体地，角点梯度为
\begin{align}
(C_x)_{i+1/2,j+1/2}&=\frac{C_{i+1,j}-C_{i,j}+C_{i+1,j+1}-C_{i,j+1}}{2\Delta x},\\
(C_y)_{i+1/2,j+1/2}&=\frac{C_{i,j+1}-C_{i,j}+C_{i+1,j+1}-C_{i+1,j}}{2\Delta y}.
\tag{7.27--7.28}
\end{align}
其余角点循环平移即可。$T_{xx}$ 与 $T_{yy}$ 由围绕单元中心的四个角点平均，$T_{xy}$ 不作此平均，因为 $x$ 面上 $\partial_yT_{xy}$ 与 $y$ 面上 $\partial_xT_{xy}$ 都正好需要角点值。该布置不是任意插值：它使每个控制体面的应力只计算一次，相邻单元使用相反符号，从而把“连续表面应力”的守恒性质保留到 MAC 离散层面。与 CSF 相比，CSS 无需先从噪声法向再求一次散度曲率，但静态压力场不再能写成 $p_2-\sigma\kappa_0C$ 的简单格点解。
CSS 的详细离散把 $C$ 的梯度放在单元角点。在角点形成 $\boldsymbol T^S_\sigma=\sigma(\boldsymbol I-\boldsymbol n\boldsymbol n)|\nabla_hC|$ 后，将 $T_{xx}$、$T_{yy}$ 以四角平均移到单元中心，而 $T_{xy}$ 保留在角点。随后在速度面上取张量散度：
\begin{align}
f_{x,i+1/2,j}&=\frac{T_{xx,i+1,j}-T_{xx,i,j}}{\Delta x}
+\frac{T_{xy,i+1/2,j+1/2}-T_{xy,i+1/2,j-1/2}}{\Delta y},\\
f_{y,i,j+1/2}&=\frac{T_{yy,i,j+1}-T_{yy,i,j}}{\Delta y}
+\frac{T_{xy,i+1/2,j+1/2}-T_{xy,i-1/2,j+1/2}}{\Delta x}.
\tag{7.32--7.33}
\end{align}
共享面的应力通量成对抵消，因此该形式具有离散守恒性。
\subsection{连续表面应力法}

% [原PDF p177]
直接把 CSF、CSS 或 PROST 的力分量加入预测速度，是最自然的实现；但基础 CSF 与 CSS 的测试结果并不理想。五点滤波给中心值权重 $1/2$，给四个最近邻各 $1/8$，仅把界面扩宽少数网格。通常有两种次序：先由原始界面算力再分配，或先平滑界面再算力。两者都等价于把奇异力铺展在界面邻带内，而 PROST 与 PCMM 不需要这种处理。
CSS 从表面应力形式出发：
\begin{equation}
  \boldsymbol f_\sigma\delta_S
  =\nabla\cdot\boldsymbol T^S_\sigma\delta_S
  =\nabla\cdot\bigl[\sigma(\boldsymbol I-\boldsymbol n\boldsymbol n)\delta_S\bigr].
  \tag{7.25}
\end{equation}
即使 $\sigma$ 随界面位置变化，该式仍然成立。无平滑的离散形式为
\begin{equation}
 \boldsymbol f_{i,j}=\nabla_h\cdot
 \left[\sigma(\boldsymbol I-\boldsymbol n_h\boldsymbol n_h)|\nabla_hC|\right].
 \tag{7.26}
\end{equation}
它避免显式计算曲率，并以应力散度形式保证动量守恒；代价是一般不存在类似式（7.10）的简单静止离散解。

% [原PDF p178]
图 7.1 表示半径为 $\varepsilon$ 的核支承域：网格点 $(i,j)$ 接收圆盘 $\Omega_\varepsilon$ 内各颜色函数值的加权贡献。离散卷积写成
\begin{equation}
\widetilde C_{ij}=A(\varepsilon)\sum_m\sum_l C_{lm}
\left[1-\frac{x_{il}^2+y_{jm}^2}{\varepsilon^2}\right]^4h^2,
\tag{7.37}
\end{equation}
其中只对支承域内的点求和，归一化常数使同一括号权重乘 $h^2$ 的总和为一。$\widetilde C$ 的梯度既可用有限差分，也可直接对核求导；后者给出的每个分量含因子 $-8A(\varepsilon)h^2$ 和核的三次幂。
具体实现中，$C$ 的梯度先在单元角点计算，例如
\begin{align}
 (\partial_xC)_{i+1/2,j+1/2}
 &=\frac{C_{i+1,j}-C_{i,j}+C_{i+1,j+1}-C_{i,j+1}}{2\Delta x},\\
 (\partial_yC)_{i+1/2,j+1/2}
 &=\frac{C_{i,j+1}-C_{i,j}+C_{i+1,j+1}-C_{i+1,j}}{2\Delta y}.
 \tag{7.27--7.28}
\end{align}
在角点构造
\begin{equation}
 \boldsymbol T^S_{\sigma,i+1/2,j+1/2}
 =\sigma(\boldsymbol I-\boldsymbol n\boldsymbol n)|\nabla_hC|,
 \tag{7.29}
\end{equation}
将对角分量平均至单元中心、非对角分量保留在角点，最后在速度面取离散散度。这样相邻控制体共享完全相反的面力，故全局动量严格守恒。

% [原PDF p179]
柱坐标轴对称界面的两条主曲率分别来自子午曲线和绕轴圆周。若 $\theta$ 是法向与 $z$ 轴夹角，原书写成
\begin{equation}
\sigma\kappa\boldsymbol n\delta_S
=\sigma\left(\kappa_1+\frac{\cos\theta}{r}\right)
\boldsymbol n\delta_S,
\tag{7.44}
\end{equation}
也可用子午面散度与径向法向分量表示。$r\to0$ 时环向项表面上含 $1/r$，但光滑轴对称界面的正则条件使分子同步趋零；数值实现必须用轴线对称极限，而不能在轴格点直接除以零。平面二维方法可负责子午曲率和 Delta 带分配，附加环向曲率必须用同一法向约定，否则球滴在轴线附近会产生错误符号或不连续压力跳跃。
柱坐标轴对称界面的两条主曲率分别来自子午曲线和绕轴圆周。若 $\theta$ 是法向与 $z$ 轴夹角，原书写成
\begin{equation}
\sigma\kappa\boldsymbol n\delta_S
=\sigma\left(\kappa_1+\frac{\cos\theta}{r}\right)
\boldsymbol n\delta_S,
\tag{7.44}
\end{equation}
也可用子午面散度与径向法向分量表示。$r\to0$ 时环向项表面上含 $1/r$，但光滑轴对称界面的正则条件使分子同步趋零；数值实现必须用轴线对称极限，而不能在轴格点直接除以零。平面二维方法可负责子午曲率和 Delta 带分配，附加环向曲率必须用同一法向约定，否则球滴在轴线附近会产生错误符号或不连续压力跳跃。
核导数可无歧义地给出平滑法向，但用核的二阶导数直接估计曲率过于不精确，所以仍采用法向散度。由于差分曲率只在比核带窄约两个网格的区域可靠，最终的近似 Delta 函数取 $6\widetilde C(1-\widetilde C)|\nabla\widetilde C|$，峰值位于带中心；若 $\widetilde C$ 沿法向单调，其法向积分等于 $\int_0^1 6C(1-C)\,dC=1$。轴对称时还必须加入绕轴方向的主曲率，而不能直接套用平面二维曲率。
\subsection{直接施加与简单平滑}

% [原PDF p180]
显式前沿中真正需要的是界面小片上的合力，而不是每个标记点的曲率值。若先在不均匀标记点上求点曲率再乘归属面积，既难保证精度，也不能保证闭合界面的总毛细力为零。把面积分转为片段边界上的积分可同时解决这两个问题：平面片各边的切向拉力恰好相消，曲面片的各边拉力因方向不同而留下法向合力。二维与三维的物理思想相同，但几何公式与数据结构不同。
CSF、CSS 以及后述 PROST 都可在预测步把表面力直接加到速度节点。对基础 CSF/CSS，这样做往往产生较大误差。一个简单改进是对颜色函数作五点平滑：
\begin{equation}
 \widetilde C_{i,j}=\frac12C_{i,j}
 +\frac18(C_{i,j-1}+C_{i,j+1}+C_{i-1,j}+C_{i+1,j}).
 \tag{7.34}
\end{equation}
更一般地，可以先由原始界面数据求力再把力分配到若干网格层，也可以先平滑界面数据再求力；两种做法都会把奇异表面力扩散到数个网格宽的带状区域。PROST 与压力修正标记法（PCMM）则不依赖这种平滑。

% [原PDF p181]
对变表面张力，式（7.48）中的第一项 $\sigma\kappa\boldsymbol n$ 沿法向，第二项 $(\partial\sigma/\partial s)\boldsymbol t$ 沿切向；后一项就是温度或浓度梯度引起的 Marangoni 剪切。用端点的 $\sigma\boldsymbol t$ 作差，可在不单独离散表面梯度的情况下同时保留二者。若前沿点不均匀，节点 $e$ 的归属段由相邻单元各半组成，端点不是原标记点而是两条单元的中点。最低阶切向对折线精确，但曲率变化快时应采用第 6 章更宽邻域的插值切向。无论阶次如何，只要共享边界使用同一个切向值，常 $\sigma$ 闭合曲线的所有节点力之和都严格为零。
对变表面张力，式（7.48）中的第一项 $\sigma\kappa\boldsymbol n$ 沿法向，第二项 $(\partial\sigma/\partial s)\boldsymbol t$ 沿切向；后一项就是温度或浓度梯度引起的 Marangoni 剪切。用端点的 $\sigma\boldsymbol t$ 作差，可在不单独离散表面梯度的情况下同时保留二者。若前沿点不均匀，节点 $e$ 的归属段由相邻单元各半组成，端点不是原标记点而是两条单元的中点。最低阶切向对折线精确，但曲率变化快时应采用第 6 章更宽邻域的插值切向。无论阶次如何，只要共享边界使用同一个切向值，常 $\sigma$ 闭合曲线的所有节点力之和都严格为零。
图 7.2 中，属于标记点 $e$ 的曲线段由相邻两条前沿单元各取一半组成，端点记为 1、2。由 $\kappa\boldsymbol n=\partial\boldsymbol t/\partial s$，段力只取决于两端切向。若 $\sigma$ 变化，$\partial_s(\sigma\boldsymbol t)$ 自动分解为法向毛细项和切向 Marangoni 项。实现时可在单元中点求切向并把合力存到节点，也可在节点求切向并把合力存到单元质心；最低阶单位切向就是端点坐标差除以线段长度。
\subsection{核平滑加权分布}

% [原PDF p182]
遍历无结构曲线网格时，可对每条单元把正切向加到起点、负切向加到终点；闭合曲线的和以望远镜形式消失。更高阶精度可用第 6 章的局部拟合式求切向。三维中，图 7.3 的曲面片 $\Delta A$ 由闭合轮廓 $C$ 围成；令 $\boldsymbol p=\boldsymbol t\times\boldsymbol n$，Stokes 定理把曲率向量的面积分化成 $\oint_C\boldsymbol p\,dl$。$\sigma\boldsymbol p$ 正是曲面片边缘单位长度上的拉力。
若以卷积平滑 Heaviside 函数，
\begin{equation}
 \widetilde H(\boldsymbol x)
 =(H*K)(\boldsymbol x)
 =\int_VH(\boldsymbol x')K(\boldsymbol x-\boldsymbol x';\varepsilon)\,d\boldsymbol x',
 \tag{7.35}
\end{equation}
其中核宽为 $\varepsilon$，且 $\varepsilon\to0$ 时 $K\to\delta_S$。卷积与微分可交换，因此“平滑阶跃函数”和“用权函数把 Delta 分布传到网格”是同一件事。书中检验采用紧支承核
\begin{equation}
 K_8(\boldsymbol x)=A(\varepsilon)
 \left(1-\frac{|\boldsymbol x|^2}{\varepsilon^2}\right)^4,
 \quad |\boldsymbol x|<\varepsilon,
 \qquad K_8=0\quad |\boldsymbol x|\ge\varepsilon,
 \tag{7.36}
\end{equation}
归一化常数 $A$ 使离散权重之和为 $1$。梯度既可对 $\widetilde C$ 作有限差分，也可直接微分核；实践表明法向可由核导数稳定获得，但曲率若直接高阶微分核会过于敏感，仍宜用
\begin{equation}
 \kappa_{i,j}=-\left[
 \frac{\widetilde n_{x,i+1,j}-\widetilde n_{x,i-1,j}}{2h}
 +\frac{\widetilde n_{y,i,j+1}-\widetilde n_{y,i,j-1}}{2h}
 \right].
 \tag{7.41}
\end{equation}
由于可靠曲率所在带宽小于 $\widetilde C$ 的平滑带，最终 Delta 近似取成更尖锐的
\begin{equation}
 \widetilde\delta_{S,i,j}=6\widetilde C_{i,j}(1-\widetilde C_{i,j})
 |\nabla\widetilde C_{i,j}|.
 \tag{7.42}
\end{equation}
当 $\widetilde C$ 沿法向单调时，其法向积分恰为 $1$。

% [原PDF p183]
当 $\sigma$ 沿界面变化时，三维合力直接写成 $\oint_C\sigma\boldsymbol p\,dl$。书中把一个顶点的归属曲面片定义为所有相邻三角形面积的三分之一。图 7.4 左图用粗线围出顶点 1 的归属区域；右图放大一个三角形 $e$。三角形质心是三个顶点坐标的算术平均，各边中点是相应两端点的平均。连接质心与边中点后，每个三角形自然分成三个顶点子区。
\subsection{轴对称界面}

% [原PDF p184]
三角形单位法向由两条边的叉积归一化得到。顶点到相邻边中点的边界拉力与邻接三角形相消，因此对顶点 1 只需积分从边中点 $x_{12}$ 到质心、再从质心到 $x_{31}$ 的两段。化简后，三角形对顶点 1 的贡献等于对边 $23$ 拉力的一半；对其余顶点循环置换即得式（7.59）。一个节点可连接任意数目的三角形，节点总力就是所有相邻面元贡献之和。高阶方法必须扩大邻域并拟合光滑曲面。
在柱坐标 $(r,\phi,z)$ 的轴对称问题中，总曲率不仅含子午面曲率，还含环向曲率。若 $\theta$ 为法向与 $z$ 轴夹角，则
\begin{equation}
 \sigma\kappa\boldsymbol n\delta_S
 =\sigma\left(\nabla_2\cdot\boldsymbol n+\frac{n_r}{r}\right)
 \boldsymbol n\delta_S,
 \qquad \nabla_2=(\partial_r,\partial_z).
 \tag{7.45}
\end{equation}
因此二维算法可用于子午面，但必须补上环向曲率项，并在轴线附近采用与正则性相容的离散。

\section{显式追踪前沿的表面张力}

% [原PDF p185]
图 7.5 展示移动标记点附近被权函数覆盖的附加固定网格点。二维分配中的 $\Delta s_l/h^2$、三维中的 $\Delta A_l/h^3$ 保证网格体积积分等于前沿总力。交错网格上，每个分量只送到对应速度分量所在的位置。编程时不必先把同一节点所有面元力合并；每个面元的三个节点贡献算出后即可立即传网格，不过这样会增加需要分配的量。直接平滑前沿力虽自然，却不是唯一选择。
当前沿由标记点和连接单元表示时，表面张力主要转化为离散几何量的计算。实际所需通常不是某点的曲率，而是一个小界面片段上的合力
\begin{equation}
 \delta\boldsymbol f_e=\int_{\Delta S}\sigma\kappa\boldsymbol n\,dS.
 \tag{7.46}
\end{equation}
“点曲率乘对应面积”的直接方法既受不均匀点距影响，也不能自动保证常表面张力闭合界面的合力为零。更好的办法是把面积分转成界面片段边界上的线积分：平面片段各边拉力相消，曲面片段则因边界拉力不能完全抵消而留下法向合力。

\subsection{二维界面}

% [原PDF p186]
这一混合构造把两类数据各自的优势保留下来：移动前沿提供较准确的几何曲率，固定网格指示函数则提供与压力差分同位置、同模板的法向 Delta 量。若二者在不同带宽上独立平滑而不再投影对齐，静止平衡的优势就会丧失。
前沿追踪也会在固定网格重建指示函数 $I$，因此可以把 $\sigma\kappa$ 与 $\boldsymbol n\delta_S$ 分开处理。用网格上的 $\nabla I$ 作为后者，可与离散压力梯度精确对齐；曲率则取自几何质量更高的前沿。Shin 等人的算法先分别把前沿合力和 $\boldsymbol n\delta_S$ 平滑成 $\boldsymbol f_{ij}$ 与 $\boldsymbol G_{ij}$，再由投影
\begin{equation}
(\sigma\kappa)_{ij}=\frac{\boldsymbol f_{ij}\cdot\boldsymbol G_{ij}}
{\boldsymbol G_{ij}\cdot\boldsymbol G_{ij}}
\tag{7.65}
\end{equation}
恢复标量毛细压力，最后乘网格梯度。
二维曲线满足 $\kappa\boldsymbol n=\partial\boldsymbol t/\partial s$，因此一段曲线上的力为
\begin{equation}
 \delta\boldsymbol f_e
 =\int_{\Delta s}\sigma\kappa\boldsymbol n\,ds
 =\sigma(\boldsymbol t_2-\boldsymbol t_1).
 \tag{7.47}
\end{equation}
无需显式求曲率，只需段两端的单位切向。表面张力变化时
\begin{equation}
 \frac{\partial(\sigma\boldsymbol t)}{\partial s}
 =\sigma\kappa\boldsymbol n+\frac{\partial\sigma}{\partial s}\boldsymbol t,
 \qquad
 \delta\boldsymbol f_e=(\sigma\boldsymbol t)_2-(\sigma\boldsymbol t)_1,
 \tag{7.48--7.49}
\end{equation}
法向毛细力与切向 Marangoni 力自然同时包含在内。线段中点的最低阶单位切向为
\begin{equation}
 \boldsymbol t=\frac{\boldsymbol x_2-\boldsymbol x_1}
 {|\boldsymbol x_2-\boldsymbol x_1|}.
 \tag{7.50}
\end{equation}
遍历无结构前沿单元时，把正切向加到起点、负切向加到终点，闭合曲线上的总和会望远镜式相消，因而严格保持零合力。

% [原PDF p187]
式（7.66）的分子是 $\sum_l(\sigma\kappa)_lw^l_{ij}(\Delta s_l/h^2)$，分母是相同几何权重之和；不能省略分母，否则标量值会随附近前沿点数改变。由合力恢复曲率时，$\boldsymbol f_e\cdot\boldsymbol f_e=(\sigma\kappa\Delta s)^2$ 只能给出绝对值。二维符号可由 $\boldsymbol f_p\times\boldsymbol t_p\cdot\boldsymbol k=f_xt_y-f_yt_x$ 判定，三维则须结合面法向取向。延拓过程只传播已有曲率值，不应跨过其他界面分支；窄间隙中若两个邻带相交，简单平均可能混合本应属于不同界面的曲率，这是选取带宽和权函数时必须防止的几何错误。
式（7.66）的分子是 $\sum_l(\sigma\kappa)_lw^l_{ij}(\Delta s_l/h^2)$，分母是相同几何权重之和；不能省略分母，否则标量值会随附近前沿点数改变。由合力恢复曲率时，$\boldsymbol f_e\cdot\boldsymbol f_e=(\sigma\kappa\Delta s)^2$ 只能给出绝对值。二维符号可由 $\boldsymbol f_p\times\boldsymbol t_p\cdot\boldsymbol k=f_xt_y-f_yt_x$ 判定，三维则须结合面法向取向。延拓过程只传播已有曲率值，不应跨过其他界面分支；窄间隙中若两个邻带相交，简单平均可能混合本应属于不同界面的曲率，这是选取带宽和权函数时必须防止的几何错误。
标量曲率不是守恒量，传递时应把“加权标量和”与“权重和”分开累积，最后相除；线段越长应具有越大权重。也可由界面小片合力的模恢复 $|\sigma\kappa\Delta s|$，再用 $\boldsymbol f\times\boldsymbol t$ 的取向确定符号。若曲率带比 $\nabla I$ 的非零带窄，可扩大权函数支承，也可反复把带内值平均到相邻空点，形成所谓 Laplacian 延拓，直至覆盖需要的网格层。
\subsection{三维界面}

三维曲面满足
\begin{equation}
 \kappa\boldsymbol n=(\boldsymbol n\times\nabla)\times\boldsymbol n,
 \tag{7.51}
\end{equation}
利用 Stokes 定理，有
\begin{equation}
 \delta\boldsymbol f_e
 =\sigma\int_{\Delta A}\kappa\boldsymbol n\,da
 =\sigma\oint_C\boldsymbol p\,dl,
 \qquad \boldsymbol p=\boldsymbol t\times\boldsymbol n.
 \tag{7.52}
\end{equation}
$\sigma\boldsymbol p$ 就是沿曲面片边缘的拉力。相邻面元共享边上的贡献大小相等、方向相反，故闭合界面总力为零；变表面张力时只需将积分改为 $\oint_C\sigma\boldsymbol p\,dl$。

% [原PDF p188]
五点平滑 CSS 的压力起伏相对温和，但图 7.6 仍能看到界面周围的尖峰；这些局部梯度在投影步中生成图 7.7 的闭合速度涡。前沿追踪若使用表 6.1 的光滑离散 Delta 权函数，寄生流通常更小，许多实算中可忽略；然而它们会限制可达密度比、黏度比和 Laplace 数。验证时应把液滴运行足够长，因为最大速度可能先振荡很久才达到统计平衡。还应同时报告压差误差和 $U_{\max}$：强黏性可以把速度压低，却不能证明曲率或毛细力正确；反过来，平均压差正确也可能掩盖界面周围正负压力尖峰。
五点平滑 CSS 的压力起伏相对温和，但图 7.6 仍能看到界面周围的尖峰；这些局部梯度在投影步中生成图 7.7 的闭合速度涡。前沿追踪若使用表 6.1 的光滑离散 Delta 权函数，寄生流通常更小，许多实算中可忽略；然而它们会限制可达密度比、黏度比和 Laplace 数。验证时应把液滴运行足够长，因为最大速度可能先振荡很久才达到统计平衡。还应同时报告压差误差和 $U_{\max}$：强黏性可以把速度压低，却不能证明曲率或毛细力正确；反过来，平均压差正确也可能掩盖界面周围正负压力尖峰。
表面张力没有均匀单相流中的对应项，所以必须单独设计验证。最基本的静态检验不仅要得到平均 Laplace 压差，还要观察界面附近的压力尖峰和速度环流。图 7.6 给出 CSS 无平滑时静止圆柱液滴的压力场，直径约占 28 个网格；图 7.7 给出相应寄生流，最大无量纲速度 $\mu U/\sigma$ 为 $O(10^{-3})$。寄生流可能经历很长瞬态，严重时足以把液滴撕裂。
对三角面元 $e=(\boldsymbol x_1,\boldsymbol x_2,\boldsymbol x_3)$，质心和边中点为
\begin{equation}
 \boldsymbol x_c=\frac{\boldsymbol x_1+\boldsymbol x_2+\boldsymbol x_3}{3},
 \qquad
 \boldsymbol x_{ij}=\frac{\boldsymbol x_i+\boldsymbol x_j}{2},
 \tag{7.54--7.55}
\end{equation}
面元法向为
\begin{equation}
 \boldsymbol n_e=\frac{\Delta\boldsymbol x_{12}\times\Delta\boldsymbol x_{13}}
 {|\Delta\boldsymbol x_{12}\times\Delta\boldsymbol x_{13}|}.
 \tag{7.56}
\end{equation}
将面元分成归属于三个顶点的子面积后，可把面元贡献直接写为
\begin{equation}
 \delta\boldsymbol f_{1e}=\frac12\boldsymbol n_e\times\Delta\boldsymbol x_{32},
 \quad
 \delta\boldsymbol f_{2e}=\frac12\boldsymbol n_e\times\Delta\boldsymbol x_{13},
 \quad
 \delta\boldsymbol f_{3e}=\frac12\boldsymbol n_e\times\Delta\boldsymbol x_{21}.
 \tag{7.59}
\end{equation}
顶点合力为所有相邻面元贡献之和。若要求更高阶，则需对更大的点邻域拟合曲面后再求切向与法向。

% [原PDF p189]
静止液滴由 Laplace 数 $\mathrm{La}=\rho_l d\sigma/\mu_l^2$、密度比和黏度比表征。表 7.1 按难度列出 F、A、B、C、D 五组算例；大 Laplace 数算例接近水中小气泡。度量取长时间极限最大速度形成的毛细数 $\mathrm{Ca}=\mu_lU/\sigma$。小 $\mathrm{La}$ 时寄生流平滑趋稳，大 $\mathrm{La}$ 时 $\mathrm{Ca}$ 强烈振荡，因此只用对称流体、低 Laplace 数作测试会显著低估算法困难。
\subsection{将前沿表面力平滑到固定网格}

% [原PDF p190]
表 7.1 中，低 $\mathrm{La}$ 时五点平滑 CSS、CSF 的寄生流均明显大于前沿追踪，CSS 略优于 CSF；提高 $\mathrm{La}$ 后，带五点平滑的液滴甚至解体，所以 C、D 两例报告的是未平滑 CSF。表 7.2 比较核 $K_8$ 的 CSF、PROST、PCMM 与高度函数：先进方法总体更稳定，HF 在充分松弛后可达到机器精度，而核平滑 CSF 的量级大体仍与前沿追踪相当。表中 ``NA'' 表示原研究未给出该组合。
在二维中，前沿片段 $l$ 的离散力 $\delta\boldsymbol f_l$ 传至固定网格点 $(i,j)$ 时写成
\begin{equation}
 \boldsymbol f_{i,j}=\sum_l\delta\boldsymbol f_l,w^l_{i,j}
 \frac{\Delta s_l}{h^2};
 \tag{7.60}
\end{equation}
三维则用面元面积并除以 $h^3$。因子 $\Delta s_l/h^d$ 保证总力守恒，$w^l_{i,j}$ 是紧支承插值权。交错网格上各力分量分配到对应速度节点。

% [原PDF p191]
动态检验先可比较上升气泡或下落液滴的稳态变形，但这种流动并不持续振荡界面。更苛刻的是小振幅毛细振荡。图 7.8 给出 $32^2$ 前沿追踪网格上初始半径扰动 $10\%$ 的液滴，在 $t\sigma/(d\mu)=9.0925$ 时的速度和压力。平面无黏波频率满足 $\omega_0^2=\sigma k^3/(\rho_1+\rho_2)$；黏性造成与 $\sqrt\nu$ 成比例、可能很显著的频移，故应使用黏性正规模或 Prosperetti 初值解比较。
另一种混合策略利用前沿提供高精度曲率，而法向从固定网格指示函数 $I$ 得到，以保持与压力梯度的一致离散：
\begin{equation}
 \boldsymbol f_{i,j}=(\sigma\kappa)_{i,j}\nabla I_{i,j}.
 \tag{7.61}
\end{equation}
先把前沿法向 Delta 量分配为
\begin{equation}
 \boldsymbol G_{i,j}=\sum_l\boldsymbol n_lw^l_{i,j}\frac{\Delta s_l}{h^2},
 \tag{7.62}
\end{equation}
并把前沿合力平滑为 $\boldsymbol f_{i,j}$，即可投影得到
\begin{equation}
 (\sigma\kappa)_{i,j}
 =\frac{\boldsymbol f_{i,j}\cdot\boldsymbol G_{i,j}}
 {\boldsymbol G_{i,j}\cdot\boldsymbol G_{i,j}}.
 \tag{7.65}
\end{equation}
也可把前沿上的标量曲率作归一化加权平均：
\begin{equation}
 (\sigma\kappa)_{i,j}
 =\frac{\sum_l(\sigma\kappa)_lw^l_{i,j}(\Delta s_l/h^2)}
 {\sum_lw^l_{i,j}(\Delta s_l/h^2)}.
 \tag{7.66}
\end{equation}
若由片段合力恢复曲率幅值，则
\begin{equation}
 (\sigma\kappa\Delta s)_e=\pm\sqrt{\boldsymbol f_e\cdot\boldsymbol f_e},
 \tag{7.69}
\end{equation}
符号必须由局部前沿取向判断。由于曲率标量的分配带常比 $\nabla I$ 的非零带窄，可通过重复邻点平均作“Laplacian 延拓”，把曲率向界面两侧扩展到所需带宽。

% [原PDF p192]
表 7.3 中的初始振幅 $a_0$、每波长网格数与 Laplace 数必须和引用研究逐项对应。若计算盒太浅，即使界面算法完美，波频也会因上下边界改变；若振幅太大，谐波和波峰--波谷不对称会污染线性误差；若初速为零却拿正规模比较，则误差来自初始条件而非表面张力离散。Prosperetti 解同时包含多个黏性衰减分量，表达式很长，原书因此指向文献而不在此抄录。实际数值曲线应同时比较振幅包络、零点时刻与相位，因为只拟合一个频率可能把衰减误差隐藏在相位补偿中。
表 7.3 中的初始振幅 $a_0$、每波长网格数与 Laplace 数必须和引用研究逐项对应。若计算盒太浅，即使界面算法完美，波频也会因上下边界改变；若振幅太大，谐波和波峰--波谷不对称会污染线性误差；若初速为零却拿正规模比较，则误差来自初始条件而非表面张力离散。Prosperetti 解同时包含多个黏性衰减分量，表达式很长，原书因此指向文献而不在此抄录。实际数值曲线应同时比较振幅包络、零点时刻与相位，因为只拟合一个频率可能把衰减误差隐藏在相位补偿中。
设界面为 $h(x,t)=a(t)\cos kx$。比较解析解时有三项容易混淆的误差：计算盒在界面上下的高度必须远大于波长，无黏有限高修正为 $\tanh(kh)-1$；$a/\lambda$ 不足够小时会出现非线性；正规模解要求界面、速度和压力都按同一模态初始化，若只给界面位移而令初速为零，则必须采用 Cortelezzi--Prosperetti 解。图 7.9 与表 7.3 在密度、黏度比均为一时比较各法的振幅与频率误差。CSS 误差降至约 $10^{-2}$ 后停滞，PCMM 先快后慢，$K_8$ 与 PROST 在细网格更好，HF 几乎各分辨率最佳。
\section{表面张力方法的检验}

% [原PDF p193]
表 7.3 以波长定义 Laplace 数，并列出每波长 8、16、32、64、128 点时的相对误差。HF 从粗网格 $0.16$ 降至 $5.5\times10^{-4}$；PROST 在 64 点时约 $1.8\times10^{-3}$。PCMM 直接由界面标记求控制体边界两端的切向拉力，不平滑地把水平投影加到 $u_{i-1/2,j}$。但界面处压力有有限跃迁，普通中心压力积分不够准确，必须按界面交点分割控制体边界。
\subsection{静止算例：寄生流}

最基本的验证是静止圆形或球形液滴：数值压力跳跃应满足 Laplace 定律，速度应保持为零。只检查压差仍不充分，因为曲率与压力梯度的微小不平衡会产生非物理环流，即寄生流（spurious/parasitic currents）。常用度量是由最大寄生速度 $U_{\max}$ 构造的毛细数
\begin{equation}
 \mathrm{Ca}_{\rm sp}=\frac{\mu U_{\max}}{\sigma},
\end{equation}
并考察它对网格分辨率、Laplace 数、密度比和黏度比的依赖。

基础 CSF/CSS 若不平滑，界面附近离散噪声可能迅速破坏液滴；五点平滑能改善稳定性，但误差仍较大。核平滑 CSF、PROST、PCMM 和高度函数（HF）通常明显更好。前沿追踪的误差与核平滑 CSF 大致同阶。对于良好平衡的 HF，静止液滴充分松弛后寄生流可降至机器精度。应注意：低寄生速度必须与正确 Laplace 压差同时满足，单独把速度耗散掉并不代表表面力离散正确。

\subsection{动态算例}

% [原PDF p194]
图 7.10 的交点 $E$ 若位于压力节点 $p_{i,j}$ 上方，右边界积分用 $DE\,p_{i,j}+EF\,p_{i,j+1}$；若位于下方，则用 $DE\,p_{i,j-1}+EF\,p_{i,j}$。该表达式替代通常的 $DF\,p_{i,j}$，把界面压力跳跃显式纳入。PCMM 已以二维前沿追踪和 VOF 重构实现，但尚无同样直接的三维实现；由于前沿追踪能较准确地求控制体边界切向，它与该框架更匹配。随后介绍的 PROST 则通过改善曲率来恢复压力--毛细力平衡。
上升气泡或下落液滴可检验表面张力、黏性应力和压力应力之间的动态平衡，但稳态变形对界面振荡的要求有限。更严格的检验是液滴、气泡或平面界面的毛细振荡。对无黏平面波，
\begin{equation}
 \omega_0^2=\frac{\sigma k^3}{\rho_1+\rho_2},
 \qquad k=\frac{2\pi}{\lambda},
\end{equation}
黏性会同时改变频率并衰减振幅，而且频移对 $\sqrt\nu$ 敏感。数值结果应与黏性正规模或由静止初速度出发的 Prosperetti 解比较。

\begin{sourceblock}
\sourcecrop{191}{78}{319}{420}{472}{0.46}
\captionof{figure}{振荡液滴在无量纲时刻 $t\sigma/(d\mu)=9.0925$ 的速度场（左）和压力场（右）。前沿追踪计算采用 $32^2$ 网格，液滴初始半径扰动为 $10\%$。图中可见界面恢复运动与跨界面压力跃迁。}
\end{sourceblock}

平面波检验还必须控制计算盒高度、初始条件和非线性振幅。盒高应远大于波长；若只初始化界面而令速度为零，不能直接拿正规模解作比较；初始振幅相对波长不够小时还会混入非线性误差。已有比较表明：CSS 在一定误差水平后收敛停滞，PCMM 先快后慢；核平滑 CSF 与 PROST 在细网格上更准确，HF 在粗网格到细网格范围内通常表现最好。

% [原PDF p195]
PROST 在界面单元周围 $3\times3$（三维为 $3\times3\times3$）邻域拟合二次零水平集。参数由标量 $a$、向量 $\boldsymbol b$ 和对称矩阵 $\boldsymbol C$ 组成；对每个试探参数，计算曲线在单元内切出的归一化面积 $v_{ij}(\boldsymbol q)$，最小化其与体积分数的平方差。标准最小二乘程序给出最优曲线，随后可解析求曲率。该法精确但每个界面邻域都要反复进行几何积分和非线性搜索，成本较高。高度函数以较便宜的一维积分替代这种拟合。
\section{更精细的表面张力方法}

\subsection{带压力修正的直接施力}

压力修正标记法（PCMM）先由显式界面计算穿过速度控制体边界的切向拉力，并把其水平或竖直投影直接施加到相应速度节点。由于压力在界面处存在有限跳跃，若仍以跨控制体两侧的节点压力作普通中心近似，会与直接表面力不匹配。PCMM 根据界面交点位于压力节点上方还是下方，对控制体边界压力积分分段。例如交点 $E$ 在节点 $p_{i,j}$ 上方时，右侧边界积分取
\begin{equation}
 \int_{DF}p\,ds=DE\,p_{i,j}+EF\,p_{i,j+1},
 \tag{7.72}
\end{equation}
若交点在下方则改用 $p_{i,j-1}$ 与 $p_{i,j}$。该法二维实现效果较好，但三维几何与压力修正显著复杂，更适合能准确获得控制体边界切向的前沿追踪框架。

\subsection{更精确曲率的 CSF：PROST}

% [原PDF p196]
图 7.11 中因 $|n_x|>|n_y|$，使用 $7\times3$ 模板求水平宽度 $x=g(y)$；若界面高度不落在目标单元内，就从相邻界面单元插值曲率。圆界面误差试验取 $R=0.15$、精确曲率约 $6.6$，在单位方形内依次使用 $32^2$ 至 $256^2$ 网格。HF 的 $L_\infty=\max|\widetilde\kappa_h-\kappa|$ 呈二阶收敛；当核宽固定为 $\varepsilon=2h$ 时，核卷积和直接有限差分误差反而近似按 $h^{-1}$ 增长。若固定物理核宽也能二阶，但参与卷积的格点数随加密急增。
PROST 的思想是：只要曲率足够准确，式（7.7）的压力--表面张力平衡就能显著抑制寄生流。它在界面单元附近的 $3\times3$ 区块内拟合二次曲线
\begin{equation}
 a+\boldsymbol b\cdot\boldsymbol x
 +\boldsymbol x\cdot\boldsymbol C\cdot\boldsymbol x=0,
 \tag{7.74}
\end{equation}
以试探曲线在各单元中围成的体积分数 $v_{ij}(\boldsymbol q)$ 与实际 $C_{ij}$ 的差为目标：
\begin{equation}
 E=\sum_{i,j}\left[v_{ij}(\boldsymbol q)-C_{ij}\right]^2.
 \tag{7.75}
\end{equation}
求得最小二乘曲线后即可解析计算曲率。三维使用 $3\times3\times3$ 邻域。PROST 精度高，但每个界面邻域都需非线性几何拟合，计算代价较大。

\subsection{由体积分数估计曲率：高度函数法}

% [原PDF p197]
图 7.12 以双对数坐标汇总三种曲率误差，并标出二阶参考斜率。实际 HF 实现若局部 VOF 数据过于不规则，会退回到穿过 VOF 重构面的最小二乘拟合。静止滴充分松弛后，HF 的寄生流可降至机器精度，毛细波结果也与表 7.3 一致。由此，数值方法已足以研究复杂多相流；后续改进仍会追求精度、效率和稳健性，但本书余下部分把重心转向用 DNS 探索具体流动物理。
高度函数法首先选择法向绝对值最大分量的坐标方向，再沿该方向对体积分数求和形成局部高度，最后对高度作中心差分。若 $|n_x|>|n_y|$，构造水平宽度 $x=g(y)$，则
\begin{equation}
 \kappa=\frac{g_{yy}}{(1+g_y^2)^{3/2}},
 \qquad
 g_y\approx\frac{x_{j+1}-x_{j-1}}{2h},
 \quad
 g_{yy}\approx\frac{x_{j+1}-2x_j+x_{j-1}}{h^2}.
 \tag{7.76}
\end{equation}
若高度函数不落在目标单元内，则由相邻界面单元的曲率插值更准确。高度是跨越 Heaviside 跳跃的一维积分，比不连续的体积分数平滑，因此曲率估计显著改善。圆界面试验显示 HF 的 $L_\infty$ 曲率误差达到二阶收敛；以固定网格宽度核平滑或直接差分法则可能不收敛。若核宽保持物理常数也能二阶收敛，但随加密涉及的网格点数迅速增加，代价高昂。

\section{数值方法小结}

直接数值模拟已经能够处理相当复杂的多相流，但表面张力仍是准确性和稳健性的核心限制。守恒、压力--毛细力平衡、曲率精度与界面拓扑处理必须共同考虑。简单 CSF 便于实现，CSS 保证动量守恒，显式前沿边界积分天然满足闭合界面零合力，而 PROST、PCMM 与高度函数则以不同方式提高静态平衡和动态振荡精度。以下各章转向这些数值方法在气泡群、雾化和撞击问题中的应用。

\chapter{弥散气泡流}

% [原PDF p198]
气泡流的预测直接关系到电站沸腾换热、冶金操作和化工气泡塔。气泡塔把气体从底部注入，气泡上升时向液体扩散并参与反应；工业容积常达数十至数百立方米，生物反应器因停留时间更长可达数千立方米。装置没有运动部件、结构简单，但流动随设备尺度显著改变，小型试验难以直接放大。工业计算因此只能求平均流场，并在混合物模型或双流体模型中用闭式关系表示未解析的气泡运动和相间力。

\section{引言}

% [原PDF p199]
多相 DNS 的目标与单相湍流 DNS 类似：在足够细的网格上解析所有相关非定常尺度，为平均方程提供漂移 Reynolds 数、速度脉动、气泡扩散以及相间作用的数据。更重要的是，它能揭示气泡主要以何种方式相互作用、是否存在优势微结构，以及是否形成远大于气泡直径的组织结构。早期计算多限于单泡、势流或点粒子；这些模型在稀薄或极限参数下有用，却不能替代中等 Reynolds 数、有限空隙率下的完整两相 Navier--Stokes 解。
气泡流广泛存在于电站沸腾换热、冶金过程和化工气泡塔中。气泡塔可用于乙烯部分氧化、异丁烯分离、重污染废水湿式氧化以及合成燃料生产：气体从底部注入，气泡上升时向液体传质并发生反应。工业装置的容积可从几十到数百立方米，生物反应器甚至达到数千立方米。它们没有运动部件、结构简单，适合大规模运行；但装置行为强烈依赖尺度，小型中试结果难以直接放大，因而数值预测十分重要。

工业尺度多相流只能依靠平均模型，从混合物模型到分别求解分散相和连续相的双流体模型不等。由于不解析每个气泡的瞬态运动，必须用闭式关系描述未解析运动以及气泡与液体之间的力。这与 Reynolds 平均单相湍流中对小尺度动量输运建模类似。直接数值模拟（DNS）则在足够细的网格上解析非定常 Navier--Stokes 方程，可为平均漂移速度、速度脉动、气泡扩散和相间作用提供基础数据，并揭示气泡是否形成特定微结构或远大于气泡尺度的组织结构。

% [原PDF p200]
除密度比、黏度比、Galileo 数和 Eötvös 数外，多泡体系必须给定空隙率，管道问题还需管流 Reynolds 数。化工文献有时以 Morton 数替代 Galileo 数，因为在给定流体和重力下前者为常数。低 Eötvös 数时气泡为球形；随 $N$ 与 $\mathrm{Eo}$ 增大，可经历椭球、裙边、摆动或螺旋轨迹，最终成为带湍流尾迹的球冠泡。摆动可能只是直线上升到螺旋上升的过渡。真实小气泡常受未知表面活性剂污染，本章则限定洁净界面。
早期研究往往限于单气泡、势流模型或点粒子模型；这些近似在特定极限下有用，但对中等 Reynolds 数、非稀疏的气泡群，需要求解完整非定常两相方程。本章讨论两个代表性问题：初始静止液体中均匀气泡群的浮升，以及竖直通道内浮升气泡与压差驱动流的耦合。

单个浮力气泡主要受四个无量纲数控制：密度比 $\rho_b/\rho_l$、黏度比 $\mu_b/\mu_l$，以及 Galileo 数 $N$ 与 Eötvös 数 $\mathrm{Eo}$。多气泡体系还需给定空隙率 $\alpha$；通道流还要给出通道 Reynolds 数等参数。低 $\mathrm{Eo}$ 对应强表面张力或小气泡，气泡近似球形。$\mathrm{Eo}$ 增大后，$N$ 决定气泡依次呈球形、椭球形、裙边形、摆动或螺旋上升，极高 $N$ 和 $\mathrm{Eo}$ 时则形成带湍流尾迹的球冠气泡。真实水--空气小气泡常受成分不明的表面活性物质污染，界面边界条件因而改变；这里仅讨论洁净气泡。

% [原PDF p201]
Ishii--Zuber 关联式汇集气泡、液滴和固体颗粒数据，通常预言漂移速度随空隙率单调降低。实验表明，即使精心选择到势流理论看似适用的参数，中等 Reynolds 数下的实际相互作用也常偏离势流。早期前沿追踪研究分别考察 $\mathrm{Re}_b\simeq1$--$2$ 与 $20$--$30$ 的二维气泡群，并做少量三维八泡计算；二维统计可收敛，但结果依赖周期域尺寸，且二维泡间作用过强，定量结论不能直接移植到三维。
\section{均匀气泡流}

均匀气泡流可在三向周期区域中模拟。一个周期单元只含一个气泡时，得到随整体上升但相对构型不变的规则阵列；增加每个周期中的气泡数后，规则阵列失稳，气泡不断近距离相互作用，但长时间平均可以达到统计稳态。计算域足够大时，每个周期内的统计量才可代表真正的均匀气泡流。

% [原PDF p202]
Bunner 与 Tryggvason 的并行三维计算在一个周期域内长期追踪超过 200 个气泡。图 8.1 左为 216 个近球形浮力泡的瞬时构型，右为 91 泡算例中局部流场；参数约为 $N=900$、$\mathrm{Eo}=1$、空隙率 $6\%$。最大算例使用 $256^3$ 均匀网格，每个泡直径约 20 点，较小系统的分辨率试验表明在该参数下已基本收敛。约十个气泡便能稳定平均上升速度和液相 Reynolds 应力，但气泡速度方差与弥散需要更大的系统。
实验关联式通常考察漂移速度随空隙率的变化。Ishii--Zuber 关联式汇集了多类气泡、液滴和固体颗粒数据，普遍预言空隙率增大时漂移速度单调下降。势流模型虽然能确认某些趋势，却通常不能描述实际中等 Reynolds 数气泡群的相互作用。前沿追踪 DNS 从二维的几十至数百个气泡逐步扩展到三维、超过两百个气泡的长期模拟。气泡平均上升 Reynolds 数定义为
\begin{equation}
 \mathrm{Re}_b=\frac{\rho_lU_bd_e}{\mu_l},
\end{equation}
其中 $U_b$ 是气泡相对液体的漂移速度，$d_e$ 是等效直径。

\begin{sourceblock}
\sourcecrop{202}{78}{461}{419}{639}{0.48}
\captionof{figure}{均匀气泡流 DNS。左：周期区域内 $216$ 个近球形浮力气泡的瞬时分布；右：另一组 $91$ 气泡计算中若干气泡周围的流线。该图说明统计均匀并不等于瞬时规则，局部尾迹相互作用十分强烈。}
\end{sourceblock}

% [原PDF p203]
两泡相互作用的取向偏好随 Reynolds 数呈连续趋势：Ladd 的 Stokes 流结果几乎无偏好，当前中等 Reynolds 数 DNS 出现水平并排倾向，势流极限则形成很强的水平层。实验也观察到同一趋势。虽然气泡弥散高度各向异性，速度脉动的近 Gaussian 分布仍支持用各向异性扩散张量描述长期位移。液相脉动按 $\alpha U_b^2$ 缩放与势流估计一致，但其各方向分量不相等，闭式模型不能只给一个标量“气泡诱导湍流强度”。此外，平均量约十泡可收敛并不意味着微结构已与周期盒无关，配对分布和低波数能谱仍需更大系统验证。
两泡相互作用的取向偏好随 Reynolds 数呈连续趋势：Ladd 的 Stokes 流结果几乎无偏好，当前中等 Reynolds 数 DNS 出现水平并排倾向，势流极限则形成很强的水平层。实验也观察到同一趋势。虽然气泡弥散高度各向异性，速度脉动的近 Gaussian 分布仍支持用各向异性扩散张量描述长期位移。液相脉动按 $\alpha U_b^2$ 缩放与势流估计一致，但其各方向分量不相等，闭式模型不能只给一个标量“气泡诱导湍流强度”。此外，平均量约十泡可收敛并不意味着微结构已与周期盒无关，配对分布和低波数能谱仍需更大系统验证。
低 Reynolds 数时自由演化泡群比规则周期阵列上升更快，中等 Reynolds 数时趋势相反。二泡主要经历尾随、接触和翻滚：后泡进入前泡尾迹，两泡接近后由竖直排列转成水平并排。这既不同于 Stokes 流中没有第三泡便保持相对取向，也不同于势流中同轴两泡互相排斥。对分布函数显示水平排列偏好随 $\mathrm{Re}_b$ 增强；速度脉动概率近似 Gaussian，竖直弥散远大于水平弥散，液相脉动约按 $\alpha U_b^2$ 缩放但同样强烈各向异性。
系统尺度检验表明，平均上升速度和液相 Reynolds 应力等量，用大约十个气泡便可得到较稳定统计；气泡速度脉动和扩散则对系统尺寸更敏感，收敛慢得多。单气泡周期阵列与自由相互作用气泡的平均速度存在明显差异：低 Reynolds 数时自由气泡群比规则阵列上升更快，而中等 Reynolds 数时趋势反转。

% [原PDF p204]
图 8.2 比较周期域内各 48 个气泡的两种计算：近球形气泡倾向组成水平“筏”，可变形气泡则更容易形成竖直气泡柱。对近球形泡，最大系统的能谱很快达到稳态，未见远大于泡径的模态继续增长。对变形泡，暂时聚集的泡群像一个大泡那样加速上升并带动羽流；球形泡在羽流边缘受到向外升力而被驱散，可变形泡的升力却指向羽流，使聚集自我加强，形成快速上升的“烟囱”。
近球形气泡的典型二体作用是“尾随--接触--翻滚”（drafting, kissing and tumbling）：后方气泡被前方尾迹吸入，二者靠近甚至接触，随后竖直气泡对翻转为并排上升。这既不同于 Stokes 流中无第三气泡作用时相对取向不变，也不同于势流所预言的同轴气泡相互排斥。气泡对分布函数显示水平并排的偏好，而且上升 Reynolds 数越大越明显；但整体空间分布仍可近似均匀。

% [原PDF p205]
把气泡初始限制在一根竖直柱内时，近球形泡立即散开，而变形泡保持柱状并比均匀泡群上升得快。更高 Reynolds 数的 48 泡计算表明，近球形泡形成水平层，可变形泡形成竖直流束；但很高 Reynolds 数下球形泡的路径不稳定也可能促成非均匀分布。形变还使配对取向从水平转向竖直，并显著增大速度脉动。图 8.2 的两帧正用于说明：空间组织不是仅由空隙率决定，形状依赖的升力是关键闭式量。
气泡速度脉动的概率分布近似 Gaussian。扩散具有强各向异性，竖直扩散显著大于水平扩散，但仍可用各向异性扩散过程建模。液相速度脉动尺度与空隙率乘以滑移速度平方 $\alpha U_b^2$ 相符，不过与势流不同，实际脉动也高度各向异性。最大规模模拟的能谱很快达到稳定，未见远大于气泡直径的模式持续增长，说明近球形、初始均匀的气泡群可在较长时间保持宏观均匀。

% [原PDF p206]
竖直通道实验早已发现：上升流的气泡往往在壁面附近形成空隙率峰，下降流则形成中心峰；气泡较大、可明显变形时，上升流的壁峰会转成中心峰。DNS 通道在流向和展向周期、壁法向两侧为固壁，重力向下，恒定压差可驱动上行或下行。随机初始气泡演化到近稳态后，壁面剪切平衡混合物重量与外加压差；通过监控总流量等平均量确认稳态，再把统计量表示成壁法向的一维剖面。
气泡可变形性会改变这一结论。近球形气泡往往形成水平“筏”，而较易变形的气泡可能聚集成竖直气泡柱或“烟囱”，并以明显更高速度上升。机制来自升力方向对形状的依赖：若局部聚集形成上升羽流，近球形气泡在羽流边缘剪切中受到指向外侧的升力，因而被驱散；变形气泡的升力可反向指向羽流，使更多气泡被吸入，产生正反馈。较高 Reynolds 数下，球形气泡还可能因路径不稳定而出现非均匀聚集，而变形气泡的形状振荡会改变配对方向与速度脉动分布。由此可见，闭式关系不能只依赖空隙率，还必须体现 Reynolds 数、变形能力和局部剪切。

\section{竖直通道中的气泡流}

% [原PDF p207]
图 8.3 给出层流上升流与下降流某时刻的气泡位置，并在通道中面叠加竖直速度等值线。近球形泡在竖直剪切中被推向液体相对它运动更快的一侧：上升流向壁，下降流离壁。气泡横迁移改变核心平均密度，直至其重量与压差精确平衡，核心剪切归零。上升流形成富泡壁层，下降流形成无泡壁层；若迁移过量，核心与壁层速度差会反向改变升力，因此该分层具有恢复机制。
竖直通道由两侧固壁限制，在展向和流向取周期边界。重力向下，气泡因浮力上升，外加压力梯度可驱动液体向上或向下。随机分布的气泡经长时间演化达到近似稳态；此时壁面剪切平衡混合物重量与外加压力梯度。对充分长、充分宽的通道，平均量主要随壁法向坐标变化，便于与一维平均模型比较，但 DNS 本身必须保持三维。

\subsection{层流中的近球形气泡}

% [原PDF p208]
图 8.4 左列把上升流、下降流的平均空隙率与简单静水平衡模型比较，右列给出平均液速及平均泡速；下降流还叠加解析速度预测。核心空隙率由混合物重量等于压差决定。上升流壁层约一泡径厚，所含泡数等于从核心移出的数量；下降流无泡壁层厚度则由补入核心的泡数决定。无泡下降壁层可直接积分稳态平行流方程，故流量预测准确；富泡上升壁层需要更复杂闭合。
层流上升流和下降流在稳态都分成两个清晰区域：约一个气泡直径厚的壁层，以及占据通道主体、近似均匀的核心区。结构由升力诱导的横向迁移形成。近球形浮升气泡处在竖直剪切流中时，升力把它推向“液体相对气泡运动更快”的一侧。因此在上升流中气泡移向壁面，在下降流中则离开壁面。

上升流的核心区因气泡被移走而平均密度增大，直至其重量恰与压力梯度平衡；下降流的气泡被移入核心区，使核心区更具浮力，最终也达到同样的静水平衡。达到平衡后核心区速度梯度为零，进一步横向迁移停止，所有显著剪切集中在壁层。上升流的富气泡壁层被压力梯度推向上方；下降流的无气泡壁层则在自重和压差共同作用下向下流动。若气泡分配偏离平衡，速度剪切的变化会把气泡推回，因此该构型具有自稳定性。

% [原PDF p209]
湍流下降流算例的摩擦 Reynolds 数为 $127.3$，无泡总体 Reynolds 数为 $3786$，采用 18--72 个近球形泡、$192\times160\times96$ 网格，壁法向节点向壁面加密。升力仍把气泡赶出壁层。即使无泡层厚度小于 50 个壁面单位，湍流也能维持；层较厚时，其外边界会因富泡核心摆动而随时间变化。图 8.5 比较有泡、无泡平均液速和平均泡速；无泡压差已调到与两相流相同的总驱动力。
核心区空隙率可由“混合物重量等于压力梯度”直接预测。上升流中壁层约一气泡直径厚，空隙率由必须从核心移出的气泡数决定；下降流中无气泡壁层厚度由必须补入核心的气泡数决定。DNS 的平均空隙率剖面与这一简单模型高度吻合。下降流壁层不含气泡，稳态平行 Navier--Stokes 方程可直接积分，速度剖面与总流量也能解析预测；上升流富气泡壁层的闭合则更复杂。

\subsection{湍流下降流}

% [原PDF p210]
加入气泡后，核心原本略有梯度的速度变成严格塞状，核心 Reynolds 应力归零，缓冲层应力也降低；但总液速改变出乎意料地小，因为主要加速仍发生在保持标准壁面律的无泡层。多分散计算显示，只要各泡仍近球形，泡径对平均速度剖面影响很小，虽然小泡的滑移速度更低。上升流形势不同：变形可能令升力反向或近零。为比较这一点，21 泡湍流上升算例采用 $256\times192\times128$ 网格，并将 $\mathrm{Eo}$ 从 $0.45$ 改到 $4.5$。
湍流下降流中，近球形气泡仍受升力驱动离壁，并形成无气泡壁层。该区域的平均速度服从标准壁面律，因此即使壁层只有几十个壁面单位厚，湍流仍可维持。气泡核心区则趋于塞状速度，核心 Reynolds 应力几乎消失，缓冲层中的 Reynolds 应力也减小。若把无气泡单相流的压力梯度调整到与两相流相同的总驱动力，加入气泡对平均液速的改变反而不大：主要变化是原本略带梯度的核心速度变得完全均匀。

气泡尺寸改变其滑移速度，但只要均保持近球形，单分散或多分散尺寸分布对平均液速剖面的影响都较小。这里关键不是每个气泡的绝对尺寸，而是形状所决定的升力方向。

% [原PDF p211]
图 8.6 左侧近球形泡形成明显壁峰，右侧可变形泡停留在通道中部；两图同时给出中面竖直速度等值线。图 8.7 的空隙率剖面显示，球形泡仍服从静水平衡模型，可变形泡则在中心最高、向壁递减。图 8.8 表明，这种相分布会显著改变流量：球形泡使平均液速大幅下降；在按混合物平均密度修正总压差后，变形泡对液速影响较小。后者因阻力大而滑移较慢，却被更快的中心液流携带，绝对上升速度反而更高。
\subsection{湍流上升流与变形效应}

在湍流上升流中，近球形气泡形成显著壁峰；更易变形的气泡则留在通道中部。近球形气泡的空隙率仍可由静水平衡模型很好预测；变形气泡的空隙率在中心最高，向壁面逐渐下降。这一差异对流量影响深刻：近球形气泡加入后平均液速显著降低，而较易变形气泡在按混合物平均密度修正总压力梯度后，对液速影响较小。变形气泡因阻力较大而滑移速度较小，却处于更快的中心液流中，所以其绝对上升速度反而可更高。

% [原PDF p212]
图 8.7 是图 8.6 两组上升流的横向平均空隙率及解析模型；核心区的基本结构由升力控制，并趋向静水平衡。均匀流研究给出的泡间动力学可用于描述核心，而壁层必须另行处理：下降流壁层无泡，动力学相对简单；上升流壁层可能聚集大量气泡，对泡数、变形能力及近壁相互作用都很敏感。第 8.4 节据此讨论 DNS 在真实物性与有限计算资源之间必须作出的取舍。
综合层流与湍流结果，通道核心结构主要由气泡升力决定，并趋向静水平衡；壁层动力学则在上升流和下降流中本质不同。下降流壁层无气泡，模型相对简单；上升流壁层富含气泡，其行为对气泡数目、变形性和相互作用十分敏感。

\section{讨论}

% [原PDF p213]
图 8.8 给出球形泡和变形泡上升算例的横向平均液速。实际水--空气密度比约 800；某些算法在此比值下失效，另一些则因变系数压力方程而显著变慢。泡内密度一方面决定浮力差，另一方面决定自身惯性。浮力通常是一阶主导因素；非定常加速更多由周围液体附加质量控制，球形泡的附加质量约为排液质量的一半。因此可在浮力项保持真实密度差、在惯性项采用审慎近似，这与 Boussinesq 论证相似。
气泡流 DNS 清楚展示了机遇与计算限制。真实装置包含数百乃至数千个气泡，并跨越很宽尺度；空气泡在水中的上升 Reynolds 数即使对小气泡也可达数百。Reynolds 数越高，解析每个气泡尾迹所需网格越密，同一计算资源下可容纳的气泡越少，因此系统尺度与分辨率相互制约。

水--空气密度比约为 $800$，大密度比会使某些算法失效，也使变系数压力方程求解显著变慢。密度差一方面决定浮力 $\Delta\rho\,g$，另一方面影响气泡惯性；通常前者更重要，因为气泡加速主要由周围液体的附加质量控制。球形气泡的附加质量约为排开液体质量的一半。因此数值上可在浮力项中保留真实密度差，而在惯性项采用经过论证的近似，这与 Boussinesq 思路相似，但适用性必须逐问题确认。

% [原PDF p214]
极大表面张力会放大寄生流，但一旦高 Reynolds 数小泡在 $\mathrm{Eo}\approx0.1$ 已保持球形，继续增大 $\sigma$ 不再改变宏观流动，故数值上无需追求更极端值。这种处理只适用于形状已饱和的气泡流。尽管系统规模、Reynolds 数、密度比和表面张力均有限制，DNS 已能为工业平均方程提供闭式依据。相关工作还包括气泡、悬浮固体颗粒等体系；随着方法和计算机进步，将能研究更复杂的分布并把微观解析结果系统地转化成平均模型。
过大的表面张力还会放大第 7 章讨论的寄生流。不过对很小、近球形气泡，一旦 $\mathrm{Eo}$ 已低到形状不再变化，进一步提高 $\sigma$ 不会改变流动；高 Reynolds 数下 $\mathrm{Eo}\approx0.1$ 往往已足以维持球形。因此可使用“足以保持正确形状”的表面张力而不必机械追求真实极大值。这样的近似只对特定气泡流成立，不能无条件移植到雾化或薄膜破裂问题。

尽管存在这些限制，DNS 已能把微观相互作用、相分布和平均闭式关系联系起来。它不仅给出关联式数据，更能解释为何球形与变形气泡会形成完全不同的空间结构，这是单纯实验平均量或经验模型难以提供的信息。

\chapter{雾化与破裂}

% [原PDF p215]
雾化流跨越从喷嘴直径到最终液滴的宽广尺度，并不断发生液柱拉伸、液膜变薄、穿孔、液丝夹断和液滴再破裂。完整 DNS 必须同时解析气液两侧旋涡与界面拓扑，而小尺度总是先受网格限制。本章先用液丝和液膜的经典失稳建立基本物理，再讨论高速射流的结构、液滴生成机制与稳定性理论，最后比较二维、三维雾化计算。这里的“原子化”并非分裂到原子尺度，而是指宏观液体被细分成大量喷雾液滴。

液体射流、液层和液滴发生大变形并最终断裂的现象遍及自然与工程。液体分裂成少量大块通常称为破裂；高速射流形成大量细小液滴云的强烈碎裂过程称为雾化。液滴尺度仍远大于原子尺度，“雾化”一词并非字面意义上的原子化，但它准确传达了液体被高度细分为喷雾的外观。

雾化器常从喷嘴高速喷出液柱，也可以产生高速液膜，或让多股液膜相互碰撞。其物理过程包含界面失稳、液膜与液丝形成、拓扑改变和二次破裂，尺度跨度极大。近年来界面算法、并行计算和自适应网格的发展，才使较完整的直接模拟成为可能。

\section{引言}

海面强风会从波峰剥离液滴，破碎波产生的气泡在自由表面爆裂后又形成细雾；这些过程影响海气质量、动量与热量交换，进而参与气候动力学。工程上则需控制喷雾性质：窄液滴粒径分布有利于均匀蒸发，大液滴可增强混合，而消防水柱又应尽量避免雾化以保持射程。燃烧应用尤其依赖初级雾化，因为把后续液滴当点粒子处理的喷雾燃烧程序，最大的不确定性正是喷嘴出口最初如何生成液滴。

% [原PDF p216]
液柱与气体同轴流过分隔板后的两相混合层，是研究喷嘴出口的标准模型。若只考虑表面张力、黏性与惯性，主要参数是两相 Reynolds 数、Weber 数、密度比、黏度比和边界层厚度比。液体可以从圆孔形成射流，也可由狭缝形成薄膜，或由两股液膜碰撞后形成更薄的片。不同装置虽几何不同，却反复出现相同的基本单元：圆柱液丝的 Plateau--Rayleigh 失稳、自由液膜回缩形成的 Taylor--Culick 边缘，以及气液剪切层的 Kelvin--Helmholtz 型放大。
完整模拟高速雾化射流需要同时覆盖喷嘴尺度、边界层、液膜、液丝和微小液滴，计算量极大。另一条研究路线是隔离基本构型，例如空气中的自由薄液膜及其端部边缘，或喷嘴分隔板后的两相混合层，分别研究失稳和破裂机制，再把所得认识用于完整雾化器。

\section{液丝、液膜与边缘破裂}

\subsection{Plateau--Rayleigh 射流失稳}

% [原PDF p217]
图 9.1 比较 Plateau--Rayleigh 射流失稳的线性阶段与非线性夹断阶段。对无黏圆柱液柱，轴对称扰动仅在波长大于周长时增长，最大增长波长约为液柱直径的若干倍；表面张力驱动液体从颈部流向鼓包，最终形成主液滴和可能的卫星滴。线性理论能给出增长率和优势波数，但不能描述颈部接近零半径后的强非线性。数值模拟若轴对称，可很经济地追踪到夹断；三维扰动则可能激发非轴对称液丝运动。
低速液柱从喷嘴流出后会周期性断成液滴。Plateau 证明，半径为 $r$ 的圆柱射流中，波长
\begin{equation}
 \lambda>2\pi r
\end{equation}
的轴对称半径扰动不稳定；Rayleigh 进一步求得最大增长率对应
$\lambda/r\simeq9.01$。其根本原因是长波扰动可在守恒体积的同时减小总表面积，释放表面能。线性阶段扰动指数增长，非线性阶段鼓包之间形成细颈，最后夹断；主液滴之外还可能产生卫星液滴。

\subsection{液膜与液丝破裂}

% [原PDF p218]
图 9.2 将悬滴实验中的细液丝夹断与数值结果比较。夹断附近出现快速收缩的颈部，局部惯性、黏性和表面张力的平衡决定最小半径随剩余时间的幂律。液丝两端的几何与远场条件会影响卫星滴，但临近奇点的局部轮廓可趋于相似。数值界面方法不能真正达到分子尺度；当最小颈宽降到几个网格时，VOF 或水平集会自动重联，前沿追踪则需显式改变连接关系。夹断后的宏观液滴运动可继续求解，但拓扑时刻带有网格依赖。
细颈最后阶段可用润滑型一维方程描述。设 $h(x,t)$ 为薄膜半厚度或液丝半径，$u(x,t)$ 为轴向平均速度。单位长度质量正比于 $h^n$，其中薄膜 $n=1$、轴对称液丝 $n=2$，质量守恒给出
\begin{equation}
 \frac{\partial h}{\partial t}
 +u\frac{\partial h}{\partial x}
 +\frac{h}{n}\frac{\partial u}{\partial x}=0.
 \tag{9.1}
\end{equation}
薄膜的动量方程可写为
\begin{equation}
 \frac{\partial u}{\partial t}+u\frac{\partial u}{\partial x}
 -\frac{4\mu}{\rho h}\frac{\partial}{\partial x}
 \left(h\frac{\partial u}{\partial x}\right)
 +\frac{\sigma}{\rho}\frac{\partial\kappa_1}{\partial x}=0,
 \tag{9.2}
\end{equation}
其中平面曲率
\begin{equation}
 \kappa_1=-\frac{h_{xx}}{(1+h_x^2)^{3/2}}.
 \tag{9.3}
\end{equation}
液丝的黏性系数和几何项不同：
\begin{equation}
 \frac{\partial u}{\partial t}+u\frac{\partial u}{\partial x}
 -\frac{3\mu}{\rho h^2}\frac{\partial}{\partial x}
 \left(h^2\frac{\partial u}{\partial x}\right)
 +\frac{\sigma}{\rho}\frac{\partial(\kappa_1+\kappa_2)}{\partial x}=0,
 \tag{9.4}
\end{equation}
环向曲率为
\begin{equation}
 \kappa_2=\frac{1}{h(1+h_x^2)^{1/2}}.
 \tag{9.5}
\end{equation}
保留精确曲率而不是只用 $1/h$，可使完整 Navier--Stokes 方程的静止液丝仍是降维方程的静止解；但夹断附近 $|h_x|$ 很大，润滑近似的严格适用范围仍需谨慎。

% [原PDF p219]
在液膜或包裹气体的薄层中，两侧界面靠近时也会重联。图 9.2 的完整模拟显示：主滴拉出长丝，细颈相继断裂并留下卫星滴。图 9.3 则示意薄层中的孔洞形成和扩展；在数值方法中，只要膜厚小于约一两个单元，两侧 VOF 重构便会合并，形成一个孔。这种自动穿孔便于继续模拟宏观回缩，却不能被解释为真实的分子破膜机制。真实液膜可能由范德华力、杂质或热涨落触发穿孔，DNS 只能把这些过程等效为拓扑规则。
惯性与表面张力主导、黏性可忽略时，临近夹断时刻 $t_c$ 的自相似尺度为
\begin{equation}
 \ell\sim\left[\frac{\sigma(t_c-t)^2}{\rho}\right]^{1/3},
 \qquad
 U\sim\left[\frac{\sigma}{\rho(t_c-t)}\right]^{1/3}.
\end{equation}
若惯性、黏性和表面张力都参与，则出现另一组自相似标度。高分辨率 VOF 模拟可以恢复这两个区域。有限时间夹断在达到分子尺度前已由连续介质方程相当准确地控制，说明分子作用通常只决定极短的最终重联阶段，而不显著改变宏观夹断动力学。

% [原PDF p220]
图 9.3 的相序列展示薄层断裂后的界面重联：初始两条界面逼近，在最薄处形成孔，孔缘卷起并向外传播。重联前后质量应守恒，且新生成边缘的曲率与速度不能产生非物理冲量。前沿追踪需删除旧面元并重新三角化，VOF 则由单元体积分数自动给出新拓扑。无论哪种方法，穿孔阈值都与网格相关；因此可以研究孔形成后的 Taylor--Culick 回缩速度和边缘失稳，却不能仅凭计算预测最初成核尺度。
薄液膜穿孔则不同。不计分子力时，处在伸长流
$u=\omega x,\ v=-\omega y$ 中的膜厚按
\begin{equation}
 h(t)\sim h_0\exp[-\omega(t-t_0)]
\end{equation}
衰减，只在无限时间趋于零；到达分子尺度 $h_m$ 的时刻含
$\omega^{-1}|\ln(h_m/h_0)|$。数值模拟若让界面在相距一个网格时自动重联，破裂时刻就会显著依赖 $\Delta x$。真实薄膜还受 van der Waals 力、表面活性剂和界面流变性质影响，可能被吸引而破裂，也可能因排斥或“膜状弹性”显著延缓排液，因此必须逐个体系建模。

% [原PDF p221]
Taylor--Culick 推导把边缘前方静止膜的表面能释放率与边缘动量增长率平衡。两张自由表面共同提供 $2\sigma$ 的拉力，所以速度中出现因子 2。边缘横截面并非静态半圆：被吸入液体在内部转向，产生涡量与回流；Oh 数降低时，惯性使颈部更薄、涡旋更强。二维模型能给出截面与回缩速度，却不能预测沿边缘方向的三维指状失稳。原书比较简化边缘方程、前沿追踪和 Gerris 结果，指出临界附近的复杂性不是单一毛细波，而是增长边缘、后方液膜和内部涡量共同作用的结果。
Taylor--Culick 推导把边缘前方静止膜的表面能释放率与边缘动量增长率平衡。两张自由表面共同提供 $2\sigma$ 的拉力，所以速度中出现因子 2。边缘横截面并非静态半圆：被吸入液体在内部转向，产生涡量与回流；Oh 数降低时，惯性使颈部更薄、涡旋更强。二维模型能给出截面与回缩速度，却不能预测沿边缘方向的三维指状失稳。原书比较简化边缘方程、前沿追踪和 Gerris 结果，指出临界附近的复杂性不是单一毛细波，而是增长边缘、后方液膜和内部涡量共同作用的结果。
自由液膜开孔后，边缘因表面张力以近恒速回缩。对膜厚 $e$、液体密度 $\rho_l$，动量平衡给出 Taylor--Culick 速度 $U_{TC}=\sqrt{2\sigma/(\rho_l e)}$。边缘不断吸收液膜质量，截面半径随传播距离增长；Ohnesorge 数 $\mathrm{Oh}=\mu_l/\sqrt{\rho_l\sigma e}$ 控制边缘是平滑、带涡还是强烈非定常。简化的随边缘移动模型可在很大尺度范围内解析其截面，而全 Navier--Stokes 计算还会显示颈部与回流。图 9.4 给出三个 Oh 数对应的边缘形状。
\subsection{Taylor--Culick 边缘}

% [原PDF p222]
图 9.4 比较不同 Ohnesorge 数的理论边缘截面；黏性越强，液膜到膨大边缘的过渡越平滑。图 9.5 是前沿追踪在无量纲时刻 $t'=8.45\times10^4$ 的边缘形状，虚线标出后续发生复杂流动的参数界限。自由液膜实验如图 9.7 所示，边缘并非总保持二维：横向扰动会产生指状结构和液滴。判定这种失稳来自孤立液环、附着边缘的毛细模态还是运动学聚焦，需要同时考虑边缘半径随时间增长以及膜颈的约束。
液膜破开后，表面张力把自由边缘向后拉，液体在端部积聚成膨大的边缘。对厚度恒定为 $e$ 的液膜，设单位边长边缘质量为 $M(t)$、回缩速度为 $V(t)$，质量和动量守恒为
\begin{equation}
 \frac{dM}{dt}=\rho Ve,
 \qquad
 \frac{d(MV)}{dt}=2\sigma.
 \tag{9.6--9.7}
\end{equation}
无黏、无气动阻力和无重力时，定速解为
\begin{equation}
 V=\left(\frac{2\sigma}{\rho e}\right)^{1/2},
 \qquad M=\rho Vet.
 \tag{9.8}
\end{equation}
特征时间 $t_c=(\rho e^3/\sigma)^{1/2}$，特征速度
$V_c=(\sigma/\rho e)^{1/2}$。厚约 $100\,\mu\mathrm m$、尺度 $10\,\mathrm{mm}$ 的飞溅液冠可在毫秒量级内回缩。

% [原PDF p223]
图 9.6 为 Gerris 在 $\mathrm{Oh}=0.0045$、$\rho_g/\rho_l=0.1$ 下计算的边缘。低黏性时，边缘内部产生多个涡量斑并相互作用；当 $\mathrm{Oh}$ 低于约 $0.07$，流动可转为复杂甚至湍动，颈部厚度显著起伏。边缘仍与液膜连接，这一点排除了把它简单视作孤立圆柱射流。若颈部真的缩到分子尺度并使边缘脱落，脱离后的液环才可独立经历 Plateau--Rayleigh 失稳；实验中这种完全脱离并非普遍前提。
黏性在足够大的控制体动量平衡中不改变 Taylor--Culick 速度，却通过
\begin{equation}
 \mathrm{Oh}=\frac{\mu}{(\rho\sigma e)^{1/2}},
 \qquad \ell_\nu=(\nu t_c)^{1/2}=\mathrm{Oh}^{1/2}e
 \tag{9.9}
\end{equation}
控制边缘形状。大 $\mathrm{Oh}$ 时边缘初期细长、厚度沿程变化；小 $\mathrm{Oh}$ 时边缘很快趋近圆形，并在其前方产生毛细波和比原膜更薄的“颈部”。低 $\mathrm{Oh}$ 的长期模拟还发现边缘内部多个涡量斑相互作用，约在 $\mathrm{Oh}\lesssim0.07$ 时转入复杂甚至湍动状态，颈部厚度强烈波动并可能触发断裂。

% [原PDF p224]
图 9.7 是自由液膜同时从正面与侧面观察的实验照片；图 9.8 则示意一种由初始变形边缘产生液滴的运动学机制。以近恒定法向速度推进的波前可能聚焦、自交并形成尖点，但计算表明，要让这种机制产生含有可观质量的液滴，初始扰动必须很大。另一种更常见的解释是：附着于液膜的膨大边缘本身发生长波失稳，形成横向手指；其增长同时受边缘半径增加和后方薄膜拉力影响。
\subsection{边缘形成液滴和指状结构}

% [原PDF p225]
对图 9.8 的示例，扰动波长与膜厚比为 $\lambda/e=130$。边缘吸收液体后半径持续变大；若在失稳尚未充分增长前达到 $R/\lambda\gtrsim1/\pi$，类似 Plateau--Rayleigh 的模态便被稳定。因此短波最容易因边缘增长而被抑制，只有足够长且初始幅值合适的扰动能发展成液滴或指状液丝。该竞争解释了实验中液滴间距并非简单等于瞬时边缘周长的固定倍数，也说明入口噪声和膜厚变化会显著改变最终破裂图样。
实验中的自由薄膜边缘常沿横向长出指状液丝并甩出液滴。候选机制之一是边缘先与液膜分离，再像孤立圆柱一样发生 Plateau--Rayleigh 失稳；但这要求连接颈缩至分子尺度，实验中似乎并不常见。另一机制是以近恒定法向速度回缩的波前发生几何聚焦和自交，形成运动学奇点；模拟表明，若要由此形成含有可观质量的液滴，初始边缘扰动必须很大，因此它也不太可能是普遍主导机制。

% [原PDF p226]
图 9.9 以分隔板后的两相混合层表示喷嘴出口：液体和气体具有各自速度与边界层厚度，界面从尾缘开始失稳。高速射流中央仍保持连续的部分称为液核，其外侧生成液膜、液丝和液滴。按喷嘴直径定义 $\mathrm{Re}_i=\rho_iU_iD/\mu_i$、$\mathrm{We}_i=\rho_iU_i^2D/\sigma$；同轴流还用动量通量比 $M=\rho_gU_g^2/(\rho_lU_l^2)$。实验中液核长度随 $M$ 增大而减小，大气速比时近似 $L/D\simeq6/\sqrt M$。
更自然的机制是仍连接液膜的膨大边缘发生类似圆柱射流的失稳。扰动波长需远大于膜厚 $e$；否则边缘半径随吸收液体快速增长，当 $R/\lambda$ 超过约 $1/\pi$ 时 Plateau--Rayleigh 模式被抑制。液膜约束、边缘增长和毛细波共同决定最终是规则液滴链还是不规则手指。

\section{高速射流}

\subsection{雾化射流的结构}

随着射流速度超过 Plateau--Rayleigh 区域，完整液柱长度先增后减，继而出现大尺度准螺旋失稳以及小尺度液丝、液膜和液滴。通常依次称为滴落区、第一风致区、第二风致区和雾化区。真正的雾化区产生远小于喷嘴直径 $D$ 的结构，是技术和数值计算最具挑战的区域。

射流中央的致密部分称为液核，外部包围片状或纤维状结构，它们向下游增长并脱落为液滴。特征尺度不仅依赖来流速度和物性，还强烈依赖喷嘴内部流动、气液边界层厚度和入口湍流。脱落液滴还可能碰撞聚并，或在强剪切和加速流中再次破裂，形成二次雾化。

% [原PDF p227]
高速雾化的候选机制包括液相入口湍流、界面线性与非线性失稳，以及喷嘴内部空化。液相湍流涡若动能足以支付新界面能，可拉出同尺度液丝；界面机制则由两相速度差放大 Kelvin--Helmholtz 波；低压区蒸汽泡或气核的增长会改变喷口内流动。图 9.10 用液相 Reynolds 数和气相 Weber 数区分这些主导区。最简单的速度跳跃模型预言 $s\propto k|U_g-U_l|$，无表面张力和黏性时短波增长率无界，数学上不适定。
基于喷嘴直径的主要无量纲数为
\begin{equation}
 \mathrm{We}_i=\frac{\rho_iU_i^2D}{\sigma},
 \qquad
 \mathrm{Re}_i=\frac{\rho_iU_iD}{\mu_i},
 \qquad i=l,g,
\end{equation}
同轴射流还使用动量通量比
\begin{equation}
 M=\frac{\rho_gU_g^2}{\rho_lU_l^2},
 \qquad r=\frac{\rho_g}{\rho_l},
 \qquad m=\frac{\mu_g}{\mu_l}.
\end{equation}
实验表明液核长度随 $M$ 增大而减小，大气速比时近似
$L/D\simeq6/\sqrt M$；单射流雾化器则有
$L/D\simeq6(\rho_l/\rho_g)^{1/2}$。以分隔板尾缘两相混合层解释时，更合适的尺度是气液两侧边界层厚度 $\delta_g,\delta_l$：
\begin{equation}
 \mathrm{We}_i^*=\frac{\rho_i(U_i-U_{\rm int})^2\delta_i}{\sigma},
 \qquad
 \mathrm{Re}_i^*=\frac{\rho_i|U_i-U_{\rm int}|\delta_i}{\mu_i}.
\end{equation}

% [原PDF p228]
图 9.10 的 Hopfinger 图并不是锐利相界，而是提示：高 $\mathrm{Re}_l$ 偏向液相湍流主导，高 $\mathrm{We}_g$ 偏向空气动力失稳主导，空化还取决于喷嘴压力史。表面张力使高波数稳定，并给出低密度比下的最大增长波长尺度；黏性边界层则把速度跃迁平滑为有限厚度。由于真实喷口两侧边界层厚度通常远小于直径，用 $D$ 定义的无量纲数并不足够，必须引入以 $\delta_g$、$\delta_l$ 和界面速度 $U_{\rm int}$ 定义的局部 Reynolds、Weber 数。
\subsection{液滴形成机制}

% [原PDF p229]
切向应力连续要求 $\mu_g|U_g-U_{\rm int}|/\delta_g=\mu_l|U_l-U_{\rm int}|/\delta_l$。若 $n=\delta_g/\delta_l$、$m=\mu_g/\mu_l$、$r=\rho_g/\rho_l$，局部参数满足 $\mathrm{Re}_g^*m^2/(\mathrm{Re}_l^*n^2r)=1$。时间发展边界层由黏性扩散给出 $\delta_i\sim\sqrt{\nu_it}$，因此 $n\sim\sqrt{m/r}$。典型气液物性使液侧速度跃迁远小于气侧，分隔板尾缘产生的大部分涡量扩散进气相。完整黏性扰动方程形成两相 Orr--Sommerfeld 特征值问题。
三类主要机制是上游液相湍流、理想化射流或液膜的界面失稳，以及喷嘴内部空化。液相湍流理论认为，尺度 $\ell$ 的液相涡具有动能约 $\rho_l\ell^3U_l^2$，把界面拉成同尺度液丝需表面能约 $\sigma\ell^2$，两者平衡给出
\begin{equation}
 \ell\simeq D\,\mathrm{We}_l^{-1}.
\end{equation}
当 $\ell\ll D$ 时还需考虑湍流能量的尺度分布。该机制要求 $\mathrm{Re}_l$ 足够大，使喷嘴内液相湍流充分发展。

% [原PDF p230]
Orr--Sommerfeld 谱除缓慢趋近无黏 Kelvin--Helmholtz 解的模态外，还含与黏性剪切层有关的 H 模态。其最大波数随 $\sqrt{\mathrm{Re}_g^*}$ 缩放，增长率随 $k_m|U_g-U_l|$ 缩放。图 9.13 将线性理论与时间型 DNS 的界面振幅比较，吻合程度还取决于混合单元中黏度采用算术平均还是调和平均；调和平均通常更接近理论。图 9.11 给出时间型混合层速度剖面，图 9.12 显示高分辨率涡量与界面形状。
稳定性理论从理想化两相混合层出发：小扰动先长成界面波，二维波的次级失稳或三维扰动的直接增长再产生液丝。空化机制则源于喷嘴上游低压区中的蒸汽泡或汽膜；它们可能在负压区成核，也可能由已有微泡长大，并在很高流速下显著改变射流形态。三种机制占据不同参数区：高液相 Reynolds 数偏向湍流主导，高气相 Weber 数偏向空气动力失稳主导。

\subsection{稳定性理论}

\subsubsection{基本 Kelvin--Helmholtz 分析}

% [原PDF p231]
图 9.11(a) 中 ``BL'' 是有限厚度边界层，``KH'' 是理想速度跳跃；(b) 给出周期时间型计算的典型界面。图 9.12(a) 用颜色表示涡量，(b) 对比不同阶段。二维时间型计算从 $128^2$ 加密到 $512^2$ 后，液丝被拉得更长，夹断时间与小滴尺度均改变，说明低分辨率的数值黏性会过早抑制界面。时间型模型适合检验增长率与机制，但因流向周期，不能直接给出从喷嘴到远场的空间液滴分布。
令界面扰动为
\begin{equation}
 h=h_0+h_1\exp(st+\mathrm i\omega t+\mathrm ikx),
 \tag{9.10}
\end{equation}
增长率 $s$ 和频率 $\omega$ 由波数 $k$ 与物性决定。若基流速度在界面发生尖锐跳跃，等密度单相极限给出
\begin{equation}
 s=\frac{k}{2}|U_g-U_l|,
 \tag{9.11}
\end{equation}
而 $\rho_g\ll\rho_l$ 的两相极限近似为
\begin{equation}
 s\simeq\left(\frac{\rho_g}{\rho_l}\right)^{1/2}
 k|U_g-U_l|.
 \tag{9.12}
\end{equation}
无黏、无表面张力模型预言 $k\to\infty$ 时增长率无界，任意短波都在任意短时间内放大，构成 Hadamard 意义下的不适定问题。表面张力和黏性会抑制高波数。低气液密度比时，表面张力给出的最大增长波长约为
$\lambda_m=3\pi D/\mathrm{We}_l$。用分段线性边界层同时保留无黏扰动方程虽可估计黏性尺度，但与实验不够一致，需要完整黏性稳定性分析。

% [原PDF p232]
图 9.13 画出算术与调和黏度平均下的振幅增长，并与 Orr--Sommerfeld 预言比较；图 9.14 展示 $512\times512$ 网格上液丝形成和伸长的三个时刻。细化网格不仅让轮廓更平滑，还延长了可解析的薄膜和液丝寿命，从而生成更小结构。若液滴谱随网格持续向小尺寸移动，说明尚未达到统计收敛。为研究实际喷口，必须改成入口--出口式的空间发展计算，并明确给定喷嘴内速度剖面或把喷嘴内部一并求解。
\subsubsection{Orr--Sommerfeld 分析}

% [原PDF p233]
图 9.15 是二维空间发展雾化射流。入口若叠加尺度接近射流半径或边界层厚度的随机小涡，下游破裂会明显加强；没有扰动时，射流可能因数值黏性表现得过于稳定。空间计算的优势是同时观察液核长度、波增长与液滴向下游输运，缺点是结果高度依赖入口湍流。仅规定均匀顶帽速度通常不够真实，因为喷嘴壁面边界层、分隔板尾迹和内部回流会决定初始涡量层。
黏性首先把速度跳跃替换为厚度 $\delta_g$、$\delta_l$ 的边界层。无表面张力梯度时，平行基流的切向应力连续：
\begin{equation}
 \left.\mu_g\frac{\partial u}{\partial z}\right|_g
 =\left.\mu_l\frac{\partial u}{\partial z}\right|_l,
 \tag{9.13}
\end{equation}
从而
\begin{equation}
 \mu_g\frac{|U_g-U_{\rm int}|}{\delta_g}
 =\mu_l\frac{|U_l-U_{\rm int}|}{\delta_l}.
 \tag{9.14}
\end{equation}
若 $n=\delta_g/\delta_l$，则无量纲参数满足
\begin{equation}
 \frac{\mathrm{Re}_g^*m^2}{\mathrm{Re}_l^*n^2r}=1.
 \tag{9.15}
\end{equation}
边界层厚度本身由黏性扩散决定；时间发展问题中
$\delta_i\sim(\nu_it)^{1/2}$，故
\begin{equation}
 n\sim(m/r)^{1/2}.
 \tag{9.16}
\end{equation}
典型气液物性下，液侧速度跳跃远小于气侧，意味着分隔板尾缘产生的大部分涡量扩散到气相。

% [原PDF p234]
图 9.16 把喷嘴内部的同向气液流也包括在二维计算中，因此不需人为在出口施加噪声，界面在尾缘附近便形成与实验相似的波。中等动量通量比下，线性理论、DNS 与实验的优势波长较一致；当 $M$ 很大并接近空气--水物性时，尾缘后几乎立即进入强非线性，线性理论随之失效。图中大气速区仍可能欠分辨，故“未发生破裂”不能自动解释成物理稳定，必须结合局部网格 Reynolds 数检查数值耗散。
线性化 Navier--Stokes 方程形成两相 Orr--Sommerfeld 特征值问题。谱中一支模态缓慢趋近无黏解，另一支与高黏性剪切层相关。后者最大增长波数随 $\sqrt{\mathrm{Re}_g^*}$ 标度，最大增长率随 $k_m|U_g-U_l|$ 标度，被称为 Hooper--Boyd/Hinch 模态（H 模态）。时间型 DNS 的小扰动增长与该理论吻合，但混合单元黏度采用算术平均还是调和平均会影响小振幅阶段；调和平均通常更接近线性理论。

% [原PDF p235]
非模态理论强调，即使所有线性特征模态在短时内并不占优，非正交扰动也能从基流抽取能量并迅速三维化；“先二维后次级失稳”的观点则把展向结构归因于液膜或液缘形成后的新基态。区分两者必须测量三维扰动何时开始指数增长，并与边缘出现时刻比较。时间型 DNS 观察到前者更早，因而排除了该算例中纯边缘毛细机制，但不能证明所有喷嘴都相同。气动力压力、局部密度跃迁下的加速度以及液膜穿孔都可能在不同阶段主导；初始展向波数、相位调制幅值和盒宽也会选择最终鱼骨间距。
非模态理论强调，即使所有线性特征模态在短时内并不占优，非正交扰动也能从基流抽取能量并迅速三维化；“先二维后次级失稳”的观点则把展向结构归因于液膜或液缘形成后的新基态。区分两者必须测量三维扰动何时开始指数增长，并与边缘出现时刻比较。时间型 DNS 观察到前者更早，因而排除了该算例中纯边缘毛细机制，但不能证明所有喷嘴都相同。气动力压力、局部密度跃迁下的加速度以及液膜穿孔都可能在不同阶段主导；初始展向波数、相位调制幅值和盒宽也会选择最终鱼骨间距。
三维化可能通过两条路线发生：非模态扰动在二维波形成前直接瞬态增长，或二维波先卷成薄膜与边缘，继而发生边缘毛细、Rayleigh--Taylor 型或气动力次级失稳。时间型三维模拟在初始界面施加展向调相扰动 $h=h_0+\varepsilon\cos[kx+a\cos(k_yy)]$，观察到完整边缘尚未形成时三维幅值已指数增长，支持气流压力驱动的直接三维机制。薄膜降到网格尺度后自动穿孔，产生主液丝两侧伸出横向细柱的“鱼骨”拓扑。
\section{雾化模拟}

% [原PDF p236]
图 9.17 展示三维 Kelvin--Helmholtz 失稳的发展：初始近二维波逐渐产生展向起伏，继而拉出膜片与液丝。图 9.18 给出 Gerris 自适应八叉树网格上的空间发展同轴射流，网格集中在界面和强涡量区；参数与图 9.16 相近，只修改了分隔板几何。自适应网格显著降低三维成本，但细化与粗化准则会影响最薄膜片的保存时间，拓扑改变必须与局部误差控制一并解释。
\subsection{二维时间发展模拟}

时间型问题在流向取周期边界，初始条件是带有限厚度气液边界层的两相剪切层。$128\times128$ 的低分辨率即可显示界面波和初始液丝；提高到 $512\times512$ 后，液丝被拉得更长，破裂时间和产生液滴的尺度都发生变化。这说明雾化对分辨率与数值黏性极为敏感，低分辨率结果只能用于识别大尺度机制，不能可靠预测液滴谱。

\subsection{二维空间发展模拟}

空间发展计算在入口给定喷嘴出口速度剖面。高速液体射流案例显示，入口是否叠加与射流半径及边界层厚度相当的随机小涡，会极大改变下游雾化量；没有入口扰动时射流过于稳定，但在极高 $\mathrm{Re}_g$ 且分辨率不足的计算中，这很可能只是数值黏性的假象。

比人为入口剖面更真实的方法是把喷嘴内部流动也纳入计算。较低 Reynolds/Weber 数的同轴气液射流模拟无需额外噪声，喷嘴出口附近便形成与实验一致的波。中等动量通量比时，线性理论、DNS 与实验相符；大 $M$、接近空气--水物性时，实验和 DNS 都显示分隔板后立即出现强非线性，线性理论因此偏离。

% [原PDF p237]
图 9.18 的三维射流形成膜片、液丝和液滴，与对应二维计算的稳定外观形成反差，原因可能是三维扰动本身提供了新的增长通道。图 9.19 则是更高速的空间发展计算：喷流头部和侧面出现鱼骨状液丝、薄膜穿孔和大量液滴。三维结构比二维液片更接近实验，但计算域、入口条件和气液物性仍受资源限制。对结果应提取液滴尺寸 PDF、速度分布与液核长度，而不能只凭瞬时图像判断成功。
\subsection{三维计算}

三维失稳有两种观点：其一认为三维扰动可由非模态瞬态增长直接放大；其二认为二维波先形成液膜与端部边缘，随后才发生次级三维失稳。后者又可由边缘毛细失稳、局部 Rayleigh--Taylor 型机制或气动力压力起伏驱动。完整三维 DNS 是区分这些机制最直接的手段。

三维时间型混合层模拟若施加相位随展向轻微调制的初始扰动
\begin{equation}
 h(x,y,0)=h_0+\varepsilon\cos[kx+a\cos(k_yy)],
\end{equation}
可观察到在完整边缘尚未形成前，三维扰动已指数增长，说明至少此时并非边缘毛细效应主导，而更像气流在液膜上施加的气动力压力。薄膜降至网格尺度后会自动穿孔，形成“鱼骨”图样：主液丝两侧伸出横向液柱。数值穿孔尺度不真实，但实验中确实观察到类似图样，提示真实液膜也可能在分子尺度穿孔后保留同类宏观拓扑。

% [原PDF p238]
图 9.19 的高速 Gerris 模拟显示，薄膜先被拉伸和穿孔，孔缘卷成细丝，细丝再经毛细失稳断为液滴。随着网格细化，新的更小液滴持续出现，粒径分布整体向小尺度移动；这意味着视觉上逼真的喷雾仍可能没有收敛。由第一性原理预测稳定液滴谱，需要更低耗散的界面输运、更严格的自适应误差估计以及跨多级网格的统计比较。雾化 DNS 的核心难题不是能否产生复杂形态，而是最小解析尺度是否已经与所报告统计量分离。
空间发展的三维同轴射流也形成液膜、液丝和液滴。较低速度案例中，自适应八叉树网格把分辨率集中在界面和涡量区；高压静止气体中的柴油射流则在更高 Reynolds 和 Weber 数下产生大量膜片、鱼骨状液丝和液滴。

\begin{sourceblock}
\sourcecrop{238}{84}{255}{408}{650}{0.60}
\captionof{figure}{Gerris 自适应网格模拟的高速雾化射流。上图为 $Ut/D\sim6$ 附近的整体发展，下图放大了射流头部和侧面的液膜、液丝与液滴。液膜穿孔后形成细丝，细丝再经毛细失稳断成液滴。}
\end{sourceblock}

这类计算已接近“由第一性原理预测液滴粒径分布”的目标，但网格收敛仍是决定性问题：细化网格后，统计分布往往整体移向更小液滴，表明先前解仍欠分辨。可靠的粒径 PDF 需要更快的计算平台、更低耗散的界面算法以及严格的自适应网格收敛研究；仅看到视觉上逼真的喷雾并不等于粒径统计已经可信。

\chapter{液滴碰撞、撞击与飞溅}

% [原PDF p239]
液滴正碰、撞击固壁或液池时会产生铺展、回弹、聚并、夹断、气泡卷吸和飞溅。燃料喷雾、喷墨、涂覆、冶金、雨蚀、农药沉积以及高速熔融物撞击都包含这些过程。惯性足够大时，甚至陨石坑的中央峰也可类比 Worthington 射流。早期高速摄影提供了经典形态，DNS 则能同时给出不可见的压力、涡量和能量转移。二维及轴对称计算发展较早，解析非轴对称手指、液冠穿孔和液滴喷射的三维计算则长期受网格量限制。

液滴碰撞和撞击因其鲜明形态成为流体力学最具代表性的现象之一。二维和轴对称模拟在两相流数值方法早期便已出现，但能解析非轴对称手指、液冠破裂和液滴喷射的三维模拟直到较近时期才成为可能；即便如此，完全复现实验尺度和真实气液物性仍非常困难。

\section{引言}

燃料液滴撞击管壁或燃烧室壁面后，可能铺展成薄膜，也可能碎裂成更细喷雾。喷墨打印、喷涂、冶金、气体喷吹过程、雨蚀和农药沉积都涉及撞击。高速液滴会损伤多相流透平；极端核事故中，熔融堆芯物质撞击安全壳也会产生高速飞溅。雨滴撞击影响土壤侵蚀和海气交换。更高能量下，陨石或金属高速撞击中的惯性远大于材料弹塑性应力，也可用类似流体动力学解释，月坑中央峰甚至可视为 Worthington 射流的遗迹。

% [原PDF p240]
Harlow 与 Shannon 的早期 MAC 计算在交错网格上用标记粒子识别自由表面，忽略黏性和表面张力；Foote 随后加入两者，并在低 Reynolds、低 Weber 数下重现实验回弹。随界面移动的有限元法适合低速润湿和固壁附近形变，固定网格 VOF/前沿追踪则更容易处理拓扑变化。低速等尺寸液滴正碰可出现反弹、完全聚并、暂时聚并后分离，以及碎成三个或更多液滴；若壁面自由滑移且无润湿，正碰的对称半域等价于单滴撞击平壁。
早期 Worthington 和后来 Edgerton 的高速摄影建立了该领域的经典图像。数值模拟的价值不仅在于重现这些形态，还在于访问全流场、能量收支和界面拓扑，从而把不同碰撞结果与控制参数联系起来。

\section{早期模拟}

Harlow 与 Shannon 最早用原始 MAC 方法模拟撞击，采用交错网格和体内标记粒子，但忽略黏性与表面张力。Foote 随后加入两者，并在低 Reynolds 数、低 Weber 数下得到与回弹液滴实验良好吻合的结果。另一条路线是随界面运动的有限元方法，适于处理润湿和固壁形变。已有工作表明，固定网格并非天然不适合撞击问题；关键在于界面输运、表面张力、压力跳跃和拓扑改变是否得到一致处理。

\section{低速撞击与碰撞}

% [原PDF p241]
图 10.1 给出两相同液滴迎头碰撞的前沿追踪结果。主要参数为 $\mathrm{Re}=\rho UD/\mu$、$\mathrm{We}=\rho U^2D/\sigma$ 和 $\mathrm{Oh}=\sqrt{\mathrm{We}}/\mathrm{Re}$；环境流体还需密度比、黏度比。两滴受压后在接触区外缘形成径向薄片，达到最大横向伸展后由毛细力回缩。守恒的前沿表面力离散确保每个闭合液滴的毛细合力为零，避免把寄生冲量误认为碰撞偏斜。图中不同帧显示压扁、回缩和可能的再次分离。
同尺寸液滴正碰在低到中等速度下可出现四种结果：完全不聚并的反弹；合并为一个液滴；短暂合并后再次分离；以及更高速度下碎成三个或更多液滴。若忽略润湿并在平壁施加自由滑移，同尺寸液滴正碰与单液滴撞击平壁具有对称等价性。干壁撞击同样可分为回弹、沉积铺展和破裂飞溅；湿壁或深液池还会形成凹坑、卷吸气泡和竖直射流。

深液池撞击的规则卷吸只出现在中间速度范围：速度太小或太大都不形成气泡。凹坑闭合时可能产生细而快的竖直射流；超过规则卷吸区后出现较粗的 Worthington 射流，并从顶端夹断出一个或多个液滴。射流高度和脱落液滴数通常随撞击速度增加。低速撞击还能形成涡环，而高速度时流动迅速转为表面飞溅，涡量反而难以深入液池。

% [原PDF p242]
代表性正碰算例取 $\mathrm{Re}=96$ 等参数，并比较界面形态随时间的演化。若简单把初始动能与扁饼状液滴新增表面能平衡，可得 $R_{\max}/D\sim\mathrm{We}^{1/2}$；考虑黏性边界层或体积耗散会给出其他指数。实验数据跨越的参数范围有限且散布较大，因此不能仅凭某一拟合断言普适标度。DNS 的优势是直接积分动能、黏性耗散与表面能，但当 $\mathrm{Re}$、$\mathrm{We}$ 增大到出现液冠或碎裂时，低速铺展理论本身已失效。
主要控制参数是
\begin{equation}
 \mathrm{We}=\frac{\rho U^2D}{\sigma},
 \qquad
 \mathrm{Re}=\frac{\rho UD}{\mu},
 \qquad
 \mathrm{Oh}=\frac{\sqrt{\mathrm{We}}}{\mathrm{Re}}
 =\frac{\mu}{(\sigma\rho D)^{1/2}},
 \tag{10.1}
\end{equation}
其中 $D$ 和 $U$ 分别为液滴直径和撞击速度，$\rho,\mu$ 表示液滴，$\rho_a,\mu_a$ 表示环境气体或载液。一般实验撞击中重力相对惯性、表面张力和黏性可忽略；研究深池凹坑与气泡卷吸时还需 Froude 数
\begin{equation}
 \mathrm{Fr}=\frac{U^2}{gD}.
 \tag{10.2}
\end{equation}
此外可能需要表面相对粗糙度、润湿参数以及液滴直径与液膜厚度之比。

% [原PDF p243]
真实碰撞前，两界面之间的气膜必须从微米厚排到约 $1$--$10\,\mathrm{nm}$ 才会破裂。典型三维自适应网格只能到约 $100\,\mu\mathrm m$，二维约 $10\,\mu\mathrm m$，离分子尺度仍差数个数量级。前沿追踪通常在指定的物理合理时刻重联，VOF 与水平集则在两界面相距一两个单元时自动合并；二者都不能预测真实排液和破膜时间。图 10.2 的三维碰撞采用人为重联后继续求解宏观拉伸与分离，必须把这一步理解成次网格模型。
前沿追踪正碰模拟采用第 7 章的守恒表面力算法，保证闭合液滴的毛细合力为零。两滴被压扁时，接触区外缘形成沿径向扩展的薄“撞击足”；Weber 数越高，它越薄。液滴达到最大铺展半径 $R_{\max}$ 后回缩并可能反弹。若简单令初始动能
$K=\pi\rho U^2D^3/12$ 与扁饼状液滴表面能
$E_S=2\pi\sigma R_{\max}^2$ 平衡，得到
\begin{equation}
 R_{\max}\sim D\,\mathrm{We}^{1/2}.
\end{equation}
其他考虑黏性耗散或边界层的理论给出不同指数，而实验数据散布和参数范围有限，结论并不统一。DNS 可系统扫描参数并直接检查能量通道，但真正的渐近大 $\mathrm{Re}$、大 $\mathrm{We}$ 区往往已进入液滴破碎和液冠飞溅，不能把低速铺展标度无限外推。

% [原PDF p244]
图 10.2 展示三维偏心液滴碰撞：冲击参数 $I=\chi/D$ 由两条平行入射轴线的距离 $\chi$ 定义。小 $I$ 接近正碰，大 $I$ 会产生擦碰；液滴短暂连接后被角动量拉成长条，再分成两个或多个子滴。图 10.3 则比较液滴撞击深液池的计算与实验。接触后先形成凹坑，毛细波沿凹坑侧壁传播并在轴线上聚焦；不同撞击速度可导致无卷吸、规则气泡卷吸或形成 Worthington 射流。
\section{更复杂的缓慢撞击}

% [原PDF p245]
图 10.3 的时序表明，凹坑闭合可夹入气泡并向上发射细射流。图 10.4 对照 MacIntyre 的自由表面气泡爆裂素描、实验和数值结果：表面膜消失后留下近似圆锥凹坑，底部曲率集中并喷出高速小液滴。虽然激励不同，液滴撞池、气泡爆裂和大振幅 Faraday 波都出现相似锥形；这提示临界局部动力学主要由惯性与毛细力决定，而对远场初始几何不甚敏感。
环境压力下的液滴碰撞通常至少发生部分聚并。两滴或液滴与液池之间的气膜需要排出并最终破裂；严格解析这一过程，网格必须细于相关分子物理尺度，约 $1$--$10\,\mathrm{nm}$。典型三维自适应计算只能达到约 $100\,\mu\mathrm m$，二维约 $10\,\mu\mathrm m$，离真实膜破裂仍差三到四个数量级。

实践中使用两种简化。前沿追踪法在预设、物理上合理的时刻重联界面；VOF 和水平集法则在两界面相距一到两个网格时自动重联。两者都不能预测微观排液与破膜时间，却常能在破膜之后给出合理的宏观运动。必须明确：这种“自动聚并”是网格模型，不是已解析的微观物理。

% [原PDF p246]
图 10.5 将临界时刻附近的凹坑轮廓按 $\ell(t)=(\sigma/\rho)^{1/3}(t-t_0)^{2/3}$ 缩放，显示近似自相似圆锥。高速撞击时，液体在接触区周围抬起为向外扩张的环形薄膜，称液冠；顶端出现手指时呈皇冠状，甩出液滴则为飞溅。图 10.6 定义薄液层撞击的初始液滴直径、膜厚和间隙，图 10.7 给出低分辨率轴对称计算。后续分析集中在薄膜撞击，因为它同时包含极早射流和较晚液冠。
偏心碰撞以冲击参数
\begin{equation}
 I=\frac{\chi}{D}
\end{equation}
描述，其中 $\chi$ 为两条入射平行轴线的间距。除正碰的四种结果外，较大 $I$ 还产生擦碰：液滴暂时聚并并被拉成长条，随后再分成两个或多个液滴。结果取决于惯性、毛细回缩、黏性耗散以及接触后角动量的共同作用。

液滴撞击深池时，接触后先形成凹坑；毛细波沿凹坑侧壁传播并在轴线上汇聚，气泡卷吸前凹坑常出现近似圆锥形状。相似的锥形也出现在自由表面气泡爆裂和大振幅 Faraday 波中。临界时刻附近，水平和竖直距离以
\begin{equation}
 \ell(t)=\left(\frac{\sigma}{\rho}\right)^{1/3}(t-t_0)^{2/3}
\end{equation}
缩放后界面轮廓趋于自相似。这反映惯性--毛细平衡，而非具体激励方式。

% [原PDF p247]
图 10.7 的参数为 $\mathrm{Re}=1000$、$\mathrm{We}=8000$、$\rho/\rho_a=500$、$\mu/\mu_a=20$、$H/D=0.15$，网格每滴径约 102 点。液滴先与液层数值重联并卷入少量气泡，随后出现径向射流，最后形成竖直液冠并从顶部脱落液滴。由于这一分辨率看不到撞击最初仍存在的薄气膜，射流出现得偏晚。把环境气体密度、黏度取得高于真实空气有助于稳定压力求解，但也可能改变早期气膜排出，参数解释时必须说明。
\section{液冠、皇冠与飞溅撞击}

撞击速度足够高时，液体在撞击点周围抬升成近似竖直、向外扩张的环形薄膜，称为液冠（corolla）。若液冠顶部形成手指，则呈皇冠状；若手指或边缘甩出液滴，便构成飞溅。深池、薄液膜和干壁上的飞溅存在明显差异，以下重点讨论数值研究最充分的薄液膜撞击。

% [原PDF p248]
图 10.8 把分辨率提高到约 $D/\Delta x=804$，显示在液滴与液膜真正接触前，交汇区已喷出极薄射流，其厚度仍只有约两个网格。图 10.9 在 $\mathrm{Re}=100$ 时射流较钝且形成较晚；图 10.10 的 $\mathrm{Re}=40$ 则不再出现射流。黏性因此控制内层厚度和早期喷射阈值。高分辨率图中的气层形状近似液滴近壁边界层，但其最小厚度仍未收敛，不能据此声称已解析纳米排液。
\subsection{薄液层上的撞击}

典型轴对称 VOF/CSF 模拟以液滴直径 $D$、膜厚 $H$ 和初始间隙定义几何。一个代表算例取
$\mathrm{Re}=1000$、$\mathrm{We}=8000$、$\rho/\rho_a=500$、
$\mu/\mu_a=20$、$H/D=0.15$。这可对应直径 $6\,\mathrm{mm}$、速度约
$8.95\,\mathrm{m/s}$ 的甘油--水液滴；为稳定计算，环境气体密度和黏度通常取得高于真实空气。

每个直径约 $102$ 个网格的中等分辨率下，液滴与液膜很快重联并卷入少量气泡，稍后才形成射流，继而发展为竖直液冠并在顶部脱落液滴。把分辨率提高到约 $804$ 点/直径后，图景显著改变：射流在极早时刻便从液滴底部与液膜交汇区喷出；在此之前，液滴与液层之间仍存活一层很薄的气膜。初生射流只有约两个网格厚，说明即便该分辨率仍可能低估更早、更薄的射流。

% [原PDF p249]
图 10.11 将早期撞击分成外部势流区与内部黏性区。假想未变形液滴与液面交点的径向位置满足 $r_J\sim\sqrt{DUt'}$，黏性层厚度按 $\sqrt{\nu t'}$ 增长并决定初生射流。匹配分析给出液冠形成判据 $\mathrm{We}^{1/2}\mathrm{Re}^{1/4}>K_c$，常数随膜厚比 $H/D$ 变化；实验在 $H/D>0.08$ 时转变值约 40。图中涡量从极窄接触区生成并随射流带出，解释了为何粗网格会完全遗漏即时飞溅。
增大液体黏性会推迟射流形成并使端部更圆。$\mathrm{Re}=100$ 时仍有较钝的射流，降至约 $40$ 时可完全不出现射流。匹配外部势流区与内部黏性区的标度理论认为，射流基部靠近未变形液滴与未变形液面假想交点，其径向位置满足
\begin{equation}
 r_J\sim\sqrt{DUt'},
\end{equation}
内部黏性层按 $\sqrt{\nu t'}$ 增长，并同时决定初生射流厚度。理论预言形成液冠的条件为
\begin{equation}
 \mathrm{We}^{1/2}\mathrm{Re}^{1/4}>K_c,
 \tag{10.3}
\end{equation}
$K_c$ 依赖 $H/D$。实验对 $H/D>0.08$ 给出的转变值约为 $40$。这一无量纲条件与干壁撞击的经验式可相互换算。

% [原PDF p250]
图 10.12 是直径约 $6.5\,\mathrm{mm}$ 的甘油滴撞击后形成螺旋液冠的实验细节；其 Reynolds 数约 1900。图 10.13 给出更接近实验物性、不同气相黏度比的轴对称计算，局部放大中的点数更多，说明自适应加密集中于早期射流。实验确实看到液滴底部立即伸出的极薄液冠，随后才发展出较粗螺旋结构。对低黏水滴，射流很快三维失稳并形成喷雾，轴对称模型只能描述最初瞬间。
实验确实观察到撞击后立即从液滴底部伸出的极薄液冠，并在稍后形成螺旋状结构。要在真实气液密度比、黏度比下解析它，需要局部极强的自适应加密。水滴黏性更低，早期液冠会很快失稳成复杂喷雾，使轴对称假设迅速失效。

\subsection{三维撞击}

% [原PDF p251]
三维计算在初始液滴上施加 $r=D/2+r_1\sin\theta\cos(N\phi)$ 的周向扰动。图 10.14 取 $N=4$；液冠初成时凹凸发生反相，顶端位置近似 $r_{\rm tip}=r_0-At\cos(N\phi)$。图 10.15 解释这种 Richtmyer--Meshkov 型机制：从重流体指向轻流体的突然加速度使速度场重排，初始波纹先反相再近线性增长。进入非线性后相邻手指拥挤、合并，优势周向波长增大。粗网格还会产生不真实的微小液滴，应从统计中剔除或通过收敛检查识别。
三维计算使用与轴对称算例相同的物理模型，但计算量大幅上升。为辨识非轴对称失稳，可在初始近球形液滴上施加小扰动
\begin{equation}
 r(\theta,\phi)=\frac D2+r_1\sin\theta\cos(N\phi),
 \tag{10.4}
\end{equation}
其中 $N$ 为周向波数。射流刚形成时，初始凹凸在液冠中发生反相，液冠顶端位置可写成
\begin{equation}
 r_{\rm tip}(\phi)=r_0-A t\cos(N\phi).
 \tag{10.5}
\end{equation}
扰动先反相、再近似线性随时间增长，是 Richtmyer--Meshkov 型过程的特征：界面突然受到从重流体指向轻流体的加速度，速度场瞬时重排，把初始波纹放大。非线性阶段相邻手指拥挤、合并，周向波长增大而手指数减少。

\begin{sourceblock}
\sourcecrop{252}{75}{430}{419}{638}{0.49}
\captionof{figure}{甘油液滴撞击薄水层的三维计算。液冠从近轴对称环形射流发展出周向手指并形成皇冠；局部放大显示早期射流和液滴脱落。参数为 $\mathrm{Re}=600$、$\mathrm{We}=443$、$\rho/\rho_a=1000$、$H/D=0.31$。}
\end{sourceblock}

% [原PDF p252]
图 10.14 的甘油滴薄水层撞击参数为 $\mathrm{Re}=600$、$\mathrm{We}=443$、$\rho/\rho_a=1000$、$H/D=0.31$，显示近轴对称液冠逐渐长出四重手指。图 10.15 用示意图表示界面突然加速后的扰动反相。早期手指和液滴在液冠高度尚小于一滴径时出现，称为即时飞溅；更晚、液冠远高于初始滴径时的破裂称为皇冠飞溅。两者在大 Reynolds、Weber 数下相隔多个数量级，不能用同一个线性阶段概括。
这种早期手指和液滴出现在液冠尚未升到一个液滴直径之前，称为“即时飞溅”（prompt splash）。较晚时刻，当液冠高度已远超初始液滴直径，出现“皇冠飞溅”（crown splash）。晚期液冠很薄，低黏性情况下连接边缘后的细颈最易穿孔；边缘可能脱离为液环，再发生 Plateau--Rayleigh 失稳。实验也观察到带孔液冠，说明这不完全是数值伪影，但模拟中的穿孔时刻仍受网格尺度控制。

% [原PDF p253]
图 10.16 左侧是与图 10.14 相近条件、但初始扰动更小的液冠；右侧为实验中的带孔液冠。晚期薄冠常在连接膨大边缘后的细颈处穿孔，边缘可脱离成液环并经历 Plateau--Rayleigh 失稳。真实实验也观察到穿孔，所以现象并非纯数值伪影，但计算中的穿孔时刻仍由网格最小膜厚触发。若液冠不穿孔，晚期手指可能来自 Taylor--Culick 边缘失稳。即时飞溅与皇冠飞溅的时间可差约 $10^2$--$10^3$，早期扰动是否保持到晚期尚无定论。视觉相似必须辅以气膜厚度、射流厚度、穿孔时刻和液滴谱的网格收敛。
如果液冠不穿孔，晚期手指更可能由第 9 章讨论的边缘失稳产生。因此液冠失稳至少有两个阶段：早期由撞击加速度引发的 Richtmyer--Meshkov 型放大主导，晚期由 Taylor--Culick 边缘及其毛细失稳主导。在大 Reynolds 数和大 Weber 数下，两阶段在时间、长度尺度上可相差数个数量级，早期扰动是否会一直“记忆”到晚期皇冠仍不清楚。低黏水滴的飞溅甚至未必可描述为某个清晰基态上的线性失稳，而可能是手指、膜片和液滴不断改变尺度的强非线性结构。

对这类问题，视觉相似不足以证明数值收敛。必须同时检查界面最小厚度、气膜分辨率、射流厚度、穿孔时刻、液滴尺寸统计以及数值耗散对这些量的影响；随着网格加密持续出现更小结构，是欠分辨雾化/飞溅的典型信号，而不是可以忽略的细节。
% ===== 译文分片 4（已合并）=====

% Tryggvason 等《气液多相流直接数值模拟》中文译稿，第 254--337 页
% 第 11 章、附录 A--D、参考文献与中文主题索引。

% [原PDF p254]
\chapter{扩展问题}
% [补译 p254]
本章的目的不是列出单流体方法的一切新应用，而是说明基本框架如何随着物理复杂度增加而扩展。新增方程必须与已有速度、压力和界面位置在同一时间精度下耦合；新增界面效应则必须进入正确的法向或切向跳跃条件。只把一个源项附加到体方程而不检查界面守恒，往往会得到网格依赖结果。反过来，若界面条件、物性插值和分布核保持一致，前几章的固定网格求解器可以在不改变总体结构的情况下容纳热、电、组分和相变等过程。真正的困难常不是多出一个未知量，而是该未知量拥有比速度场更薄的边界层，或者引入新的未解析微观尺度。

新问题通常由新的场方程、新的界面效应或两者共同构成。增加场方程看似是最简单的推广，但它会带来新的时间尺度和长度尺度。例如液体中的组分扩散层可能远薄于热边界层和动量边界层，所以原本足以解析速度的网格未必能给出可靠传质率。传热是最直接的例子：在求温度对流--扩散方程时，若两相导热率差别很大，界面附近需要更细的网格，而且未定义在同一点的导热系数宜像大黏度比问题那样使用调和平均，以保持法向热通量的一致性。


多相流直接数值模拟正在迅速发展。除了继续改进数值技术并更深入地研究前述问题，研究者也越来越多地转向包含复杂物理的新问题。本章简要介绍若干扩展方向，目的不是穷尽单流体表述的所有新应用，而是为读者进入相关文献提供线索。

\section{附加场与界面物理}

广义地说，新物理由新的场方程、新的界面效应或两者共同组成。增加一个场方程看似简单，却常引入新的时间与长度尺度。例如液体中的传质边界层，可能远薄于热边界层和动量边界层，其分辨率要求会比不含传质的同一流动严格得多。

最简单的扩展是传热：求解温度的对流-扩散方程。跨界面导热系数差别很大时，需要细网格，并且像处理大黏度比一样，在导热系数没有直接定义的位置采用调和平均通常更合适。下面重点讨论比单纯传热更复杂的效应。

% [原PDF p255]
\subsection{热毛细运动}
% [补译 p255]
表面张力通常随温度升高而降低，因此界面温度不均匀会产生切向 Marangoni 应力，并由黏性作用传给两侧流体。浸在温度梯度中的气泡或液滴一般向热端迁移。在地面条件下这种效应常被浮力掩盖，在微重力环境却可能占主导；它既可帮助从凝固前的熔体中移走气泡，也会使气泡聚集在两相换热器热壁上而增强隔热。数值实现中把温度插值到界面，以实验或给定的 $\sigma(T)$ 求表面张力，再计算曲面梯度即可。


在两流体传热中加入温度依赖的表面张力，是最直接的扩展。表面张力通常随温度升高而减小；界面上的非均匀表面张力产生切向应力，再通过黏性作用传给两侧流体。因此，处在温度梯度中的气泡或液滴通常向高温区迁移。在地面重力条件下热毛细力常被浮力掩盖，但在微重力环境中可以占主导。空间材料加工可利用它在凝固前从熔体中移除气泡；空间两相换热器设计也必须考虑热壁上聚集气泡的隔热作用。

单流体算法推广到变表面张力并不困难：把温度插值到界面，由给定的 $\sigma(T)$ 计算表面张力，再把
\begin{equation*}
 \bm f_\sigma=\sigma\kappa\bm n+\nabla_S\sigma,
 \qquad \nabla_S\sigma=\frac{\mathrm d\sigma}{\mathrm dT}\nabla_ST
\end{equation*}
加入界面力。界面追踪、VOF 等方法均已用于二维和三维单液滴及多液滴热毛细迁移。

\subsection{电流体动力学}

对含互不相溶界面的多流体系统施加电场，会因自由电荷向界面迁移以及两相介电性质不匹配而产生界面力。理想介电体或理想导体中该力可为纯法向；一般漏电介质中还存在切向 Maxwell 应力，进而通过黏性应力驱动流动。

Taylor 漏电介质近似下，首先求解电势的变系数 Poisson 方程，再由 Maxwell 应力张量的散度计算体力或界面跳跃：
\begin{equation*}
 \nabla\cdot(\epsilon\nabla\phi)=0,
 \qquad \bm E=-\nabla\phi,
 \qquad \bm T_E=\epsilon\left(\bm E\bm E-\frac12|\bm E|^2\bm I\right).
\end{equation*}
% [原PDF p256]
% [补译 p256]
\begin{quote}\small 图 11.1　电场作用下通道内液滴悬浮体系的一个瞬时画面。通道上下为刚性壁，水平方向周期；液滴电导率低于连续相且剪切较弱，异号端部电荷的吸引使液滴形成贯穿通道高度的长链。经 Fernandez（2008b）许可转载。\end{quote}
Fernandez 的研究后来推广到更宽参数范围和三维液滴。Hua 等用界面追踪法研究电场诱导的单滴变形；Tomar 等的 VOF/水平集方法强调，当两相电性质相差很大时必须以调和平均插值，否则电通量与 Maxwell 应力都会失真。这些方法也已用于电场影响膜态沸腾的问题。电解质界面还存在电毛细效应，其数值处理与温度依赖表面张力很相似。

电性质相差很大时，对介电常数或电导率的界面插值也宜采用与扩散通量一致的调和平均。模拟表明，在通道内电场作用下，液滴端部异号电荷相互吸引，可形成跨越通道高度的长链。电场还可显著改变液滴变形和膜态沸腾。对电解质界面，电场还会改变表面张力；数值处理与温度依赖表面张力相似。

% [原PDF p257]
\subsection{传质与化学反应}
% [补译 p257]
气--液或液--液体系中的反应可以发生在任一体相内，也可以局限于界面；它们还可能依赖悬浮催化颗粒，产物甚至会构成新相。完整过程一般同时包含跨界面传质、表面吸附与解吸、组分在主体相中的扩散，以及反应物向催化颗粒内部的渗透。在 Fischer--Tropsch 合成和硝基芳烃加氢中，解吸与颗粒内渗透相对快，液相中的气体扩散往往成为限速步骤。反应集中在气泡邻域时，结果对尾迹混合高度敏感，因而 DNS 特别适合揭示控制机理。


许多多相流装置的最终目标是化学反应：液体燃料雾化是为了燃烧，气泡柱本质上是化学反应器。完整模型除流动外还要包含组分传质与反应。反应可发生在任一体相或界面上，也可能依赖悬浮催化颗粒，反应产物甚至会形成新的相。常见环节包括跨界面传质、界面吸附/解吸、体相扩散以及反应物向催化颗粒内部渗透。

在 Fischer--Tropsch 合成和硝基芳烃催化加氢等气泡柱过程中，解吸和进入催化颗粒可能很快，而气体在液相中的扩散决定总速率。反应集中在气泡附近并高度依赖局部混合时，DNS 尤其有价值。单个固定气泡的研究表明，尾迹细节会显著改变反应产物；大而可变形的气泡若形成闭合尾迹，会长时间困住不希望出现的中间产物。多气泡自由运动时，这种敏感性仍然存在，而且参数依赖更复杂。

% [原PDF p258]
% [补译 p258]
\begin{quote}\small 图 11.2　硝基芳烃催化加氢模拟。图（a）、（b）为氢浓度，图（c）、（d）为羟胺浓度；（a）、（c）的反应相对传质较慢，（b）、（d）的反应相对较快。引自 Radl 等（2008）。\end{quote}
羟胺形成很快，而最终生成芳胺的步骤较慢，因此闭合尾迹中容易积聚中间产物。尾迹里的羟胺若接触不到新的溶解氢，就不能继续反应。反应较慢时主体液相仍保存相当多未反应氢；反应较快时氢在气泡附近即被消耗，主体浓度几乎为零，中间体分布也随之改变。显著尾迹存在时，数值选择性与忽略局部混合的简单膜理论明显不同。

硝基芳烃加氢的主要困难是羟胺中间体生成很快，而继续加氢为芳胺较慢，容易造成羟胺积累。羟胺必须与含溶解氢的液体充分混合；气泡后方缺氢的闭合尾迹尤其不利。反应相对较慢时，主体液相仍含较多未反应氢；反应较快时，氢在气泡附近就被消耗，主体浓度接近零，同时中间体分布也发生显著改变。由于传质边界层比流动结构更薄，相关 DNS 常对组分方程使用比流场更细的网格。

% [原PDF p259]
% [补译 p259]
早期上升气泡实验测得的速度更接近固体球，而不是洁净气泡的 Hadamard--Rybczynski 解，原因通常是微量表面活性剂使界面近似不可动。Frumkin--Levich 模型令 $\sigma$ 依赖界面浓度 $\Gamma$；$\Gamma$ 增大通常降低表面张力，非均匀浓度产生的 Marangoni 力倾向于使分布变均匀，而流动对流又可能与其抗衡。对不溶性活性剂，忽略表面扩散时每个物质面元携带的总量守恒；可溶性情形还必须耦合吸附、解吸和液相扩散。

即使没有化学反应，传质也会反过来强烈影响流动。早期气泡上升实验的速度更接近固体球而不是洁净气泡的解析结果，原因通常是表面活性剂污染使界面近似不可动。Frumkin--Levich 模型令表面张力依赖界面活性剂浓度 $\Gamma$；浓度增大会降低 $\sigma$，非均匀 $\Gamma$ 产生 Marangoni 力，使界面输运趋于均匀。流动对流又可能抵消这一趋势，形成非均匀稳态分布。

不溶性表面活性剂始终停留在界面上；忽略表面扩散时，一个物质面元所携带的总量守恒，因此浓度与面元面积成反比。一般情况下还要处理界面吸附、解吸以及液相中的扩散，这构成更困难的全耦合传质问题。

\subsection{沸腾}

许多重要系统包含相变。沸腾是从固体表面移除热量的高效方式，广泛用于发电和热表面冷却；但巨大的汽液体积变化和高温也会放大设计或运行错误的后果。

沸腾通常从热壁成核，也可从液体内部悬浮核开始。成核沸腾中，蒸气泡在细小凹穴等明确成核点形成并脱离。热流更高时，气泡合并成覆盖壁面的蒸气膜，形成膜态沸腾；蒸气泡会从膜层脱落，而界面蒸发继续补充蒸气膜。蒸气层隔热，通常降低传热并抬高壁温，因此膜态沸腾一般不希望出现。静止液池中的过程称池沸腾，流动液体中的过程称流动沸腾或强制对流沸腾。降压也能使液体汽化，但通常单独归入空化；冷凝则是相反的热力学过程。混合液体的不同沸点和气相中的不可凝气体都会改变相变。

% [原PDF p260]
% [补译 p260]
% [覆盖增强 p260]
池沸腾指初始近静止液池中的相变，流动沸腾或强制对流沸腾则有显著来流。膜态沸腾时气泡从蒸气层脱落，而界面蒸发不断补充薄层；由于蒸气导热差，这层膜会降低热流并抬高固壁温度。Welch 的移动三角网格、Son 与 Dhir 的贴体坐标只能跟随有限变形，后来固定网格界面方法才允许界面断裂、合并和长时间演化。现有膜态沸腾算法包括界面追踪、VOF、水平集和耦合水平集/VOF，而成核沸腾还必须给出壁面成核与接触线闭合。

降压同样会使液体汽化，但一般另称空化；冷凝则是在移走热量时蒸气转为液体的逆过程。混合液体的不同沸点和气相中不可凝气体都会改变相变。当蒸气分压只占总压力很小部分、液体转化引起的局部膨胀在动力学上可忽略时，通常称为蒸发而非沸腾。早期移动网格计算只能跟随半球形气泡或中等变形很短时间；界面追踪、水平集、VOF 及耦合方法解除这一限制，使膜态沸腾可长期计算，但成核点和微液膜仍使成核沸腾更困难。

早期计算常假定气泡保持半球形，或用移动贴体网格只能跟踪有限变形。界面追踪、水平集、VOF 及耦合水平集/VOF 消除了这一限制，已能可靠计算膜态沸腾。成核沸腾还涉及成核点、固体导热、微小初生泡和微液膜，进展相对较慢。

% [原PDF p261]
% [补译 p261]
相变计算除质量和动量守恒外，还必须求两相温度、规定相界温度、计算液体和蒸气之间的质量转化率，并在不可压缩两相模型中表示体积生成。式（11.1）把潜热写成集中于界面的源项。它可以平滑到固定网格上直接求解，也可修改界面附近离散并显式施加界面温度；汽液密度比很大时后一做法通常更可靠。局部热力学平衡允许令温度跨界面连续，并近似等于系统压力下的饱和温度。法向探针在两侧离界面约 $1.5$ 个网格距处插值温度，再与已知界面温度构造单边梯度。

除质量和动量方程外，沸腾必须求解能量方程，规定界面温度，计算液-汽转化率，并计入相变体积源。能量方程的单场形式为
\begin{equation}
 \frac{\partial(\rho cT)}{\partial t}+\nabla\cdot(\rho cT\bm u)
 =\nabla\cdot(k\nabla T)-\dot q\,\delta(n). \tag{11.1}
\end{equation}
$\dot q$ 是界面单位面积潜热源。可把它平滑到固定网格上直接求解，也可修改界面附近离散，显式强制界面温度；汽液密度比大时，后者往往更合适。

多数情形可假设界面局部热力学平衡，温度跨界面连续，并近似等于系统压力下的饱和温度。界面热通量跳跃为
\begin{equation}
 \dot q=\left.k_v\frac{\partial T}{\partial n}\right|_v
 -\left.k_l\frac{\partial T}{\partial n}\right|_l, \tag{11.2}
\end{equation}
温度梯度可用“法向探针”求得：沿界面法向在两侧取约 $1.5$ 个网格间距处插值温度，再与已知界面温度构造单侧梯度。质量相变率为
\begin{equation*}
 \dot m=\frac{\dot q}{h_{lv}},
\end{equation*}
其中 $h_{lv}$ 为汽化潜热。

% [原PDF p262]
% [补译 p262]
潜热通量除以汽化潜热给出单位面积质量通量 $\dot m=\dot q/h_{lv}$。虽然每一相内部都可保持常密度、无散，合成速度却在界面上有散度集中项。Rankine--Hugoniot 条件说明两侧相对于移动界面的质量通量必须相等；消去界面速度后得到法向速度跳跃 $u_v-u_l=\dot m(1/\rho_v-1/\rho_l)$，再乘界面分布即得式（11.6）。界面推进速度不能简单取某一侧流速，而应取两侧平均并减去相变造成的相对运动；在 $\rho_v\ll\rho_l$ 时该修正尤其显著。

两相内部密度分别取常数，体积膨胀只发生在界面。写
\begin{equation*}
 \bm u=\bm u_vH+\bm u_l(1-H),
\end{equation*}
并令两侧延拓速度各自无散，则
\begin{equation}
 \nabla\cdot\bm u=(u_v-u_l)\delta(n). \tag{11.3}
\end{equation}
Rankine--Hugoniot 条件为
\begin{equation}
 \rho_v(u_v-V)=\rho_l(u_l-V)=\dot m. \tag{11.4}
\end{equation}
消去界面速度 $V$ 得单位界面面积的体积生成率
\begin{equation}
 u_v-u_l=\dot m\left(\frac1{\rho_v}-\frac1{\rho_l}\right), \tag{11.5}
\end{equation}
因此单流体速度满足
\begin{equation}
 \nabla\cdot\bm u=\dot m\left(\frac1{\rho_v}-\frac1{\rho_l}\right)\delta(n), \tag{11.6}
\end{equation}
界面法向速度为
\begin{equation}
 V=\frac12(u_v+u_l)-\frac{\dot m}{2}\left(\frac1{\rho_v}+\frac1{\rho_l}\right). \tag{11.7}
\end{equation}
第一项是两侧流体平均法向速度，第二项是相变导致的界面相对运动。

相内温度造成的密度变化可用 Boussinesq 近似产生浮力。若界面温度近似常数，热毛细效应通常很小；需要时也可按前述方法加入。黏性耗散和界面伸展产生的能量通常远小于潜热，常被忽略。

% [原PDF p263]
% [补译 p263]
\begin{quote}\small 图 11.3　二维膜态沸腾模拟。热壁上方覆盖薄蒸气膜；界面逐渐失稳，蒸气泡周期性脱落并上升。左列表示相界面，右列表示温度场。引自 Esmaeeli 与 Tryggvason（2004b）。\end{quote}
算例在水平方向使用周期边界，底部热壁无滑移，顶部开放。最初仅受轻微扰动的蒸气膜在持续蒸发作用下形成长波、颈缩和气泡脱落；脱落留下的空缺又由界面蒸发补充。由于最薄蒸气层控制总热阻，竖向采用非均匀加密网格，把分辨率集中在壁面附近。

二维膜态沸腾模拟中，静止液池位于热水平壁上方，壁面被薄蒸气膜覆盖。蒸发使汽液界面失稳，气泡周期性从膜层脱落、上升并穿出自由表面。水平周期、底部无滑移、顶部开放；为分辨薄膜，竖向网格可非均匀加密。壁面平均传热常用 Nusselt 数
\begin{equation*}
 Nu=\left.\frac{\lambda}{T_w-T_{\mathrm{sat}}}\frac{\partial T}{\partial y}\right|_{y=0}
\end{equation*}
表示，并随无量纲壁面过热度 Jacob 数变化。二维 DNS 与已有三维实验关联式总体符合良好。

% [原PDF p264]
% [补译 p264]
\begin{quote}\small 图 11.4　平均 Nusselt 数 $Nu$ 随无量纲壁面过热度 Jacob 数 $Ja$ 的变化，二维 DNS 与理论及实验关联式比较。\end{quote}
底壁局部热流经空间和时间平均得到图中的 $Nu$。Jacob 数衡量液体显热相对于汽化潜热的大小，壁面过热改变蒸发速度、蒸气膜厚度和气泡脱落周期，因此平均传热并非单纯导热。尽管二维气泡几何不同于实验的三维体系，所得无量纲传热量级与趋势仍较一致。成核沸腾还需给定成核点活化规律、求解固体导热、处理远小于网格的初生泡，并以亚网格模型描述接触线后方微液膜的蒸发。

成核沸腾还需解决四个核心问题：成核点的位置和活化规律取决于表面粗糙度及过热度；气泡生成会改变固体内部温度，原则上必须求固体能量方程；初生泡尺度极小，即便 AMR 也难完整分辨，以有限尺寸气泡重启成核只是近似；气泡扩张时接触线后方留下极薄微液膜，其蒸发可能显著影响初期生长，通常只能用薄膜亚网格模型表示。尽管如此，多个成核点的三维模拟已经显示，DNS 有望显著提升对复杂沸腾系统的认识和预测能力。

% [原PDF p265]
\subsection{空化}
% [补译 p265]
\begin{quote}\small 图 11.5　四个成核点上成核沸腾模拟的三个时刻。气泡在热壁上生长、相互作用并脱离，同时显示固体和流体中的温度场。引自 Shin 等（2005b）。\end{quote}
成核位置由表面微凹穴、润湿性和局部过热共同决定，并不是连续介质方程自动给出的量。气泡生长还会冷却成核点附近的固体，所以原则上必须联立固体能量方程。初生泡及接触线后的微液膜远小于主体网格；以有限可分辨气泡重新启动成核会引入人为尺度，而微液膜蒸发又可能支配早期增长。这些困难解释了为什么成核沸腾的分辨模拟晚于膜态沸腾。


空化可视为降压触发的特殊沸腾，但常使用更简单模型：泡内压力取系统温度下的饱和蒸气压，或按理想气体定律变化。已有模拟研究壁面附近汽泡在剪切中的坍缩、黏性对单泡坍缩的影响以及含大量气泡液体中的压力波传播。

% [原PDF p266]
% [补译 p266]
\begin{quote}\small 图 11.6　压力波向含大量可空化气泡的区域传播的一个瞬时画面。上边界压力突然升高后，上层气泡先坍缩，继而触发下方气泡，形成向下传播的波前。\end{quote}
空化模型常把泡内压力固定为系统温度下的饱和蒸气压，或者把不可凝气体按理想气体处理，使压力随泡体积变化。壁面附近气泡在剪切中的坍缩、黏性对单泡坍缩的影响以及气泡云中的压力波均已得到研究。波速有限并非数值耗散的产物，而是每一层气泡从受压到坍缩都需要有限时间；计算得到的传播速度和波形与平均理论预测相符。

若区域初始压力等于蒸气压，顶部压力突然升高，最上方气泡先坍缩并提高局部压力，继而触发更深位置气泡坍缩。由于每个气泡坍缩需要有限时间，压力扰动以有限速度向下传播，所得传播速度和激波结构与理论预测相符。

% [原PDF p267]
\section{浸入边界}
% [补译 p267]
Peskin 的原始浸入边界法为心脏运动建模而提出：代表心壁的一维弹性纤维束以 Lagrangian 标记点离散，弹性力通过光滑离散 $\delta$ 核分配到均匀 Eulerian 流体网格，流体速度再用同一核插回标记点推动纤维。流体密度和黏度在全域保持不变，因而其结构与界面追踪法相似而物性跳跃更简单。心壁纤维有限伸长，标记点通常不会无限聚集或稀疏，所以不像强变形流体界面那样频繁重网格化。该思想后来从二维发展到完整三维，并广泛用于生物流体问题。


把静止固体周围的流动称作多相流稍显宽泛，但近年来，规则固定网格处理复杂固壁重新受到重视，其发展与流体界面方法十分平行。下面先介绍 Peskin 针对弹性移动边界的浸入边界法，再讨论任意固体边界。

\subsection{Peskin 浸入边界法}

原始浸入边界法为心脏运动建模而设计，后来从二维推广到全三维。它与界面追踪法相似，但流体密度和黏度处处为常数；代表心壁的弹性固体由一维纤维束组成，纤维以标记点离散。有限伸长使标记点通常不需重分布，因此其几何处理比会强烈变形并需要重网格化的流体界面简单。

早期实现多用同位网格和一阶时间推进，后来发展了高阶积分。Peskin 方法推动了众多生物流体模拟。虽然均匀流体和小变形边界缺少多流体中的许多困难，但其核心思想深刻影响了现代单场方法；就历史影响而言，仅次于 MAC 和 VOF。

% [原PDF p268]
\subsection{固体边界}
% [补译 p268]
固定 Cartesian 网格处理固壁的早期做法，是把落入固体的单元直接遮蔽。对曲线壁面若简单取最近网格面，会形成低分辨率阶梯边界；贴体和非结构网格因而曾占主导。随着规则网格的高效率、局部加密能力和浸入边界思想发展，任意固壁方法重新兴起。这里的 $I(\bm x)$ 是纯几何指示函数，而非物性场：流体内为一，固体内为零。全域速度 $\bm u=I\bm u_f+(1-I)\bm u_S$ 把满足 Navier--Stokes 方程的流体解与给定固体速度拼接起来。


早期 CFD 就曾在矩形结构网格上屏蔽固体单元：先处理与网格对齐的矩形壁面，再用最邻近网格面逼近曲线边界。由于分辨率低、局部加密能力有限，这一路线一度让位于贴体和非结构网格。随着规则网格的高效性与复杂几何处理优势重新受到重视，固定 Cartesian 网格上的浸入固体方法再次兴起。

在含任意固体的区域定义指示函数
\begin{equation}
 I(\bm x)=\begin{cases}1,&\bm x\text{ 在流体内},\\0,&\bm x\text{ 在固体内},\end{cases} \tag{11.8}
\end{equation}
全域速度写为
\begin{equation}
 \bm u(\bm x)=I\bm u_f+(1-I)\bm u_S, \tag{11.9}
\end{equation}
其中 $\bm u_f$ 满足均匀流体 Navier--Stokes 方程
\begin{equation}
 \frac{\partial\bm u_f}{\partial t}+\nabla_h\cdot(\bm u_f\bm u_f)
 =-\frac{\nabla_hp}{\rho}+\nu\nabla_h^2\bm u_f, \tag{11.10}
\end{equation}
$\bm u_S$ 为给定固体速度。

% [原PDF p269]
% [补译 p269]
\begin{quote}\small 图 11.7　$Re=60$ 时圆柱周围的稳态速度。圆柱覆盖在固定 Cartesian 网格上，固体内部速度每个时间步被重置。\end{quote}
投影步骤先忽略固体得到预测速度，再以压力使最终全域速度无散。展开 $\nabla\cdot[I\bm u_f+(1-I)\bm u_S]$ 时，$\nabla I$ 在固壁产生集中项；当固体内速度为常量且法向速度跨壁连续时，这些奇异贡献相互抵消，因此仍可使用与无固体区域相同的压力 Poisson 方程。最简单实现只需在时间步结束时把固体点速度重置为 $\bm u_S$，即可得到合理的圆柱绕流。

一阶投影法的预测速度为
\begin{equation}
 \bm u_f^*=\bm u^n+\Delta t[-\nabla_h\cdot(\bm u^n\bm u^n)+\nu\nabla_h^2\bm u^n], \tag{11.11}
\end{equation}
再由
\begin{equation}
 \bm u_f^{n+1}=\bm u_f^*-\Delta t\nabla p \tag{11.12}
\end{equation}
修正。最终全域速度
\begin{equation}
 \bm u^{n+1}=I\bm u_f^{n+1}+(1-I)\bm u_S^{n+1} \tag{11.13}
\end{equation}
应无散。若固体内速度为常数且壁面法向速度连续，$\nabla I$ 产生的奇异项相互抵消，仍可使用与无固体区域相同的压力方程
\begin{align}
 \nabla\cdot\bm u^{n+1}
 &=\nabla\cdot(I\bm u_f^{n+1})+
   \nabla\cdot[(1-I)\bm u_S^{n+1}] \notag\\
 &=I\nabla\cdot\bm u_f^{n+1}+(1-I)\nabla\cdot\bm u_S^{n+1}
  +(\bm u_f^{n+1}-\bm u_S^{n+1})\cdot\nabla I. \tag{11.14}
\end{align}
\begin{equation}
 \nabla^2p=\frac{\rho}{\Delta t}\nabla_h\cdot\bm u^*. \tag{11.15}
\end{equation}
所以最简单实现只需在每步末把固体内速度重置为 $\bm u_S^{n+1}$。圆柱绕流等问题可用这种方法直接计算。

% [原PDF p270]
% [补译 p270]
% [覆盖增强 p270]
更一般地，可以把固体约束看成时间推进后的投影：先在包括固体的整个矩形域上计算一个临时速度，再在固体内部用目标速度替换它。替换所造成的动量变化除以时间步，就是等效的浸入体力。与显式反馈相比，这种直接强制没有额外松弛参数，也不会因反馈增益选择不当而振荡。然而，压力投影仍通过整个域传播约束，所以固壁附近的速度插值、指示函数宽度和压力边界的一致性决定最终精度，不能仅凭固体点“看起来静止”判断方法已经收敛。

指示函数若只在很窄的带内平滑，网格看到的壁面会呈锯齿状，速度和压力产生格点振荡；平滑带太宽则会模糊真实壁面位置。Goldstein 等提出反馈体力，使固体区计算速度逐渐趋于目标值。若目标速度已知，直接在每步把固体速度重置为 $\bm u_S$，等价于施加一个恰能完成这次修正的离散体力，所以不必另外积分反馈方程。该做法结构简单，但几何仅通过扩散指示函数进入，壁面附近误差通常使整体收敛阶不高于一阶。

实践中常把 $I$ 在固壁附近数个单元内平滑。过窄会造成锯齿速度，过宽会损失边界精度；非对称平滑有时效果更好。Goldstein 等提出反馈体力，使固体区计算速度趋于目标值；但若时间离散为
\begin{equation}
 \bm u^{n+1}=\bm u^n+\Delta t\left[-\nabla_h\cdot(\bm u\bm u)+\nu\nabla_h^2\bm u-\frac{\nabla_hp}{\rho}-\bm f_S\right], \tag{11.16}
\end{equation}
而目标 $\bm u_S$ 已知，则直接重置速度已经隐式包含所需 $\bm f_S$，无须显式计算反馈力。

% [原PDF p271]
% [补译 p271]
\begin{quote}\small 图 11.8　固体边界附近速度的插值。先沿壁面法向找到与邻近网格线的交点，在交点处由流体节点双线性插值，再结合壁面规定速度求靠壁节点值。\end{quote}
Fadlun 方法对所有邻点都在流体内的节点使用标准 Navier--Stokes 更新，对紧邻固体的流体节点则用壁面速度与更远流体值插值。沿坐标线处理便宜，却会因曲壁方向产生选择歧义；法向构造更几何一致。代价是靠壁节点的守恒方程被插值条件替代，不再严格局部守恒。幽灵单元方法反其道而行：保留流体节点方程，求固体侧外推值，使几何壁面上的插值准确满足边界条件。

直接平滑浸入边界通常至多一阶收敛。为获得二阶精度，应更准确地处理近壁节点。Fadlun 方法先正常更新所有邻点均在流体内的节点，再把与固体相邻的流体节点速度，由壁面目标速度和更远流体节点线性插值得到。若沿网格线插值方向含糊，可过待定点作固壁法线，在法线与最近网格线交点处先双线性插值，再沿法线插到待定点。这样能改善精度，但丢弃近壁节点的 Navier--Stokes 离散，不再严格局部守恒。另一思路是反过来设定固体内幽灵值，使插值得到的壁面速度准确，从而保留流体节点方程。

% [原PDF p272]
% [补译 p272]
自由运动刚体可以先赋予真实密度和很大的有效黏度，与流体一起推进，然后把其内部速度场投影到刚体运动空间。固体区域线动量给出质心平移速度，关于质心的角动量给出角速度，最终速度写成 $\bm u_{Sc}+\bm r\times\bm\Omega_S$。这种投影消除固体内部形变，并把流体作用传递给颗粒。颗粒间隙小于网格时，分辨流体压力不足以阻止几何穿透，必须加入短程接触或未分辨润滑力；弹性体则需求解相应的应变--应力关系，不能使用刚体投影。

自由运动刚体可赋予真实密度和很大的有效黏度；流场更新后，把每个刚体内部速度投影为刚体运动：
\begin{equation}
 \bm u_S^{n+1}=\bm u_{Sc}^{n+1}+\bm r\times\bm\Omega_S^{n+1},
 \qquad \bm r=\bm x-\bm x_c^{n+1}. \tag{11.17}
\end{equation}
平移与转动速度由固体区域内线动量和角动量平均确定：
\begin{equation}
 M_S\bm u_{Sc}^{n+1}=\int(1-I)\rho\bm u\,\mathrm dV,
 \qquad \bm I_S\bm\Omega_S^{n+1}=\int(1-I)\bm r\times\rho\bm u\,\mathrm dV. \tag{11.18}
\end{equation}
不可变形固体相互碰撞或撞壁时，流体压力不足以阻止几何穿透，必须加入短程接触力。其他方案包括以零形变率约束刚体运动、按固体平均加速度预测速度等；弹性体则需求解相应应变-应力关系。

% [原PDF p273]
\subsection{凝固}
% [补译 p273]
凝固控制岩浆与海冰等自然过程，也决定金属、聚合物、混凝土和玻璃的加工品质。液体通常先在容器壁或固体晶核上形成固相；显热与潜热继续被移走后，固相区域逐渐扩张。合金各组分的凝固温度不同，界面前还会形成溶质边界层，最终组织记录了冷却历史。纯物质至少需求温度，二元合金还需求溶质浓度；固液体积变化常可忽略，把全域视为不可压缩。若包含熔体流动，则还要解 Navier--Stokes 方程，并以浸入固体方式处理已凝固区域。


凝固对岩浆与熔岩、海冰和冰川等自然过程，以及金属、聚合物、混凝土和玻璃等材料加工都极为重要。液体冷却后通常从容器壁或固体晶核开始凝固；继续移除显热和潜热，固相区域逐渐长大，直至熔体完全凝固。数值模拟的核心是准确预测温度场和相界面速度。

凝固前沿常会失稳；合金还包含不同凝固温度的组分，最终形成成分不同的复杂微观组织。工艺加热/冷却速率和历史会强烈影响组织和材料性能，非均匀凝固还可产生残余应力、斑点和孔洞。高温、腐蚀环境和有限光学通道使瞬态温度与溶质分布很难实验测量，因此数值模拟尤其重要。

纯物质凝固需求能量方程，二元合金还要解溶质输运。凝固体积变化通常较小，可把固液全域视为不可压缩；若流动可忽略，能量与组分方程即可描述演化，若存在熔体对流，则还要解完整 Navier--Stokes 方程，并按前述浸入固体方法处理固相。熔体热膨胀导致的自然对流可用 Boussinesq 项 $\beta(T-T_0)$ 表示。

% [原PDF p274]
% [补译 p274]
局部热力学平衡使温度跨固液界面连续。最简单时界面温度等于平衡熔点 $T_m$；曲率通过 Gibbs--Thomson 项改变熔点，液侧溶质浓度通过液相线斜率改变平衡温度，有限界面动力学则产生与法向速度成正比的过冷。表面能和动力学迁移率一般随晶向变化，数值研究常用解析各向异性函数代替逐材料实验数据，而枝晶形态会强烈依赖其形式。温度更新后，由两侧热通量跳跃求潜热释放率和界面速度；由于固体与熔体物性通常比汽液物性接近，这一步往往比沸腾问题容易。

局部热力学平衡下，温度跨固液界面连续。界面温度由 Gibbs--Thomson 条件给出，最简单时等于平衡熔点 $T_m$；考虑曲率、溶质与界面动力学后常写为
\begin{equation}
 T_S=T_m-T_m\frac{\sigma\kappa}{L}+mC_L-\frac{V}{\eta}. \tag{11.19}
\end{equation}
其中 $\sigma$ 为界面能，$\kappa$ 为曲率，$L$ 为潜热，$m$ 为液相线斜率，$C_L$ 为界面液侧溶质浓度，$V$ 为法向凝固速度，$\eta$ 为动力学迁移率。界面能和迁移率常随晶向各向异性变化，微观组织对各向异性形式非常敏感。

能量、溶质和流动方程可在同一规则网格上统一求解。更新温度后，由界面两侧热通量跳跃求潜热释放和 $V$。固液物性通常比汽液物性接近，因此界面温度比沸腾问题更容易准确计算。

% [原PDF p275]
% [补译 p275]
铸件宏观尺度只关心各处冷却和完全凝固时间时，焓法最常用。它把潜热吸收到状态变量中，将固液转变表示为恒温或窄温区的糊状区；含流动时又把糊状区当作孔隙率从液相的一变到固相的零的多孔介质。该方法适合整体热史，却不解析真实微结构，并依赖经验阻力。若前沿稳定且近似平面，例如热液滴撞击冷壁后的快速凝固，显式跟踪界面更直接。微结构源于形态不稳定：小凸起伸入更冷或溶质更少的熔体后增长更快，最终形成胞状或枝晶结构。

在铸件整体尺度，只关注冷却与凝固时间而不分辨微观组织时，焓法最常用。它把液固转变表示为等温或窄温区的“糊状区”，焓在液固值之间连续变化；含流动时把糊状区视作孔隙率由液相的 1 变到固相的 0 的多孔介质。这一方法依赖经验糊状区模型。若凝固前沿稳定且近似平面，则可显式追踪界面，无需糊状区经验输入，例如热液滴撞击冷壁后的快速凝固。

连续介质尺度下最困难的是微观组织形成。平面或球形前沿上的凸起伸入更冷或溶质更少、因而更有利于凝固的熔体，增长速度超过周围区域，最终形成胞状或枝晶结构。合金中早、晚凝固固相成分不同，凝固历史被保存在最终组织中；纯物质也可暂时形成复杂枝晶，但完全凝固后几何痕迹消失。

% [原PDF p276]
% [补译 p276]
贴体有限元、水平集、相场和界面追踪都已用于大振幅胞状或枝晶增长。无流动纯物质的三维枝晶如今已相对常规，而合金还要处理界面分凝和更薄的溶质层。实验已证实熔体流动显著改变微结构。耦合计算中，来流把相界释放的潜热扫向下游，使迎风尖端热边界层更薄、增长更快，背风枝臂与下游侧枝受到抑制。三维流体可绕过侧枝，尾迹明显小于二维。代表性计算在 $256^3$ 网格上从轻微扰动的小球形晶核开始，还复现了来流减小迎风尖端半径的实验趋势。

水平集、相场、界面追踪与贴体有限元都已用于大振幅胞状/枝晶生长。无流动纯物质三维枝晶已较常规，合金凝固仍更复杂。实验确认熔体流动会显著影响微观组织；耦合模拟显示，来流把界面释放的潜热从上游带到下游，使迎风枝晶尖端热边界层更薄、增长更快，背风枝臂受到抑制，侧枝也呈明显不对称。三维流体可绕过侧枝，尾迹结构与二维有本质差别。来流还会减小迎风尖端半径，与实验相符。

% [原PDF p277]
% [补译 p277]
\begin{quote}\small 图 11.9　三维枝晶增长的大振幅阶段。液相深色区域表示欠冷液体，浅灰区域温度较高；同时显示穿过域中心平面的温度和速度矢量。引自 Al-Rawahi 与 Tryggvason（2004）。\end{quote}
工业铸造计算已经用于工艺规划，却常以经验关系代替未分辨的枝晶、偏析和糊状区输运，所以尚不能可靠预测最终微结构。界面分辨模拟的长期目标，是建立从工艺参数到组织统计之间的可验证关系，再把这种认识压缩成宏观求解器可用的亚网格模型。这种跨尺度传递若能实现，可能显著改变材料加工设计。

工业铸造模拟已经广泛使用，但宏观模型常依赖未分辨过程的经验闭合，不能准确预测微观组织。若能由分辨模拟建立可靠亚网格模型，并将微观组织机理传递到工业尺度，将可能深刻改变材料加工。

\section{多尺度问题}

即使只研究可由连续介质理论描述的系统，真实多相流的尺度范围仍限制 DNS。最小主导尺度通常至少需要 $O(10)$ 个网格点。当前三维计算可在每方向数百点上推进数千时间步，$1000^3$ 级别也逐渐可行，这使可分辨长度尺度跨度超过两个数量级；然而薄膜、细丝和微滴经常远小于主导流动尺度。

% [原PDF p278]
% [补译 p278]
% [覆盖增强 p278]
当前三维计算通常能在每个方向数百个节点上推进数千时间步，超过 $1000^3$ 的网格也逐渐可行，因此可同时解析两个数量级以上的长度尺度。但界面流动很容易在主导尺度之下生成局部缺陷。正碰两滴若高度对称，可以把分辨率集中到轴线薄膜而勉强覆盖巨大尺度比；多滴任意运动时薄膜位置、方向和数量都不断变化，网格层级与时间步限制会急剧增长。因而多尺度模型不是单纯为了省算力，而是目前计算能力下使某些物理问题可解的必要组成。

多相流的“主导小尺度”通常由表面张力与惯性或黏性的平衡决定，例如射流破碎后的平均液滴尺寸；若用正确的局部尺度定义，$We$、$Ca$、$Oh$ 或 $Re$ 在此常为 $O(1)$。但流动还会产生远小于它的薄膜、细丝与微滴。两颗直径约 $500\,\mu\mathrm m$ 的液滴碰撞时，夹带气膜要排液到几百埃厚才会破裂，尺度跨度可达 $10^4$。自适应网格可把节点集中到薄膜，却显著增加算法复杂度和计算时间；对多个任意运动液滴，完整分辨目前仍基本不可行。

主导小尺度通常由表面张力同惯性或黏性的平衡决定，对射流破碎而言它控制平均液滴尺度，相应局部 $We$、$Ca$、$Oh$ 或 $Re$ 常为 $O(1)$。但两颗直径约 $500\,\mu\mathrm m$ 的液滴碰撞时，夹带气膜必须排液到几百埃厚才会破裂，尺度比可达 $O(10^4)$。AMR 能把网格集中于薄膜，却会显著增加算法复杂度和代价；对多个任意运动液滴，完全分辨目前仍基本不可行。

% [原PDF p279]
% [补译 p279]
\begin{quote}\small 图 11.10　多相流多尺度问题及 E 与 Engquist（2003）的分类。极薄膜、小液滴等是需要局部建模的“孤立缺陷”（A 型）；由微观模型建立宏观本构关系属于 B 型。\end{quote}
点粒子近似是亚网格思想的成熟例子：颗粒远小于最小流动尺度且 Stokes 数很小时，主体流仍完全分辨，而颗粒仅保留位置、速度与点力。低颗粒 Reynolds 数下，流固相互作用可用解析关系给出。已有气泡浆态反应器计算在分辨大气泡的同时把催化颗粒视为点，也有雾化计算把低于网格尺度的小滴转换成 Lagrangian 点滴。薄膜也可由润滑理论描述，但与大尺度界面模拟动态耦合仍不普遍。

许多极小结构的几何和动力学相对简单：黏性占优，强表面张力使形状规则，可用解析或准解析模型代替网格分辨。典型例子是点粒子近似：当液滴远小于最小流动尺度且 Stokes 数远小于 1 时，主体流完全分辨，而微粒视为点；低粒子 Reynolds 数下可用解析阻力，几乎不需经验调节。已有气泡浆态反应器模拟把小催化颗粒视为点粒子而分辨大气泡，也有雾化模拟把极小液滴转为 Lagrangian 点滴。

薄膜同样适合降维模型。薄膜方程已有成熟理论，但与大尺度界面模拟动态耦合仍不普遍；液滴碰撞排液、成核沸腾微液膜和热颗粒上的膜态沸腾是较早实例。% [原PDF p280]
% [补译 p280]
这些推广共同检验数值方法对守恒、跳跃条件和尺度分离的处理能力。

从研究方法看，界面分辨 DNS 同时承担两种任务。其一是在给定尺度范围内直接预测瞬态流动、界面拓扑与传热传质；其二是作为“数值实验”，为更大尺度模型提供闭合、临界条件和误差可量化的数据。后一任务要求模拟不仅在视觉上合理，还要对网格、界面厚度、物性比和小尺度替代模型做系统验证。本章列举的热毛细、电流体、反应、沸腾、固体边界和凝固问题，虽然物理各异，却都沿用同一原则：体相守恒方程与界面跳跃条件必须作为一个整体离散。未来进展将同时依赖更强计算能力、适应极端尺度比的算法，以及能把分辨结果可靠传递到工程尺度的多尺度理论。

当尺度继续缩小时，经典连续介质假设最终失效。此时可以完全转向分子动力学或耗散粒子动力学，也可使用相场等把微观效应平均进自由能与输运关系的修正连续模型；分子--连续区域耦合已经有初步尝试。E 与 Engquist 将多尺度问题分为两类：A 型处理薄膜、小滴等局部孤立缺陷，需要在大尺度计算中嵌入局部模型；B 型由微观或分辨模型提取宏观本构规律。前述碰撞薄膜与点粒子属于 A 型，第 8 章从大量分辨气泡研究群体动力学并构造闭合则属于 B 型。

尺度继续减小时，经典连续介质假设失效，需要分子动力学、耗散粒子动力学，或含小尺度效应的相场/修正连续模型，并可能与宏观求解器区域耦合。

按 E 与 Engquist 的分类，多尺度问题可分为 A 类“孤立缺陷”和 B 类“由微观模型建立本构”。薄膜和微滴的局部亚网格处理属于 A 类；从大量分辨气泡的群体动力学提炼宏观闭合关系则属于 B 类。实际多相模拟常同时包含两类问题。

\section{小结}

多相流贯穿能源、制造、化工、污水处理和农业等领域。其 DNS 虽起源于 CFD 早期，却在计算能力和算法近年成熟后才开始真正改变研究方式。应用范围扩大后，附加场、复杂界面物理、相变、固体边界和多尺度闭合已成为不可回避的核心。本章只勾勒了若干方向，但足以说明：多相流 DNS 已取得巨大进展，同时仍是一个充满基础问题和重要应用机会的活跃领域。

\appendix
% [原PDF p281]
\chapter{界面的描述与定义}
% [补译 p281]
准确跟踪分隔不同流体或不同相的可变形界面，是可靠预测多相流的核心。本附录建立描述界面几何和运动所需的基本记号。三维曲面几何本可引出完整的微分几何体系，但本书只使用较初等的概念，不引入 Christoffel 符号；论述主要遵循 Weatherburn。三维界面是分隔两相的曲面，二维界面则是一条曲线，故先从曲线开始。以单调参数 $u$ 表示曲线 $\bm x(u)=(x(u),y(u))$。导数 $\mathrm d\bm x/\mathrm du$ 沿切向，除以 $\mathrm ds/\mathrm du$ 后得到单位切向；弧长定义为 $\mathrm ds^2=\mathrm dx^2+\mathrm dy^2$。


准确跟踪分隔不同流体或相的可变形界面，需要一套描述几何与运动的数学记号。三维曲面几何可以非常复杂，但本书只需基本结果，无须引入 Christoffel 符号。以下先讨论二维曲线，再推广到三维曲面。

\section{二维几何}

二维曲线用单调参数 $u$ 表示：
\begin{equation}
 \bm x(u)=(x(u),y(u)). \tag{A.1}
\end{equation}
$\mathrm d\bm x/\mathrm du$ 切于曲线。弧长满足 $\mathrm ds^2=\mathrm dx^2+\mathrm dy^2$。记
\begin{equation*}
 x_u=\frac{\mathrm dx}{\mathrm du},\quad x_{uu}=\frac{\mathrm d^2x}{\mathrm du^2},\quad
 y_u=\frac{\mathrm dy}{\mathrm du},\quad y_{uu}=\frac{\mathrm d^2y}{\mathrm du^2},
\end{equation*}
则单位切向和法向为
\begin{equation}
 \bm t=\frac{(x_u,y_u)}{(x_u^2+y_u^2)^{1/2}},\qquad
 \bm n=\frac{(-y_u,x_u)}{(x_u^2+y_u^2)^{1/2}}. \tag{A.2}
\end{equation}
按此定向，曲率
\begin{equation}
 \kappa=\bm n\cdot\frac{\mathrm d\bm t}{\mathrm ds}
 =\frac{x_uy_{uu}-y_ux_{uu}}{(x_u^2+y_u^2)^{3/2}}. \tag{A.3}
\end{equation}

% [原PDF p282]
% [补译 p282]
法向取向使 $(\bm t,\bm n)$ 构成右手基。曲率是单位切向对弧长的变化率在法向上的投影，因此参数 $u$ 无论怎样单调重标定，都得到同一几何量。高度函数 $y=h(x)$ 对应取 $u=x$，于是 $x_u=1$、$x_{uu}=0$，一般参数式直接化为式（A.4）和（A.5）。椭圆 $x=a\cos u$、$y=b\sin u$ 的外法向与原先右手取向相反，切向也相应反向；曲率为负是法向约定的结果。当 $a=b$ 时恢复圆的 $-1/b$。

若曲线是高度函数 $y=h(x)$，取 $u=x$，则
\begin{equation}
 \kappa=\frac{h_{xx}}{(1+h_x^2)^{3/2}},\qquad
 \bm t=\frac{(1,h_x)}{(1+h_x^2)^{1/2}},\qquad
 \bm n=\frac{(-h_x,1)}{(1+h_x^2)^{1/2}}. \tag{A.4--A.5}
\end{equation}

椭圆可参数化为
\begin{equation}
 x=a\cos u,\qquad y=b\sin u. \tag{A.6}
\end{equation}
若法向取向外，
\begin{equation}
 \bm t=\frac{(-a\sin u,b\cos u)}{(a^2\sin^2u+b^2\cos^2u)^{1/2}},\quad
 \bm n=\frac{(b\cos u,a\sin u)}{(a^2\sin^2u+b^2\cos^2u)^{1/2}}, \tag{A.7}
\end{equation}
曲率为
\begin{equation}
 \kappa=-\frac{ab}{(a^2\sin^2u+b^2\cos^2u)^{3/2}}. \tag{A.8}
\end{equation}
$a=b$ 时恢复圆的 $\kappa=-1/b$。用 Monge 形式 $y=\pm b(1-x^2/a^2)^{1/2}$ 也可得到等价切向、法向和曲率表达式。

% [原PDF p283]
% [补译 p283]
同一椭圆也可写为 Monge 形式
\begin{equation}
 y=\pm b\left(1-\frac{x^2}{a^2}\right)^{1/2}. \tag{A.9}
\end{equation}
对椭圆上半部，Monge 表示给出的切向、法向和曲率为
\begin{align}
 \bm t&=\frac{(-a\sqrt{a^2-x^2},\,bx)}{\sqrt{a^4+(b^2-a^2)x^2}},\qquad
 \bm n=\frac{(bx,\,a\sqrt{a^2-x^2})}{\sqrt{a^4+(b^2-a^2)x^2}}, \tag{A.10}\\
 \kappa&=-\frac{a^4b}{[a^4+(b^2-a^2)x^2]^{3/2}}. \tag{A.11}
\end{align}

\section{三维几何}
% [覆盖增强 p283]
$E,F,G$ 构成第一基本形式，只依赖曲面内在度量；$L,M,N$ 构成第二基本形式，依赖法向并刻画曲面嵌入空间后的弯曲。恒等式 $H^2=EG-F^2$ 表明，$H\,\mathrm du\mathrm dv$ 正是参数小矩形映射到真实曲面后的面积。坐标正交只要求 $F=0$，正交归一还要求 $E=G=1$。任意曲面一般不能在有限区域上保持这样的坐标，但在单个点附近总可用局部正交归一基理解公式；数值计算则通常保留一般参数形式，避免因坐标选择而丢失几何不变性。

三维曲面需要两个参数 $u,v$。任意曲面内位移是两条坐标切向 $\bm x_u,\bm x_v$ 的线性组合；第一基本形式中的 $E,F,G$ 给出长度和夹角，$H=|\bm x_u\times\bm x_v|$ 给出面积伸缩。第二基本形式的 $L,M,N$ 则描述曲面法向沿坐标方向的变化。只有在特殊点或特殊曲面上才可把坐标全局选成正交归一。


三维界面以两个曲面坐标表示：
\begin{equation}
 \bm x(u,v)=(x(u,v),y(u,v),z(u,v)). \tag{A.12}
\end{equation}
曲面内任意微小位移
\begin{equation}
 \mathrm d\bm x=\bm x_u\,\mathrm du+\bm x_v\,\mathrm dv, \tag{A.13}
\end{equation}
其第一基本形式为
\begin{equation}
 \mathrm ds^2=E\,\mathrm du^2+2F\,\mathrm du\mathrm dv+G\,\mathrm dv^2, \tag{A.14}
\end{equation}
其中
\begin{equation}
 E=\bm x_u\cdot\bm x_u,\quad F=\bm x_u\cdot\bm x_v,\quad
 G=\bm x_v\cdot\bm x_v,\quad H=|\bm x_u\times\bm x_v|,\quad H^2=EG-F^2. \tag{A.15}
\end{equation}
第二基本形式系数为
\begin{equation}
 L=\bm n\cdot\bm x_{uu},\qquad M=\bm n\cdot\bm x_{uv},\qquad N=\bm n\cdot\bm x_{vv}. \tag{A.16}
\end{equation}
$F=0$ 时坐标切向正交；再有 $E=G=1$ 时为正交归一，但任意曲面一般不能全局选取这种坐标。

% [原PDF p284]
% [补译 p284]
\begin{quote}\small 图 A.1　曲面坐标。坐标线 $u=\mathrm{const}$ 与 $v=\mathrm{const}$ 的切向为 $\bm x_v$ 和 $\bm x_u$，$\bm n$ 为曲面法向，$\bm t$ 与 $\bm p$ 位于切平面。\end{quote}
曲面内不同曲线一般有不同曲率，而且曲线自身的法向通常并不等于曲面法向。主方向的定义要求沿该方向位移时，曲面法向变化与位移平行，仿照平面曲线写成
\begin{equation}
 (\kappa\bm x_u+\bm n_u)\,\mathrm du+(\kappa\bm x_v+\bm n_v)\,\mathrm dv=0. \tag{A.17}
\end{equation}
分别与 $\bm x_u$、$\bm x_v$ 点乘得到式（A.18）；非平凡方向存在的条件给出关于 $\kappa$ 的二次方程。两个根是唯一的主曲率，其和足以描述本书所需的平均曲率。

寻找法向变化与曲面位移满足 $\mathrm d\bm n=-\kappa\,\mathrm d\bm x$ 的方向，可得
\begin{equation}
 (\kappa E-L)\,\mathrm du+(\kappa F-M)\,\mathrm dv=0,
\quad
 (\kappa F-M)\,\mathrm du+(\kappa G-N)\,\mathrm dv=0. \tag{A.18}
\end{equation}
非零解要求
\begin{equation}
 H^2\kappa^2-(EN-2FM+GL)\kappa+(LN-M^2)=0, \tag{A.19}
\end{equation}
其两根为主曲率 $\kappa_1$、$\kappa_2$。本书采用二者之和作为平均曲率：
\begin{equation}
 \kappa=\kappa_1+\kappa_2=\frac{EN-2FM+GL}{H^2}. \tag{A.20--A.21}
\end{equation}
有些文献把二者平均值称平均曲率，相差因子 $1/2$。Gauss 曲率为主曲率乘积
\begin{equation}
 \kappa_G=\kappa_1\kappa_2=\frac{LN-M^2}{H^2}. \tag{A.22}
\end{equation}
主方向可由
\begin{equation}
 \frac{\mathrm du}{\mathrm dv}=-\frac{\kappa_iF-M}{\kappa_iE-L},\qquad i=1,2 \tag{A.23}
\end{equation}
求得，但多数数值算法只需平均曲率。

若曲面为 $z=z(x,y)$，则 $\bm x_x=(1,0,z_x)$、$\bm x_y=(0,1,z_y)$，
\begin{equation*}
 \bm n=\frac{(-z_x,-z_y,1)}{(1+z_x^2+z_y^2)^{1/2}},
\end{equation*}
并可把上述基本形式代入得到高度函数的平均曲率。

% [原PDF p285]
\section{轴对称几何}
% [补译 p285]
本书把两个主曲率之和称为平均曲率；有些文献把二者平均值称平均曲率，故定义相差因子二。Gauss 曲率是主曲率乘积，主方向可由式（A.23）求出，但多相流算法通常只需曲率和。对高度曲面 $z=z(x,y)$，有 $E=1+z_x^2$、$F=z_xz_y$、$G=1+z_y^2$、$H^2=1+z_x^2+z_y^2$，并可得到
\[
 \kappa=\frac{z_{xx}(1+z_y^2)-2z_xz_yz_{xy}+z_{yy}(1+z_x^2)}{(1+z_x^2+z_y^2)^{3/2}}.
\]
轴对称曲面取 $u=r,v=\phi$。总曲率由子午截面曲率与环向曲率相加；环向曲率半径 $R=r/\sin\theta$ 是沿法向到旋转轴的距离。球面两项相等，立即验证式（A.24）。


在柱坐标 $(r,z,\phi)$ 中，轴对称界面可写成 $z=z(r)$。子午面曲线曲率为 $\kappa_1$，法向与 $z$ 轴夹角为 $\theta$，法向延长线到旋转轴的距离为 $R=r/\sin\theta$。总平均曲率是子午曲率与环向曲率之和：
\begin{equation}
 \kappa=\kappa_1+\frac1R. \tag{A.24}
\end{equation}
对球面，两项相等并给出正确结果。

% [原PDF p286]
\section{曲面上的微分与积分}
% [补译 p286]
曲面梯度的方向必须垂直于标量等值线。若等值线上位移为 $(\delta u,\delta v)$，则
\begin{equation}
 \varphi_u\delta u+\varphi_v\delta v=0,
 \qquad \frac{\delta u}{\delta v}=-\frac{\varphi_v}{\varphi_u}. \tag{A.25}
\end{equation}
再要求梯度方向对应的曲面位移与等值线位移在三维空间正交，并用第一基本形式消去内积，便得到式（A.26）。分母 $H^2$ 保证 $\mathrm d\varphi=\nabla_S\varphi\cdot\mathrm d\bm x$。向量场曲面散度的式（A.27）以相同度量构造；在局部正交归一坐标下，$F=0$、$E=G=H=1$，它退化为沿两条单位切向的普通方向导数之和。


标量 $\varphi(u,v)$ 的曲面梯度必须垂直于其等值线。由基本形式可得
\begin{equation}
 \nabla_S\varphi=\frac1{H^2}\left[
 \bm x_u(G\varphi_u-F\varphi_v)+\bm x_v(E\varphi_v-F\varphi_u)
 \right]. \tag{A.26}
\end{equation}
向量 $\bm F$ 的曲面散度为
\begin{equation}
 \nabla_S\cdot\bm F=\frac1{H^2}\left[
 \bm x_u\cdot(G\bm F_u-F\bm F_v)+
 \bm x_v\cdot(E\bm F_v-F\bm F_u)\right]. \tag{A.27}
\end{equation}
局部正交归一时，它退化为熟悉的切向导数之和。

% [原PDF p287]
% [补译 p287]
% [覆盖增强 p287]
投影张量 $\bm I_S=\bm I-\bm n\bm n$ 把任意空间向量投到切平面。因而曲面梯度可先在邻域内任意光滑延拓标量，再用 $\bm I_S$ 去掉法向分量；结果与延拓方式无关，只要界面上的标量值相同。法向是单位向量，$\bm n\cdot\bm n=1$ 的导数说明其沿任何方向的变化都垂直于自身。若把法向沿自身方向保持常值延拓，则法向导数为零，普通散度与曲面散度相等，均为 $-\kappa$。这解释了水平集法为何常从归一化梯度的散度直接计算曲率。

把曲面散度作用于法向，先得到
\begin{equation}
 \nabla_S\cdot\bm n=\frac{1}{H^2}\!\left[
 \bm x_u\cdot(G\bm n_u-F\bm n_v)+
 \bm x_v\cdot(E\bm n_v-F\bm n_u)\right]. \tag{A.28}
\end{equation}
对恒等式 $\bm n\cdot\bm x_u=\bm n\cdot\bm x_v=0$ 分别求导，可把法向导数与第二基本形式系数联系起来，得到式（A.29），继而证明 $\nabla_S\cdot\bm n=-\kappa$。任意向量可分解为法向部分 $(\bm n\cdot\bm F)\bm n$ 和切向部分 $(\bm I-\bm n\bm n)\cdot\bm F$；因此 $\nabla_S=(\bm I-\bm n\bm n)\cdot\nabla$。适当延拓法向后，普通散度也给出 $-\kappa$。

利用
\begin{equation}
 \bm n_u\cdot\bm x_u=-L,\quad
 \bm n_v\cdot\bm x_u=\bm n_u\cdot\bm x_v=-M,\quad
 \bm n_v\cdot\bm x_v=-N, \tag{A.29}
\end{equation}
可得
\begin{equation}
 \nabla_S\cdot\bm n=-\frac{EN-2FM+GL}{H^2}=-\kappa. \tag{A.30}
\end{equation}
任意向量分解为法向与切向部分：
\begin{equation}
 \bm F_S=\bm F-(\bm n\cdot\bm F)\bm n=(\bm I-\bm n\bm n)\cdot\bm F. \tag{A.31}
\end{equation}
因此 $\nabla_S=(\bm I-\bm n\bm n)\cdot\nabla$，并且对恰当延拓的法向场
\begin{equation}
 \nabla\cdot\bm n=-\kappa. \tag{A.32}
\end{equation}
对圆 $f=x^2+y^2-R^2=0$，$\bm n=\nabla f/|\nabla f|=(x,y)/\sqrt{x^2+y^2}$，直接计算得
\begin{equation}
 \nabla\cdot\bm n=\frac{y^2}{(x^2+y^2)^{3/2}}+
 \frac{x^2}{(x^2+y^2)^{3/2}}=\frac1R. \tag{A.33}
\end{equation}
圆 $f=x^2+y^2-R^2=0$ 上，$\bm n=(x,y)/(x^2+y^2)^{1/2}$，直接取散度即得 $1/R$，符号由法向约定决定。还常用恒等式
\begin{equation}
 (\bm n\times\nabla)\times\bm n=\kappa\bm n,
 \qquad \nabla_S\cdot\bm I_S=\kappa\bm n. \tag{A.34--A.35}
\end{equation}

% [原PDF p288]
% [补译 p288]
% [覆盖增强 p288]
恒等式（A.34）可用 Levi--Civita 符号展开证明：双重叉积留下 $-\bm n(\nabla\cdot\bm n)$，而由单位法向得到的另一项为零，于是按本书符号约定等于 $\kappa\bm n$。式（A.35）是同一几何事实的张量形式，其中 $\bm I_S$ 为曲面单位投影。曲面散度定理与平面 Green 公式的差别正是 $-\int_A\kappa\bm F\cdot\bm n\,\mathrm dA$；曲面弯曲时，沿边界的余法向通量不仅反映曲面内散度，还必须补偿法向随位置的变化。该项在推导界面应力和曲面输运守恒式时不可省略。

符号与本书法向约定一致。利用指标记号还可证明 $(\bm n\times\nabla)\times\bm n=\kappa\bm n$ 和 $\nabla_S\cdot\bm I_S=\kappa\bm n$。曲面散度定理把面积积分转成边界曲线积分，并多出曲率项。令 $\bm F=\bm c$ 为任意常向量，式（A.36）化为
由任意性可从式（A.37）得到式（A.38）。

曲面散度定理为
\begin{equation}
 \int_A\nabla_S\cdot\bm F\,\mathrm dA
 =\oint_C\bm F\cdot\bm p\,\mathrm dl-
 \int_A\kappa\bm F\cdot\bm n\,\mathrm dA, \tag{A.36}
\end{equation}
其中 $C=\partial A$，$\bm p=\bm t\times\bm n$ 位于曲面切平面内且垂直于 $C$。令 $\bm F$ 为任意常向量，得到
\begin{equation}
 \bm c\cdot\left(\oint_C\bm p\,\mathrm dl-\int_A\kappa\bm n\,\mathrm dA\right)=0. \tag{A.37}
\end{equation}
\begin{equation}
 \oint_C\bm p\,\mathrm dl=\int_A\kappa\bm n\,\mathrm dA, \tag{A.38}
\end{equation}
进而
\begin{equation}
 \kappa\bm n=\lim_{\delta A\to0}\frac1{\delta A}\oint_{\partial(\delta A)}\bm p\,\mathrm dl. \tag{A.39}
\end{equation}
% [原PDF p289]
% [补译 p289]
把式（A.38）用于围绕一点的无穷小曲面元，得到曲率向量的边界积分定义（A.39），这正是许多离散表面张力算法的几何基础。对闭合曲面，边界曲线消失，所以 $\int_A\kappa\bm n\,\mathrm dA=0$；常表面张力在闭合界面上的总合力因而为零。

对闭合曲面，
\begin{equation}
 \int_A\kappa\bm n\,\mathrm dA=\bm0, \tag{A.40}
\end{equation}
这保证常表面张力在闭合界面上的总合力为零。将曲面散度定理应用于 $\bm F\times\bm c$，其中 $\bm c$ 为任意常向量，可进一步得到
\begin{equation}
 \int_A\nabla_S\times\bm F\,\mathrm dA
 =\oint_C\bm p\times\bm F\,\mathrm dl-
 \int_A\kappa\bm n\times\bm F\,\mathrm dA. \tag{A.41}
\end{equation}
选择不同的 $\bm F$ 还可导出其他曲面积分恒等式。

% [原PDF p290]
\chapter{集中在界面上的分布}
% [补译 p290]
阶跃函数 $H$ 在指定流体内为一、外部为零，是单流体表述的基本对象。二维时可把它写为区域 $V$ 上两个一维 Dirac 分布乘积的积分。若观察点位于闭合轮廓内，积分过程中必有 $(x',y')=(x,y)$，结果为一；若在外部则永不重合，结果为零。为简化记号这里只推导二维式，三维推广完全相同。界面恰是 $\nabla H$ 非零之处。


单流体表述中的阶跃函数 $H$ 在流体 1 内为 1、其他区域为 0。二维中可写成
\begin{equation}
 H(x,y)=\int_V\delta(x-x')\delta(y-y')\,\mathrm dV'. \tag{B.1}
\end{equation}
把对未加撇坐标的梯度移入积分得
\begin{equation}
 \nabla H=\int_V\nabla[\delta(x-x')\delta(y-y')]\,\mathrm dV'. \tag{B.2}
\end{equation}
利用 $\nabla=-\nabla'$，再用散度定理，就把区域积分化为沿界面的分布积分。
对非撇变量求梯度，再换成对积分变量求导并使用散度定理：
\begin{align}
 \nabla H&=-\int_V\nabla'[\delta(x-x')\delta(y-y')]\,\mathrm dV' \tag{B.3}\\
 &=-\oint_S\delta(x-x')\delta(y-y')\bm n'\,\mathrm dS'. \tag{B.4}
\end{align}
% [原PDF p291]
% [补译 p291]
虽然推导先假定标记区域有限、边界闭合，但分布只在观察点附近有贡献，所以也可把闭合积分换成局部界面段积分：
\begin{equation}
 \nabla H=-\int_S\delta(x-x')\delta(y-y')\bm n'\,\mathrm dS'. \tag{B.5}
\end{equation}
在局部切向--法向坐标中，二维分布写成 $\delta(s)\delta(n)$；沿切向积分后只剩 $-\delta(n)\bm n$。因此定义界面分布 $\delta_S(\bm x)=\delta(n)$，它把任意体积分转换为界面积分，并满足 $\nabla H=-\delta_S\bm n$。数值计算必须把 Dirac 分布平滑到有限网格宽度；平滑 $H$ 的离散导数与平滑 $\delta_S$ 应保持一致，才能避免界面力和物性跳跃错位。

在界面局部切向-法向坐标中
\begin{equation}
 \delta(x-x')\delta(y-y')=\delta(s)\delta(n), \tag{B.6}
\end{equation}
沿界面积分后得到
\begin{equation}
 \nabla H=-\delta(n)\bm n. \tag{B.7}
\end{equation}
定义界面分布
\begin{equation}
 \delta_S(\bm x)=\delta(n), \tag{B.8}
\end{equation}
它把体积分转化为界面积分：
\begin{equation}
 \int_V\delta_S(\bm x)f(\bm x)\,\mathrm dV=\int_Sf(\bm x)\,\mathrm dS,
 \qquad \nabla H=-\delta_S\bm n. \tag{B.9--B.10}
\end{equation}

数值计算中用有限宽光滑核近似 Dirac 分布。可从两类极限理解：
\begin{equation}
 \delta(x)=\lim_{\varepsilon\to0}\frac{\varepsilon}{\pi(x^2+\varepsilon^2)},
 \qquad
 \delta(x)=\lim_{\varepsilon\to0}\frac{e^{-x^2/(4\varepsilon)}}{2\sqrt{\pi\varepsilon}}. \tag{B.11}
\end{equation}
固定网格上在有限 $\varepsilon$ 停止极限过程，使界面扩展为数个网格宽。因为 $\delta$ 是阶跃函数的导数，可对平滑 $H$ 离散求导得到 $\delta$，或积分离散 $\delta$ 得到 $H$；两者应保持离散一致。

% [原PDF p292]
\section{一个简单例子}
% [补译 p292]
\begin{quote}\small 图 B.1　当宽度参数 $\varepsilon\to0$ 时，光滑尖峰函数趋于 Dirac $\delta$ 分布。图 B.2　用跨越数个网格距的平滑阶跃函数表示圆形区域。\end{quote}
固定网格上不能真正令 $\varepsilon$ 为零，而是在有限宽度停止极限过程，正如有限差分从导数定义在有限步长处截断一样。因为 $\delta$ 是阶跃函数的导数，可以先选择平滑 $H$ 再离散求导，也可先选平滑 $\delta$ 再积分恢复 $H$。简单例子考虑常速 $U_0$ 输运标量 $f$，在 $x_s$ 以速率 $q$ 注入。界面两侧分别满足式（B.12）的齐次对流方程，并由跳跃 $[U_0f]=q$ 耦合。


考虑均匀速度 $U_0$ 输运标量 $f$，来流侧初始 $f=0$，在 $x_s$ 以速率 $q$ 注入。分区形式在 $x_s$ 两侧满足
\begin{equation}
 \frac{\partial f}{\partial t}+U_0\frac{\partial f}{\partial x}=0,
 \qquad x<x_s\ \text{或}\ x>x_s, \tag{B.12}
\end{equation}
并以跳跃 $[U_0f]=q$ 耦合。解为左侧 0、右侧 $q/U_0$。

% [原PDF p293]
% [补译 p293]
\begin{quote}\small 图 B.3　一维对流中不连续解的平滑近似；给出理想跳跃和两种不同宽度的过渡层。\end{quote}
分区问题的解显然是注入点左侧 $f=0$、右侧 $f=q/U_0$。把跳跃写进全域方程，得到奇异源形式（B.13）。为证明等价性，按式（B.14）将解分解为两侧光滑函数与互补阶跃函数的乘积。这里 $H_1$ 在 $x_s$ 左侧为一、右侧为零，故 $\partial_xH_2=-\partial_xH_1=\delta(x-x_s)$。

等价的单场方程是
\begin{equation}
 \frac{\partial f}{\partial t}+U_0\frac{\partial f}{\partial x}=q\delta(x-x_s). \tag{B.13}
\end{equation}
写
\begin{equation}
 f=f_1H_1(x-x_s)+f_2H_2(x-x_s),\qquad H_2=1-H_1, \tag{B.14}
\end{equation}
代入得
% [原PDF p294]
% [补译 p294]
把分解式代入含奇异源的全域对流方程，普通函数项分别给出两侧的齐次方程，而乘 $\delta$ 的项给出
\[
 U_0(f_2-f_1)\delta=q\delta.
\]
因此恢复 $q=U_0(f_2-f_1)=[U_0f]$，与原分区跳跃条件完全相同。这个例子说明，集中在界面上的源项不是经验修正，而是把体相方程与界面跳跃条件合写成一个分布方程。当 $\delta$ 用有限宽平滑函数近似时，理想不连续相应变成有限宽过渡区；只要 $H$ 与 $\delta$ 的离散定义一致，积分跳跃仍可保持。

\begin{equation}
 H_1(f_{1t}+U_0f_{1x})+H_2(f_{2t}+U_0f_{2x})
 +U_0(f_2-f_1)\delta=q\delta. \tag{B.15}
\end{equation}
体相方程与跳跃条件同时恢复，证明两种表述等价。把 $\delta$ 替换为平滑核时，理想跳跃变成有限宽过渡层。

% [原PDF p295]
\chapter{立方体切割算法}
% [补译 p295]
立方体切割问题求平面在直角六面体内截出的体积，是几何 VOF 重构的基本运算。平面 $\bm m\cdot\bm x=\alpha$ 的法向指向所求体积之外，先假定三个 $m_i$ 为正。直接问题由已知法向和截距求体积，逆问题则由体积分数反求截距。式（C.2）可用容斥原理解读：$\alpha^3/(6m_1m_2m_3)$ 是原点侧大四面体；第一组 $F_3$ 项扣除越过单元面的部分，第二组再加回被重复扣除的对角区域。$F_n(z)$ 在 $z>0$ 时等于 $z^n$，否则为零，使所有几何分段统一写进一个公式。二维对应式只需扣除越界三角形。


考虑平面切割边长为 $\Delta x_1,\Delta x_2,\Delta x_3$、体积 $V_0=\Delta x_1\Delta x_2\Delta x_3$ 的直角六面体。平面方程为
\begin{equation}
 \bm m\cdot\bm x=m_1x_1+m_2x_2+m_3x_3=\alpha, \tag{C.1}
\end{equation}
法向 $\bm m$ 指向待计算体积之外，先假定 $m_i>0$。直接问题是在已知 $\bm m,\alpha$ 时计算平面下方体积；逆问题是由体积分数反求 $\alpha$。

定义 $F_n(z)=z^n$（$z>0$）且 $F_n(z)=0$（$z\le0$），以及 $\alpha_{\max}=\sum_{i=1}^3m_i\Delta x_i$。三维切割体积为
\begin{equation}
 V=\frac1{6m_1m_2m_3}\left[
 \alpha^3-\sum_{i=1}^3F_3(\alpha-m_i\Delta x_i)
 +\sum_{i=1}^3F_3(\alpha-\alpha_{\max}+m_i\Delta x_i)
 \right]. \tag{C.2}
\end{equation}
几何上，第一项是原点侧直四面体体积；第一组和项扣除平面越过各单元面后伸出单元的四面体；第二组和项把被重复扣除的对角区域加回。二维面积相应为
\begin{equation}
 A=\frac1{2m_1m_2}\left[
 \alpha^2-\sum_{i=1}^2F_2(\alpha-m_i\Delta x_i)
 \right]. \tag{C.3}
\end{equation}

归一化为 $\Delta x_i=1$、$\sum m_i=1$ 后，$V(\alpha)$ 连续、单调一一对应且一阶导连续；$V,\alpha\in[0,1]$，关于 $(1/2,1/2)$ 奇对称，所以只需求 $[0,1/2]$；公式在指标置换下不变，故可排序 $m_1\le m_2\le m_3$；$m_1\to0$ 时三维式光滑退化为二维式。

% [原PDF p296]
\section{二维问题}
% [补译 p296]
\begin{quote}\small 图 C.1　“切割体积”为直角六面体内部且位于平面 $ABCD$ 下方的灰色区域。\end{quote}
归一化为 $\Delta x_i=1$、$\sum m_i=1$ 后，$V(\alpha)$ 具有四项关键性质：它连续、单调且一一对应，第一导数也连续；$V$ 与 $\alpha$ 都在 $[0,1]$，并关于 $(1/2,1/2)$ 奇对称，故只需处理前半区间；交换坐标指标不改变关系，所以可排序二维 $m_1\le m_2$、三维 $m_1\le m_2\le m_3$；最后，当三维某个系数趋零时，体积公式光滑退化为二维面积公式。二维直接式分为三角形段和梯形段，逆式分别为平方根和线性关系。


令 $m_1+m_2=1$、$m_1\le m_2$、$A_1=m_1/(2m_2)$。直接关系为
\begin{equation*}
 A=\begin{cases}
 \alpha^2/(2m_1m_2),&0\le\alpha<m_1,\\
 \alpha/m_2-A_1,&m_1\le\alpha\le1/2,
 \end{cases}
\end{equation*}
逆关系为
\begin{equation*}
 \alpha=\begin{cases}
 \sqrt{2m_1m_2A},&0\le A<A_1,\\
 m_2A+m_1/2,&A_1\le A\le1/2.
 \end{cases}
\end{equation*}

% [原PDF p297]
\section{三维问题}
% [补译 p297]
三维令 $m_{12}=m_1+m_2$、$\widehat m=\min(m_{12},m_3)$，按平面依次越过三个坐标方向把 $\alpha$ 区间分成四段。第一段仅切出原点四面体，第二段减去越过第一面的四面体，第三段再处理第二个方向。第四段取决于 $m_3<m_{12}$ 还是 $m_{12}<m_3$：前者仍为三次式，后者几何已退化成截面积恒定的棱柱，体积对 $\alpha$ 线性。逆问题的前两段可直接开立方根或解二次方程；后两段需解三次多项式，并只保留落在当前几何区间且使 $P(\alpha)$ 单调增加的实根。


令 $m_1+m_2+m_3=1$、$m_1\le m_2\le m_3$，$m_{12}=m_1+m_2$，$\widehat m=\min(m_{12},m_3)$。直接问题按 $\alpha$ 分成四段：
\begin{align}
 V&=\frac{\alpha^3}{6m_1m_2m_3},&&0\le\alpha<m_1,\\
 V&=\frac{\alpha(\alpha-m_1)}{2m_2m_3}+V_1,&&m_1\le\alpha<m_2,\\
 V&=\frac{\alpha^2(3m_{12}-\alpha)+m_1^2(m_1-3\alpha)+m_2^2(m_2-3\alpha)}{6m_1m_2m_3},
 &&m_2\le\alpha<\widehat m.
\end{align}
第四段又依 $m_3<m_{12}$ 或 $m_{12}<m_3$ 分别为
\begin{equation*}
 V=\frac{\alpha^2(3-2\alpha)+\sum_{i=1}^3m_i^2(m_i-3\alpha)}{6m_1m_2m_3},
 \quad m_3\le\alpha\le\frac12,
\end{equation*}
或
\begin{equation*}
 V=\frac{2\alpha-m_{12}}{2m_3},
 \quad m_{12}\le\alpha\le\frac12.
\end{equation*}

% [原PDF p298]
% [补译 p298]
分段端点为 $V_1=m_1^2/[\max(6m_2m_3,\varepsilon)]$、$V_2=V_1+(m_2-m_1)/(2m_3)$，另有 $V_{31}$ 与 $V_{32}$ 分别对应两种第四区间。$\max$ 中的 $\varepsilon$ 只用于避免退化几何下分母为零，正确极限仍是 $m_1,m_2\to0$ 时 $V_1\to0$。三次式可按 Abramowitz--Stegun 的三角公式求三个实根；判别式为负时，只有满足单调性和区间约束的根可用。推广到矩形网格时除以单元体积得到式（C.4）。若某 $m_i<0$，用 $x_i'=\Delta x_i-x_i$ 镜像到正系数标准情形；对偏移六面体则先用 $x_i'=x_i-x_{0i}$ 平移局部原点。

逆问题前两段显式可解：
\begin{equation*}
 \alpha=(6m_1m_2m_3V)^{1/3},\qquad 0\le V<V_1,
\end{equation*}
\begin{equation*}
 \alpha=\frac12\left[m_1+\sqrt{m_1^2+8m_2m_3(V-V_1)}\right],
 \qquad V_1\le V<V_2.
\end{equation*}
后续区间需解三次多项式 $P(\alpha)=0$；三个根中只取使 $P(\alpha)$ 随 $\alpha$ 单调增加、且落在当前分段区间的根。为避免 $m_1,m_2\to0$ 时分母为零，可令
\begin{equation*}
 V_1=\frac{m_1^2}{\max(6m_2m_3,\varepsilon)},
\end{equation*}
$\varepsilon$ 是仅用于退化几何的极小正数。三次方程可用标准三角解析式或稳定的区间 Newton/二分法求解；单调性使后者在数值实现中很可靠。

对一般矩形网格，除以 $V_0$ 得体积分数
\begin{equation}
 C=\frac1{6\prod_{i=1}^3(m_i\Delta x_i)}\left[
 \alpha^3-\sum_{i=1}^3F_3(\alpha-m_i\Delta x_i)
 +\sum_{i=1}^3F_3(\alpha-\alpha_{\max}+m_i\Delta x_i)
 \right]. \tag{C.4}
\end{equation}
把平面方程归一化为 $\alpha_{\max}=1$ 即回到上述标准形式。若某 $m_i<0$，用镜像变换 $x_i'=\Delta x_i-x_i$ 映射到正系数情形，求解后再镜像回原位置。若要计算同一平面在偏移直角六面体内切出的通量体积，先用 $x_i'=x_i-x_{0i}$ 把局部原点平移到该六面体顶点即可。

% [原PDF p299]
\chapter{液膜动力学：线性理论}
% [补译 p299]
本附录汇总上下叠置液层的线性化理论。近静止气液层可产生重力波、毛细波和 Rayleigh--Taylor 不稳定性；两层有切向速度差时则出现 Kelvin--Helmholtz 不稳定性。设相 $n=1,2$ 的密度、黏度和基态水平速度分别为 $\rho_n,\mu_n,U_n$，相 2 位于 $y>0$，相 1 位于 $y<0$，未扰动界面为 $h_0=0$。无黏近似允许 $U_1\ne U_2$ 的速度跳跃。小扰动采用正规模 $q'=q_1(y)e^{ik(x-ct)+i\varphi}$；复波速 $c=c_r+ic_i$ 的实部给频率 $\omega=kc_r$，虚部给增长率 $\gamma=kc_i$。


本附录汇总上下叠置液层的线性理论。近静止气-液层会产生重力波、毛细波和 Rayleigh--Taylor 不稳定性；两层存在剪切时则出现 Kelvin--Helmholtz 不稳定性。

\section{流动构型}

考虑两个水平叠置流体层。相 $n=1,2$ 的密度、黏度分别为 $\rho_n,\mu_n$，基态水平速度为 $U_n$。二维界面 $y=h(x,t)$，未扰动时 $h_0=0$；相 2 位于 $y>0$，相 1 位于 $y<0$。无黏近似允许 $U_1\ne U_2$ 的速度跳跃。

波长为 $\lambda$ 的小扰动写成正规模
\begin{equation}
 q'(x,y,t)=q_1(y)\exp\{ik(x-ct)+i\varphi\}, \tag{D.1}
\end{equation}
其中 $k=2\pi/\lambda$，$c=c_r+ic_i$ 为复波速，增长率 $\gamma=kc_i$，角频率 $\omega=kc_r$。

\section{无黏结果}

\subsection{一般色散关系}

令 $\mu=0$，Euler 方程为
\begin{equation}
 \rho\frac{D\bm u}{Dt}=\rho\bm g-\nabla p, \tag{D.2}
\end{equation}
并满足 $\nabla\cdot\bm u=0$ 与 $D\rho/Dt=0$。相 $n$ 的速度分解为
\begin{equation}
 \bm u=(U_n,0)+(u',v'). \tag{D.3}
\end{equation}
% [原PDF p300]
% [补译 p300]
\begin{quote}\small 图 D.1　本节使用的三种基流：（a）理想尖锐速度跳跃；（b）早期分析采用的两侧分段线性边界层；（c）用于线性理论比较和完整模拟的光滑边界层。\end{quote}
无黏时从 Euler 方程、不可压条件、密度随质点守恒以及界面运动学条件出发。把速度写成基态 $(U_n,0)$ 与小扰动 $(u',v')$ 之和，并围绕定常基流保留一阶项，得到式（D.4）--（D.8）。界面扰动 $h'(x,t)$ 是一条线的位置变化，不依赖 $y$；其物质速度必须等于界面处流体法向速度。正规模把 $\partial_x$、$\partial_t$ 分别替换为 $ik$ 和 $-ikc$。

围绕基态线性化后
\begin{align}
 \rho_0(\partial_tu'+U\partial_xu')&=-\partial_xp', \tag{D.4}\\
 \rho_0(\partial_tv'+U\partial_xv')&=-\partial_yp'-\rho'g, \tag{D.5}\\
 \partial_xu'+\partial_yv'&=0, \tag{D.6}\\
 \partial_t\rho'+U\partial_x\rho'&=-v'\partial_y\rho_0, \tag{D.7}\\
 v'_S&=\partial_th'+U\partial_xh'. \tag{D.8}
\end{align}
% [原PDF p301]
% [补译 p301]
% [覆盖增强 p301]
式（D.9）在界面之外因 $D\rho_0=0$、$DU=0$ 而大幅简化，但界面上的密度与速度跳跃不能直接丢掉，它们通过应力和运动学条件进入色散关系。指数解上、下两侧的衰减方向相反：相 2 中为 $e^{-ky}$，相 1 中为 $e^{ky}$。共同常数 $A$ 保证两侧由式（D.10）得到同一 $h'$。将这些解代回线性动量方程求压力，再在 $y=0$ 施加含重力和表面张力的法向应力跳跃，就得到式（D.15）；根号内若为负，波速出现虚部，扰动随时间指数增长。

消去压力、水平速度和密度扰动后，垂向速度幅值满足
\begin{equation}
 -k^3\rho_0(U-c)v_1+D[\rho_0k(U-c)Dv_1-\rho_0k(DU)v_1]
 =gk\frac{D\rho_0}{U-c}v_1. \tag{D.9}
\end{equation}
运动学条件写成
\begin{equation}
 h'=-\frac{i v'_S}{k(U-c)}. \tag{D.10}
\end{equation}
远离界面，基态分片常数使方程退化为 $(D^2-k^2)v_1=0$。要求 $|y|\to\infty$ 时衰减，并保证两侧给出同一界面位移，便得到指数解（D.12）、（D.13）。把压力幅值和线性曲率代入法向应力跳跃，最终得到二次色散关系（D.15）和复波速（D.16）。

正规模下 $\partial_x\to ik$、$\partial_t\to-ikc$、$D=\mathrm d/\mathrm dy$。远离界面有
\begin{equation}
 (D^2-k^2)v_1(y)=0. \tag{D.11}
\end{equation}
满足 $|y|\to\infty$ 衰减和界面法向位移唯一性的解是
\begin{equation}
 v_1=A k(U_2-c)e^{-ky},\quad y>0;\qquad
 v_1=A k(U_1-c)e^{ky},\quad y<0. \tag{D.12--D.13}
\end{equation}
线性曲率为
\begin{equation}
 \kappa=\frac{\partial^2h'}{\partial x^2}. \tag{D.14}
\end{equation}
代入法向应力跳跃得到一般色散关系
\begin{equation}
 \rho_2k^2(U_2-c)^2+\rho_1k^2(U_1-c)^2
 =gk(\rho_1-\rho_2)+\sigma k^3. \tag{D.15}
\end{equation}
记 $\rho_p=\rho_1+\rho_2$、$\rho_m=\rho_1-\rho_2$，解得
\begin{equation}
 c(k)=\frac{\rho_1U_1+\rho_2U_2}{\rho_p}
 \pm\left[
 \frac{\rho_mg}{\rho_pk}+\frac{\sigma k}{\rho_p}
 -\frac{\rho_1\rho_2}{\rho_p^2}(U_1-U_2)^2
 \right]^{1/2}. \tag{D.16}
\end{equation}

% [原PDF p302]
\subsection{毛细-重力波}
% [补译 p302]
无平均水平流时，色散关系只含重力项与毛细项。小波数对应长波，重力占主导；大波数对应短波，表面张力占主导。空气--水转折尺度即毛细长度 $\ell_c=\sqrt{\sigma/(\rho_{\mathrm{water}}g)}\approx2.7\,\mathrm{mm}$。若忽略重力与毛细力，速度跳跃使任意波数扰动按 $e^{\gamma t}$ 增长，形成 Kelvin--Helmholtz 不稳定性。气液密度比很大时，增长率按 $k\sqrt{\rho_g/\rho_l}|U_l-U_g|$ 缩放。加入重力、黏性或表面张力后，高波数不再无限增长，而会出现最大增长波数和截止波数。


无平均水平运动时 $U_1=U_2=0$，
\begin{equation}
 c(k)=\left(\frac{\rho_mg}{\rho_pk}+\frac{\sigma k}{\rho_p}\right)^{1/2}. \tag{D.17}
\end{equation}
小 $k$、长波由重力主导；大 $k$、短波由表面张力主导。空气-水体系的毛细长度
\begin{equation*}
 \ell_c=\sqrt{\frac{\sigma}{\rho_{\mathrm{water}}g}}
\end{equation*}
约为 $2.7\,\mathrm{mm}$。

\subsection{Kelvin--Helmholtz 不稳定性}

忽略重力和毛细力，频率与增长率为
\begin{equation}
 \omega(k)=\frac{k}{\rho_p}(\rho_1U_1+\rho_2U_2), \tag{D.18}
\end{equation}
\begin{equation}
 \gamma(k)=\frac{k}{\rho_p}(\rho_1\rho_2)^{1/2}|U_1-U_2|. \tag{D.19}
\end{equation}
扰动按 $e^{\gamma t}$ 增长。气液密度比很大时
\begin{equation}
 \gamma(k)\simeq k\left(\frac{\rho_2}{\rho_1}\right)^{1/2}|U_1-U_2|. \tag{D.20}
\end{equation}
重力、黏性或表面张力会使增长率在有限波数达到最大并在截止波数后变为零。仅有表面张力时
\begin{equation}
 \gamma(k)=\left[
 \frac{\rho_1\rho_2k^2}{\rho_p^2}(U_1-U_2)^2-
 \frac{\sigma k^3}{\rho_p}
 \right]^{1/2}, \tag{D.21}
\end{equation}
不稳定范围为
\begin{equation}
 k<k_c=\frac{\rho_1\rho_2(U_1-U_2)^2}{\rho_p\sigma}. \tag{D.22}
\end{equation}
以波长定义
\begin{equation}
 We(\lambda)=\frac{\rho_2\lambda(U_1-U_2)^2}{\sigma}, \tag{D.23}
\end{equation}
则 $We>2\pi$ 时失稳。最大增长波数 $k_m=2k_c/3$，对应 $We=3\pi$。若以预定尺度 $D$ 定义 $We_D$，最危险波长为
\begin{equation}
 \lambda_m=3\pi D\,We_D^{-1}. \tag{D.24}
\end{equation}

% [原PDF p303]
\subsection{无黏框架中有限厚边界层的影响}
% [补译 p303]
只保留表面张力时，不稳定要求 $k<k_c=\rho_1\rho_2(U_1-U_2)^2/(\rho_p\sigma)$。以波长定义 Weber 数，失稳条件是 $We(\lambda)>2\pi$；最大增长位于 $k_m=2k_c/3$，即 $We=3\pi$，用预定尺度 $D$ 表示则 $\lambda_m=3\pi D/We_D$。真实黏性把尖锐速度跳跃扩散成厚度 $\delta_g,\delta_l$ 的剪切层。若仍在无黏扰动方程中使用分段线性基流，可以解析得到单峰增长率和式（D.25）、（D.26）的尺度律；但这种“黏性基流加无黏扰动”的混合近似即使在高 Reynolds 数下也不能定量替代完整黏性稳定性理论。


黏性扩散把理想速度跳跃变成气侧厚度 $\delta_g$、液侧厚度 $\delta_l$ 的连续剪切层。可先仍用 Euler 方程、只把基态改成分段线性 $U(y)$；所得增长率仍只有一个峰。对 $U_g\gg U_l$、$\rho_g\ll\rho_l$、$\delta_l=0$，近似给出
\begin{equation}
 k_m\sim1.5\left(\frac{\rho_g}{\rho_l}\right)^{1/2}\frac1{\delta_g},
 \qquad
 \gamma(k_m)\sim\frac{\rho_gU_g}{\rho_l\delta_g}, \tag{D.25}
\end{equation}
以及
\begin{equation}
 \lambda_m\sim\frac{4\pi}{3}\left(\frac{\rho_l}{\rho_g}\right)^{1/2}\delta_g. \tag{D.26}
\end{equation}
但即使 Reynolds 数很大，这种“有边界层基态、无黏扰动方程”的混合近似也不能定量替代完整黏性理论。

% [原PDF p304]
\section{Kelvin--Helmholtz 不稳定性的黏性理论}
% [补译 p304]
一致的黏性理论必须线性化 Navier--Stokes 方程，得到两相耦合的 Orr--Sommerfeld 特征值问题。以气相速度 $U_g$ 和边界层厚度 $\delta_g$ 无量纲化，控制参数为 $Re=U_g\delta_g/\nu_g$ 与 $We=\rho_gU_g^2\delta_g/\sigma$。流函数保证不可压条件自动满足，扰动取 $\psi_n=e^{ik(x-ct)}\phi_n(y)$。对每个实波数 $k$，复特征值 $c$ 同时决定相速度与增长率。两相常微分方程含四阶黏性算子 $(D^2-k^2)^2$，液相方程通过密度比 $r=\rho_g/\rho_l$ 和黏度比 $m=\mu_g/\mu_l$ 与气相方程缩放。界面处首先要求法向和切向速度连续，得到式（D.32）、（D.33）。


一致处理必须线性化 Navier--Stokes 方程，得到 Orr--Sommerfeld 特征值问题。用气相速度 $U_g$ 和气相边界层厚度 $\delta_g$ 无量纲化，
\begin{equation*}
 Re=\frac{U_g\delta_g}{\nu_g},\qquad
 We=\frac{\rho_gU_g^2\delta_g}{\sigma}.
\end{equation*}
用流函数表示速度
\begin{equation}
 u=\frac{\partial\psi}{\partial y},\qquad
 v=-\frac{\partial\psi}{\partial x}, \tag{D.27}
\end{equation}
并取
\begin{equation}
 \psi_1=e^{ik(x-ct)}\phi_1(y),\quad y<0;\qquad
 \psi_2=e^{ik(x-ct)}\phi_2(y),\quad y>0. \tag{D.28--D.29}
\end{equation}
得到以实波数 $k$ 为参数、复波速 $c$ 为特征值的耦合常微分方程。令 $r=\rho_g/\rho_l$、$m=\mu_g/\mu_l$，则
\begin{align}
 (U_2-c)(D^2-k^2)\phi_2-(D^2U_2)\phi_2
 &=\frac1{ikRe}(D^2-k^2)^2\phi_2, \tag{D.30}\\
 (U_1-c)(D^2-k^2)\phi_1-(D^2U_1)\phi_1
 &=\frac{r}{mikRe}(D^2-k^2)^2\phi_1. \tag{D.31}
\end{align}
界面速度连续给出
\begin{equation}
 \phi_1=\phi_2, \tag{D.32}
\end{equation}
\begin{equation}
 D\phi_1+\frac{DU_1}{c}\phi_1
 =D\phi_2+\frac{DU_2}{c}\phi_2. \tag{D.33}
\end{equation}
% [原PDF p305]
% [补译 p305]
法向与切向应力连续还给出
\begin{align}
 -\frac{k^2}{cWe}\phi_2={}&-\frac1r(cD\phi_1+\phi_1DU_1)
 -\frac{1}{mikRe}(D^3-3k^2D)\phi_1 \\
 &+(cD\phi_2+\phi_2DU_2)+\frac{1}{ikRe}(D^3-3k^2D)\phi_2, \tag{D.34}\\
 \frac1m\left(D^2+k^2+\frac{D^2U_2}{c}\right)\phi_2={}&
 \left(D^2+k^2+\frac{D^2U_1}{c}\right)\phi_1. \tag{D.35}
\end{align}
毛细项通过 $We^{-1}$ 进入法向应力条件，黏性项含三阶导数。式（D.30）--（D.35）可用 Chebyshev 多项式配置离散为广义特征值问题，逐个 $k$ 求 $c(k)$，再由 $kc_r$ 和 $kc_i$ 计算频率与增长率；第 9 章将这些结果与完整气液剪切层模拟比较。

另有法向和切向应力连续条件，包含 $We^{-1}$ 毛细项及三阶导数黏性项；与式（D.30）--（D.33）共同构成完整特征值问题。实际可用 Chebyshev 谱配置离散，求得 $c(k)$，再由 $kc_r$ 与 $kc_i$ 计算频率和增长率。完整黏性结果用于第 9 章气液剪切层和雾化分析。

\chapter*{参考文献}

\addcontentsline{toc}{chapter}{参考文献}

\markboth{参考文献}{参考文献}

\begingroup\small

\setlength{\parindent}{0pt}

\setlength{\parskip}{0.45em}

下列题录完整保留原书参考文献的作者、年份、英文题名及出版信息。

% [原PDF p306]

\par\hangindent=2em Abramowitz, M. and Stegun, I. (eds.) (1964). Handbook of Mathematical Functions with Formulas, Graphs, and Mathematical Tables. National Bureau of Standards Applied Mathematics Series 55. US Government Printing Office, Washington, DC.

\par\hangindent=2em Afkhami, S., Zaleski, S., and Bussmann, M. (2009). A mesh-dependent model for applying dynamic contact angles to VOF simulations. J. Comput. Phys., 228: 5370--5389.

\par\hangindent=2em Agarwal, D., Welch, S.W.J., Biswas, G., and Durst, F. (2004). Planar simulation of bubble growth in film boiling in near-critical water using a variant of the VOF method. ASME J. Heat Transfer, 126: 329--338.

\par\hangindent=2em Agresar, G., Linderman, J.J., Tryggvason, G., and Powell. K.G. (1998). An adaptive, Cartesian, front-tracking method for the motion, deformation and adhesion of circulating cells. J. Comput. Phys., 43: 346--380.

\par\hangindent=2em Al-Rawahi, N. and Tryggvason, G. (2002). Numerical simulation of dendritic solidification with convection: two-dimensional geometry. J. Comput. Phys., 180: 471--496.

\par\hangindent=2em Al-Rawahi, N. and Tryggvason, G. (2004). Numerical simulation of dendritic solidification with convection: Three-dimensional flow. J. Comput. Phys., 194: 677--696.

\par\hangindent=2em Alexiades, V. and Solomon, A.D. (1993). Mathematical Modeling of Melting and Freezing Processes. Hemisphere.

\par\hangindent=2em Amsden, A.A. (1993). KIVA-3: a KIVA program with block-structured mesh for complex geometries. Technical report, LA-12503-MS, Los Alamos National Lab., NM.

% [原PDF p307]

\par\hangindent=2em Antal, S.P., Lahey, R.T., and Flaherty, J.E. (1991). Analysis of phase distribution in fully developed laminar bubbly two-phase flows. Int. J. Multiphase Flow, 15: 635--652.

\par\hangindent=2em Aris, R. (1962). Vectors, Tensors and Basic Equations of Fluid Mechanics. Dover.

\par\hangindent=2em Arp, P.A., Foister, R.T., and Mason, S.G. (1980). Some electrohydrodynamic effects in fluid dispersions. Adv. Colloid Interface Sci., 12: 295--356.

\par\hangindent=2em Ashgriz, N. and Poo, J.Y. (1991). FLAIR: flux line-segment model for advection and interface reconstruction. J. Comput. Phys., 93: 449--468.

\par\hangindent=2em Ashgriz, N., Shirani, E., and Mostaghini, J. (2005). Interface pressure calculation based on conservation of momentum for front capturing methods. J. Comput. Phys., 203: 154--175.

\par\hangindent=2em Aubert, F., Aulisa, E., Manservisi, S., and Scardovelli, R. (2006). Interface tracking with dynamically-redistributed surface markers in unstructured quadrangular grids. Comput. Fluids, 35: 1332--1343.

\par\hangindent=2em Aulisa, E., Manservisi, S., and Scardovelli, R. (2003a). A mixed markers and volume-of-fluid method for the reconstruction and advection of interfaces in two-phase and free-boundary flows. J. Comput. Phys., 188: 611--639.

\par\hangindent=2em Aulisa, E., Manservisi, S., Scardovelli, R., and Zaleski, S. (2003b). A geometrical area-preserving volume of fluid advection method. J. Comput. Phys., 192: 355-- 364.

\par\hangindent=2em Aulisa, E., Manservisi, S., and Scardovelli, R. (2004). A surface marker algorithm coupled to an area-preserving marker redistribution method for threedimensional interface tracking. J. Comput. Phys., 197: 555--584.

\par\hangindent=2em Aulisa, E., Manservisi, S., Scardovelli, R., and Zaleski, S. (2007). Interface reconstruction with least-squares fit and split advection in three-dimensional Cartesian geometry. J. Comput. Phys., 225: 2301--2319.

\par\hangindent=2em Azpitarte, O.E. and Buscaglia, G.C. (2003). Analytical and numerical evaluation of two-fluid model solutions for laminar fully developed bubbly two-phase flows. Chem. Eng. Sci., 58: 3765--3776.

\par\hangindent=2em Baker, G.R., Meiron, D.I., and Orszag, S.A. (1980). Vortex simulation of the Rayleigh--Taylor instability. Phys. Fluids, 23: 1485--1490.

\par\hangindent=2em Baker, G.R., Meiron, D.I., and Orszag, S.A. (1982). Generalized vortex methods for free surface flows problems. J. Fluid Mech., 123: 477.

\par\hangindent=2em Bargteil, A.W., Goktekin, T.G., O’Brien, J.F., and Strain, J.A. (2006). A semi- Lagrangian contouring method for fluid simulation. ACM Trans. Graphics (TOG), 25(1): 19--38.

\par\hangindent=2em Bassano, E. and Castagnolo, D. (2003). Marangoni migration of a methanol drop in cyclohexane matrix in a closed cavity. Microgravity Sci. Technol., XIV(1): 20--33.

% [原PDF p308]

\par\hangindent=2em Batchelor, G.K. (1970). An Introduction to Fluid Dynamics. Cambridge University Press.

\par\hangindent=2em Bayvel, L. and Orzechowski, Z. (1993). Liquid Atomization. Taylor and Francis.

\par\hangindent=2em Beckermann, C., Diepers, H.-J., Steinbach, I., Karma, A., and Tong, X. (1999). Modeling melt convection in phase-field simulations of solidification. J. Comput. Phys., 154: 468--496.

\par\hangindent=2em Bell, J.B., Colella, P., and Glaz, H.M. (1989). A second-order projection method for the incompressible Navier--Stokes equations. J. Comput. Phys., 85: 257--283.

\par\hangindent=2em Ben Rayana, F., Cartellier, A., and Hopfinger, E. (2006). Assisted atomization of a liquid layer: investigation of the parameters affecting the mean drop size prediction. In Proc. ICLASS 2006, Aug. 27--Sept.1, Kyoto Japan. Academic Publication and Printings. (ISBN 4-9902774-1-4).

\par\hangindent=2em Berenson, P.J. (1961). Film boiling heat transfer from a horizontal surface. J. Heat Transfer, 83: 351--358.

\par\hangindent=2em Bierbrauer, F. (1995). Water drop impact on galvanised steel surfaces. In Proc. The Splash and Free Surface Workshop (ed. J.-L. Liow and P. Schwarz), The University of Melbourne, Melbourne, Australia.

\par\hangindent=2em Birdsall, C.K. and Langdon, A.B. (1985). Plasma Physics via Computer Simulation. McGraw-Hill.

\par\hangindent=2em Birkhoff, G. (1954). Taylor instability and laminar mixing. Rep. LA-1862; appendices in Rep. LA-1927, Los Alamos Scientific Laboratory (unpublished).

\par\hangindent=2em Blake, J.R. and Gibson, D.C. (1981). Growth and collapse of a vapour cavity near a free surface. J. Fluid Mech., 111: 123--140.

\par\hangindent=2em Blake, T.D. and Shikmurzaev, Y. (2002). Dynamic wetting by liquids of different viscosity. J. Colloid Interface Sci., 253: 196--202.

\par\hangindent=2em Boeck, T. and Zaleski, S. (2005). Viscous versus inviscid instability of two-phase mixing layers with continuous velocity profile. Phys. Fluids, 17: 032106.

\par\hangindent=2em Boeck, T., Li, J., Lo'pez-Page's, E., Yecko, P., and Zaleski, S. (2007). Ligament formation in sheared liquid--gas layers. Theor. Comput. Fluid Dyn., 21: 59--76.

\par\hangindent=2em Boettinger, W.J., Warren, J.A., Beckermann, C., and Karma, A. (2002). Phase field simulations of solidification. Annu. Rev. Mater. Res., 32: 163--194.

\par\hangindent=2em Bonn, D., Eggers, J., Indekeu, J., Meunier, J., and Rolley, E. (2009). Wetting and spreading. Rev. Mod. Phys., 81: 739--805.

\par\hangindent=2em Boris, J.P. and Book, D.L. (1973). Flux-corrected transport. I. SHASTA, a fluid transport algorithm that works. J. Comput. Phys., 11: 38--69.

\par\hangindent=2em Bothe, D., Koebe, M., Wielage, K., and Warnecke, H.J. (2003). FEDSM2003- 45155: VOF-simulations of mass transfer from single bubbles and bubble chains rising in aqueous solutions. In Proc. 2003 ASME Joint U.S.--European Fluids Eng. Conf., Honolulu, USA.

% [原PDF p309]

\par\hangindent=2em Bothe, D., Schmidtke, M., and Warnecke, H.-J. (2006). VOF-simulation of the lift force for single bubbles in a simple shear flow. Chem. Eng. Technol., 29: 1048--1053.

\par\hangindent=2em Boulton-Stone, J.M. and Blake, J.R. (1993). Gas-bubbles bursting at a free-surface. J. Fluid Mech., 254: 437--466.

\par\hangindent=2em Brackbill, J.U. and Ruppel, H.M. (1986). FLIP: a method for adaptively zoned, particle-in-cell calculations in two dimensions. J. Comput. Phys., 65: 314--343.

\par\hangindent=2em Brackbill, J.U., Kothe, D.B., and Zemach, C. (1992). A continuum method for modeling surface tension. J. Comput. Phys., 100: 335--354.

\par\hangindent=2em Bradley, S.G. and Stow, C.D. (1978). Collision between liquid drops. Phil. Trans. R. Soc. Lond. A, 287: 635.

\par\hangindent=2em Brady, J.F. and Bossis, G. (1988). Stokesian dynamics. Annu. Rev. Fluid Mech., 20: 111--157.

\par\hangindent=2em Brenner, M. and Gueyffier, D. (1999). On the bursting of viscous films. Phys. Fluids, 11: 737--739.

\par\hangindent=2em Briggs, W.L. (1987). A Multigrid Tutorial. SIAM Books, Philadelphia.

\par\hangindent=2em Bro¨der, D. and Sommerfeld, M. (2007). Planar shadow image velocimetry for the analysis of the hydrodynamics in bubbly flows. Meas. Sci. Technol., 18: 2513-- 2528.

\par\hangindent=2em Bunner, B. and Tryggvason, G. (2002a). Dynamics of homogeneous bubbly flows. Part 1. Rise velocity and microstructure of the bubbles. J. Fluid Mech., 466: 17--52.

\par\hangindent=2em Bunner, B. and Tryggvason, G. (2002b). Dynamics of homogeneous bubbly flows. Part 2. Velocity fluctuations. J. Fluid Mech., 466: 53--84.

\par\hangindent=2em Bunner, B. and Tryggvason, G. (2003). Effect of bubble deformation on the stability and properties of bubbly flows. J. Fluid Mech., 495: 77--118.

\par\hangindent=2em Carey, V.P. (2007). Liquid Vapor Phase Change Phenomena: An Introduction to the Thermophysics of Vaporization and Condensation Processes in Heat Transfer Equipment, 2nd edition, Taylor \& Francis.

\par\hangindent=2em Cartellier, A. and Rivie`re, N. (2001). Bubble-induced agitation and microstructure in uniform bubbly flows at small to moderate particle. Phys. Fluids, 13: 2165-- 2181.

\par\hangindent=2em Ceniceros, H.D. and Roma, A.M. (2005). A multi-phase flow method with a fast, geometry-based fluid indicator. J. Comput. Phys., 205: 391--400.

\par\hangindent=2em Cha, P.-H., Yeon, D.-H., and Yoon, J.-K. (2005). Phase-field model for multicomponent alloy solidification. J. Cryst. Growth, 274: 281--293.

\par\hangindent=2em Chahine, G.L. and Duraiswami, R. (1992). Dynamic interactions in a multibubble cloud. ASME J. Fluids Eng., 114(4): 680--686.

\par\hangindent=2em Chan, R.K.-C. and Street, R.L. (1970). A computer study of finite-amplitude water waves. J. Comput. Phys., 6: 68--94.

% [原PDF p310]

\par\hangindent=2em Chandra, S. and Avedisian, C.T. (1991). On the collision of a droplet with a solid surface. Proc. R. Soc. London, Ser. A: Math. Phys. Sci., 432(1884): 13--41.

\par\hangindent=2em Chandrasekhar, S. (1961). Hydrodynamic and Hydromagnetic Stability. Oxford University Press.

\par\hangindent=2em Chang, Y.C., Hou, T.Y., Merriman, B., and Osher, S. (1996). Temporal evolution of periodic disturbances in two-layer Couette flow. J. Comput. Phys., 124: 449-- 464.

\par\hangindent=2em Chapman, R.B. and Plesset, M.S. (1972). Nonlinear effects in the collapse of a nearly spherical cavity in a liquid. Trans. ASME, J. Basic Eng., 94: 142.

\par\hangindent=2em Che, J., Ceccio, S.L., and Tryggvason, G. (2004). Computations of structures formed by the solidification of impinging molten metal drops. J. Appl. Math. Model., 28: 127--144.

\par\hangindent=2em Chen, S., Johnson, D.B., Raad, P.E., and Fadda, D. (1997a). The surface marker and micro cell method. Int. J. Numer. Meth. Fluids, 25: 749--778.

\par\hangindent=2em Chen, S., Merriman, B., Osher, S., and Smereka, P. (1997b). A simple level set method for solving Stefan problems. J. Comput. Phys., 135: 8--29.

\par\hangindent=2em Chern, I.-L., Glimm, J., McBryan, O., Plohr, B., and Yaniv, S. (1986). Front tracking for gas dynamics. J. Comput. Phys., 62: 83--110.

\par\hangindent=2em Choi, H.G. and Joseph, D.D. (2001). Fluidization by lift of 300 circular particles in plane Poiseuille flow by direct numerical simulation. J. Fluid Mech., 438: 101--128.

\par\hangindent=2em Chorin, A.J. (1968). Numerical solutions of the Navier--Stokes equations. Math. Comput., 22: 745--762.

\par\hangindent=2em Chorin, A.J. (1967). A numerical method for solving incompressible viscous flow problems. J. Comput. Phys., 2: 12.

\par\hangindent=2em Christiansen, J.P. (1973). Numerical simulation of hydrodynamics by the method of point vortices. J. Comput. Phys., 13: 363--379.

\par\hangindent=2em Clift, R., Grace, J.R., and Weber, M.E. (1978). Bubbles, Drops, and Particles. Academic Press.

\par\hangindent=2em Coriell, S.R. and McFadden, G.B. (1993). Morphological stability. In Handbook of Crystal Growth, vol. 1B (ed. D.T.J. Hurle), Elsevier, pp. 785--857.

\par\hangindent=2em Cortelezzi, L. and Prosperetti, A. (1981). Small amplitude waves on the surface of a layer of a viscous liquid. Q. Appl. Math., 38: 375--388.

\par\hangindent=2em Cortez, R. and Minion, M. (2000). The blob projection method for immersed boundary problems. J. Comput. Phys., 161: 428--453.

\par\hangindent=2em Coward, A.V., Renardy, Y.Y., Renardy, M., and Richards, J.R. (1997). Temporal evolution of periodic disturbances in two-layer Couette flow. J. Comput. Phys., 132: 346--361.

% [原PDF p311]

\par\hangindent=2em Crapper, G.D., Dombrowski, N., Jepson, W.P., and Pyott, G.A.D. (1973). A note on the growth of Kelvin--Helmholtz waves on thin liquid sheets. J. Fluid. Mech., 57: 671--672.

\par\hangindent=2em Cristini, V., Bławzdziewicz, J., and Loewenberg, M. (1998). Drop breakup in threedimensional viscous flows. Phys. Fluids., 10: 1781--1784.

\par\hangindent=2em Cummins, S.J., Francois, M.M., and Kothe, D.B. (2005). Estimating curvature from volume fractions. Comput. Struct., 83: 425--434.

\par\hangindent=2em Daly, B.J. (1969a). A technique for including surface tension effects in hydrodynamic calculations. J. Comput. Phys., 4: 97--117.

\par\hangindent=2em Daly, B.J. (1969b). Numerical study of the effect of surface tension on interface instability. Phys. Fluids, 12: 1340--1354.

\par\hangindent=2em Daly, B.J. and Pracht, W.E. (1968). Numerical study of density-current surges. Phys. Fluids, 11: 15--30.

\par\hangindent=2em Dandy, D.S. and Leal, G.L. (1989). Buoyancy-driven motion of a deformable drop through a quiescent liquid at intermediate Reynolds numbers. J. Fluid Mech., 208: 161--192.

\par\hangindent=2em Davidson, M.R. and Rudman, M. (2002). Volume-of-fluid calculation of heat or mass transfer across deforming interfaces in two-fluid flow. Numer. Heat Transfer Part B Fund., 41: 291--308.

\par\hangindent=2em Davies, J.T. and Rideal, E.K. (1966). Interfacial Phenomena. Academic Press.

\par\hangindent=2em Davis, R.H., Schonberg, J.A., and Rallison, J.M. (1989). The lubrication force between two viscous drops. Phys. Fluids A: Fluid Dyn., 1: 77--81.

\par\hangindent=2em de Gennes, P.G. (1985). Wetting: statics and dynamics. Rev. Mod. Phys., 57: 827--863.

\par\hangindent=2em de Gennes, P.G., Quere, D., and Brochard, F. (2003). Capillarity and Wetting Phenomena: Drops, Bubbles, Pearls, Waves. Springer Verlag.

\par\hangindent=2em de Josselin de Jong, G. (1960). Singularity distribution for the analysis of multiplefluid flow through porous media. J. Geophys. Res., 65: 3739--3758.

\par\hangindent=2em DeBar, R. (1974). Fundamentals of the KRAKEN code. Technical Report UCIR- 760, LLNL.

\par\hangindent=2em Deckwer, W.-D. (1992). Bubble Column Reactors. Wiley.

\par\hangindent=2em Deen, N.G., van Sint Annaland, M., and Kuipers, J.A.M. (2004). Multi-scale modeling of dispersed gas--liquid two-phase flow. Chem. Eng. Sci., 59: 1853--1861.

\par\hangindent=2em DeGregoria, A.J. and Schwartz, L.W. (1985). Finger breakup in Hele--Shaw cells. Phys. Fluids, 28: 2313.

\par\hangindent=2em Delale, C.F., Nas, S., and Tryggvason, G. (2005). Direct numerical simulations of shock propagation in bubbly liquids. Phys. Fluids, 17: 121705.

\par\hangindent=2em Dhir, V.K. (1998). Boiling heat transfer. Annu. Rev. Fluid Mech., 30: 365--401.

\par\hangindent=2em Ding, H., Spelt, P.D.M., and Shu, C. (2007). Diffuse interface model for incompressible two-phase flows with large density ratios. J. Comput. Phys., 226: 2078--2095.

% [原PDF p312]

\par\hangindent=2em Dombrowski, N. and Johns, W.R. (1963). The aerodynamic instability and disintegration of viscous liquid sheets. Chem. Eng. Sci., 18: 203--214.

\par\hangindent=2em Drew, D.A. (1983). Mathematical modeling of two-phase flow. Annu. Rev. Fluid Mech., 15: 261--291.

\par\hangindent=2em Drew, D.A. and Passman, S.L. (1999). Theory of Multicomponent Fluids. Springer.

\par\hangindent=2em Drumright-Clarke, M.A. and Renardy, Y. (2004). The effect of insoluble surfactant at dilute concentration on drop breakup under shear with inertia. Phys. Fluids, 16: 14--21.

\par\hangindent=2em Druzhinin, O.A. and Elghobashi, S.E. (2001). Direct numerical simulation of a three-dimensional spatially-developing bubble-laden mixing layer with two-way coupling. J. Fluid Mech., 429: 23--61.

\par\hangindent=2em Du, J., Fix, B., Glimm, J., Jia, X., Li, X., Li, Y., and Wu, L. (2006). A simple package for front tracking. J. Comput. Phys., 213: 613--628.

\par\hangindent=2em Duchemin, L., Popinet, S., Josserand, C., and Zaleski, S. (2002). Jet formation in bubbles bursting at a free surface. Phys. Fluids, 14(9): 3000--3008.

\par\hangindent=2em Dussan, V.E.B. (1979). On the spreading of liquid on solid surfaces: static and dynamic contact lines. Annu. Rev. Fluid Mech., 11: 371--400.

\par\hangindent=2em E, W. and Engquist, B. (2003). The heterogeneous multiscale methods. Commun. Math. Sci., 1: 87--133.

\par\hangindent=2em Edgerton, H. (1987). Stopping Time. Harry N. Abrams, New York.

\par\hangindent=2em Edwards, D.A., Brenner, H., and Wasan, D.T. (1961). Interfacial Transport Processes and Rheology. Butterworth.

\par\hangindent=2em Eggers, J. (1993). Universal pinching of 3D axisymmetric free-surface flow. Phys. Rev. Lett., 71: 3458--3461.

\par\hangindent=2em Eggers, J. and Dupont, T.D. (1994). Drop formation in a one-dimensional approximation of the Navier--Stokes equation. J. Fluid Mech., 262: 205--221.

\par\hangindent=2em Eggers, J. and Villermaux, E. (2008). Physics of liquid jets. Rep. Prog. Phys., 71: 036601.

\par\hangindent=2em Ellingsen, K. and Risso, F. (2001). On the rise of an ellipsoidal bubble in water: oscillatory paths and liquid-induced velocity. J. Fluid Mech., 440: 235--268.

\par\hangindent=2em Ervin, E.A. and Tryggvason, G. (1997). The rise of bubbles in a vertical shear flow. ASME J. Fluid Eng., 119: 443--449.

\par\hangindent=2em Esmaeeli, A. and Tryggvason, G. (1996). An inverse energy cascade in twodimensional, low Reynolds number bubbly flows. J. Fluid Mech., 314: 315--330.

\par\hangindent=2em Esmaeeli, A. and Tryggvason, G. (1998). Direct numerical simulations of bubbly flows. Part I. Low Reynolds number arrays. J. Fluid Mech., 377: 313--345.

\par\hangindent=2em Esmaeeli, A. and Tryggvason, G. (1999). Direct numerical simulations of bubbly flows. Part II. Moderate Reynolds number arrays. J. Fluid Mech., 385: 325--358.

\par\hangindent=2em Esmaeeli, A. and Tryggvason, G. (2004a). A front tracking method for computations of boiling in complex geometries. Int. J. Multiphase Flow, 30: 1037--1050.

% [原PDF p313]

\par\hangindent=2em Esmaeeli, A. and Tryggvason, G. (2004b). Computations of film boiling. Part I: numerical method. Int. J. Heat Mass Transfer, 47: 5451--5461.

\par\hangindent=2em Esmaeeli, A. and Tryggvason, G. (2004c). Computations of film boiling. Part II: multi-mode film boiling. Int. J. Heat Mass Transfer, 47: 5463--5476.

\par\hangindent=2em Esmaeeli, A. and Tryggvason, G. (2005). A direct numerical simulation study of the buoyant rise of bubbles at O(100) Reynolds number. Phys. Fluids, 17: 093303.

\par\hangindent=2em Fadlun, E.A., Verzicco, R., Orlandi, P., and Mohd-Yusof, J. (2000). Combined immersed-boundary finite-difference methods for three-dimensional complex flow simulations. J. Comput. Phys., 161: 35--60.

\par\hangindent=2em Farmer, E.E. (1973). Relative detachability of soil particles by simulated rainfall. Soil Sci. Soc. Am. J., 37: 629--633.

\par\hangindent=2em Fauci, L.J. and Dillon, R. (2006). Biofluidmechanics of reproduction. Annu. Rev. Fluid Mech., 38: 371--394.

\par\hangindent=2em Fedkiw, R., Aslam, T., Merriman, B., and Osher, S. (1999). A non-oscillatory Eulerian approach to interfaces in multimaterial flows (the ghost fluid method). J. Comput. Phys., 152: 457--492.

\par\hangindent=2em Feng, J., Hu, H.H., and Joseph, D.D. (1994). Direct simulation of initial value problems for the motion of solid bodies in a Newtonian fluid. Part 1. Sedimation. J. Fluid Mech., 261: 95--134.

\par\hangindent=2em Feng, J., Hu, H.H., and Joseph, D.D. (1995). Direct simulation of initial value problems for the motion of solid bodies in a Newtonian fluid. Part 2. Couette and Poiseuille flows. J. Fluid Mech., 277: 271--301.

\par\hangindent=2em Fernandez, A. (2008a). Response of an emulsion of leaky dielectric drops immersed in a simple shear flow: drops more conductive than the suspending fluid. Phys. Fluids, 20: 043303.

\par\hangindent=2em Fernandez, A. (2008b). Response of an emulsion of leaky dielectric drops immersed in a simple shear flow: drops less conductive than the suspending fluid. Phys. Fluids, 20: 043304.

\par\hangindent=2em Fernandez, A., Che, J., Ceccio, S.L., and Tryggvason, G. (2005). The effects of electrostatic forces on the distribution of drops in a channel flow -- twodimensional oblate drops. Phys. Fluids, 17: 093302.

\par\hangindent=2em Foote, G.B. (1973). A numerical method for studying liquid drop behavior: simple oscillations. J. Comput. Phys., 11: 507--530.

\par\hangindent=2em Foote, G.B. (1975). The water drop rebound problem: dynamics of collision. J. Atmos. Sci., 32: 390--402.

\par\hangindent=2em Fortes, A., Joseph, D.D., and Lundgren, T. (1987). Nonlinear mechanics of fluidization of beds of spherical particles. J. Fluid Mech., 177: 467--483.

% [原PDF p314]

\par\hangindent=2em Foster, N. and Fedkiw, R. (2001). Practical animation of liquids. In Proc. 28th Annual Conf. on Computer Graphics and Interactive Techniques. ACM Press, New York, NY, pp. 23--30.

\par\hangindent=2em Francois, M.M., Cummins, S.J., Dendy, E.D., Kothe, D.B., Sicilian, J.M., and Williams, M.W. (2006). A balanced-force algorithm for continuous and sharp interfacial surface tension models within a volume tracking framework. J. Comput. Phys., 213(1): 141--173.

\par\hangindent=2em Frumkin, A. and Levich, V.G. (1947). O vliyanii poverkhnostno-aktivnikh veshestv na dvizhenie na granitse zhidkikh sred. Zh. Fiz. Khim., 21: 1183--1204.

\par\hangindent=2em Fukai, J., Shiiba, Y., Yamamoto, T., Miyatake, O., Poulikakos, D., Megaridis, C.M., and Zhao, Z. (1995). Wetting effects on the spreading of a liquid droplet colliding with a flat surface: experiment and modeling. Phys. Fluids, 7: 236--247.

\par\hangindent=2em Furusaki, S., Fan, L.-S., and Garside, J. (2001). The Expanding World of Chemical Engineering (2nd edition). Taylor and Francis.

\par\hangindent=2em Fuster, D., Agbaglah, G., Josserand, C., Popinet, S., and Zaleski, S. (2009a). Numerical simulation of droplets, bubbles and waves: state of the art. Fluid Dyn. Res. 41: 065001.

\par\hangindent=2em Fuster, D., Bague', A., Boeck, T., Le Moyne, L., Leboissetier, A., Popinet, S., Ray, P., Scardovelli, R., and Zaleski, S. (2009b). Simulation of primary atomization with an octree adaptive mesh refinement and VOF method. Int. J. Multiphase Flow, 35: 550--565.

\par\hangindent=2em Gao, P., Yin, Z., and Hu, W. (2008). Thermocapillary motion of droplets at large Marangoni numbers. Adv. Space Res., 41(12): 2101--2106.

\par\hangindent=2em Ge, Y. and Fan, L.-S. (2006). Three-dimensional direct numerical simulation for film-boiling contact of moving particle and liquid droplet. Phys. Fluids, 18: 117104.

\par\hangindent=2em Gerlach, D., Tomar, G., Biswas, G., and Durst, F. (2006). Comparison of volumeof-fluid methods for surface tension-dominant two-phase flows. Int. J. Heat Mass Transfer, 49: 740--754.

\par\hangindent=2em Gibou, F., Chen, L., Nguyen, D., and Banerjee, S. (2007). A level set based sharp interface method for the multiphase incompressible Navier--Stokes equations with phase change. J. Comput. Phys., 222: 536--555.

\par\hangindent=2em Gilmanov, A., Sotiropoulos, F., and Balaras, E. (2003). A general reconstruction algorithm for simulating flows with complex 3D immersed boundaries on Cartesian grids. J. Comput. Phys., 191: 660--669.

\par\hangindent=2em Glicksman, M.E., Coriell, S.R., and McFadden, G.B. (1986). Interaction of flows with the crystal--melt interface. Annu. Rev. Fluid Mech., 18: 307--335.

\par\hangindent=2em Glimm, J. and McBryan, O. (1985). A computational model for interfaces. Adv. Appl. Math., 6: 422--435.

\par\hangindent=2em Glimm, J., Marchesin, D., and McBryan, O. (1981). A numerical method for two phase flow with an unstable interface. J. Comput. Phys, 39: 179--200.

% [原PDF p315]

\par\hangindent=2em Glimm, J., Grove, J.W., Li, X.L., Oh, W., and Sharp, D.H. (2001). A critical analysis of Rayleigh--Taylor growth rates. J. Comput. Phys., 169: 652--677.

\par\hangindent=2em Glowinski, R., Pan, T.W., Helsa, T.I., Joseph, D.D., and Periaux, J. (2001). A fictitious domain approach to the direct numerical simulation of incompressible viscous flow past moving rigid bodies: application to particulate flow. J. Comput. Phys., 169: 363--426.

\par\hangindent=2em Goldstein, D., Handler, R., and Sirovich, L. (1993). Modeling a no-slip flow boundary with an external force field. J. Comput. Phys., 105: 354--366.

\par\hangindent=2em Gordillo, J.M. and Perez-Saborid, M. (2005). On the first wind atomization regime. J. Fluid. Mech, 541: 1--20.

\par\hangindent=2em Guet, S., Ooms, G., Oliemans, R.V.A., and Mudde, R.F. (2004). Bubble size effect on low liquid input drift-flux parameters. Chem. Eng. Sci., 59: 3315--3329.

\par\hangindent=2em Guet, S., Ooms, G., and Oliemans, R.V.A. (2005). Simplified two-fluid model for gas-lift efficiency predictions. AIChe J., 51: 1885--1896.

\par\hangindent=2em Gueyffier, D. and Zaleski, S. (1998). Formation de digitations lors de l’impact d’une goutte sur un film liquide. C. R. Acad. Sci. Paris se'r. II b, 326: 839--844.

\par\hangindent=2em Gueyffier, D., Nadim, A., Li, J., Scardovelli, R., and Zaleski, S. (1999). Volume of fluid interface tracking with smoothed surface stress methods for threedimensional flows. J. Comput. Phys., 152: 423--456.

\par\hangindent=2em Hadamard, J. (1911). Mouvement permanent lent d’une sphere liquide et visqueuse dans un liquide visqueux. C. R. Acad. Sci. Paris, 152: 1735--1738.

\par\hangindent=2em Hammerlin, G. and Hoffmann, K.-H. (1991). Numerical Mathematics. Springer.

\par\hangindent=2em Harlow, F.H. (1964). The particle-in-cell computing method for fluid dynamics. Methods Comput. Phys., 3: 319--343.

\par\hangindent=2em Harlow, F.H. and Fromm, J.E. (1965). Computer experiments in fluid dynamics. Sci. Am., 212: 104.

\par\hangindent=2em Harlow, F.H. and Shannon, J.P. (1967). The splash of a liquid drop. J. Appl. Phys., 38: 3855--3866.

\par\hangindent=2em Harlow, F.H. and Welch, J.E. (1965). Numerical calculation of time-dependent viscous incompressible flow of fluid with a free surface. Phys. Fluids, 8: 2182-- 2189.

\par\hangindent=2em Harlow, F.H. and Welch, J.E. (1966). Numerical study of large-amplitude freesurface motions. Phys. Fluids, 9: 842--851.

\par\hangindent=2em Harvie, D.J.E. and Fletcher, D.F. (2000). A new volume of fluid advection algorithm: the Stream scheme. J. Comput. Phys., 162: 1.

\par\hangindent=2em Hervet, H. and de Gennes, P.G. (1984). Dynamique du mouillage: films pre'curseurs sur solide ‘sec’. C. R. Acad. Sci. Paris, Ser. II, 299: 499--503.

\par\hangindent=2em Hinch, E.J. (1984). A note on the mechanism of the instability at the interface between two shearing fluids. J. Fluid Mech., 144: 463--465.

% [原PDF p316]

\par\hangindent=2em Hirt, C.W. and Nichols, B.D. (1981). Volume of fluid (VOF) method for the dynamics of free boundaries. J. Comput. Phys., 39: 201--226.

\par\hangindent=2em Hirt, C.W., Cook, J.L., and Butler, T.D. (1970). A Lagrangian method for calculating the dynamics of an incompressible fluid with a free surface. J. Comput. Phys., 5: 103--124.

\par\hangindent=2em Hockney, R.W. and Eastwood, J.W. (1981). Computer Simulation Using Particles. McGraw-Hill.

\par\hangindent=2em Hooper, A.P. and Boyd, W.G.C. (1983). Shear-flow instability at the interface between two viscous fluids. J. Fluid Mech., 128: 507--528.

\par\hangindent=2em Hou, T.Y., Lowengrub, J.S., and Shelley, M.J. (2001). Boundary integral methods for multicomponent fluids and multiphase materials. J. Comput. Phys., 169: 302--362.

\par\hangindent=2em Hu, H.H. (1996). Direct simulation of flows of solid--liquid mixtures. Int. J. Multiphase Flow, 22: 335--352.

\par\hangindent=2em Hu, H.H., Patankar, N.A., and Zhu, M.Y. (2001). Direct numerical simulations of fluid-solid systems using arbitrary-Lagrangian--Eularian techniques. J. Comput. Phys., 169: 427--462.

\par\hangindent=2em Hua, J. and Lou, J. (2007). Numerical simulation of bubble rising in viscous liquid. J. Comput. Phys., 222: 769--795.

\par\hangindent=2em Hua, J., Lim, L.K., and Wang, C.-H. (2008). Numerical simulation of deformation/motion of a drop suspended in viscous liquids under influence of steady electric fields. Phys. Fluids, 20: 113302.

\par\hangindent=2em Huh, C. and Scriven, L. (1971). Hydrodynamic model of steady movement of a solid/liquid/fluid contact line. J. Coll. Interf. Sci, 35: 85.

\par\hangindent=2em Ishii, M. (1987). Interfacial Area Modeling, Vol. 3, McGraw-Hill, Chapter 3, pp. 31--62.

\par\hangindent=2em Ishii, M. and Zuber, N. (1979). Drag coefficient and relative velocity in bubbly, droplet or particulate flows. AIChE J., 25: 843--855.

\par\hangindent=2em Issa, R.I. (1985). Solution of the implicitly discretised fluid flow equations by operator-splitting. J. Comput. Phys., 62: 40--65.

\par\hangindent=2em Jacqmin, D. (1999). Calculation of two-phase Navier--Stokes flows using phasefield modeling. J. Comput. Phys., 155: 96--127.

\par\hangindent=2em Jacqmin, D. (2000). Contact-line dynamics of a diffuse fluid interface. J. Fluid Mech., 402: 57--88.

\par\hangindent=2em James, A.J. and Lowengrub, J. (2004). A surfactant-conserving volume-of-fluid method for interfacial flows with insoluble surfactant. J. Comput. Phys., 201: 685--722.

\par\hangindent=2em Jamet, D., Lebaigue, O., Coutris, N., and Delhaye, J.M. (2001). Feasibility of using the second gradient theory for the direct numerical simulations of liquid--vapor flows with phase-change. J. Comput. Phys., 169: 624--651.

% [原PDF p317]

\par\hangindent=2em Jan, Y.-J. (1994). Computational studies of bubble dynamics. PhD thesis, The University of Michigan.

\par\hangindent=2em Jeong, J.-H., Goldenfield, N., and Dantzig, J.A. (2001). Phase field model for three-dimensional dendritic growth with fluid flow. Phys. Rev. E, 64: 041602.

\par\hangindent=2em Johansson, B.C.V. (1993). Boundary conditions for open boundaries for the incompressible Navier--Stokes equation. J. Comput. Phys., 105: 233--251.

\par\hangindent=2em Johnson, A. A. and Tezduyar, T.E. (1997). 3D simulation of fluid--particle interactions with the number of particles reaching 100. Comput. Methods Appl. Mech. Eng., 145: 301--321.

\par\hangindent=2em Josserand, C. and Zaleski, S. (2003). Droplet splashing on a thin liquid film. Phys. Fluids, 15: 1650--1657.

\par\hangindent=2em Juric, D. and Tryggvason, G. (1998). Computations of boiling flows. Int. J. Multiphase Flow, 24: 387--410.

\par\hangindent=2em Kadioglu, S.Y. and Sussman, M. (2008). Adaptive solution techniques for simulating underwater explosions and implosions. J. Comput. Phys., 227: 2083--2104.

\par\hangindent=2em Kajishima, T. and Takiguchi, S. (2002). Interaction between particle clusters and fluid turbulence. Int. J. Heat Fluid Flow, 23: 639--646.

\par\hangindent=2em Kanai, A. and Miyata, H. (2001). Direct numerical simulation of wall turbulent flows with microbubbles. Int. J. Numer. Methods Fluids, 35: 593--615.

\par\hangindent=2em Kang, I.S. and Leal, L.G. (1987). Numerical solution of axisymmetric, unsteady free-boundary problems at finite Reynolds number. I. Finite-difference scheme and its applications to the deformation of a bubble in a uniaxial straining flow. Phys. Fluids, 30: 1929--1940.

\par\hangindent=2em Karma, A. and Rappel, W.J. (1997). Phase-field simulation of three-dimensional dendrites: is microscopic solvability theory correct? J. Cryst. Growth, 174: 54--64.

\par\hangindent=2em Karma, A. and Rappel, W.-J. (1998). Quantitative phase-field modeling of dendritic growth in two and three dimensions. Phys. Rev. E, 57: 4323--4349.

\par\hangindent=2em Kashinsky, O.N. and Randin, V.V. (1999). Downward bubbly gas--liquid flow in a vertical pipe. Int. J. Multiphase Flow, 25: 109--138.

\par\hangindent=2em Kawamura, T. and Kodama, Y. (2002). Numerical simulation method to resolve interactions between bubbles and turbulence. Int. J. Heat Fluid Flow, 23: 627-- 638.

\par\hangindent=2em Keller, J.B. and Miksis, M.J. (1983). Surface tension driven flows. SIAM J. Appl. Math., 43(2): 268--277.

\par\hangindent=2em Kermeen, R.W. (1956). Water tunnel tests of NACA 4412 and Walchner profile 7 hydrofoils in non-cavitating and cavitating flows. Technical report, California Institute of Technology Hydro. Lab. Rep. 47-5.

% [原PDF p318]

\par\hangindent=2em Khinast, J.G., Koynov, A.A., and Leib, T.M. (2003). Reactive mass transfer at gas-- liquid interfaces: impact of micro-scale fluid dynamics on yield and selectivity of liquid-phase cyclohexane oxidation. Chem. Eng. Sci., 58: 3961--3971.

\par\hangindent=2em Kim, J. and Moin, P. (1985). Application of a fractional step method to incompressible Navier--Stokes equations. J. Comput. Phys., 59: 308--323.

\par\hangindent=2em Kim, J., Kim, D., and Choi, H. (2001). An immersed-boundary finite-volume method for simulations of flow in complex geometries. J. Comput. Phys., 171: 132--150.

\par\hangindent=2em Kitagawa, A., Sugiyama, K., and Murai, Y. (2004). Experimental detection of bubble--bubble interactions in a wall-sliding bubble swarm. Int. J. Multiphase Flow, 30: 1213--1234.

\par\hangindent=2em Klimenko, V.V. (1981). Film boiling on a horizontal plate -- new correlation. Int. J. Heat Mass Transfer, 24: 69--79.

\par\hangindent=2em Kobayashi, R. (1992). Simulations of three-dimensional dendrites. In Pattern Formation in Complex Dissipative Systems (ed. S. Kai). World Scientific, Singapore, pp. 121--128.

\par\hangindent=2em Kobayashi, R. (1993). Modeling and numerical simulations of dendritic crystalgrowth. Phys. D, 63: 410--423.

\par\hangindent=2em Koplik, J., Banavar, J.R., and Willemsen, J.F. (1988). Molecular dynamics of Poiseuille flow and moving contact lines. Phys. Rev. Lett., 60: 1282--1285.

\par\hangindent=2em Koynov, A., Khinast, J.G., and Tryggvason, G. (2005). Mass transfer and chemical reactions in bubble swarms with dynamic interfaces. AIChE. J., 51: 2786--2800.

\par\hangindent=2em Kuo, T.C., Pan, C., and Chieng, C.C. (1997). Eulerian--Lagrangian computations on phase distribution of two-phase bubbly flows. Int. J. Numer. Methods Fluids, 24: 579--593.

\par\hangindent=2em Ladd, A.J.C. (1993). Dynamical simulations of sedimenting spheres. Phys. Fluids A, 5: 299--310.

\par\hangindent=2em Lafaurie, B., Nardone, C., Scardovelli, R., Zaleski, S., and Zanetti, G. (1994). Modelling merging and fragmentation in multiphase flows with SURFER. J. Comput. Phys., 113: 134--147.

\par\hangindent=2em Lai, M.-C. and Peskin, C.S. (2000). An immersed boundary method with formal second order accuracy and reduced numerical viscosity. J. Comput. Phys., 160: 705--719.

\par\hangindent=2em Lance, M. and Bataille, J. (1991). Turbulence in the liquid phase of a uniform bubbly air--water flow. J. Fluid Mech., 222: 95--118.

\par\hangindent=2em Lasheras, J.C. and Hopfinger, E.J. (2000). Liquid jet instability and atomization in a coaxial gas stream. Annu. Rev. Fluid Mech., 32: 275--308.

\par\hangindent=2em Leboissetier, A. (2002). Simulation nume'rique directe de l’atomisation primaire d’un jet liquide a` haute vitesse. Technical report, Universite' de Paris 6. The`se de doctorat.

% [原PDF p319]

\par\hangindent=2em Lee, J. and Pozrikidis, C. (2006). Effect of surfactants on the deformation of drops and bubbles in Navier--Stokes flow. Comput. Fluids, 35: 43--60.

\par\hangindent=2em Lee, L. and LeVeque, R.J. (2003). An immersed interface method for incompressible Navier--Stokes equations. SIAM J. Sci. Comput., 25: 832--856. .

\par\hangindent=2em Lee, R.C. and Nydahl, J.E. (1989). Numerical calculations of bubble growth in nucleate boiling from inception through departure. J. Heat Transfer, 111: 474-- 479.

\par\hangindent=2em Lee, Y.-W., Ananth, R., and Gill, W.N. (1993). Selection of length scale in unconstrained dendritic growth with convection in the melt. J. Cryst. Growth, 132: 226--230.

\par\hangindent=2em Lefebvre, A.H. (1989). Atomization and Sprays. Taylor and Francis.

\par\hangindent=2em Lenard, P. (1892). U¨ ber die Electricita¨t der Wasserfa¨lle. Ann. Phys., 46: 584--636.

\par\hangindent=2em Leonard, B.P. (1979). A stable and accurate convective modelling procedure based on quadratic upstream interpolation. Comput. Methods Appl. Mech. Eng., 19: 59--98.

\par\hangindent=2em Li, J. (1995). Calcul d’interface affine par morceaux (Piecewise linear interface calculation). C. R. Acad. Sci. Paris, Se'r. IIb (Paris), 320: 391--396.

\par\hangindent=2em Li, J. (1996). Re'solution nume'rique de l’e'quation de Navier--Stokes avec reconnection d’interfaces. Me'thode de suivi de volume et application a` l’atomisation. Technical report, Universite' de Paris 6. The`se de doctorat.

\par\hangindent=2em Liow, J.L., Morton, D.E., Guerra, E., and Gray, N.B. (1996). Dynamics of splash formation in gas injected systems. In The Howard Worner Int. Symp. on Injection in Pyrometallurgy, Melbourne, pp. 175--190.

\par\hangindent=2em Lister, J. and Zhang, W.W. (1999). Similarity solutions for van der Waals rupture of a thin film on a solid substrate. Phys. Fluids, 11: 2454--2462.

\par\hangindent=2em Liu, H., Krishnan, S., Marella, S., and Udaykumar, H.S. (2005). Sharp interface Cartesian grid method II: a technique for simulating droplet interactions with surfaces of arbitrary shape. J. Comput. Phys., 210: 32--54.

\par\hangindent=2em Liu, T.J. (1997). Investigation of the wall shear stress in vertical bubbly flow under different bubble size conditions. Int. J. Multiphase Flow, 23: 1085--1109.

\par\hangindent=2em Liu, T.J. and Bankoff, S.G. (1993). Structure of air-water bubbly flow in a vertical pipe -- I. Liquid mean velocity and turbulence measurements. Int. J. Heat Mass Transfer, 36: 1049--1060.

\par\hangindent=2em Loewenberg, M. and Hinch, E.J. (1996). Numerical simulation of a concentrated emulsion in shear flow. J. Fluid Mech., 321: 395--419.

\par\hangindent=2em Longuet-Higgins, M.S. and Cokelet, E.D. (1976). The deformation of steep surface waves on water. Proc. R. Soc. London Ser. A, 358: 1.

\par\hangindent=2em Lo'pez, J., Hernandez, J., Go'mez, P., and Faura, F. (2004). A volume of fluid method based on multidimensional advection and spline interface reconstruction. J. Comput. Phys., 195: 718--742.

% [原PDF p320]

\par\hangindent=2em Lopez De Bertodano, M., Lahey, R.T., Jr., and Jones, O.C. (1987). Development of a k--ε model for bubbly two-phase flow. J. Fluids Eng., 13: 327--343.

\par\hangindent=2em Lopez De Bertodano, M., Lahey, R.T., Jr., and Jones, O.C. (1994). Phase distribution in bubbly two-phase flow in vertical ducts. Int. J. Multiphase Flow, 20: 805--818.

\par\hangindent=2em Louge, M.Y. (1994). Computer simulations of rapid granular flows of spheres interacting with a flat, frictional boundary. Phys. Fluids, 6: 2253--2269.

\par\hangindent=2em Lu, J. and Tryggvason, G. (2006). Numerical study of turbulent bubbly downflows in a vertical channel. Phys. Fluids, 18: 103302.

\par\hangindent=2em Lu, J. and Tryggvason, G. (2007). Effect of bubble size in turbulent bubbly downflow in a vertical channel. Chem. Eng. Sci., 62: 3008--3018.

\par\hangindent=2em Lu, J. and Tryggvason, G. (2008). Effect of bubble deformability in turbulent bubbly upflow in a vertical channel. Phys. Fluids, 20: 040701.

\par\hangindent=2em Lu, J., Fernandez, A., and Tryggvason, G. (2005a). The effect of bubbles on the wall shear in a turbulent channel flow. Phys. Fluids, 17: 095102.

\par\hangindent=2em Lu, J., Biswas, S., and Tryggvason, G. (2006). A DNS study of laminar bubbly flows in a vertical channel. Int. J. Multiphase Flow, 32: 643--660.

\par\hangindent=2em Lu, Y., Beckermann, C., and Ramirez, J.C. (2005b). Three-dimensional phase field simulations of the effect of convection on free dendritic growth. J. Cryst. Growth, 280: 320--334.

\par\hangindent=2em Luo, R., Pan, X.H., and Yang, X.Y. (2003). Laminar light particle and liquid twophase flows in a vertical pipe. Int. J. Multiphase Flow, 29: 603--620.

\par\hangindent=2em MacIntyre, F. (1972). Flow patterns in breaking bubbles. J. Geophys. Res., 77(27): 5211--5228.

\par\hangindent=2em MacNeice, P., Olson, K.M., Mobarry, C., de Fainchtein, R., and Packer, C. (2000). PARAMESH: a parallel adaptive mesh refinement community toolkit. Comput. Phys. Commun., 126: 330--354.

\par\hangindent=2em Marella, S., Krishnan, S., Liu, H., and Udaykumar, H.S. (2005). Sharp interface Cartesian grid method I: an easily implemented technique for 3D moving boundary computations. J. Comput. Phys., 214: 1--31.

\par\hangindent=2em Marmottant, P. and Villermaux, E. (2001). Ligament mediated drop formation. Phys. Fluids, 13: S7.

\par\hangindent=2em Marmottant, P. and Villermaux, E. (2002). Atomisation primaire dans les jets coaxiaux. Combustion (Rev. des Sci. Tech. Combustion), 2: 89--126.

\par\hangindent=2em Martin, D.F. (1998). An adaptive cell-centered projection method for the incompressible Euler equations. PhD thesis, University of California, Berkeley.

\par\hangindent=2em Meiburg, E. and Homsy, G.M. (1988). Nonlinear unstable viscous fingers in Hele-- Shaw flows. II. Numerical simulation. Phys. Fluids, 31: 429.

% [原PDF p321]

\par\hangindent=2em Melcher, J.R. and Taylor, G.I. (1969). Electrohydrodynamics: a review of the role of interfacial shear stresses. Annu. Rev. Fluid Mech., 1: 111--146.

\par\hangindent=2em Menchaca-Rocha, A., Huidobro, F., Martinez-Davalos, A., Michaelian, K., Perez, A., Rodriguez, V., and Carjan, N. (2000). Coalescence and fragmentation of colliding mercury drops. J. Fluid Mech., 346: 291--318.

\par\hangindent=2em Miller, G.H. and Colella, P. (2002). A conservative three-dimensional Eulerian method for coupled solid--fluid shock capturing. J. Comput. Phys., 183: 26--82.

\par\hangindent=2em Mitchell, T.M. and Hammitt, F.H. (1973). Asymmetric cavitation bubble collapse. Trans. ASME, J. Fluids Eng., 95: 29.

\par\hangindent=2em Mittal, R. and Iaccarino, G. (2005). Immersed boundary methods. Annu. Rev. Fluid Mech., 37: 239--261.

\par\hangindent=2em Morton, D., Rudman, M., and Jong-Leng, L. (2000). An investigation of the flow regimes resulting from splashing drops. Phys. Fluids, 12(4): 747--763.

\par\hangindent=2em Mukherjee, A. and Dhir, V.K. (2004). Study of lateral merger of vapor bubbles during nucleate pool boiling. J. Heat Transfer, 126: 1023--1039.

\par\hangindent=2em Mukherjee, A. and Kandlikar, S.G. (2007). Numerical study of single bubbles with dynamic contact angle during nucleate pool boiling. Int. J. Heat Mass Transfer, 50: 127--138.

\par\hangindent=2em Muradoglu, M. and Kayaalp, A.D. (2006). An auxiliary grid method for computations of multiphase flows in complex geometries. J. Comput. Phys., 214: 858--877.

\par\hangindent=2em Muradoglu, M. and Tryggvason, G. (2008). A front-tracking method for computation of interfacial flows with soluble surfactants. J. Comput. Phys., 227: 2238--2262.

\par\hangindent=2em Nakamura, T. and Yabe, T. (1999). Cubic interpolated propagation scheme for solving the hyper-dimensional Vlasov--Poisson equation in phase space. Comput. Phys. Commun., 120: 122.

\par\hangindent=2em Nas, S. and Tryggvason, G. (2003). Thermocapillary interaction of two bubbles or drops. Int. J. Multiphase Flow, 29: 1117--1135.

\par\hangindent=2em Nas, S., Muradoglu, M., and Tryggvason, G. (2006). Pattern formation of drops in thermocapillary migration. Int. J. Heat Mass Transfer, 49: 2265--2276.

\par\hangindent=2em Navier, C.L. (1823). Sur les lois du mouvement des fluides. Me'm. Acado. Sci. France, 6: 389--440 (appeared in 1827).

\par\hangindent=2em Nestler, B. (2005). A 3D parallel simulator for crystal growth and solidification in complex alloy systems. J. Cryst. Growth, 275: 273--278.

\par\hangindent=2em Nie, X.B., Chen, S.Y., E, W.N., and Robbins, M.O. (2004). A continuum and molecular dynamics hybrid method for micro-and nano-fluid flow. J. Fluid Mech., 500: 55--64.

% [原PDF p322]

\par\hangindent=2em Nikolopoulos, N., Theodorakakos, A., and Bergeles, G. (2007). Three-dimensional numerical investigation of a droplet impinging normally onto a wall film. J. Comput. Phys., 227: 1428--1469.

\par\hangindent=2em Nobari, M.R. and Tryggvason, G. (1996). Numerical simulations of threedimensional drop collisions. AIAA J., 34: 750--755.

\par\hangindent=2em Nobari, M.R., Jan, Y.-J., and Tryggvason, G. (1996). Head-on collision of drops -- a numerical investigation. Phys. Fluids, 8: 29--42.

\par\hangindent=2em Noh, W.F. and Woodward, P. (1976). SLIC (simple line interface calculation). In Proceedings, Fifth International Conference on Fluid Dynamics (eds. A.I. van de Vooren and P.J. Zandbergen), Lecture Notes in Physics, volume 59. Berlin, Springer, Berlin, pp. 330--340.

\par\hangindent=2em Nordmark, H.O. (1991). High-order vortex methods with regridding. J. Comput. Phys., 97: 366.

\par\hangindent=2em Oguz, H. and Prosperetti, A. (1990). Bubble entrainment by the impact of drops on liquid surfaces. J. Fluid Mech., 219: 143--179.

\par\hangindent=2em Oguz, H. and Prosperetti, A. (1993). Dynamics of bubble growth and detachment from a needle. J. Fluid Mech., 257: 111--145.

\par\hangindent=2em O’Keefe, J.D. and Ahrens, T.J. (1986). Oblique impact: a process for obtaining meteorite samples from other planets. Science, 234(4774): 346.

\par\hangindent=2em Oran, E.S. and Boris, J.P. (1987). Numerical Simulation of Reactive Flow. Elsevier.

\par\hangindent=2em Oron, A., Davis, S.H., and Bankoff, S.G. (1997). Long-scale evolution of thin films. Rev. Mod. Phys., 69: 931--980.

\par\hangindent=2em Osher, S. and Fedkiw, R. (2002). Level Set Methods and Dynamic Implicit Surfaces. Springer.

\par\hangindent=2em Osher, S. and Fedkiw, R.P. (2001). Level set methods: an overview and some recent results. J. Comput. Phys., 169: 463--502.

\par\hangindent=2em Osher, S. and Sethian, J. (1988). Fronts propagating with curvature-dependent speed: algorithms based on Hamilton--Jacobi formulations. J. Comput. Phys., 79: 12--49.

\par\hangindent=2em Osher, S. and Shu, C.-W. (1991). High-order essentially nonsocillatory schemes for Hamilton--Jacobi equations. SIAM J. Numer. Anal., 28: 907--922.

\par\hangindent=2em Pan, T.W., Joseph, D.D., Bai, R., Glowinski, R., and Sarin, V. (2002). Fluidization of 1204 spheres: simulation and experiment. J. Fluid Mech., 451: 169--191.

\par\hangindent=2em Pasandideh-Fard, M., Chandra, S., and Mostaghimi, J. (2002). A three-dimensional model of droplet impact and solidification. Int. J. Heat Mass Transfer, 45: 2229-- 2242.

\par\hangindent=2em Patankar, S.V. (1980). Numerical Heat Transfer and Fluid Flow. Taylor and Francis.

\par\hangindent=2em Peng, D., Merriman, B., Zhao, H., Osher, S., and Kang, M. (1999). A PDE based fast local level set method. J. Comput. Phys., 155: 410--438.

% [原PDF p323]

\par\hangindent=2em Peregrine, D.H., Shoker, G., and Symon, A. (1990). The bifurcation of liquid bridges. J. Fluid. Mech., 212: 25--39.

\par\hangindent=2em Peskin, C.S. (1977). Numerical analysis of blood flow in the heart. J. Comput. Phys., 25: 220.

\par\hangindent=2em Peskin, C.S. (2002). The immersed boundary method. Acta Numer., 11: 479--517.

\par\hangindent=2em Peskin, C.S. and McQueen, D.M. (1989). A three-dimensional computational method for blood flow in the heart I. Immersed elastic fibers in a viscous incompressible fluid. J. Comput. Phys., 81: 372--405.

\par\hangindent=2em Peskin, C.S. and Printz, B.F. (1993). Improved volume conservation in the computation of flows with immersed boundaries. J. Comput. Phys., 105: 33--46.

\par\hangindent=2em Pilliod, J.E., Jr. and Puckett, E.G. (2004). Second-order accurate volume-of-fluid algorithms for tracking material interfaces. J. Comput. Phys., 199: 465--502.

\par\hangindent=2em Plapp, M. (2007). Three-dimensional phase-field simulations of directional solidification. J. Cryst. Growth, 303: 49--57.

\par\hangindent=2em Plateau, J. (1873). Statique Expe'rimentale et The'orique des Liquide. Gauthier- Villars.

\par\hangindent=2em Pomeau, Y. (2002). Recent progress in contact line dynamics: a review. C. R. Acad. Sci. Paris, Se'r. IIb, (Paris). Me'c., 330: 207--222.

\par\hangindent=2em Popinet, S. (2003). Gerris: a tree-based adaptive solver for the incompressible Euler equations in complex geometries. J. Comput. Phys., 190: 572--600.

\par\hangindent=2em Popinet, S. (2008). The Gerris flow solver. http://gfs.sf.net.

\par\hangindent=2em Popinet, S. (2009). An accurate adaptive solver for surface-tension-driven interfacial flows. J. Comput. Phys., 228: 5838--5866.

\par\hangindent=2em Popinet, S. and Zaleski, S. (1999). A front-tracking algorithm for accurate representation of surface tension. Int. J. Numer. Methods Fluids, 30: 775--793.

\par\hangindent=2em Popinet, S. and Zaleski, S. (2002). Bubble collapse near a solid boundary: a numerical study of the influence of viscosity. J. Fluid Mech., 464: 137--163.

\par\hangindent=2em Po¨schel, T. and Schwage, T. (2005). Computational Granular Dynamics: Models and Algorithms. Springer.

\par\hangindent=2em Powell, K. (1998). Solution of the Euler equations on solution-adaptive Cartesian grids. In Computational Fluid Dynamics Reviews, (eds. M. Hafez, and K. Oshima), volume 1. World Scientific, pp. 65--92.

\par\hangindent=2em Pozrikidis, C. (1992). Integral and Singularity Methods for Linearized Viscous Flow. Cambridge University Press.

\par\hangindent=2em Pozrikidis, C. (2001). Interfacial dynamics for Stokes flow. J. Comput. Phys., 169: 250--301.

\par\hangindent=2em Press, W.H., Flannery, B.P., Teukolsky, S.A., and Vetterling, W.T. (1993). Numerical Recipes in C: The Art of Scientific Computing. Cambridge University Press.

\par\hangindent=2em Prosperetti, A. (1981). Motion of two superposed viscous fluids. Phys. Fluids, 24: 1217--1223.

% [原PDF p324]

\par\hangindent=2em Prosperetti, A. (2007). Averaged Equations for Multiphase Flow. Cambridge University Press, pp. 237--281.

\par\hangindent=2em Prosperetti, A. and Og˜uz, H.N. (1993). The impact of drops on liquid surfaces and the underwater noise of rain. Annu. Rev. Fluid Mech., 25: 577--602.

\par\hangindent=2em Prosperetti, A. and Tryggvason, G. (2007). Computational Methods for Multiphase Flow. Cambridge University Press.

\par\hangindent=2em Puckett, E.G. (1991). A volume-of-fluid interface tracking algorithm with applications to computing shock wave refraction. In Proceedings of the 4th Int. Symp. on Computational Fluid Dynamics, Davis, CA (ed. H. Dwyer) pp. 933--938.

\par\hangindent=2em Puckett, E.G., Almgren, A.S., Bell, J.B., Marcus, D.L., and Rider, W.J. (1997). A high-order projection method for tracking fluid interfaces in variable density incompressible flows. J. Comput. Phys., 130: 269--282.

\par\hangindent=2em Pumphrey, H.C. and Elmore, P.A. (1990). The entrainment of bubbles by drop impacts. J. Fluid Mech., 220: 539--567.

\par\hangindent=2em Purvis, R. and Smith, F.T. (2005). Droplet impact on water layers: post-impact analysis and computations. Phil. Trans. R. Soc. A, 363: 1209--1221.

\par\hangindent=2em Quirk, J.J. (1994). An alternative to unstructured grids for computing gas dynamic flows around arbitrarily complex two-dimensional bodies. Comput. Fluids, 23: 125--142.

\par\hangindent=2em Radl, S., Koynov, A., Tryggvason, G., and Khinast, J.G. (2008). DNS-based prediction of the selectivity of fast multiphase reactions: hydrogenation of nitroarenes. Chem. Eng. Sci., 63: 3279--3291.

\par\hangindent=2em Rallison, J.M. and Acrivos, A. (1978). A numerical study of the deformation and burst of a viscous drop in an extensional flow. J. Fluid Mech., 89: 191--200.

\par\hangindent=2em Ranade, V. (2001). Computational Flow Modeling for Chemical Reactor Engineering. Academic Press.

\par\hangindent=2em Rayleigh, J.W. Strutt, Baron. (1879). On the capillary phenomena of jets. Proc. R. Soc. London, 29: 71--97.

\par\hangindent=2em Rayleigh, J.W. Strutt, Baron. (1900). Scientific Papers. Cambridge University Press.

\par\hangindent=2em Raynal, L. (1997). PhD thesis, Universite Joseph Fourier, Grenoble.

\par\hangindent=2em Rein, M. (1993). Phenomena of liquid drop impact on solid and liquid surfaces. Fluid Dyn. Res., 12(2): 61--93.

\par\hangindent=2em Renardy, Y. and Renardy, M. (2002). PROST: a parabolic reconstruction of surface tension for the volume-of-fluid method. J. Comput. Phys., 183: 400--421.

\par\hangindent=2em Rhie, C.M. and Chow, W.L. (1983). Numerical study of the turbulent flow past an airfoil with trailing edge separation. AIAA J., 21(11): 1525--1532.

\par\hangindent=2em Richtmyer, R.D. and Morton, K.W. (1967). Difference Methods for Initial Value Problems. Interscience, New York.

% [原PDF p325]

\par\hangindent=2em Richtmyer, R.D. (1960). Taylor instability in shock acceleration of compressible fluids. Commun. Pure Appl. Math., 13: 297--319.

\par\hangindent=2em Rider, W.J. and Kothe, D.B. (1998). Reconstructing volume tracking. J. Comput. Phys., 141: 112--152.

\par\hangindent=2em Rieber, M. and Frohn, A. (1999). A numerical study on the mechanism of splashing. Int. J. Heat Fluid Flow, 20(5): 455--461.

\par\hangindent=2em Rioboo, R., Bauthier, C., Conti, J., Voue', M., and De Coninck, J. (2003). Experimental investigation of splash and crown formation during single drop impact on wetted surfaces. Exp. Fluids, 35(6): 648--652.

\par\hangindent=2em Robinson, P.B., Blake, J.R., Kodama, T., Shima, A., and Tomita, Y. (2001). Interaction of cavitation bubbles with a free surface. J. Appl. Phys., 89(12): 8225--8237.

\par\hangindent=2em Roisman, I.V., Rioboo, R., and Tropea, C. (2002). Normal impact of a liquid drop on a dry surface: model for spreading and receding. Proc. R. Soc. Lond. Ser. A, 458(2022): 1411--1430.

\par\hangindent=2em Roisman, I.V., Horvat, K., and Tropea, C. (2006). Spray impact: rim transverse instability initiating fingering and splash, and description of a secondary spray. Phys. Fluids, 18(10): 102104--102104.

\par\hangindent=2em Rothman, D.H. and Zaleski, S. (1997). Lattice-Gas Cellular Automata: Simple Models of Complex Hydrodynamics. Cambridge University Press.

\par\hangindent=2em Rudman, M. (1997). Volume-tracking methods for interfacial flow calculations. Int. J. Numer. Methods Fluids, 24: 671--691.

\par\hangindent=2em Rybczynski, W. (1911). U¨ ber die fortschreitende Bewegung einer flu¨ssigen Kugel in einem za¨hen Medium. Bull. Acad. Sci. Crac., A: 40--46.

\par\hangindent=2em Ryskin, G. and Leal, L.G. (1984). Numerical solution of free-boundary problems in fluid mechanics. Part 2. Buoyancy-driven motion of a gas bubble through a quiescent liquid. J. Fluid Mech., 148: 19--35.

\par\hangindent=2em Saiki, E.M. and Biringen, S. (1996). Numerical simulation of a cylinder in uniform flow: application of a virtual boundary method. J. Comput. Phys., 123: 450--465.

\par\hangindent=2em Sainsbury, G.L. and Cheeseman, I.C. (1950). Effect of rain in calming the sea. Nature, 166(4210): 79.

\par\hangindent=2em Sangani, A.S. and Didwania, A.K. (1993). Dynamic simulations of flows of bubbly liquids at large Reynolds numbers. J. Fluid Mech., 250: 307--337.

\par\hangindent=2em Sankaranarayanan, K., Shan, X., Kevrekidis, I.G., and Sundaresan, S. (2002). Analysis of drag and virtual mass forces in bubbly suspensions using an implicit formulation of the lattice Boltzmann method. J. Fluid Mech., 452: 61--96.

\par\hangindent=2em Sankaranarayanan, K., Kevrekidis, I.G., Sundaresan, S., Lu, J., and Tryggvason, G. (2003). A comparative study of lattice Boltzmann and front-tracking finitedifference methods for bubble simulations. Int. J. Multiphase Flow, 29: 109-- 116.

% [原PDF p326]

\par\hangindent=2em Saville, D.A. (1997). Electrohydrodynamics: the Taylor--Melcher leaky dielectric model. Annu. Rev. Fluid Mech., 29: 27--64.

\par\hangindent=2em Scardovelli, R. and Zaleski, S. (2000). Analytical relations connection linear interfaces and volume fractions in rectangular grids. J. Comput. Phys., 164: 228--237.

\par\hangindent=2em Scardovelli, R. and Zaleski, S. (2003). Interface reconstruction with least-square fit and split Lagrangian--Eulerian advection. Int J. Numer. Methods Fluids, 41: 251--274.

\par\hangindent=2em Schlick, T. (2002). Molecular Modeling and Simulation. Springer.

\par\hangindent=2em Schmidt, A. (1996). Computations of three dimensional dendrites with finite elements. J. Comput. Phys., 126: 293--312.

\par\hangindent=2em Schultz, W.W., Huh, J., and Griffin, O.M. (1994). Potential-energy in steep and breaking waves. J. Fluid Mech, 278: 201--228.

\par\hangindent=2em Scriven, L.E. (1960). The Marangoni effects. Chem. Eng. Sci., 12: 98--108.

\par\hangindent=2em Serizawa, A., Kataoka, I., and Michiyoshi, I. (1975a). Turbulence structure of air-- water bubbly flow -- II. Local properties. Int. J. Multiphase Flow, 2: 235--246.

\par\hangindent=2em Serizawa, A., Kataoka, I., and Michiyoshi, I. (1975b). Turbulence structure of air--water bubbly flow -- III. Transport properties. Int. J. Multiphase Flow, 2: 247--259.

\par\hangindent=2em Sethian, J.A. (2001). Evolution, implementation, and application of level set and fast marching methods for advancing fronts. J. Comput. Phys., 169: 503--555.

\par\hangindent=2em Sethian, J.A. and Strain, J. (1992). Crystal growth and dendritic solidification. J. Comput. Phys., 98: 231--253.

\par\hangindent=2em Shan, X.W. and Chen, H.D. (1993). Lattice Boltzmann model for simulationg flows with multiple phases and components. Phys. Rev. E, 47: 1815--1819.

\par\hangindent=2em Shikmurzaev, Y.D. (1997). Moving contact lines in liquid/liquid/solid systems. J. Fluid Mech., 334: 211--249.

\par\hangindent=2em Shin, S. and Juric, D. (2002). Modeling three-dimensional multiphase flow using a level contour reconstruction method for front tracking without connectivity. J. Comput. Phys, 180: 427--470.

\par\hangindent=2em Shin, S., Abdel-Khalik, S.I., Daru, V., and Juric, D. (2005a). Accurate representation of surface tension using the level contour reconstruction method. J. Comput. Phys., 203: 493--516.

\par\hangindent=2em Shin, S., Abdel-Khalik, S.I., and Juric, D. (2005b). Direct three-dimensional numerical simulation of nucleate boiling using the level contour reconstruction method. Int. J. Multiphase Flow, 31: 1231--1242.

\par\hangindent=2em Shin, S.W. and Juric, D. (2000). Computation of microstructure in solidification with fluid convection. J. Mech. Behav. Mater., 11: 313.

\par\hangindent=2em Shopov, P.J., Minev, P.D., Bazhekov, I.B., and Zapryanov, Z.D. (1990). Interaction of a deformable bubble with a rigid wall at moderate Reynolds numbers. J. Fluid Mech., 219: 241--271.

% [原PDF p327]

\par\hangindent=2em Shu, C.W. and Osher, S. (1989). Efficient implementation of essentially nonoscillatory shock capturing schemes, II. J. Comput. Phys., 83: 32--78.

\par\hangindent=2em Shyy, W. and Rao, M.M. (1994). Enthalpy based formulations for phase-change problems with application to g-jitter. Microgravity Sci. Technol., 7: 41--49.

\par\hangindent=2em Shyy, W., Udaykumar, H.S., Rao, M., and Smith, R. (1996). Computational Fluid Dynamics with Moving Boundaries. Taylor and Francis.

\par\hangindent=2em Smereka, P. (1993). On the motion of bubbles in a periodic box. J. Fluid Mech., 254: 79--112.

\par\hangindent=2em So, S., Morikita, H., Takagi, S., and Matsumoto, Y. (2002). Laser doppler velocimetry measurement of turbulent bubbly channel flow. Exp. Fluids, 33: 135-- 142.

\par\hangindent=2em Son, G. and Dhir, V.K. (1997). Numerical simulation of saturated film boiling on a horizontal surface. J. Heat Transfer, 119: 525--533.

\par\hangindent=2em Son, G. and Dhir, V.K. (1998). Numerical simulation of film boiling near critical pressures with a level set method. J. Heat Transfer, 120: 183--192.

\par\hangindent=2em Son, G., Dhir, V.K., and Ramanujapu, N. (1999). Dynamics and heat transfer associated with a single bubble during nucleate boiling on a horizontal surface. J. Heat Transfer, 121: 623--631.

\par\hangindent=2em Son, G., Ramanujapu, N., and Dhir, V.K. (2002). Numerical simulation of bubble merger process on a single nucleation site during pool nucleate boiling. ASME J. Heat Transfer, 124: 51--62.

\par\hangindent=2em Song, M. and Tryggvason, G. (1999). The formation of a thick border on an initially stationary fluid sheet. Phys. Fluids, 11: 2487--2493.

\par\hangindent=2em Song, Q., Luo, R., Yang, X.Y., and Wang, Z. (2001). Phase distributions for upward laminar dilute bubbly flows with non-uniform bubble sizes in a vertical pipe. Int. J. Multiphase Flow, 27: 379--390.

\par\hangindent=2em Stow, C.D. and Hadfield, M.G. (1981). An experimental investigation of fluid flow resulting from the impact of a water drop with an unyielding dry surface. Proc. R. Soc. Lond. Ser. A, 373(1755): 419--441.

\par\hangindent=2em Subramanian, R.S. and Balasubramaniam, R. (2001). The Motion of Bubbles and Drops in Reduced Gravity. Cambridge University Press.

\par\hangindent=2em Sussman, M. (2003). A second order coupled level set and volume-of-fluid method for computing growth and collapse of vapor bubbles. J. Comput. Phys., 187: 110--136.

\par\hangindent=2em Sussman, M. and Fatemi, E. (1999). An efficient, interface preserving level set redistancing algorithm and its application to interfacial incompressible fluid flow. SIAM J. Sci. Comput., 20(4): 1165--1191.

\par\hangindent=2em Sussman, M. and Puckett, E.G. (2000). A coupled level set and volume-of-fluid method for computing 3D and axisymmetric incompressible two-phase flows. J. Comput. Phys., 162: 301--337.

% [原PDF p328]

\par\hangindent=2em Sussman, M. and Smereka, P. (1997). Axisymmetric free boundary problems. J. Fluid Mech., 341: 269--294.

\par\hangindent=2em Sussman, M., Smereka, P., and Osher, S. (1994). A level set approach for computing solutions to incompressible two-phase flows. J. Comput. Phys., 114: 146--159.

\par\hangindent=2em Sussman, M., Fatemi, E., Smereka, P., and Osher, S. (1998). An improved level set method for incompressible two-phase flows. Comput. Fluids, 27: 663--680.

\par\hangindent=2em Sussman, M., Almgren, A.S., Bell, J.B., Colella, P., Howell, L.H., and Welcome, M.L. (1999). An adaptive level set approach for incompressible two-phase flows. J. Comput. Phys., 148: 81--124.

\par\hangindent=2em Swaminathan, C.R. and Voller, V.R. (1993). On the enthalpy method. Int. J. Numer. Methods Heat Fluid Flow, 3: 233--244.

\par\hangindent=2em Swarztrauber, P.N. and Sweet, R.A. (1979). Algorithm 541: efficient Fortran subprograms for the solution of separable elliptic partial differential equations [D3]. ACM Trans. Math. Software, 5(3): 352--364.

\par\hangindent=2em Sweby, P.K. (1984). High-resolution schemes using flux limiters for hyperbolic conservation-laws. SIAM J. Numer. Anal., 21: 995--1011.

\par\hangindent=2em Takada, N., Misawa, M., Tomiyama, A., and Hosokawai, S. (2001). Simulation of bubble motion under gravity by lattice Boltzmann method. J. Nucl. Sci. Technol., 38: 330--341.

\par\hangindent=2em Takagi, S., Matsumoto, Y., and Huang, H. (1997). Numerical analysis of a single rising bubble using boundary-fitted coordinate system. JSME Int. J., Ser. B, 40: 42--50.

\par\hangindent=2em Takahira, H., Horiuchi, T., and Banerjee, S. (2004). An improved threedimensional level set method for gas--liquid two-phase flows. J. Fluids Eng., 126: 578.

\par\hangindent=2em Takewaki, H. and Yabe, T. (1987). The cubic-interpolated pseudo particle (CIP) method: application to nonlinear and multi-dimensional hyperbolic equations. J. Comput. Phys., 70: 355--372.

\par\hangindent=2em Takewaki, H., Nishiguchi, A., and Yabe, T. (1985). Cubic interpolated pseudoparticle method (CIP) for solving hyperbolic-type equations. J. Comput. Phys., 61: 261--268.

\par\hangindent=2em Tan, L. and Zabaras, N. (2007). A level set simulation of dendritic solidification of multi-component alloys. J. Comput. Phys., 221: 9--40.

\par\hangindent=2em Tasoglu, S., Demirci, U., and Muradoglu, M. (2008). The effect of soluble surfactant on the transient motion of a buoyancy-driven bubble. Phys. Fluids, 20: 040805.

\par\hangindent=2em Taylor, G.I. (1963). Generation of Ripples by Wind Blowing over a Viscous Liquid. Cambridge University Press, Cambridge, UK, pp. 244--254.

% [原PDF p329]

\par\hangindent=2em Taylor, G.I. (1966). Studies in electrohydrodynamics. I. The circulation produced in a drop by electrical field. Proc. R. Soc. Lond. Ser. A, 291: 159--166.

\par\hangindent=2em Thompson, J.F., Soni, B.K., and Weatherill, N.P. (1998). Handbook of Grid Generation. CRC Press.

\par\hangindent=2em Thomson, J.J. and Newall, H.F. (1885). On the formation of vortex rings by drops falling into liquids, and some allied phenomena. Proc. R. Soc. Lond., 39: 417-- 436.

\par\hangindent=2em Thoroddsen, S.T. (2002). The ejecta sheet generated by the impact of a drop. J. Fluid Mech., 451: 373--381.

\par\hangindent=2em Tomar, G., Biswas, G., Sharma, A., and Agrawal, A. (2005). Numerical simulation of bubble growth in film boiling using a coupled level-set and volume-of-fluid method. Phys. Fluids, 17: 112103.

\par\hangindent=2em Tomar, G., Gerlach, D., Biswas, G., Alleborn, N., Sharma, A., Durst, F., S.W.J. Welch, and Delgado, A. (2007). Two-phase electrohydrodynamic simulations using a volume-of-fluid approach. J. Comput. Phys., 227: 1267--1285.

\par\hangindent=2em Tomar, G., Biswas, G., Sharma, A., and Welch, S.W.J. (2008). Multi-mode analysis of bubble growth in saturated film boiling. Phys. Fluids, 20: 092101.

\par\hangindent=2em Tomar, G., Biswas, G., Sharma, A., and Welch, S.W.J. (2009). Influence of electric field on saturated film boiling. Phys. Fluids, 21: 032107.

\par\hangindent=2em Tomar, G., Fuster, D., Zaleski, S., and Popinet, S. (2010). Multiscale simulations of primary atomization using Gerris. Comput. Fluids, 39(10): 1864--1874.

\par\hangindent=2em Tomiyama, A., Zun, I., Sou, A., and Sakaguchi, T. (1993). Numerical analysis of bubble motion with the VOF method. Nucl. Eng. Des., 141: 69--82.

\par\hangindent=2em Tomiyama, A., Tamai, H., Zun, I., and Hosokawa, S. (2002). Transverse migration of single bubbles in simple shear flows. Chem. Eng. Sci., 57: 1849--1858.

\par\hangindent=2em Tonhardt, R. and Amberg, G. (1998). Phase-field simulations of dendritic growth in a shear flow. J. Cryst. Growth, 194: 406--425.

\par\hangindent=2em Torres, D.J and Trujillo, M.F. (2006). KIVA-4: an unstructured ALE code for compressible gas flow with sprays. J. Comput. Phys., 219(2): 943--975.

\par\hangindent=2em Tryggvason, G. and Aref, H. (1983). Numerical experiments on Hele--Shaw flow with a sharp interface. J. Fluid Mech., 136: 1--30.

\par\hangindent=2em Tryggvason, G. and Aref, H. (1985). Finger interaction mechanisms in stratified Hele--Shaw flow. J. Fluid Mech., 154: 284--301.

\par\hangindent=2em Tsimplis, M. and Thorpe, S.A. (1989). Wave damping by rain. Nature, 342(6252): 893--895.

\par\hangindent=2em Udaykumar, H.S., Kan, H.C., Shyy, W., and Tran-Son-Tay, R. (1997). Multiphase dynamics in arbitrary geometries on fixed Cartesian grids. J. Comput. Phys., 137: 366--405.

% [原PDF p330]

\par\hangindent=2em Udaykumar, H.S., Mittal, R., and Shyy, W. (1999). Computation of solid--liquid phase fronts in the sharp interface limit on fixed grids. J. Comput. Phys., 153: 535--574.

\par\hangindent=2em Udaykumar, H.S., Mittal, R., Rampunggoon, P., and Khanna, A. (2001). A sharp interface Cartesian grid method for simulating flows with complex moving boundaries. J. Comput. Phys., 174: 345--380.

\par\hangindent=2em Ungar, L.H. and Brown, R.A. (1985). Cellular interface morphologies in directional solidification. IV. The formation of deep cells. Phys. Rev. B, 31: 5931-- 5940.

\par\hangindent=2em Unverdi, S.O. and Tryggvason, G. (1992). A front-tracking method for viscous, incompressible, multi-fluid flows. J. Comput. Phys., 100: 25--37.

\par\hangindent=2em van der Vorst, H.A. (1992). Bi-CGSTAB: a fast and smoothly converging variant of Bi-CG for the solution of nonsymmetric linear systems. SIAM J. Sci. Stat. Comput., 13: 631.

\par\hangindent=2em van der Vorst, H.A. (2003). Iterative Krylov Methods for Large Linear Systems. Cambridge University Press.

\par\hangindent=2em van Leer, B. (1979). Towards the ultimate conservative difference scheme. V. A second-order sequel to Godunov’s method. J. Comput. Phys., 32: 101--136.

\par\hangindent=2em van Sint Annaland, M., Dijkhuizen, W., Deen, N.G., and Kuipers, J.A.M. (2006). Numerical simulation of gas bubbles behaviour using a 3D front tracking method. AIChE J., 52: 99--110.

\par\hangindent=2em Verschueren, M., van de Vosse, F.N., and Meijer, H.E.H. (2001). Diffuse-interface modelling of thermo-capillary flow instabilities in a Hele--Shaw cell. J. Fluid Mech., 434: 153--166.

\par\hangindent=2em Viecelli, J.A. (1969). A method for including arbitrary external boundaries in the MAC incompressible fluid computing technique. J. Comput. Phys., 4: 543--551.

\par\hangindent=2em Villermaux, E. (1998). On the role of viscosity in shear instabilities. Phys. Fluids, 10(2): 368--373.

\par\hangindent=2em Villermaux, E. and Clanet, C. (2002). Life of a flapping liquid sheet. J. Fluid Mech., 462: 341--363.

\par\hangindent=2em Villermaux, E., Marmottant, Ph., and Duplat, J. (2004). Ligament-mediated spray formation. Phys. Rev. Lett., 92(7): 074501.

\par\hangindent=2em Vinje, T. and Brevig, P. (1981). Numerical simulations of breaking waves. Adv. Water Resources, 4: 77--82.

\par\hangindent=2em Wang, S.K., Lee, S.J., Jones, O.C., Jr. and Lahey, R.T. Jr. (1987). 3-D turbulence structure and phase distribution in bubbly two-phase flows. Int. J. Multiphase Flow, 13: 327--343.

\par\hangindent=2em Wang, Y., Lu, X., Zhuang, L., Tang, Z., and Hu, W. (2004). Numerical simulation of drop Marangoni migration under microgravity. Acta Astronaut., 54: 325--335.

% [原PDF p331]

\par\hangindent=2em Wanga, J., Lua, P., Wanga, Z., Chao, C., and Maoa, Z.-S. (2008). Numerical simulation of unsteady mass transfer by the level set method. Chem. Eng. Sci., 63: 3141--3151.

\par\hangindent=2em Weatherburn, C.E. (1927). Differential Geometry of Three Dimensions, Vol. I. Cambridge University Press, Cambridge.

\par\hangindent=2em Welch, S.W.J. (1995). Local simulation of two-phase flows including interface tracking with mass transfer. J. Comput. Phys., 121: 142--154.

\par\hangindent=2em Welch, S.W.J. and Biswas, G. (2007). Direct simulation of film boiling including electrohydrodynamic forces. Phys. Fluids, 19: 012106.

\par\hangindent=2em Welch, S.W.J. and Wilson, J. (2000). A volume of fluid based method for fluid flows with phase change. J. Comput. Phys., 160: 662--682.

\par\hangindent=2em Werder, T., Walther, J.H., and Koumoutsakos, P. (2005). Hybrid atomisticcontinuum method for the simulation of dense fluid flow. J. Comput. Phys., 205: 373--390.

\par\hangindent=2em Wesseling, P. (1992). An Introduction to Multigrid Methods. Wiley, New York.

\par\hangindent=2em Wesseling, P. (2001). Principles of Computational Fluid Dynamics. Springer- Verlag Berlin.

\par\hangindent=2em Wheeler, A.A., Murray, B.T., and Schaefer, R.J. (1993). Computations of dendrites using a phase-field model. Phys. D, 66: 243--262.

\par\hangindent=2em Wheeler, A.A., Ahmad, N.A., Boettinger, W.J., Braun, R.J., McFadden, G.B., and Murray, B.T. (1995). Recent development in phase-field models of solidification. Adv. Space Res., 16: 163--172.

\par\hangindent=2em Worthington, A.M. (1908). A Study of Splashes. Longmans, Green and Co., London.

\par\hangindent=2em Xue, M., Xu, H.B., Liu, Y.M., and Yue, D.K.P. (2001). Computations of fully nonlinear three-dimensional wave--wave and wave--body interactions. Part 1. Dynamics of steep three-dimensional waves. J. Fluid Mech., 438: 11--39.

\par\hangindent=2em Yabe, T., Xiao, F., and Utsumi, T. (2001). The constrained interpolation profile (CIP) method for multi-phase analysis. J. Comput. Phys., 169: 556--593.

\par\hangindent=2em Yanenko, N.N. (1971). The Method of Fractional Steps for Solving Multidimensional Problems of Mathematical Physics in Several Variables. Springer-Verlag, Berlin.

\par\hangindent=2em Yang, C. and Maoa, Z.-S. (2005). Numerical simulation of interphase mass transfer with the level set approach. Chem. Eng. Sci., 60: 2643--2660.

\par\hangindent=2em Yang, X., Feng, J.J., Liu, C., and Shen, J. (2006a). Numerical simulations of jet pinching-off and drop formation using an energetic variational phase-field method. J. Comput. Phys., 218(1): 417--428.

\par\hangindent=2em Yang, X., James, A.J., Lowengrub, J., Zheng, X., and Cristini, V. (2006b). An adaptive coupled level-set/volume-of-fluid interface capturing method for unstructured triangular grids. J. Comput. Phys., 217: 364--394.

% [原PDF p332]

\par\hangindent=2em Yang, Y. and Udaykumar, H.S. (2005). Sharp interface Cartesian grid method III: solidification of pure materials and binary solutions. J. Comput. Phys., 210: 55--74.

\par\hangindent=2em Yarin, A.L. (2006). Drop impact dynamics: splashing, spreading, receding, bouncing. Annu. Rev. Fluid Mech., 38: 159--192.

\par\hangindent=2em Yarin, A.L. and Weiss, D.A. (1995). Impact of drops on solid surfaces: self-similar capillary waves, and splashing as a new type of kinematic discontinuity. J. Fluid Mech., 283: 141--173.

\par\hangindent=2em Yecko, P. and Zaleski, S. (2005). Transient growth in two-phase mixing layers. J. Fluid Mech., 528: 43--52.

\par\hangindent=2em Yecko, P., Zaleski, S., and Fullana, J.-M. (2002). Viscous modes in two-phase mixing layers. Phys.Fluids, 14: 4115--4122.

\par\hangindent=2em Youngren, G.K. and Acrivos, A. (1976). On the shape of a gas bubble in a viscous extensional flow. J. Fluid Mech., 76: 433.

\par\hangindent=2em Youngs, D.L. (1982). Time dependent multimaterial flow with large fluid distortion. In Numerical Methods for Fluid Dynamics (eds. K.M. Morton and M.J. Baines) Academic Press, New York, pp. 27--39.

\par\hangindent=2em Youngs, D.L. (1984). An interface tracking method for a 3D Eulerian hydrodynamics code. Technical report, AWRE. Technical Report 44/92/35.

\par\hangindent=2em Yu, D. and Tryggvason, G. (1990). The free surface signature of unsteady, twodimensional vortex flows. J. Fluid Mech., 218: 547--572.

\par\hangindent=2em Yu, P.-W., Ceccio, S.L., and Tryggvason, G. (1995). The collapse of a cavitation bubble in shear flows -- a numerical study. Phys. Fluids, 7: 2608--2616.

\par\hangindent=2em Zaleski, S., Zanetti, G., Scardovelli, R., et al. (1992--2008). SURFER code. The code is covered by the GNU Public Licence and may be downloaded from \url{http://www.lmm.jussieu.fr/~zaleski}.

\par\hangindent=2em Zeff, B.W., Kleber, B., Fineberg, J., and Lathrop, D.P. (2000). Singularity dynamics in curvature collapse and jet eruption on a fluid surface. Nature, 403: 401--404.

\par\hangindent=2em Zenit, R., Koch, D.L., and Sangani, A.S. (2001). Measurements of the average properties of a suspension of bubbles rising in a vertical channel. J. Fluid Mech., 429: 307--342.

\par\hangindent=2em Zhang, D.Z. and Prosperetti, A. (1994). Ensemble phase-averaged equations for bubbly flows. Phys. Fluids, 6: 2956--2970.

\par\hangindent=2em Zhang, J., Eckmann, D.M., and Ayyaswamy, P.S. (2006). A front tracking method for a deformable intravascular bubble in a tube with soluble surfactant transport. J. Comput. Phys., 214: 366--396.

\par\hangindent=2em Zhou, H. and Pozrikidis, C. (1993). The flow of ordered and random suspensions of two-dimensional drops in a channel. J. Fluid Mech., 255: 103--127.

\par\hangindent=2em Zinchenko, A.Z. and Davis, R.H. (2000). An efficient algorithm for hydrodynamical interaction of many deformable drops. J. Comput. Phys, 157: 539--587.

\endgroup

% reference_entries=450

\chapter*{中英双语索引}
\addcontentsline{toc}{chapter}{中英双语索引}
\markboth{中英双语索引}{中英双语索引}
\begingroup\small
\setlength{\parindent}{0pt}
\setlength{\parskip}{0.25em}
% [原PDF p333]
附加质量（\textit{added mass}），203\par
平流项（\textit{advection term}）\par
\quad 定义（\textit{definition}），51\par
\quad 离散化（\textit{discretization}），59, 61\par
自适应网格加密（AMR）（\textit{AMR (adaptive mesh refinement)}），19, 22, 71, 72, 225, 227, 232\par
解析解（\textit{analytical solutions}），1, 180, 181, 198, 201, 214, 285, 287, 292\par
雾化（\textit{atomization}），3, 4, 19, 72, 204, 205, 211, 212, 214, 215, 219, 222, 226, 228\par
\quad 同向伴流（\textit{co-flowing}），4, 215\par
\quad 喷嘴的影响（\textit{effect of nozzle}），214, 219, 222, 223\par
不良单元（\textit{bad elements}）\par
\quad 前沿追踪中（\textit{in front tracking}），145\par
Berenson 关联式（\textit{Berenson’s correlation}），252\par
双共轭梯度稳定法（\textit{Bi-Conjugate Gradient STABilized method}），68\par
BiCGSTAB 方法（\textit{BiCGSTAB method}），68\par
沸腾（\textit{boiling}），248\par
\quad 池式（\textit{pool}），251\par
\quad 对流（\textit{convection}），249\par
\quad 膜态（\textit{film}），248, 249, 251, 252\par
\quad 流动（\textit{flow}），249\par
\quad 成核（\textit{nucleate}），254\par
\quad 池式（\textit{pool}），249\par
Bond 数（\textit{Bond number}），43\par
边界条件（\textit{boundary condition}），29, 30, 37, 40, 50, 67, 69, 71, 141, 189\par
\quad 自由表面条件（\textit{for free surface}），40\par
\quad 无滑移（\textit{no slip}），37, 69\par
\quad 出流（\textit{outflow}），29, 30, 66, 70\par
\quad 部分滑移（\textit{partial slip}），47\par
\quad 周期条件（\textit{periodic}），30, 71, 146, 189, 191, 195, 219\par
边界积分法（\textit{boundary integral method}），8--11\par
液滴破碎（\textit{breakup of droplet}）\par
\quad 在线性剪切流中（\textit{in linear shear flow}），9, 10\par
\quad 冲击或碰撞下（\textit{under impact or collision}），228, 230\par
边缘液环破碎（\textit{breakup of rim}），205, 210, 211\par
细液丝破碎（\textit{breakup of thin threads}），158, 206--208\par
气泡柱（\textit{bubble columns}），19, 187\par
气泡（\textit{bubbles}）\par
\quad 上升（\textit{rising}），2, 4, 179, 189, 193, 196\par
\quad 小尺度（\textit{small}），203\par
气泡流（\textit{bubbly flows}）\par
均匀气泡流（\textit{homogeneous}），189--191\par
浮力（\textit{buoyancy}），203\par
毛细数（\textit{capillary number}），43\par
CFL（Courant--Friedrichs--Lewy）条件（\textit{CFL (Courant, Friedrichs, and Lewy) condition}），53, 54, 109\par
VOF 方法中的 CFL 数（\textit{CFL number in VOF method}），109, 123--125\par
CIP（约束插值传播）方法（\textit{CIP (constrained interpolated propagation) method}），14, 75, 77, 91--93\par
封闭关系（\textit{closure relations}），187, 203\par
CLSVOF（耦合水平集/VOF）方法（\textit{CLSVOF (coupled level-set volume-of-fluid) method}），128\par
聚并（\textit{coalescence}），8, 158, 159, 230, 232, 233\par
相干结构（\textit{coherent structures}），223\par
同位网格（\textit{collocated grid}），56, 58, 73\par
颜色函数（\textit{color function}），76--78, 80, 82, 95--98, 128, 129, 131, 161, 167, 184\par
燃烧（\textit{combustion}），4, 205, 228\par
复杂物理（\textit{complex physics}），17, 19\par
撞击坑（\textit{craters}），5, 229--231, 233--235\par
冠状飞溅（\textit{crown splash}），见“液滴撞击”\par
CSF（连续表面力）方法（\textit{CSF (continuous surface force) method}），13, 14, 161, 162\par
CSS（连续表面应力）方法（\textit{CSS (continuous surface stress) method}），161, 164--166, 169, 177--179, 181\par
曲率（\textit{curvature}）\par
由高度函数计算（\textit{computed from height function}），184--186, 271\par
前沿追踪中的计算（\textit{computing in front tracking}），169--171, 175, 176\par
VOF 中的计算（\textit{computing in VOF}），163, 164, 167, 168, 184--186\par
定义（\textit{definition}），32\par
表达式（\textit{expression for}），34, 35, 276\par
Gauss 曲率（\textit{Gauss}），274\par
轴对称几何中（\textit{in axisymmetric geometry}），275\par
平均曲率（\textit{mean}）\par
定义（\textit{definition}），33, 34, 169, 273, 274, 277\par
总曲率（\textit{total}），274\par
数据结构（\textit{data structure}）\par
前沿追踪中（\textit{in front tracking}），85, 86, 133, 134\par
枝晶（\textit{dendrites}），264, 265\par
枝晶结构（\textit{dendritic structure}），6, 264\par
密度比（\textit{density ratio}），202\par
% [原PDF p334]
分离压力（\textit{disjoining pressure}），44, 45, 48\par
跳跃条件（\textit{jump conditions}），37, 39, 40\par
直接数值模拟（DNS）（\textit{DNS (direct numerical simulation)}），2, 3, 18, 19, 188, 195, 201, 203, 251\par
尾随、接触与翻滚（\textit{drafting, kissing, and tumbling}），18, 192\par
滴（\textit{drop}），见“液滴”\par
液滴形成（\textit{droplet formation}），213, 215, 225\par
液滴撞击（\textit{droplet impact}）\par
\quad 花冠状液片（\textit{corolla}），213, 235, 236, 238--242\par
\quad 冠状飞溅（\textit{crown splash}），232, 235, 240--242\par
\quad 射流形成（\textit{jet formation}），236\par
\quad 即时飞溅（\textit{prompt splash}），240\par
Eötvös 数（\textit{Eo¨tvo¨s number}），43, 188\par
\quad 小尺度（\textit{small}），203\par
电流体动力学（\textit{electrohydrodynamics}），244\par
能量守恒（\textit{energy conservation}），25\par
ENO（本质无振荡）格式（\textit{ENO (essentially non-oscillating) scheme}），60, 61, 88\par
《幻想曲》（迪士尼动画）（\textit{Fantasia (Disney animation)}），229\par
Faraday 波（\textit{Faraday waves}），234\par
有限体积法（\textit{Finite-volume method}），26, 50, 55, 86\par
鱼骨形态（\textit{fishbone patterns}），226\par
漂浮碎片与喷射碎片（\textit{floatsam and jetsam}），12, 83, 116\par
通量修正输运法（\textit{flux-corrected transport method}），80\par
Fourier 定律（\textit{Fourier’s law}），26\par
自由表面（\textit{free surface}），1, 8, 11\par
\quad 其边界条件（\textit{boundary condition at}），40, 45\par
\quad 其上的气泡破裂（\textit{bubble bursting at}），204, 234, 235\par
界面捕捉法（\textit{front capturing method}），91\par
前沿--网格通信（\textit{front-grid communications}），145--148\par
\quad 权重函数（\textit{weighting functions}），148\par
前沿追踪/VOF 混合法（\textit{front-tracking volume-of-fluid method}），129\par
Froude 数（\textit{Froude number}），44\par
Galileo 数（\textit{Galileo number}），43, 188\par
气液流（\textit{gas–liquid flows}），1, 6, 189, 291\par
Gerris 程序（\textit{Gerris code}），61\par
Gibbs--Thomson 条件（\textit{Gibbs–Thompson condition}），263\par
GMRES 方法（\textit{GMRES method}），68\par
Godunov 方法（\textit{Godunov method}），80\par
网格生成（\textit{grid generation}），17\par
H 模式（\textit{H-mode}），219\par
Hamaker 常数（\textit{Hamaker’s constant}），44, 45\par
传热（\textit{heat transfer}），243\par
Heaviside（阶跃）函数（\textit{Heaviside or step function}），35, 41, 76, 77, 279\par
高度函数法（\textit{height-function method}），179, 181, 184\par
人的因素（\textit{human aspect}），20\par
混合方法（\textit{hybrid methods}），128\par
\quad CLSVOF，见“CLSVOF（耦合水平集/VOF）方法”（\textit{CLSVOF, see CLSVOF (coupled level-set volume-of-fluid) method}）\par
\quad 前沿追踪 VOF，见“前沿追踪/VOF 混合法”（\textit{front-tracking VOF, see front-tracking volume-of-fluid method}）\par
浸入边界法（\textit{immersed-boundary method}），14, 86, 149\par
指示函数（\textit{indicator function}），68, 76, 77\par
工业应用（\textit{industrial applications}），1, 3, 187\par
惯性（\textit{inertia}），203\par
界面（\textit{interface}）\par
\quad 几何描述（\textit{geometrical description}），30, 31, 47, 270\par
分子间作用力（\textit{intermolecular forces}），21, 22, 44, 45, 48, 208\par
跳跃记号（\textit{jump notation}），38\par
运动学奇性（\textit{kinematic singularity}），213\par
运动黏度（\textit{kinematic viscosity}），29\par
Krylov 方法（\textit{Krylov method}），68\par
Lagrangian 网格（\textit{Lagrangian grids}），17\par
Laplace 数（\textit{Laplace number}），43\par
Laplace 定律（\textit{Laplace’s law}），45, 163, 177\par
大尺度结构（\textit{large-scale structures}）\par
雾化中（\textit{in atomization}），214\par
气泡流中（\textit{in bubbly flows}），192\par
格子 Boltzmann 法（LBM）（\textit{LBM (lattice–Boltzmann method)}），18\par
水平集法（\textit{level-set method}），14, 15, 75, 86, 87, 89, 90, 96, 128, 175, 232, 264\par
升力（\textit{lift force}），193, 196, 197, 199\par
链表（\textit{linked list}），139--141, 143, 154\par
液体燃料（\textit{liquid fuel}），205\par
生命系统（\textit{living systems}），1, 6\par
低 Reynolds 数（\textit{low Reynolds numbers}），53, 190\par
润滑近似（\textit{lubrication approximation}），206, 207, 210\par
MAC（标记--单元）法（\textit{MAC (marker-and-cell) method}），11--14, 51, 56, 59, 85, 109, 114, 121, 122, 125, 163, 229\par
标记函数（\textit{marker function}），12--16, 34, 68, 74--77, 80--82, 84, 86--88, 92--94, 97, 115, 133, 139, 150, 152--158, 161, 162, 174, 175\par
标记粒子（前沿追踪）（\textit{marker particle (in front tracking)}），14, 85, 86, 133, 151, 152\par
标记粒子（体域内）（\textit{marker particle (in volume)}），12, 229\par
质量守恒（\textit{mass conservation}），14, 22--24, 26, 28, 37, 41, 89, 90, 97, 112, 123, 136, 143, 206, 249\par
VOF 中（\textit{in VOF}），97, 112, 123\par
平均曲率（\textit{mean curvature}），见“曲率”\par
微观结构（\textit{microstructure}），6, 20, 188, 262--265\par
混合层（\textit{mixing layer}），205, 215, 216, 220, 291\par
稳定性理论（\textit{stability theory}），216\par
Kelvin--Helmholtz 不稳定性（\textit{Kelvin–Helmholtz}），216, 288\par
Orr--Sommerfeld 理论（\textit{Orr–Sommerfeld}），218\par
瞬态增长理论（\textit{transient-growth theory}），224\par
动量守恒（\textit{momentum conservation}），24\par
含界面时（\textit{with interfaces}），39\par
月球撞击坑（\textit{moon craters}），228\par
中央峰（\textit{central peak}），229\par
Morton 数（\textit{Morton number}），43, 188\par
多重网格方法（\textit{multigrid methods}），67\par
多重网格求解器（\textit{multigrid solvers}），67\par
Navier--Stokes 方程（\textit{Navier–Stokes equations}），11, 21\par
定义（\textit{definition}），25\par
Newton 流体（\textit{Newtonian fluids}），22, 26, 28, 29, 40, 42\par
定义（\textit{definition}），24\par
非 Newton 流体（\textit{non-Newtonian fluids}），20\par
非均匀表面张力（\textit{nonuniform surface tension}），244\par
成核沸腾（\textit{nucleate boiling}），253\par
Nusselt 数（\textit{Nusselt number}），252\par
海洋（\textit{oceans}），1, 3, 5, 204\par
Ohnesorge 数（\textit{Ohnesorge number}），43, 210, 231\par
单流体方法（\textit{one-fluid approach}），15--17, 21, 41, 42, 54, 161\par
% [原PDF p335]
概率密度函数（PDF）（\textit{PDF (probability distribution functions)}），226\par
相变（\textit{phase change}），5, 19, 37--39, 41, 91, 248\par
相场法（\textit{phase-field method}），15, 48, 90, 91, 264, 265\par
Plateau--Rayleigh 射流不稳定性（毛细不稳定性）（\textit{Plateau--Rayleigh jet instability (capillary instability)}），4, 205, 206, 213, 214, 242\par
Poisson 方程（\textit{Poisson equation}），52, 64\par
势流（\textit{potential flow}）\par
\quad 液滴撞击中（\textit{in droplet impact}），237\par
势流模拟（\textit{potential flow simulations}）\par
\quad 气泡流中（\textit{in bubbly flows}），188, 190, 192\par
压力方程（\textit{pressure equation}），52, 64\par
压力求解器（\textit{pressure solver}），66\par
投影法（\textit{projection method}），12, 38, 51--53, 59\par
PROST（表面张力适当表示）方法（\textit{PROST (proper representation of surface tension) method}），165, 166, 179, 181, 183, 184\par
QUICK（二次迎风对流运动学插值）格式（\textit{QUICK (quadratic upstream interpolation for convective kinematics) scheme}），60, 61\par
雨（\textit{rain}），1, 4, 5, 228\par
Rayleigh--Taylor 不稳定性（\textit{Rayleigh--Taylor instability}），8, 11, 12, 85, 91\par
\quad 雾化中（\textit{in atomization}），224\par
重分布（\textit{redistribution}）\par
\quad VOF 中（\textit{in VOF}），120, 122, 130\par
\quad 前沿追踪标记点的（\textit{of markers in front tracking}），137\par
重新初始化（\textit{reinitialization}）\par
\quad 水平集法中（\textit{in level-set method}），14, 88, 89\par
Reynolds 数（\textit{Reynolds number}），42\par
Richtmyer--Meshkov 不稳定性（\textit{Richtmyer--Meshkov instability}）\par
\quad 液滴撞击中（\textit{in droplet impact}），241, 242\par
边缘液环（\textit{rim}），205, 209, 210, 212--214\par
Saffman--Taylor 不稳定性（\textit{Saffman--Taylor instability}），10\par
自相似解（\textit{self-similar solutions}），48, 208, 234, 235\par
尖锐界面极限（\textit{sharp interface limit}），15, 16, 21, 22, 34, 37, 42, 44, 73, 80, 81, 91, 93\par
奇异项（\textit{singular terms}），41\par
SLIC（简单线界面计算）方法（\textit{SLIC (simple line interface calculation) method}），82--84, 96, 116\par
SOR（逐次超松弛）方法（\textit{SOR (successive over-relaxation) method}），66--68\par
球体（\textit{spheres}）\par
\quad 下落（\textit{falling}），18\par
飞溅（\textit{splash}），4, 5, 11, 13, 210, 211, 213, 228, 230, 231, 240\par
\quad 即时飞溅（\textit{prompt}），240\par
喷雾形成（\textit{spray formation}），204\par
交错网格（MAC 网格）（\textit{staggered grid (MAC grid)}），12, 54, 56--58, 65, 69, 74, 114, 121, 125, 146, 153, 174\par
阶跃函数（\textit{step function}）\par
\quad 前沿追踪中的近似（\textit{approximation in front tracking}），84, 166\par
曲面几何（\textit{surface geometry}）\par
\quad 曲率（\textit{curvature}），见“曲率”\par
\quad 主方向（\textit{principal directions}），274\par
表面张力（\textit{surface tension}），203\par
\quad 轴对称情形（\textit{axisymmetry}），168\par
\quad 由标记函数计算（\textit{computed from marker functions}），161--167, 184, 185\par
\quad 前沿追踪中的计算（\textit{computed in front tracking}），168, 169, 171--177\par
\quad 定义（\textit{definition}），22, 38\par
\quad 力的定义（\textit{force (definition)}），42\par
\quad 表面张力张量（\textit{surface tension tensor}），38, 39, 164\par
表面活性剂（\textit{surfactants}）\par
\quad 气泡动力学中（\textit{in bubble dynamics}），189\par
悬浮体系（\textit{suspensions}），6, 8, 10\par
Taylor--Culick 边缘液环（\textit{Taylor--Culick rim}），见“边缘液环”\par
热毛细运动（\textit{thermocapillary motion}），244\par
薄膜（\textit{thin films}），19, 31, 44, 48, 206--208, 228\par
\quad 薄膜上的飞溅（\textit{splashing on}），235\par
时间积分（\textit{time integration}），51, 53\par
拓扑改变（\textit{topology change}），8, 44, 158, 159, 230, 232, 236\par
\quad 前沿追踪中（\textit{in front tracking}），87, 134, 158, 159\par
\quad VOF 中（\textit{in VOF}），96\par
单位法向（\textit{unit normal}）\par
\quad 定义（\textit{definition}），31--33\par
非结构网格（\textit{unstructured grids}），18, 85\par
\quad 前沿追踪法中（\textit{in front-tracking method}），135, 143\par
\quad VOF 中（\textit{in VOF}），131\par
黏性项（\textit{viscous terms}）\par
\quad 离散化（\textit{discretization}），61, 64\par
\quad 调和平均（\textit{harmonic mean}），63, 64, 221\par
VOF（流体体积）法（\textit{VOF (volume-of-fluid) method}），12--15, 72, 81, 83, 84, 88--90, 94--98, 103, 124, 128, 129, 150, 161, 163, 165, 166, 169, 175, 182--184, 236, 242, 249\par
\quad 平流（\textit{advection}），108\par
\quad Eulerian 隐式（\textit{Eulerian implicit}），115, 116\par
\quad Lagrangian 显式（\textit{Lagrangian explicit}），115, 116\par
\quad 非分裂（\textit{unsplit}），116\par
\quad 法向估计（\textit{normal evaluation}）\par
\qquad 中心列法（\textit{centered columns method}），100\par
\qquad ELVIRA（\textit{ELVIRA}），101--104, 107, 108, 123--125, 127, 128\par
\qquad 最小二乘拟合法（\textit{least-squares fit method}），103\par
\qquad Youngs 法（\textit{Youngs' method}），99, 103, 104, 107, 108\par
\quad 重构（\textit{reconstruction}），12, 83, 84, 95--99, 101, 103, 104, 106--108, 112, 116, 122--129, 183\par
\quad 细丝状碎屑（\textit{wisps}），116\par
体积变化（\textit{volume change}），248, 262, 263\par
体积膨胀（\textit{volume expansion}），250, 251\par
体积分数（\textit{volume fraction}）\par
\quad 前沿追踪中（\textit{in front tracking}），155\par
界面体积源（\textit{volume source at the interface}），41\par
盒中涡测试（\textit{vortex-in-a-box test}），123, 125--127, 131, 132, 151, 152\par
壁面层（\textit{wall layer}）\par
\quad 气泡流中（\textit{in bubbly flow}），196, 201\par
Weber 数（\textit{Weber number}），43\par
加权 ENO 格式（\textit{weighted ENO scheme}），128\par
润湿（\textit{wetting}），46, 48, 229, 230\par
细丝状碎屑（\textit{wisps}），116\par
Worthington 射流（\textit{Worthington jet}），229\par
\endgroup
% index_entries=232
\end{document}
